\documentclass{article}
\usepackage[margin=0.6in]{geometry}
\usepackage{graphicx} % Required for inserting images
\usepackage{amsmath}
\usepackage{amssymb}
\usepackage{pdflscape}
\usepackage{longtable}
\usepackage{array}
\usepackage{booktabs}
\usepackage{makecell}
\newtheorem{definition}{Definition}

\setlength{\tabcolsep}{3pt}
\renewcommand{\arraystretch}{1.12}




\title{PGD attack}
\author{yifeisun }
\date{May 2026}

\begin{document}

\maketitle

\section{Planned Experiments for Cross-Framework PDE Solver--Model Attacks}

This section summarizes the next set of experiments for studying adversarial attacks on PDE surrogate models. The main goal is to compare numerical solvers and neural surrogate models implemented in JAX and PyTorch, under different attack losses, different optimization strategies, and different ways of measuring runtime and GPU memory usage.

The experiments focus on two PDE settings:
\[
\text{1D Burgers equation},
\qquad
\text{2D Stokes / Navier--Stokes equation}.
\]
For each PDE case, both the numerical solver and the surrogate model have JAX and PyTorch implementations. The central question is how the framework choice, the loss definition, and the attack optimization method affect the true adversarial error:
\[
\mathcal{L}^{\mathrm{true}}(\delta)
=
\left\|
f(x+\delta) - g(x+\delta)
\right\|^2,
\]
where \(f\) denotes the surrogate model, \(g\) denotes the numerical solver, \(x\) is the original input, and \(\delta\) is the adversarial perturbation.

\subsection{Cross-Framework Solver--Model PGD Attack Experiments}

The first direction is to systematically compare four combinations of solver and surrogate model implementations:
\[
\begin{aligned}
&\text{Solver}_{\mathrm{JAX}} + \text{Model}_{\mathrm{JAX}}, \\
&\text{Solver}_{\mathrm{JAX}} + \text{Model}_{\mathrm{PyTorch}}, \\
&\text{Solver}_{\mathrm{PyTorch}} + \text{Model}_{\mathrm{JAX}}, \\
&\text{Solver}_{\mathrm{PyTorch}} + \text{Model}_{\mathrm{PyTorch}}.
\end{aligned}
\]
These four combinations should be tested separately on both the 1D Burgers case and the 2D Stokes / Navier--Stokes case. Therefore, the full experimental matrix contains
\[
2 \times 4 = 8
\]
main solver--model attack settings.

For each setting, the attack input is
\[
x_{\mathrm{adv}} = x + \delta.
\]
The perturbation \(\delta\) is updated using a PGD-type method:
\[
\delta_{k+1}
=
\Pi_{\mathcal{B}_{\epsilon}}
\left(
\delta_k
+
\alpha \cdot \operatorname{UpdateDirection}
\left(
\nabla_{\delta}\mathcal{L}(x+\delta_k)
\right)
\right),
\]
where \(\Pi_{\mathcal{B}_{\epsilon}}\) denotes projection onto the admissible perturbation set, such as
\[
\|\delta\|_{\infty} \leq \epsilon
\]
or
\[
\|\delta\|_2 \leq \epsilon.
\]

The goal is not only to measure the final attack strength, but also to record runtime and GPU memory usage during every attack step. For each PGD step, the following quantities should be recorded:
\[
\begin{aligned}
&\text{forward inference time}, \\
&\text{loss evaluation time}, \\
&\text{backward / gradient computation time}, \\
&\text{projection time}, \\
&\text{total PGD step time}, \\
&\text{allocated GPU memory}, \\
&\text{reserved GPU memory}, \\
&\text{peak GPU memory}.
\end{aligned}
\]

A particularly important issue is the overhead caused by converting tensors between JAX and PyTorch. In mixed-framework settings, the forward pass may require conversions such as
\[
\text{PyTorch tensor}
\rightarrow
\text{JAX array}
\rightarrow
\text{PyTorch tensor},
\]
and the backward pass may require gradient transfer between the two frameworks. Therefore, the following conversion costs should be measured separately:
\[
\text{PyTorch} \rightarrow \text{JAX},
\qquad
\text{JAX} \rightarrow \text{PyTorch}.
\]
The cost should be measured both in terms of runtime and GPU memory usage.

The most important mixed-framework case is
\[
\text{Solver}_{\mathrm{JAX}} + \text{Model}_{\mathrm{PyTorch}}.
\]
In this case, the solver output is computed by JAX, while the neural surrogate output is computed by PyTorch. A typical forward computation has the form
\[
x_{\mathrm{torch}}
\rightarrow
x_{\mathrm{jax}}
\rightarrow
g_{\mathrm{jax}}(x)
\rightarrow
g_{\mathrm{torch}}(x),
\]
while the model prediction is
\[
f_{\mathrm{torch}}(x_{\mathrm{torch}}).
\]
The loss can then be evaluated on the PyTorch side as
\[
\mathcal{L}
=
\left\|
f_{\mathrm{torch}}(x)
-
g_{\mathrm{torch}}(x)
\right\|^2.
\]
During backpropagation, the gradient must be passed through the JAX solver and then returned to PyTorch. This is the main place where tensor conversion and cross-framework gradient propagation may introduce extra overhead.

The final result of this part should include a table of the following form:
\[
\begin{array}{c|c|c|c|c|c}
\text{PDE Case}
&
\text{Solver}
&
\text{Model}
&
\text{Avg. Step Time}
&
\text{Peak Memory}
&
\text{Final True Loss}
\\
\hline
\text{Burgers}
&
\text{JAX}
&
\text{JAX}
&
\cdots
&
\cdots
&
\cdots
\\
\text{Burgers}
&
\text{JAX}
&
\text{PyTorch}
&
\cdots
&
\cdots
&
\cdots
\\
\text{Burgers}
&
\text{PyTorch}
&
\text{JAX}
&
\cdots
&
\cdots
&
\cdots
\\
\text{Burgers}
&
\text{PyTorch}
&
\text{PyTorch}
&
\cdots
&
\cdots
&
\cdots
\\
\text{Stokes / NS}
&
\text{JAX}
&
\text{JAX}
&
\cdots
&
\cdots
&
\cdots
\\
\text{Stokes / NS}
&
\text{JAX}
&
\text{PyTorch}
&
\cdots
&
\cdots
&
\cdots
\\
\text{Stokes / NS}
&
\text{PyTorch}
&
\text{JAX}
&
\cdots
&
\cdots
&
\cdots
\\
\text{Stokes / NS}
&
\text{PyTorch}
&
\text{PyTorch}
&
\cdots
&
\cdots
&
\cdots
\end{array}
\]

\subsection{Comparison of Three Attack Losses}

The second direction is to compare three different adversarial loss functions. The key difference among them is how the ground-truth solution is treated.

Let
\[
f(x)
\]
be the surrogate model prediction, and let
\[
g(x)
\]
be the numerical solver output. The perturbed input is
\[
x_{\delta} = x+\delta.
\]

\subsubsection{Loss 1: Surrogate-Only Output Change}

The first loss only measures the change of the surrogate model output:
\[
\mathcal{L}_1(\delta)
=
\left\|
f(x+\delta)-f(x)
\right\|^2.
\]
This loss does not require the numerical solver. It only asks how much the model output changes when the input is perturbed.

By local linearization,
\[
f(x+\delta)
\approx
f(x) + J_f(x)\delta,
\]
so
\[
f(x+\delta)-f(x)
\approx
J_f(x)\delta.
\]
Therefore,
\[
\mathcal{L}_1(\delta)
\approx
\left\|
J_f(x)\delta
\right\|^2
=
\delta^\top J_f(x)^\top J_f(x)\delta.
\]
This is approximately a quadratic form. Under a norm constraint such as \(\|\delta\|_2=1\), the maximization problem is closely related to the largest singular value of \(J_f(x)\):
\[
\max_{\|\delta\|_2=1}
\left\|
J_f(x)\delta
\right\|
=
\sigma_{\max}(J_f(x)).
\]
Thus, this loss is the most suitable one for power iteration or generalized power iteration methods.

\subsubsection{Loss 2: Fixed Ground Truth}

The second loss assumes that a ground-truth output \(y\) is known and fixed:
\[
\mathcal{L}_2(\delta)
=
\left\|
f(x+\delta)-y
\right\|^2.
\]
Here \(y\) can be the solver output at the clean input:
\[
y = g(x),
\]
or it can be a precomputed ground-truth label from the dataset.

Using the local approximation
\[
f(x+\delta)
\approx
f(x)+J_f(x)\delta,
\]
we obtain
\[
\mathcal{L}_2(\delta)
\approx
\left\|
f(x)-y+J_f(x)\delta
\right\|^2.
\]
Expanding this expression gives
\[
\mathcal{L}_2(\delta)
\approx
\left\|
f(x)-y
\right\|^2
+
2(f(x)-y)^\top J_f(x)\delta
+
\delta^\top J_f(x)^\top J_f(x)\delta.
\]
This objective contains a linear term:
\[
2(f(x)-y)^\top J_f(x)\delta.
\]
Therefore, it is not a pure Rayleigh quotient and is not exactly equivalent to a largest singular vector problem. PGD is more natural for this loss because PGD directly optimizes the full nonlinear objective.

A related version is to recompute the solver output at the perturbed input but stop its gradient:
\[
\mathcal{L}_{2,\mathrm{sg}}(\delta)
=
\left\|
f(x+\delta)
-
\operatorname{stopgrad}(g(x+\delta))
\right\|^2.
\]
In this case, the numerical value of the target changes with \(x+\delta\), but the gradient does not flow through the solver.

\subsubsection{Loss 3: Fully Coupled Solver--Model Loss}

The third loss is the true solver--model discrepancy at the perturbed input:
\[
\mathcal{L}_3(\delta)
=
\left\|
f(x+\delta)-g(x+\delta)
\right\|^2.
\]
This is the most faithful objective for attacking the surrogate model relative to the numerical solver.

Using local linearization,
\[
f(x+\delta)
\approx
f(x)+J_f(x)\delta,
\]
and
\[
g(x+\delta)
\approx
g(x)+J_g(x)\delta.
\]
Then
\[
f(x+\delta)-g(x+\delta)
\approx
f(x)-g(x)
+
\left(J_f(x)-J_g(x)\right)\delta.
\]
Therefore,
\[
\mathcal{L}_3(\delta)
\approx
\left\|
f(x)-g(x)
+
\left(J_f(x)-J_g(x)\right)\delta
\right\|^2.
\]
Expanding this expression gives a linear term and a quadratic term:
\[
2(f(x)-g(x))^\top
\left(J_f(x)-J_g(x)\right)\delta
\]
and
\[
\delta^\top
\left(J_f(x)-J_g(x)\right)^\top
\left(J_f(x)-J_g(x)\right)
\delta.
\]
Thus, this loss is also not a pure quadratic Rayleigh quotient. However, it directly optimizes the true quantity of interest:
\[
\left\|
f(x+\delta)-g(x+\delta)
\right\|^2.
\]

Empirically, this loss is expected to produce the fastest growth in true solver--model error because the optimization objective and the final evaluation metric are aligned.

\subsubsection{Unified True-Loss Evaluation}

Although the three optimization losses are different, all attacks must be evaluated using the same true loss:
\[
\boxed{
\mathcal{L}^{\mathrm{true}}(\delta)
=
\left\|
f(x+\delta)-g(x+\delta)
\right\|^2
}
\]
This is important because \(\mathcal{L}_1\) and \(\mathcal{L}_2\) are only surrogate attack objectives. They should not be used as the final measure of attack strength.

For every attack method, after obtaining the adversarial perturbation \(\delta\), the final evaluation should recompute
\[
x_{\mathrm{adv}} = x+\delta,
\]
\[
f_{\mathrm{adv}} = f(x_{\mathrm{adv}}),
\]
\[
g_{\mathrm{adv}} = g(x_{\mathrm{adv}}),
\]
and
\[
\mathcal{L}^{\mathrm{true}}
=
\left\|
f_{\mathrm{adv}}-g_{\mathrm{adv}}
\right\|^2.
\]

The comparison table should have the following form:
\[
\begin{array}{c|c|c|c|c}
\text{Optimized Loss}
&
\text{Uses Solver?}
&
\text{Solver Gradient?}
&
\text{True Loss Growth}
&
\text{Final True Loss}
\\
\hline
\left\|f(x+\delta)-f(x)\right\|^2
&
\text{No}
&
\text{No}
&
\cdots
&
\cdots
\\
\left\|f(x+\delta)-y\right\|^2
&
\text{Optional}
&
\text{No}
&
\cdots
&
\cdots
\\
\left\|f(x+\delta)-g(x+\delta)\right\|^2
&
\text{Yes}
&
\text{Yes}
&
\cdots
&
\cdots
\end{array}
\]

\subsection{Adding Input Perturbation Size into the Objective}

The third direction is to modify the attack objective so that the perturbation size is included directly in the loss, rather than only being controlled by PGD projection.

The standard projected attack solves
\[
\max_{\delta}
\mathcal{L}_{\mathrm{out}}(\delta)
\quad
\text{subject to}
\quad
\|\delta\| \leq \epsilon.
\]
In PGD, this constraint is enforced by projection:
\[
\delta_{k+1}
=
\Pi_{\|\delta\|\leq \epsilon}
\left(
\delta_k + \alpha \nabla_\delta \mathcal{L}_{\mathrm{out}}(\delta_k)
\right).
\]
In this setup, the perturbation size is controlled by the projection step, but it is not explicitly included in the objective.

The proposed experiment is to directly include the input perturbation size in the objective. The goal is to encourage both large output error and small input perturbation.

\subsubsection{Penalty-Type Objective}

The first modified objective is a penalty-type objective:
\[
\max_\delta
\left[
\mathcal{L}_{\mathrm{out}}(\delta)
-
\lambda
\mathcal{L}_{\mathrm{in}}(\delta)
\right],
\]
where
\[
\mathcal{L}_{\mathrm{out}}(\delta)
\]
is the output discrepancy and
\[
\mathcal{L}_{\mathrm{in}}(\delta)
\]
measures the input perturbation size.

For example, using the true solver--model loss,
\[
\mathcal{L}_{\mathrm{out}}(\delta)
=
\left\|
f(x+\delta)-g(x+\delta)
\right\|_p^p,
\]
and
\[
\mathcal{L}_{\mathrm{in}}(\delta)
=
\|\delta\|_q^r.
\]
The full objective becomes
\[
\max_\delta
\left[
\left\|
f(x+\delta)-g(x+\delta)
\right\|_p^p
-
\lambda
\|\delta\|_q^r
\right].
\]
The coefficient \(\lambda>0\) controls the trade-off between attack strength and perturbation size. A larger \(\lambda\) penalizes large perturbations more strongly, while a smaller \(\lambda\) makes the objective closer to the ordinary attack loss.

Possible values include
\[
\lambda
\in
\{10^{-4},10^{-3},10^{-2},10^{-1},1,10\},
\]
although the final values should be adjusted according to the numerical scale of the PDE loss.

For this experiment, the following quantities should be recorded:
\[
\mathcal{L}_{\mathrm{out}}(\delta),
\qquad
\|\delta\|,
\qquad
\mathcal{L}_{\mathrm{out}}(\delta)-\lambda\|\delta\|,
\qquad
\mathcal{L}^{\mathrm{true}}(\delta).
\]
The main questions are:
\[
\begin{aligned}
&\text{Does the penalty objective produce smaller perturbations?}\\
&\text{Does it still produce large true solver--model error?}\\
&\text{Is it more stable than pure projected PGD?}\\
&\text{How sensitive is the result to }\lambda?
\end{aligned}
\]

\subsubsection{Ratio-Type Objective}

The second modified objective is a ratio-type objective:
\[
\max_\delta
\frac{
\mathcal{L}_{\mathrm{out}}(\delta)
}{
\mathcal{L}_{\mathrm{in}}(\delta)+\varepsilon
}.
\]
For example,
\[
\max_\delta
\frac{
\left\|
f(x+\delta)-g(x+\delta)
\right\|_p
}{
\|\delta\|_q+\varepsilon
}.
\]
A squared version is
\[
\max_\delta
\frac{
\left\|
f(x+\delta)-g(x+\delta)
\right\|_p^2
}{
\|\delta\|_q^2+\varepsilon
}.
\]

This objective is similar to an operator norm or local Lipschitz ratio:
\[
\|A\|
=
\max_{\delta\neq 0}
\frac{\|A\delta\|}{\|\delta\|}.
\]
However, in the present problem,
\[
f(x+\delta)-g(x+\delta)
\]
is generally nonlinear. Therefore, the ratio objective may be unstable if the denominator becomes too small. To avoid this issue, a stabilizing constant \(\varepsilon\) must be added:
\[
\frac{
\mathcal{L}_{\mathrm{out}}(\delta)
}{
\|\delta\|_q^r+\varepsilon
}.
\]

This experiment should record:
\[
\text{ratio objective},
\qquad
\mathcal{L}_{\mathrm{out}}(\delta),
\qquad
\|\delta\|,
\qquad
\mathcal{L}^{\mathrm{true}}(\delta).
\]
A large ratio value is not necessarily meaningful if the actual true loss is small. Therefore, the ratio objective must always be interpreted together with
\[
\mathcal{L}^{\mathrm{true}}(\delta)
\]
and
\[
\|\delta\|.
\]

\subsubsection{Norm Choices}

The norm used for the input perturbation should also be treated as an experimental variable. Possible choices include
\[
\|\delta\|_2,
\qquad
\|\delta\|_\infty,
\qquad
\|\delta\|_1.
\]
The output discrepancy can also be measured in several ways:
\[
\left\|
f-g
\right\|_2^2,
\qquad
\left\|
f-g
\right\|_2,
\qquad
\frac{
\left\|
f-g
\right\|_2
}{
\left\|
g
\right\|_2+\varepsilon
}.
\]
For PDE solutions, it may also be useful to use a spatial average:
\[
\frac{1}{N}
\left\|
f-g
\right\|_2^2.
\]
If the output is a full space-time trajectory, different output losses can be compared:
\[
\text{final-time loss},
\qquad
\text{all-time trajectory loss},
\qquad
\text{energy-spectrum loss},
\qquad
\text{vorticity loss}.
\]

\subsection{Power-Iteration-Like Attacks and Generalized Power Iteration}

The fourth direction is to replace standard PGD updates with power-iteration-like updates. This is motivated by the fact that, for a pure quadratic loss, the optimal direction is often related to the leading singular vector of a Jacobian.

The standard PGD update is
\[
\delta_{k+1}
=
\Pi
\left(
\delta_k
+
\alpha \nabla_\delta \mathcal{L}(\delta_k)
\right).
\]
This update accumulates the current gradient into the existing perturbation direction.

For a quadratic loss such as
\[
\mathcal{L}(\delta)
=
\|J\delta\|^2
=
\delta^\top J^\top J\delta,
\]
the gradient is
\[
\nabla_\delta \mathcal{L}(\delta)
=
2J^\top J\delta.
\]
Therefore, the power iteration update is
\[
v_{k+1}
=
\frac{
J^\top Jv_k
}{
\|J^\top Jv_k\|
}.
\]
This suggests a more aggressive update:
\[
\delta_{k+1}
=
\epsilon
\frac{
\nabla_\delta \mathcal{L}(\delta_k)
}{
\|\nabla_\delta \mathcal{L}(\delta_k)\|+\varepsilon
}.
\]
Unlike PGD, this update directly resets the perturbation direction to the current gradient direction.

\subsubsection{Power Iteration for the Surrogate-Only Loss}

The most natural case for power iteration is
\[
\mathcal{L}_1(\delta)
=
\left\|
f(x+\delta)-f(x)
\right\|^2.
\]
By local linearization,
\[
\mathcal{L}_1(\delta)
\approx
\left\|
J_f(x)\delta
\right\|^2.
\]
Thus, the leading right singular vector of \(J_f(x)\) gives the locally most sensitive input direction.

This can be implemented without explicitly forming the Jacobian. One only needs Jacobian-vector products and vector-Jacobian products:
\[
v
\mapsto
J_f(x)v,
\]
and
\[
w
\mapsto
J_f(x)^\top w.
\]
Then one step of power iteration is
\[
v_k
\rightarrow
J_f(x)v_k
\rightarrow
J_f(x)^\top J_f(x)v_k
\rightarrow
v_{k+1}.
\]
This is especially suitable for JAX, where JVP and VJP are directly available, but it can also be implemented in PyTorch using autograd-based JVP/VJP tools.

\subsubsection{Heuristic Power-Like Updates for the Fixed-Target and Solver-Coupled Losses}

For the fixed-target loss
\[
\mathcal{L}_2(\delta)
=
\left\|
f(x+\delta)-y
\right\|^2,
\]
the local expansion contains a linear term:
\[
2(f(x)-y)^\top J_f(x)\delta.
\]
For the solver-coupled loss
\[
\mathcal{L}_3(\delta)
=
\left\|
f(x+\delta)-g(x+\delta)
\right\|^2,
\]
the local expansion contains
\[
2(f(x)-g(x))^\top
\left(J_f(x)-J_g(x)\right)\delta.
\]
Thus, these two losses are not pure Rayleigh quotient problems. Power iteration is no longer theoretically exact. However, it is still useful to test power-like updates as heuristic attack methods.

One possible method is the gradient-reset update:
\[
\delta_{k+1}
=
\epsilon
\frac{
\nabla_\delta \mathcal{L}(\delta_k)
}{
\|\nabla_\delta \mathcal{L}(\delta_k)\|+\varepsilon
}.
\]
For an \(L_\infty\) perturbation constraint, the corresponding update is
\[
\delta_{k+1}
=
\epsilon
\cdot
\operatorname{sign}
\left(
\nabla_\delta \mathcal{L}(\delta_k)
\right).
\]
This is similar to repeatedly choosing the current steepest direction, rather than accumulating gradients as in PGD.

Another possible method is to define the solver--model error operator
\[
h(x) = f(x)-g(x).
\]
Then
\[
\mathcal{L}_3(\delta)
=
\|h(x+\delta)\|^2.
\]
Local linearization gives
\[
h(x+\delta)
\approx
h(x)+J_h(x)\delta,
\]
where
\[
J_h(x)=J_f(x)-J_g(x).
\]
Ignoring the residual term \(h(x)\), one can run power iteration on
\[
J_h(x)^\top J_h(x).
\]
The update becomes
\[
v_{k+1}
=
\frac{
J_h(x)^\top J_h(x)v_k
}{
\|J_h(x)^\top J_h(x)v_k\|
}.
\]
This method finds a locally sensitive direction of the solver--model error operator, although it does not exactly optimize the full nonlinear loss.

\subsubsection{Comparison with PGD}

The following attack methods should be compared:
\[
\begin{aligned}
&\text{standard PGD},\\
&\text{gradient-reset power-like update},\\
&\text{JVP/VJP-based power iteration},\\
&\text{generalized power iteration}.
\end{aligned}
\]
For each method, the following quantities should be recorded:
\[
\text{convergence speed},
\qquad
\text{final true loss},
\qquad
\text{average step time},
\qquad
\text{peak GPU memory}.
\]

The expected behavior is:
\[
\begin{aligned}
&\mathcal{L}_1 \text{ should be the most suitable for power iteration},\\
&\mathcal{L}_2 \text{ may be more stable under PGD},\\
&\mathcal{L}_3 \text{ should best match the true attack objective},\\
&\text{gradient-reset updates may grow quickly but may also oscillate},\\
&\text{JVP/VJP power iteration may reduce iteration count but increase per-step cost}.
\end{aligned}
\]

The final comparison table should have the form
\[
\begin{array}{c|c|c|c|c|c}
\text{Loss}
&
\text{Attack Method}
&
\text{Iterations to Threshold}
&
\text{Final True Loss}
&
\text{Avg. Step Time}
&
\text{Peak Memory}
\\
\hline
\mathcal{L}_1
&
\text{PGD}
&
\cdots
&
\cdots
&
\cdots
&
\cdots
\\
\mathcal{L}_1
&
\text{Power-like}
&
\cdots
&
\cdots
&
\cdots
&
\cdots
\\
\mathcal{L}_2
&
\text{PGD}
&
\cdots
&
\cdots
&
\cdots
&
\cdots
\\
\mathcal{L}_2
&
\text{Power-like}
&
\cdots
&
\cdots
&
\cdots
&
\cdots
\\
\mathcal{L}_3
&
\text{PGD}
&
\cdots
&
\cdots
&
\cdots
&
\cdots
\\
\mathcal{L}_3
&
\text{Power-like}
&
\cdots
&
\cdots
&
\cdots
&
\cdots
\end{array}
\]

\subsection{Overall Experimental Structure}

The full experiment can be organized hierarchically as
\[
\text{PDE Case}
\rightarrow
\text{Solver--Model Framework Combination}
\rightarrow
\text{Loss Type}
\rightarrow
\text{Attack Method}
\rightarrow
\text{Metrics}.
\]
More explicitly:
\[
\begin{aligned}
&\text{PDE Case:}
&&
\text{1D Burgers},\quad \text{2D Stokes / Navier--Stokes},
\\
&\text{Framework Combination:}
&&
(\text{JAX},\text{JAX}),\quad
(\text{JAX},\text{PyTorch}),\quad
(\text{PyTorch},\text{JAX}),\quad
(\text{PyTorch},\text{PyTorch}),
\\
&\text{Loss Type:}
&&
\mathcal{L}_1,\quad
\mathcal{L}_2,\quad
\mathcal{L}_3,
\\
&\text{Attack Method:}
&&
\text{PGD},\quad
\text{penalty objective},\quad
\text{ratio objective},\quad
\text{power-like update},
\\
&\text{Metrics:}
&&
\text{optimized loss},\quad
\text{true loss},\quad
\|\delta\|,\quad
\text{runtime},\quad
\text{GPU memory}.
\end{aligned}
\]

The most important evaluation metric is always the true solver--model discrepancy:
\[
\boxed{
\mathcal{L}^{\mathrm{true}}(\delta)
=
\left\|
f(x+\delta)-g(x+\delta)
\right\|^2
}
\]
Different attack losses and optimization methods may be used to generate \(\delta\), but the final comparison must always be based on this true loss. This guarantees that all methods are compared according to the same physical and numerical meaning: the adversarially induced error between the neural surrogate model and the numerical PDE solver.

\subsection{Structure-Aware and Alignment-Aware Solver--Model Losses}

The fifth direction is to replace the purely pointwise \(L^2\) solver--model discrepancy with a more structure-aware loss. The current true loss is usually written as
\[
\mathcal{L}^{\mathrm{true}}(\delta)
=
\left\|
f(x+\delta)-g(x+\delta)
\right\|_2^2.
\]
This loss compares the two output fields point by point. However, a pointwise \(L^2\) loss may overestimate the true structural difference between two outputs. For example, if one output is only slightly shifted, stretched, compressed, rotated, or locally warped relative to the other, the pointwise \(L^2\) loss can become very large even though the two fields still have similar global shapes or textures.

Therefore, the fifth experimental direction is to design losses that still compare
\[
f(x+\delta)
\quad\text{and}\quad
g(x+\delta),
\]
but no longer rely only on direct pointwise subtraction. The purpose is to measure whether the adversarial perturbation creates a real structural difference, rather than only producing a small spatial misalignment or a simple intensity rescaling.

There are two main types of transformations to consider.

First, one can allow limited spatial alignment before computing the loss. This means that the spatial coordinates of one output field can be transformed within a controlled range. Possible transformations include
\[
\text{translation},\quad
\text{rotation},\quad
\text{stretching},\quad
\text{compression},\quad
\text{local warping}.
\]
In one-dimensional time-series settings, this is similar in spirit to dynamic time warping (DTW), where two signals are aligned by allowing local stretching and compression along the time axis. In two-dimensional PDE outputs, the corresponding idea is to allow limited spatial deformations before comparing the model output and the solver output.

A general form of this alignment-aware loss is
\[
\mathcal{L}_{\mathrm{align}}(\delta)
=
\min_{T \in \mathcal{T}}
\left\|
f(x+\delta)\circ T
-
g(x+\delta)
\right\|^2,
\]
where \(T\) is a spatial transformation and \(\mathcal{T}\) is the allowed transformation class. The set \(\mathcal{T}\) must be carefully restricted. For example, one may allow only translations up to a fixed number of grid points, rotations up to a small angle, or local deformations with bounded smoothness and bounded displacement. If \(\mathcal{T}\) is too flexible, then almost any output can be transformed to match another output, and the loss may become meaningless.

Second, one can allow limited value-space transformations while keeping the spatial grid fixed. In this case, the spatial points are not moved, but the field values are allowed to undergo simple global transformations, such as scaling or shifting:
\[
u \mapsto au+b.
\]
This is useful because two outputs may have very similar spatial structure but differ by an overall amplitude shift or scaling factor. Possible losses include correlation-based losses, SSIM-type losses, normalized \(L^2\) losses, or feature-based losses.

A simple value-alignment loss can be written as
\[
\mathcal{L}_{\mathrm{value}}(\delta)
=
\min_{a,b}
\left\|
a f(x+\delta) + b
-
g(x+\delta)
\right\|^2,
\]
where the parameters \(a\) and \(b\) must also be constrained. For example, \(a\) should stay close to \(1\), and \(b\) should stay close to \(0\). Otherwise, the value transformation may become too flexible and hide real differences between the model output and the solver output.

More generally, the final structure-aware loss may combine pointwise error, spatially aligned error, value-aligned error, and feature-level error:
\[
\mathcal{L}_{\mathrm{struct}}(\delta)
=
\lambda_1
\mathcal{L}_{\mathrm{pointwise}}(\delta)
+
\lambda_2
\mathcal{L}_{\mathrm{align}}(\delta)
+
\lambda_3
\mathcal{L}_{\mathrm{value}}(\delta)
+
\lambda_4
\mathcal{L}_{\mathrm{feature}}(\delta).
\]
Here \(\mathcal{L}_{\mathrm{feature}}\) may be a perceptual loss, a correlation-based loss, an SSIM-type loss, or a loss computed from physically meaningful features such as gradients, vorticity, energy spectrum, or dominant spatial patterns.

The key principle is that the loss should allow small, controlled transformations that remove unimportant misalignment, but it should not allow arbitrary transformations. Otherwise, the loss may collapse and fail to measure real solver--model discrepancy. Therefore, this direction should test both the type of allowed transformation and the allowed transformation magnitude:
\[
\text{allowed transformation type},
\qquad
\text{allowed transformation range}.
\]

This fifth direction is still a loss-design experiment, similar to the comparison amongst the
\[
\mathcal{L}_1,\quad
\mathcal{L}_2,\quad
\mathcal{L}_3.
\]
However, it focuses specifically on replacing the pointwise \(L^2\) form of
\[
\left\|
f(x+\delta)-g(x+\delta)
\right\|^2
\]
with losses that better capture structural, geometric, and perceptual differences between the two PDE output fields.

\section{Cross-Framework PGD Attacks for PDE Solver--Model Comparisons}

This section describes the first major experimental direction: comparing PGD attacks across different combinations of numerical PDE solvers and neural surrogate models implemented in JAX and PyTorch. The goal is to evaluate not only the final adversarial error, but also the computational cost, GPU memory usage, and tensor-conversion overhead caused by mixing JAX and PyTorch in the same attack pipeline.

The experiments focus on two PDE cases:
\[
\text{1D Burgers},
\qquad
\text{2D Stokes / Navier--Stokes}.
\]
For each PDE case, there are two main computational components:
\[
\text{numerical solver},
\qquad
\text{surrogate model}.
\]
Both the numerical solver and the surrogate model have JAX and PyTorch implementations. Therefore, for each PDE case, four solver--model framework combinations need to be tested:
\[
\begin{aligned}
&\text{Solver}_{\mathrm{JAX}} + \text{Model}_{\mathrm{JAX}},\\
&\text{Solver}_{\mathrm{JAX}} + \text{Model}_{\mathrm{PyTorch}},\\
&\text{Solver}_{\mathrm{PyTorch}} + \text{Model}_{\mathrm{JAX}},\\
&\text{Solver}_{\mathrm{PyTorch}} + \text{Model}_{\mathrm{PyTorch}}.
\end{aligned}
\]
Thus, the full experimental matrix contains
\[
2 \text{ PDE cases} \times 4 \text{ framework combinations}
=
8 \text{ main attack experiments}.
\]

\subsection{Experimental Cases and Framework Combinations}

The two PDE cases are:
\[
\text{1D Burgers},
\qquad
\text{2D Stokes / Navier--Stokes}.
\]
For each case, the numerical solver is denoted by
\[
g(x),
\]
and the surrogate neural model is denoted by
\[
f(x).
\]
The input is an initial condition or forcing field
\[
x,
\]
and the adversarial perturbation is
\[
\delta.
\]
The adversarial input is
\[
x_{\mathrm{adv}} = x+\delta.
\]

The four framework combinations are summarized as follows:
\[
\begin{array}{c|c|c}
\text{PDE Case} & \text{Solver Backend} & \text{Model Backend} \\
\hline
\text{Burgers} & \text{JAX} & \text{JAX} \\
\text{Burgers} & \text{JAX} & \text{PyTorch} \\
\text{Burgers} & \text{PyTorch} & \text{JAX} \\
\text{Burgers} & \text{PyTorch} & \text{PyTorch} \\
\text{Stokes / NS} & \text{JAX} & \text{JAX} \\
\text{Stokes / NS} & \text{JAX} & \text{PyTorch} \\
\text{Stokes / NS} & \text{PyTorch} & \text{JAX} \\
\text{Stokes / NS} & \text{PyTorch} & \text{PyTorch}
\end{array}
\]

The purpose of this comparison is to determine whether different framework combinations produce different attack trajectories, different final adversarial examples, different true solver--model errors, and different computational costs.

\subsection{PGD Attack Formulation}

Each experiment uses a PGD attack. The adversarial perturbation is updated iteratively. Given the current perturbation \(\delta_k\), the next perturbation is computed as
\[
\delta_{k+1}
=
\Pi_{\mathcal{B}_{\epsilon}}
\left(
\delta_k
+
\alpha \cdot
\operatorname{UpdateDirection}
\left(
\nabla_{\delta}\mathcal{L}(x+\delta_k)
\right)
\right),
\]
where \(\alpha\) is the PGD step size, \(\mathcal{L}\) is the attack loss, and
\[
\Pi_{\mathcal{B}_{\epsilon}}
\]
is the projection onto the admissible perturbation set.

For an \(L^{\infty}\)-bounded attack, the constraint is
\[
\|\delta\|_{\infty} \leq \epsilon.
\]
For an \(L^2\)-bounded attack, the constraint is
\[
\|\delta\|_2 \leq \epsilon.
\]

The main output of the attack is not only the final perturbation, but the full attack trajectory:
\[
\delta_0,\delta_1,\dots,\delta_K,
\]
where \(K\) is the number of PGD steps.

During the attack, both the optimized attack loss and the true solver--model loss should be recorded. The true loss is always defined as
\[
\mathcal{L}^{\mathrm{true}}(\delta)
=
\left\|
f(x+\delta)-g(x+\delta)
\right\|^2.
\]
This quantity is the final evaluation metric, regardless of which loss is used internally to update \(\delta\).

\subsection{Runtime Metrics}

For every PGD step, the runtime should be decomposed into several components. The following quantities should be recorded:
\[
\begin{aligned}
&\text{forward time},\\
&\text{loss evaluation time},\\
&\text{backward / gradient time},\\
&\text{projection time},\\
&\text{optimizer / update time},\\
&\text{total step time}.
\end{aligned}
\]
The total runtime of the full 100-step attack should also be recorded:
\[
\text{total attack time}.
\]

A central purpose of this experiment is to test the following observations:
\[
\begin{aligned}
&\text{JAX may have a large first-step compilation cost},\\
&\text{JAX may have faster forward and backward passes after compilation},\\
&\text{JAX may use more GPU memory than PyTorch},\\
&\text{mixed JAX--PyTorch pipelines may introduce additional conversion overhead}.
\end{aligned}
\]

Therefore, the timing analysis should distinguish between:
\[
\text{first-step JAX compilation time}
\]
and
\[
\text{post-compilation per-step runtime}.
\]
For mixed-framework experiments, tensor conversion time should be recorded separately from the solver and model computation time.

\subsection{GPU Memory Metrics}

For every PGD step, GPU memory usage should also be recorded. The required memory metrics include
\[
\begin{aligned}
&\text{allocated GPU memory},\\
&\text{reserved GPU memory},\\
&\text{peak GPU memory}.
\end{aligned}
\]

For PyTorch, these quantities can be recorded using functions such as
\[
\texttt{torch.cuda.memory\_allocated()},
\]
\[
\texttt{torch.cuda.memory\_reserved()},
\]
and
\[
\texttt{torch.cuda.max\_memory\_allocated()}.
\]

For JAX, GPU memory should be measured using available JAX-side device memory information or an external GPU memory sampler. The important requirement is that the memory measurement should reflect the true memory usage of each individual experiment.

The four framework combinations should not be loaded and executed in the same process. Otherwise, GPU memory measurements may be contaminated by previously loaded models, cached arrays, or compiled programs. The recommended procedure is:
\[
\begin{aligned}
&\text{run one framework combination},\\
&\text{save its metrics and outputs},\\
&\text{clear memory and terminate the process},\\
&\text{run the next framework combination},\\
&\text{aggregate results only after all runs are finished}.
\end{aligned}
\]

Thus, the four combinations should be executed as separate runs:
\[
\text{JAX--JAX},
\qquad
\text{JAX--PyTorch},
\qquad
\text{PyTorch--JAX},
\qquad
\text{PyTorch--PyTorch}.
\]

\subsection{Mixed JAX--PyTorch Framework Interface}

The most technically important part of this task is the mixed-framework interface between JAX and PyTorch. The two mixed cases are:
\[
\text{Solver}_{\mathrm{JAX}} + \text{Model}_{\mathrm{PyTorch}},
\]
and
\[
\text{Solver}_{\mathrm{PyTorch}} + \text{Model}_{\mathrm{JAX}}.
\]

The most important case is
\[
\text{Solver}_{\mathrm{JAX}} + \text{Model}_{\mathrm{PyTorch}}.
\]
In this case, the forward computation has the form
\[
x_{\mathrm{torch}}
\rightarrow
x_{\mathrm{jax}}
\rightarrow
g_{\mathrm{jax}}(x)
\rightarrow
g_{\mathrm{torch}}(x).
\]
At the same time, the PyTorch model computes
\[
f_{\mathrm{torch}}(x_{\mathrm{torch}}).
\]
The loss is then evaluated on the PyTorch side:
\[
\mathcal{L}
=
\left\|
f_{\mathrm{torch}}(x_{\mathrm{torch}})
-
g_{\mathrm{torch}}(x_{\mathrm{torch}})
\right\|.
\]

During backpropagation, the gradient from the PyTorch loss must be passed back to the JAX solver output, then propagated through the JAX solver using a JAX VJP or a custom backward rule:
\[
\frac{\partial g_{\mathrm{JAX}}(x)}{\partial x}.
\]
The resulting input gradient must then be converted back to a PyTorch tensor so that PGD can update the PyTorch-side perturbation.

The mixed-framework interface must therefore support both forward and backward conversions:
\[
\text{PyTorch tensor} \rightarrow \text{JAX array},
\]
\[
\text{JAX array} \rightarrow \text{PyTorch tensor},
\]
\[
\text{PyTorch gradient} \rightarrow \text{JAX cotangent},
\]
\[
\text{JAX input gradient} \rightarrow \text{PyTorch gradient}.
\]

The following conversion costs should be recorded separately:
\[
\begin{aligned}
&\text{forward conversion time: PyTorch} \rightarrow \text{JAX},\\
&\text{forward conversion time: JAX} \rightarrow \text{PyTorch},\\
&\text{backward conversion time: PyTorch} \rightarrow \text{JAX},\\
&\text{backward conversion time: JAX} \rightarrow \text{PyTorch},\\
&\text{JAX VJP time}.
\end{aligned}
\]

\subsection{Recommended Mixed-Framework Control Strategy}

For engineering simplicity, mixed-framework experiments should use PyTorch as the main control framework whenever at least one component is implemented in PyTorch. Under this strategy:
\[
\begin{aligned}
&\text{JAX--JAX}: \text{run fully in JAX},\\
&\text{PyTorch--PyTorch}: \text{run fully in PyTorch},\\
&\text{JAX--PyTorch}: \text{wrap the JAX component as a PyTorch-callable function},\\
&\text{PyTorch--JAX}: \text{wrap the JAX component as a PyTorch-callable function and keep PGD in PyTorch}.
\end{aligned}
\]

For the JAX solver and PyTorch model case, the JAX solver should be wrapped as a PyTorch autograd function. From the PyTorch side, the wrapped solver behaves like
\[
g_{\mathrm{wrapped}}(x_{\mathrm{torch}}),
\]
while internally it performs
\[
x_{\mathrm{torch}}
\rightarrow
x_{\mathrm{jax}}
\rightarrow
g_{\mathrm{jax}}(x_{\mathrm{jax}})
\rightarrow
g_{\mathrm{torch}}.
\]

In the backward pass, the wrapper performs
\[
\bar{g}_{\mathrm{torch}}
\rightarrow
\bar{g}_{\mathrm{jax}}
\rightarrow
\bar{x}_{\mathrm{jax}}
\rightarrow
\bar{x}_{\mathrm{torch}}.
\]

This wrapper should explicitly record:
\[
\begin{aligned}
&\text{torch-to-jax forward conversion time},\\
&\text{jax solver forward time},\\
&\text{jax-to-torch forward conversion time},\\
&\text{torch-to-jax gradient conversion time},\\
&\text{jax VJP time},\\
&\text{jax-to-torch gradient conversion time}.
\end{aligned}
\]

\subsection{Attack Hyperparameter Pool}

The PGD hyperparameters should be selected from previously used experimental settings. The goal is not to introduce unrelated new hyperparameters, but to reuse reasonable values from the existing attack code.

\subsubsection{Hyperparameters for 1D Burgers}

For the 1D Burgers case, the following parameter sets are candidates.

The first candidate is an \(L^{\infty}\)-bounded attack:
\[
\text{norm}=\infty,
\qquad
\text{index}=0,
\qquad
\text{steps}=100,
\qquad
\epsilon=0.5,
\qquad
\alpha=0.05.
\]

The second candidate is an \(L^2\)-bounded attack:
\[
\text{norm}=2,
\qquad
\epsilon=10,
\qquad
\alpha=0.05,
\qquad
\text{steps}=100.
\]

The third candidate is:
\[
\text{norm}=2,
\qquad
\text{steps}=100,
\qquad
\epsilon=20,
\qquad
\alpha=1.
\]

The fourth candidate is:
\[
\text{norm}=2,
\qquad
\alpha=2,
\qquad
\epsilon=100,
\qquad
\text{steps}=100,
\qquad
\text{index}=0,1,2.
\]

The fifth candidate is:
\[
\text{norm}=2,
\qquad
\alpha=0.06,
\qquad
\epsilon=20,
\qquad
\text{steps}=100,
\qquad
\text{index}=0.
\]

A stable first-choice setting for Burgers is:
\[
\boxed{
\text{norm}=2,
\qquad
\epsilon=20,
\qquad
\alpha=0.06,
\qquad
\text{steps}=100,
\qquad
\text{index}=0.
}
\]
An additional \(L^{\infty}\) comparison can use:
\[
\boxed{
\text{norm}=\infty,
\qquad
\epsilon=0.5,
\qquad
\alpha=0.05,
\qquad
\text{steps}=100,
\qquad
\text{index}=0.
}
\]

\subsubsection{Hyperparameters for 2D Stokes / Navier--Stokes}

For the 2D Stokes / Navier--Stokes case, the first candidate is:
\[
\text{norm}=2,
\qquad
\epsilon=32768,
\qquad
\alpha=5,
\qquad
\text{steps}=100,
\qquad
\text{index}=1.
\]

The second candidate is:
\[
\text{norm}=\infty,
\qquad
\epsilon=1,
\qquad
\alpha=0.01,
\qquad
\text{steps}=100,
\qquad
\text{index}=2.
\]

A first-choice setting for the NS / Stokes case is:
\[
\boxed{
\text{norm}=2,
\qquad
\epsilon=32768,
\qquad
\alpha=5,
\qquad
\text{steps}=100,
\qquad
\text{index}=1.
}
\]
An additional \(L^{\infty}\) comparison can use:
\[
\boxed{
\text{norm}=\infty,
\qquad
\epsilon=1,
\qquad
\alpha=0.01,
\qquad
\text{steps}=100,
\qquad
\text{index}=2.
}
\]

\subsection{Code Implementation Requirements}

The implementation should reuse the existing project structure as much as possible. The current codebase already contains components for:
\[
\begin{aligned}
&\text{PyTorch model attack},\\
&\text{PyTorch solver attack},\\
&\text{PGD attack logic},\\
&\text{visualization},\\
&\text{loss curve saving and plotting}.
\end{aligned}
\]

The new code should mainly introduce a unified solver--model interface and a cross-framework attack runner. The implementation should support command-line options:
\[
\texttt{--solver\_backend} \in \{\texttt{jax},\texttt{torch}\},
\]
\[
\texttt{--model\_backend} \in \{\texttt{jax},\texttt{torch}\}.
\]

The full set of supported combinations should be:
\[
\begin{aligned}
&\texttt{--solver\_backend jax --model\_backend jax},\\
&\texttt{--solver\_backend jax --model\_backend torch},\\
&\texttt{--solver\_backend torch --model\_backend jax},\\
&\texttt{--solver\_backend torch --model\_backend torch}.
\end{aligned}
\]

The core attack runner should perform the following operations:
\[
\begin{aligned}
&\text{load input sample},\\
&\text{load solver},\\
&\text{load model},\\
&\text{construct framework-compatible solver--model pair},\\
&\text{run PGD attack},\\
&\text{record runtime and memory metrics},\\
&\text{save all results},\\
&\text{generate plots or save data for plotting}.
\end{aligned}
\]

The code modification should be limited to the parts needed for this task. The existing model architectures, numerical solvers, checkpoint loading logic, and data loading logic should not be changed unless the change is necessary for the backend interface.

\subsection{Result Saving Structure}

Each experiment should save its outputs in a separate directory. A recommended directory structure is:
\[
\begin{aligned}
&\texttt{results/burgers/solver\_jax\_model\_jax/},\\
&\texttt{results/burgers/solver\_jax\_model\_torch/},\\
&\texttt{results/burgers/solver\_torch\_model\_jax/},\\
&\texttt{results/burgers/solver\_torch\_model\_torch/},\\
&\texttt{results/ns/solver\_jax\_model\_jax/},\\
&\texttt{results/ns/solver\_jax\_model\_torch/},\\
&\texttt{results/ns/solver\_torch\_model\_jax/},\\
&\texttt{results/ns/solver\_torch\_model\_torch/}.
\end{aligned}
\]

Each experiment directory should contain at least:
\[
\begin{aligned}
&\texttt{metrics.csv},\\
&\texttt{config.json},\\
&\texttt{attack\_loss.npy},\\
&\texttt{true\_loss.npy},\\
&\texttt{delta\_norm.npy},\\
&\texttt{step\_time.npy},\\
&\texttt{forward\_time.npy},\\
&\texttt{backward\_time.npy},\\
&\texttt{projection\_time.npy},\\
&\texttt{memory\_allocated.npy},\\
&\texttt{memory\_reserved.npy},\\
&\texttt{memory\_peak.npy},\\
&\texttt{x\_clean.npy},\\
&\texttt{x\_adv.npy},\\
&\texttt{model\_output\_clean.npy},\\
&\texttt{solver\_output\_clean.npy},\\
&\texttt{model\_output\_adv.npy},\\
&\texttt{solver\_output\_adv.npy},\\
&\texttt{difference\_adv.npy}.
\end{aligned}
\]

For mixed-framework experiments, the following additional files should be saved:
\[
\begin{aligned}
&\texttt{conversion\_forward\_torch\_to\_jax.npy},\\
&\texttt{conversion\_forward\_jax\_to\_torch.npy},\\
&\texttt{conversion\_backward\_torch\_to\_jax.npy},\\
&\texttt{conversion\_backward\_jax\_to\_torch.npy},\\
&\texttt{jax\_vjp\_time.npy}.
\end{aligned}
\]

\subsection{Visualization Requirements}

The visualization should compare the four framework combinations under the same PDE case and the same attack hyperparameters.

The following curves should be plotted:
\[
\begin{aligned}
&\text{attack loss curve},\\
&\text{true loss curve},\\
&\text{perturbation norm curve},\\
&\text{step time curve},\\
&\text{forward time curve},\\
&\text{backward time curve},\\
&\text{memory usage curve},\\
&\text{peak memory curve}.
\end{aligned}
\]

The four combinations
\[
\text{JAX--JAX},
\qquad
\text{JAX--PyTorch},
\qquad
\text{PyTorch--JAX},
\qquad
\text{PyTorch--PyTorch}
\]
should be plotted on the same figure whenever possible.

The comparison should answer:
\[
\begin{aligned}
&\text{Do the four combinations produce the same loss values?}\\
&\text{Do they have the same loss growth rate?}\\
&\text{Do they reach the same final true loss?}\\
&\text{Are mixed-framework attacks slower?}\\
&\text{Do mixed-framework attacks use more memory?}\\
&\text{Is the first JAX step much slower because of compilation?}\\
&\text{Are later JAX steps faster than PyTorch steps?}
\end{aligned}
\]

\subsubsection{Visualization for 1D Burgers}

For 1D Burgers, the main visualization should use line plots. The following quantities should be plotted:
\[
\begin{aligned}
&\text{clean input},\\
&\text{adversarial input},\\
&\text{clean model output},\\
&\text{clean solver output},\\
&\text{adversarial model output},\\
&\text{adversarial solver output},\\
&\text{adversarial difference}.
\end{aligned}
\]
The adversarial difference is
\[
f(x+\delta)-g(x+\delta).
\]

It is also useful to generate a GIF showing the PGD trajectory:
\[
k=0,1,2,\dots,100.
\]
The GIF can show:
\[
\begin{aligned}
&x+\delta_k,\\
&f(x+\delta_k),\\
&g(x+\delta_k),\\
&f(x+\delta_k)-g(x+\delta_k).
\end{aligned}
\]

\subsubsection{Visualization for 2D Stokes / Navier--Stokes}

For 2D Stokes / Navier--Stokes, the main visualization should use heatmaps. The following fields should be plotted:
\[
\begin{aligned}
&\text{clean input heatmap},\\
&\text{adversarial input heatmap},\\
&\text{clean model output heatmap},\\
&\text{clean solver output heatmap},\\
&\text{adversarial model output heatmap},\\
&\text{adversarial solver output heatmap},\\
&\text{adversarial difference heatmap}.
\end{aligned}
\]
The adversarial difference heatmap is
\[
f(x+\delta)-g(x+\delta).
\]

The visualization should also compare differences among framework combinations, such as
\[
x_{\mathrm{adv}}^{\mathrm{JAX-JAX}}
-
x_{\mathrm{adv}}^{\mathrm{Torch-Torch}},
\]
\[
f_{\mathrm{adv}}^{\mathrm{JAX-JAX}}
-
f_{\mathrm{adv}}^{\mathrm{Torch-Torch}},
\]
and
\[
g_{\mathrm{adv}}^{\mathrm{JAX-JAX}}
-
g_{\mathrm{adv}}^{\mathrm{Torch-Torch}}.
\]
These difference plots are important for checking whether the same attack parameters produce numerically similar adversarial inputs and outputs across framework combinations.

A GIF can also be generated to show the heatmap evolution during the PGD attack.

\subsection{Summary Tables}

The final summary table should include the following columns:
\[
\begin{array}{c|c|c|c|c|c|c|c|c|c|c|c}
\text{PDE Case}
&
\text{Solver}
&
\text{Model}
&
\text{Norm}
&
\epsilon
&
\alpha
&
\text{Steps}
&
\text{Index}
&
\text{Avg. Step Time}
&
\text{Total Time}
&
\text{Peak Memory}
&
\text{Final True Loss}
\\
\hline
\cdots & \cdots & \cdots & \cdots & \cdots & \cdots & \cdots & \cdots & \cdots & \cdots & \cdots & \cdots
\end{array}
\]

For mixed-framework experiments, additional columns should be included:
\[
\begin{array}{c|c|c|c|c}
\text{Torch} \rightarrow \text{JAX Forward}
&
\text{JAX} \rightarrow \text{Torch Forward}
&
\text{Torch} \rightarrow \text{JAX Backward}
&
\text{JAX} \rightarrow \text{Torch Backward}
&
\text{JAX VJP Time}
\\
\hline
\cdots & \cdots & \cdots & \cdots & \cdots
\end{array}
\]

The table should make it possible to compare:
\[
\begin{aligned}
&\text{final true loss},\\
&\text{average per-step runtime},\\
&\text{total attack runtime},\\
&\text{peak GPU memory},\\
&\text{conversion overhead},\\
&\text{numerical consistency among framework combinations}.
\end{aligned}
\]

\subsection{Command-Line Interface}

The implementation should provide a command-line interface that allows each experiment to be run separately. This is important for accurate memory measurements.

For the Burgers case, the four commands can be:

\begin{verbatim}
python run_cross_framework_attack.py \
  --case burgers \
  --solver_backend jax \
  --model_backend jax \
  --norm 2 \
  --epsilon 20 \
  --alpha 0.06 \
  --steps 100 \
  --index 0 \
  --output_dir results/burgers/solver_jax_model_jax
\end{verbatim}

\begin{verbatim}
python run_cross_framework_attack.py \
  --case burgers \
  --solver_backend jax \
  --model_backend torch \
  --norm 2 \
  --epsilon 20 \
  --alpha 0.06 \
  --steps 100 \
  --index 0 \
  --output_dir results/burgers/solver_jax_model_torch
\end{verbatim}

\begin{verbatim}
python run_cross_framework_attack.py \
  --case burgers \
  --solver_backend torch \
  --model_backend jax \
  --norm 2 \
  --epsilon 20 \
  --alpha 0.06 \
  --steps 100 \
  --index 0 \
  --output_dir results/burgers/solver_torch_model_jax
\end{verbatim}

\begin{verbatim}
python run_cross_framework_attack.py \
  --case burgers \
  --solver_backend torch \
  --model_backend torch \
  --norm 2 \
  --epsilon 20 \
  --alpha 0.06 \
  --steps 100 \
  --index 0 \
  --output_dir results/burgers/solver_torch_model_torch
\end{verbatim}

For the 2D Stokes / Navier--Stokes case, the four commands can be:

\begin{verbatim}
python run_cross_framework_attack.py \
  --case ns \
  --solver_backend jax \
  --model_backend jax \
  --norm 2 \
  --epsilon 32768 \
  --alpha 5 \
  --steps 100 \
  --index 1 \
  --output_dir results/ns/solver_jax_model_jax
\end{verbatim}

\begin{verbatim}
python run_cross_framework_attack.py \
  --case ns \
  --solver_backend jax \
  --model_backend torch \
  --norm 2 \
  --epsilon 32768 \
  --alpha 5 \
  --steps 100 \
  --index 1 \
  --output_dir results/ns/solver_jax_model_torch
\end{verbatim}

\begin{verbatim}
python run_cross_framework_attack.py \
  --case ns \
  --solver_backend torch \
  --model_backend jax \
  --norm 2 \
  --epsilon 32768 \
  --alpha 5 \
  --steps 100 \
  --index 1 \
  --output_dir results/ns/solver_torch_model_jax
\end{verbatim}

\begin{verbatim}
python run_cross_framework_attack.py \
  --case ns \
  --solver_backend torch \
  --model_backend torch \
  --norm 2 \
  --epsilon 32768 \
  --alpha 5 \
  --steps 100 \
  --index 1 \
  --output_dir results/ns/solver_torch_model_torch
\end{verbatim}

After all four runs are finished for one PDE case, the comparison plots can be generated by:

\begin{verbatim}
python plot_cross_framework_attack_results.py \
  --case burgers \
  --results_root results/burgers \
  --output_dir figures/burgers_cross_framework
\end{verbatim}

and

\begin{verbatim}
python plot_cross_framework_attack_results.py \
  --case ns \
  --results_root results/ns \
  --output_dir figures/ns_cross_framework
\end{verbatim}

\subsection{Code Search and Modification Constraints}

When locating files, functions, or references in the repository, the search commands should use \texttt{grep} rather than \texttt{rg}. For example:
\begin{verbatim}
grep -R "def pgd" -n .
grep -R "attack" -n .
grep -R "torch.cuda.memory" -n .
grep -R "jax" -n .
grep -R "solver" -n .
\end{verbatim}

The code modification should be limited to the parts directly required by this experiment. The main additions should be:
\[
\begin{aligned}
&\text{a unified cross-framework attack runner},\\
&\text{a solver/model backend selector},\\
&\text{JAX--PyTorch tensor conversion utilities},\\
&\text{mixed-framework autograd wrappers},\\
&\text{per-step timing recorders},\\
&\text{per-step memory recorders},\\
&\text{result-saving utilities},\\
&\text{cross-framework plotting scripts}.
\end{aligned}
\]

The following parts should remain unchanged unless absolutely necessary:
\[
\begin{aligned}
&\text{existing model architecture},\\
&\text{existing numerical solver implementation},\\
&\text{existing data loading logic},\\
&\text{existing checkpoint loading logic},\\
&\text{existing visualization functions},\\
&\text{existing PGD logic outside the backend interface}.
\end{aligned}
\]

\subsection{Expected Final Outputs}

The final outputs of this experimental direction should include:
\[
\begin{aligned}
&\text{attack results for all four framework combinations for each PDE case},\\
&\text{per-step loss, true loss, runtime, and memory files},\\
&\text{single-run output visualizations},\\
&\text{four-combination comparison plots},\\
&\text{1D Burgers line plots and optional GIFs},\\
&\text{2D Stokes / NS heatmaps and optional GIFs},\\
&\text{difference plots across framework combinations},\\
&\text{summary tables},\\
&\text{directly executable Linux commands},\\
&\text{explicit records of JAX--PyTorch tensor conversion overhead}.
\end{aligned}
\]

The experiment should answer the following questions:
\[
\begin{aligned}
&\text{Under the same PDE case, input sample, and PGD parameters,}\\
&\text{do different JAX/PyTorch solver--model combinations produce different adversarial inputs?}\\
&\text{Do they produce different attack loss and true loss curves?}\\
&\text{Is JAX faster after compilation?}\\
&\text{Does JAX use more GPU memory?}\\
&\text{Does JAX--PyTorch conversion become a runtime bottleneck?}\\
&\text{Does mixed-framework backward propagation significantly slow down PGD?}\\
&\text{Are the final adversarial examples numerically consistent across frameworks?}\\
&\text{Is the JAX solver + PyTorch model pipeline worth keeping,}\\
&\text{or should the implementation be converted to a pure JAX or pure PyTorch pipeline?}
\end{aligned}
\]

The main conclusion should be based on both numerical attack strength and system-level cost:
\[
\text{final true loss},
\qquad
\text{attack trajectory},
\qquad
\text{total runtime},
\qquad
\text{peak GPU memory},
\qquad
\text{cross-framework conversion overhead}.
\]


\section{Task Direction 2: Systematic Comparison of Attack Losses, PGD, Power Iteration, and Solver-Gradient Control}

\subsection{Core Problem Setup}

This task studies adversarial perturbations for PDE surrogate models.

The fixed computational setting is:

\[
\text{JAX numerical solver } g(x),
\qquad
\text{PyTorch FNO surrogate model } f(x).
\]

The clean input is \(x\). The adversarial perturbation is \(\delta\). The adversarial input is

\[
x_\delta = x+\delta.
\]

The true solver-model discrepancy that we ultimately care about is

\[
L_{\mathrm{true}}(\delta)
=
L_3(\delta)
=
\|f(x+\delta)-g(x+\delta)\|_2.
\]

Here:

\[
f(x+\delta)
\]

is the PyTorch FNO model output, and

\[
g(x+\delta)
\]

is the JAX numerical solver output at the same perturbed input.

The goal is not only to produce a large final attack loss, but also to understand why the attack works, which loss is actually being optimized, and whether the optimized direction is related to the leading singular vector direction of the local Jacobian.

The experiments should answer the following questions:

\begin{enumerate}
    \item For the model self-sensitivity loss \(L_1\), is power iteration more suitable than PGD?
    \item For \(L_2\), \(L_{2,\mathrm{sg}}\), and \(L_3\), is PGD more suitable than power iteration / power-like update?
    \item If PGD is used to optimize different surrogate losses, which optimized loss produces the largest true loss growth?
    \item Does the PGD perturbation direction align with the leading right singular vector direction of the local model Jacobian?
    \item How much does solver forward information help if the solver gradient is detached?
    \item How much does full solver backward propagation help compared with detached solver forward?
    \item For recurrent 2D experiments, how do different frame-level solver modes affect attack strength?
\end{enumerate}

All scripts should be generated first. The experiments should not be run immediately, because the full attack experiments may take a long time.

\subsection{Existing Code Should Be Reused}

This task should not be implemented from scratch.

The new code should reuse the existing JAX-solver-plus-PyTorch-FNO attack pipeline already present in the codebase.

For 1D Burgers, first inspect:

\begin{verbatim}
1D_Burgers/
  perturbation_methods/
    runner_*.py
    test_*.py
    new_*.py
    ...
\end{verbatim}

Use the following commands to locate the existing attack code:

\begin{verbatim}
grep -R "perturbation" -n 1D_Burgers
grep -R "PGD" -n 1D_Burgers
grep -R "pgd" -n 1D_Burgers
grep -R "jax" -n 1D_Burgers
grep -R "torch" -n 1D_Burgers
grep -R "solver" -n 1D_Burgers
grep -R "FNO" -n 1D_Burgers
\end{verbatim}

For 2D NS / Stokes, inspect:

\begin{verbatim}
2D_NS_FNO*/
  continuous_f_attack/
  perturbation_methods/
    pgd_attack*.py
    ...
\end{verbatim}

Use:

\begin{verbatim}
grep -R "continuous_f_attack" -n .
grep -R "perturbation_methods" -n 2D_NS*
grep -R "pgd_attack" -n 2D_NS*
grep -R "jax" -n 2D_NS*
grep -R "torch" -n 2D_NS*
grep -R "solver" -n 2D_NS*
grep -R "FNO" -n 2D_NS*
\end{verbatim}

The trained PyTorch FNO model should also be reused. Do not retrain the model.

Locate checkpoints and model-loading code with:

\begin{verbatim}
grep -R "load_state_dict" -n .
grep -R "torch.load" -n .
grep -R ".pt" -n .
grep -R ".pth" -n .
grep -R "checkpoint" -n .
grep -R "ckpt" -n .
grep -R "FNO" -n .
\end{verbatim}

The JAX solver should come from the already installed solver package and the existing code path. Locate existing solver usage with:

\begin{verbatim}
grep -R "x2nox" -n .
grep -R "jax" -n .
grep -R "vmap" -n .
grep -R "jit" -n .
grep -R "solver" -n .
\end{verbatim}

The new implementation should mainly add:

\begin{itemize}
    \item loss selector;
    \item attack method selector;
    \item stop-gradient solver option;
    \item frame-mode option for recurrent 2D experiments;
    \item metric recorder;
    \item singular-vector diagnostics;
    \item plotting scripts;
    \item final summary table generator.
\end{itemize}

The new implementation should not rewrite:

\begin{itemize}
    \item the numerical solver;
    \item the PyTorch FNO architecture;
    \item the checkpoint-loading logic;
    \item the dataset-loading logic;
    \item the normalization and denormalization logic;
    \item the existing basic visualization utilities.
\end{itemize}

\subsection{Attack Losses}

Let

\[
f_\delta = f(x+\delta),
\qquad
g_\delta = g(x+\delta),
\qquad
f_0 = f(x),
\qquad
g_0 = g(x).
\]

Four losses should be considered:

\[
L_1,
\qquad
L_2,
\qquad
L_{2,\mathrm{sg}},
\qquad
L_3.
\]

\subsubsection{Loss 1: Model Self-Sensitivity Loss}

The first loss is

\[
L_1(\delta)
=
\|f(x+\delta)-f(x)\|_2.
\]

This loss does not use the solver output during optimization. It only measures how much the PyTorch FNO model changes when the input is perturbed.

Using local linearization,

\[
f(x+\delta)
\approx
f(x)+J_f(x)\delta.
\]

Therefore,

\[
f(x+\delta)-f(x)
\approx
J_f(x)\delta.
\]

Hence,

\[
L_1(\delta)
\approx
\|J_f(x)\delta\|_2.
\]

For the squared loss,

\[
L_1^2(\delta)
\approx
\delta^\top J_f(x)^\top J_f(x)\delta.
\]

Thus, under an \(L_2\)-norm perturbation constraint, this is locally a leading singular-vector problem:

\[
\max_{\|\delta\|_2 \leq \epsilon}
\|J_f(x)\delta\|_2.
\]

The optimal local direction is the leading right singular vector of \(J_f(x)\), equivalently the leading eigenvector of

\[
J_f(x)^\top J_f(x).
\]

This is the loss where power iteration should be most meaningful.

\subsubsection{Loss 2: Fixed Solver Target Loss}

The second loss is

\[
L_2(\delta)
=
\|f(x+\delta)-g(x)\|_2.
\]

Here the solver output is computed only once at the clean input:

\[
y = g(x),
\]

and then treated as a fixed target during the attack.

The local expansion is

\[
f(x+\delta)-g(x)
\approx
f(x)-g(x)+J_f(x)\delta.
\]

Therefore,

\[
L_2^2(\delta)
\approx
\|f(x)-g(x)\|_2^2
+
2(f(x)-g(x))^\top J_f(x)\delta
+
\delta^\top J_f(x)^\top J_f(x)\delta.
\]

This loss contains a linear term:

\[
2(f(x)-g(x))^\top J_f(x)\delta.
\]

Therefore, \(L_2\) is not a pure Rayleigh quotient problem. Power iteration is only a heuristic baseline for this loss. PGD may be more suitable.

\subsubsection{Loss 2 Stop-Gradient Variant: Moving Solver Target without Solver Backward}

The stop-gradient version is

\[
L_{2,\mathrm{sg}}(\delta)
=
\|f(x+\delta)-\operatorname{stopgrad}(g(x+\delta))\|_2.
\]

This loss uses the perturbed solver output in the forward pass:

\[
g(x+\delta),
\]

but does not propagate gradient through the solver.

That means:

\begin{itemize}
    \item the solver output changes with the current adversarial input;
    \item the loss value uses the true perturbed solver output;
    \item the backward pass only differentiates through the PyTorch FNO model;
    \item the solver part is detached.
\end{itemize}

This loss is between Loss 2 and Loss 3:

\[
L_2:
\text{ fixed solver target},
\]

\[
L_{2,\mathrm{sg}}:
\text{ moving solver target, solver gradient detached},
\]

\[
L_3:
\text{ moving solver target, solver gradient included}.
\]

The key question is whether using the updated solver target in the forward pass is already enough to increase the true loss, even without solver backward propagation.

\subsubsection{Loss 3: True Solver-Model Error Loss}

The third loss is the true attack target:

\[
L_3(\delta)
=
\|f(x+\delta)-g(x+\delta)\|_2.
\]

Using local linearization,

\[
f(x+\delta)
\approx
f(x)+J_f(x)\delta,
\]

\[
g(x+\delta)
\approx
g(x)+J_g(x)\delta.
\]

Therefore,

\[
f(x+\delta)-g(x+\delta)
\approx
f(x)-g(x)+(J_f(x)-J_g(x))\delta.
\]

The squared loss satisfies

\[
L_3^2(\delta)
\approx
\|f(x)-g(x)\|_2^2
+
2(f(x)-g(x))^\top (J_f(x)-J_g(x))\delta
+
\delta^\top
(J_f(x)-J_g(x))^\top
(J_f(x)-J_g(x))
\delta.
\]

This also contains a linear term. It is not a pure Rayleigh quotient problem.

However, it directly optimizes the true solver-model discrepancy:

\[
\|f(x+\delta)-g(x+\delta)\|_2.
\]

Therefore, it is expected to be the most faithful attack loss for the final evaluation.

\subsection{Important Constraint: Do Not Explicitly Construct the Jacobian}

The Jacobians

\[
J_f(x),
\qquad
J_g(x),
\qquad
J_f(x)-J_g(x)
\]

must not be explicitly constructed.

The input and output dimensions are too large. Explicitly forming the Jacobian will likely cause GPU memory overflow.

All Jacobian-related quantities should be computed implicitly using JVP, VJP, or autograd vector products.

For Loss 1, the local power iteration should use:

\[
v_k
\longmapsto
J_f(x)v_k
\longmapsto
J_f(x)^\top J_f(x)v_k
\longmapsto
v_{k+1}.
\]

The update is

\[
v_{k+1}
=
\frac{
J_f(x)^\top J_f(x)v_k
}{
\|J_f(x)^\top J_f(x)v_k\|_2 + 10^{-12}
}.
\]

The estimated singular value is

\[
\hat{\sigma}_k
=
\frac{
\|J_f(x)v_k\|_2
}{
\|v_k\|_2 + 10^{-12}
}.
\]

The Rayleigh quotient is

\[
R_k
=
\frac{
v_k^\top J_f(x)^\top J_f(x)v_k
}{
v_k^\top v_k + 10^{-12}
}
=
\frac{
\|J_f(x)v_k\|_2^2
}{
\|v_k\|_2^2 + 10^{-12}
}.
\]

Thus,

\[
\hat{\sigma}_k
=
\sqrt{R_k}.
\]

For PGD, the perturbation gradient should be computed by standard autograd, but the full Jacobian should never be materialized.

\subsection{Two Main Experimental Goals}

This task has two main experiments.

\subsection{Experiment A: PGD versus Power Iteration / Power-Like Optimization}

\subsubsection{Goal of Experiment A}

Experiment A compares two optimization methods:

\[
\text{PGD}
\qquad
\text{versus}
\qquad
\text{power iteration / power-like update}.
\]

This comparison should be performed for each loss:

\[
L_1,
\qquad
L_2,
\qquad
L_{2,\mathrm{sg}},
\qquad
L_3.
\]

Experiment A has two equally important goals.

\subsubsection{Goal A1: Verify Whether Power Iteration is More Suitable for Loss 1}

For Loss 1,

\[
L_1(\delta)
=
\|f(x+\delta)-f(x)\|_2,
\]

the local squared objective is approximately

\[
L_1^2(\delta)
\approx
\delta^\top J_f(x)^\top J_f(x)\delta.
\]

Therefore, Loss 1 is locally close to a Rayleigh quotient / leading singular-vector problem.

The experiment should test whether power iteration finds a stronger attack direction than PGD for this loss.

The key diagnostics are:

\[
L_1(\delta_k),
\]

\[
L_{\mathrm{true}}(\delta_k),
\]

\[
\hat{\sigma}_k,
\]

\[
R_k,
\]

and

\[
\left|
\cos
\left(
\frac{\delta_{\mathrm{PGD}}}{\|\delta_{\mathrm{PGD}}\|_2+10^{-12}},
v_{\max}
\right)
\right|.
\]

Here \(v_{\max}\) is the estimated leading right singular vector of \(J_f(x)\).

This comparison answers whether PGD is actually finding the same direction as the dominant singular vector direction.

\subsubsection{Goal A2: Verify Whether PGD is More Suitable for Loss 2, Stop-Gradient Loss 2, and Loss 3}

For Loss 2,

\[
L_2(\delta)
=
\|f(x+\delta)-g(x)\|_2,
\]

the local squared-loss expansion contains a linear term:

\[
2(f(x)-g(x))^\top J_f(x)\delta.
\]

For the stop-gradient loss,

\[
L_{2,\mathrm{sg}}(\delta)
=
\|f(x+\delta)-\operatorname{stopgrad}(g(x+\delta))\|_2,
\]

the target changes in the forward pass because \(g(x+\delta)\) changes with \(\delta\), but the solver gradient is detached.

For Loss 3,

\[
L_3(\delta)
=
\|f(x+\delta)-g(x+\delta)\|_2,
\]

the local squared-loss expansion contains the linear term

\[
2(f(x)-g(x))^\top (J_f(x)-J_g(x))\delta.
\]

Thus, \(L_2\), \(L_{2,\mathrm{sg}}\), and \(L_3\) are not pure Rayleigh quotient problems.

Power iteration is therefore only a heuristic local sensitivity baseline for these losses.

Experiment A should test whether PGD gives:

\[
\text{larger optimized loss growth},
\]

\[
\text{larger true loss growth},
\]

\[
\text{larger final true loss},
\]

and

\[
\text{more stable per-step improvement}
\]

than power iteration / power-like update for \(L_2\), \(L_{2,\mathrm{sg}}\), and \(L_3\).

The important comparisons are:

\[
\text{PGD on } L_2
\quad
\text{vs.}
\quad
\text{power-like update on } L_2,
\]

\[
\text{PGD on } L_{2,\mathrm{sg}}
\quad
\text{vs.}
\quad
\text{power-like update on } L_{2,\mathrm{sg}},
\]

\[
\text{PGD on } L_3
\quad
\text{vs.}
\quad
\text{power-like update on } L_3.
\]

The expected result is:

\[
\text{Power iteration may be more suitable for } L_1,
\]

while

\[
\text{PGD may be more suitable for } L_2,\ L_{2,\mathrm{sg}},\ \text{and } L_3.
\]

The conclusion for \(L_2\), \(L_{2,\mathrm{sg}}\), and \(L_3\) should not be based only on the optimized loss. The main evidence should be the true loss:

\[
L_{\mathrm{true}}(\delta)
=
\|f(x+\delta)-g(x+\delta)\|_2.
\]

\subsubsection{Summary of Experiment A}

\[
\begin{array}{c|c|c|c}
\text{Loss} & \text{Local structure} & \text{Comparison} & \text{Main hypothesis} \\
\hline
L_1
&
\delta^\top J_f^\top J_f\delta
&
\text{PGD vs. power iteration}
&
\text{power iteration may be more suitable}
\\
L_2
&
\text{linear term}+\delta^\top J_f^\top J_f\delta
&
\text{PGD vs. power-like}
&
\text{PGD may be more suitable}
\\
L_{2,\mathrm{sg}}
&
\text{moving detached solver target}
&
\text{PGD vs. power-like}
&
\text{PGD may be more suitable}
\\
L_3
&
\text{linear term}+\delta^\top(J_f-J_g)^\top(J_f-J_g)\delta
&
\text{PGD vs. power-like}
&
\text{PGD may be more suitable}
\end{array}
\]

\subsection{Experiment B: Which Optimized Loss Produces the Largest True Loss Growth?}

\subsubsection{Goal of Experiment B}

Experiment B fixes the optimizer to PGD.

Then it compares four optimized losses:

\[
L_1,
\qquad
L_2,
\qquad
L_{2,\mathrm{sg}},
\qquad
L_3.
\]

The central question is:

\[
\text{Which optimized loss produces the largest growth of the true solver-model loss?}
\]

The true loss is always

\[
L_{\mathrm{true}}(\delta)
=
L_3(\delta)
=
\|f(x+\delta)-g(x+\delta)\|_2.
\]

\subsubsection{Both Surrogate Loss and True Loss Must Be Recorded}

For every PGD run in Experiment B, two types of losses must be recorded at every iteration.

First, record the loss that is actually being optimized:

\[
L_{\mathrm{opt}}(\delta_k).
\]

This is the surrogate, fake, approximate, or selected attack loss. Depending on the run,

\[
L_{\mathrm{opt}}
\in
\{L_1,\ L_2,\ L_{2,\mathrm{sg}},\ L_3\}.
\]

Second, always record the true solver-model loss:

\[
L_{\mathrm{true}}(\delta_k)
=
\|f(x+\delta_k)-g(x+\delta_k)\|_2.
\]

This distinction is essential.

For example, if the attack optimizes \(L_1\), then at each step record both:

\[
L_{\mathrm{opt}}(\delta_k)
=
L_1(\delta_k)
=
\|f(x+\delta_k)-f(x)\|_2,
\]

and

\[
L_{\mathrm{true}}(\delta_k)
=
L_3(\delta_k)
=
\|f(x+\delta_k)-g(x+\delta_k)\|_2.
\]

If the attack optimizes \(L_2\), record both:

\[
L_{\mathrm{opt}}(\delta_k)
=
L_2(\delta_k)
=
\|f(x+\delta_k)-g(x)\|_2,
\]

and

\[
L_{\mathrm{true}}(\delta_k)
=
L_3(\delta_k)
=
\|f(x+\delta_k)-g(x+\delta_k)\|_2.
\]

If the attack optimizes \(L_{2,\mathrm{sg}}\), record both:

\[
L_{\mathrm{opt}}(\delta_k)
=
L_{2,\mathrm{sg}}(\delta_k)
=
\|f(x+\delta_k)-\operatorname{stopgrad}(g(x+\delta_k))\|_2,
\]

and

\[
L_{\mathrm{true}}(\delta_k)
=
L_3(\delta_k)
=
\|f(x+\delta_k)-g(x+\delta_k)\|_2.
\]

If the attack optimizes \(L_3\), then

\[
L_{\mathrm{opt}}(\delta_k)
=
L_{\mathrm{true}}(\delta_k)
=
L_3(\delta_k).
\]

Even in this case, both columns should still be saved separately in the result files for consistency.

\subsubsection{Main Evidence for Experiment B}

The most important plot is the true loss curve:

\[
L_{\mathrm{true}}(\delta_k)
=
\|f(x+\delta_k)-g(x+\delta_k)\|_2.
\]

The main comparison is:

\[
\text{PGD optimizing } L_1
\quad
\text{vs.}
\quad
\text{PGD optimizing } L_2
\quad
\text{vs.}
\quad
\text{PGD optimizing } L_{2,\mathrm{sg}}
\quad
\text{vs.}
\quad
\text{PGD optimizing } L_3.
\]

The expected result is:

\[
\text{Directly optimizing } L_3
\text{ should usually produce the largest true loss growth.}
\]

The stop-gradient loss may produce an intermediate result:

\[
L_2
\quad
\longrightarrow
\quad
L_{2,\mathrm{sg}}
\quad
\longrightarrow
\quad
L_3.
\]

However, this is an empirical hypothesis. The experiment may confirm or refute it.

\subsection{Experiment C: Frame-Mode Control for 2D Recurrent Experiments}

\subsubsection{Motivation}

For 1D Burgers, the output is usually a single final frame, so testing \(L_2\), \(L_{2,\mathrm{sg}}\), and \(L_3\) is relatively direct.

For 2D recurrent FNO / NS-style models, the model may use a sequence of frames. In this case, solver outputs may enter different frames in different ways.

Therefore, each frame should be controlled by a tag.

\subsubsection{Frame Tags}

Each frame has one of three possible tags:

\[
\texttt{a},
\qquad
\texttt{d},
\qquad
\texttt{w}.
\]

The meaning is:

\[
\texttt{a}
=
\text{absent solver frame}.
\]

This means the solver frame is not used in the attack computation for that frame.

\[
\texttt{d}
=
\text{detached solver frame}.
\]

This means the solver frame is used in the forward pass, but detached in the backward pass.

\[
\texttt{w}
=
\text{with-gradient solver frame}.
\]

This means the solver frame is used in the forward pass and gradient is propagated through it in the backward pass.

For a recurrent window of length 10, the first round should only test:

\[
\texttt{aaaaaaaaaa},
\qquad
\texttt{dddddddddd},
\qquad
\texttt{wwwwwwwwww}.
\]

Do not test all mixed combinations at this stage. The number of combinations is too large.

The comparison is:

\[
\texttt{all-a}
\quad
\text{vs.}
\quad
\texttt{all-d}
\quad
\text{vs.}
\quad
\texttt{all-w}.
\]

This tests:

\begin{itemize}
    \item whether solver forward values help;
    \item whether solver backward gradients help;
    \item whether full solver-model recurrent coupling gives stronger attacks.
\end{itemize}

\subsection{Fixed Main Parameter Settings}

The first round of experiments should use fixed parameters. Do not perform a large hyperparameter search at this stage.

\subsubsection{1D Burgers Main \(L_2\)-Norm Parameters}

Use:

\[
\text{norm}=2,
\qquad
\epsilon=20,
\qquad
\alpha=0.06,
\qquad
\text{steps}=100,
\qquad
\text{index}=0.
\]

The command-line arguments are:

\begin{verbatim}
--norm 2
--epsilon 20
--alpha 0.06
--steps 100
--index 0
\end{verbatim}

\subsubsection{1D Burgers Additional \(L_\infty\)-Norm Check}

Use:

\[
\text{norm}=\infty,
\qquad
\epsilon=0.5,
\qquad
\alpha=0.05,
\qquad
\text{steps}=100,
\qquad
\text{index}=0.
\]

The command-line arguments are:

\begin{verbatim}
--norm inf
--epsilon 0.5
--alpha 0.05
--steps 100
--index 0
\end{verbatim}

\subsubsection{2D NS / Stokes Main \(L_2\)-Norm Parameters}

Use:

\[
\text{norm}=2,
\qquad
\epsilon=32768,
\qquad
\alpha=5,
\qquad
\text{steps}=100,
\qquad
\text{index}=1.
\]

The command-line arguments are:

\begin{verbatim}
--norm 2
--epsilon 32768
--alpha 5
--steps 100
--index 1
\end{verbatim}

\subsubsection{2D NS / Stokes Additional \(L_\infty\)-Norm Check}

Use:

\[
\text{norm}=\infty,
\qquad
\epsilon=1,
\qquad
\alpha=0.01,
\qquad
\text{steps}=100,
\qquad
\text{index}=2.
\]

The command-line arguments are:

\begin{verbatim}
--norm inf
--epsilon 1
--alpha 0.01
--steps 100
--index 2
\end{verbatim}

\subsubsection{2D Recurrent Frame-Mode Parameters}

For the frame-mode test, use the main 2D \(L_2\)-norm parameters:

\[
\text{norm}=2,
\qquad
\epsilon=32768,
\qquad
\alpha=5,
\qquad
\text{steps}=100,
\qquad
\text{index}=1.
\]

Only test:

\[
\texttt{aaaaaaaaaa},
\qquad
\texttt{dddddddddd},
\qquad
\texttt{wwwwwwwwww}.
\]

The command-line arguments are:

\begin{verbatim}
--norm 2
--epsilon 32768
--alpha 5
--steps 100
--index 1
--frame_mode aaaaaaaaaa
\end{verbatim}

\begin{verbatim}
--norm 2
--epsilon 32768
--alpha 5
--steps 100
--index 1
--frame_mode dddddddddd
\end{verbatim}

\begin{verbatim}
--norm 2
--epsilon 32768
--alpha 5
--steps 100
--index 1
--frame_mode wwwwwwwwww
\end{verbatim}

\subsection{Metrics to Record at Every Attack Step}

For every run, record the following values at each iteration \(k\).

\subsubsection{Basic Variables}

Record:

\[
k,
\qquad
\delta_k,
\qquad
x_k=x+\delta_k.
\]

Record perturbation norms:

\[
\|\delta_k\|_2,
\qquad
\|\delta_k\|_\infty,
\qquad
\|\delta_k\|_p.
\]

Record budget usage:

\[
\frac{\|\delta_k\|_p}{\epsilon}.
\]

\subsubsection{Loss Variables}

Record all of the following, even if only one is being optimized:

\[
L_1(\delta_k)
=
\|f(x+\delta_k)-f(x)\|_2,
\]

\[
L_2(\delta_k)
=
\|f(x+\delta_k)-g(x)\|_2,
\]

\[
L_{2,\mathrm{sg}}(\delta_k)
=
\|f(x+\delta_k)-\operatorname{stopgrad}(g(x+\delta_k))\|_2,
\]

\[
L_3(\delta_k)
=
\|f(x+\delta_k)-g(x+\delta_k)\|_2.
\]

Also record:

\[
L_{\mathrm{opt}}(\delta_k),
\]

where \(L_{\mathrm{opt}}\) is the actual optimized loss in the current run.

Always record:

\[
L_{\mathrm{true}}(\delta_k)
=
L_3(\delta_k).
\]

This is especially important for Experiment B, because the optimized loss may be a surrogate loss, while the final evaluation must always use the true solver-model loss.

Also record squared losses:

\[
L_1^2,
\qquad
L_2^2,
\qquad
L_{2,\mathrm{sg}}^2,
\qquad
L_3^2.
\]

\subsubsection{Loss Growth Variables}

Record initial and final values:

\[
L_{\mathrm{opt}}(\delta_0),
\qquad
L_{\mathrm{opt}}(\delta_T),
\]

\[
L_{\mathrm{true}}(\delta_0),
\qquad
L_{\mathrm{true}}(\delta_T).
\]

Record absolute growth:

\[
\Delta L_{\mathrm{opt}}
=
L_{\mathrm{opt}}(\delta_T)
-
L_{\mathrm{opt}}(\delta_0),
\]

\[
\Delta L_{\mathrm{true}}
=
L_{\mathrm{true}}(\delta_T)
-
L_{\mathrm{true}}(\delta_0).
\]

Record growth ratios:

\[
\text{growth ratio}_{\mathrm{opt}}
=
\frac{
L_{\mathrm{opt}}(\delta_T)
}{
L_{\mathrm{opt}}(\delta_0)+10^{-12}
},
\]

\[
\text{growth ratio}_{\mathrm{true}}
=
\frac{
L_{\mathrm{true}}(\delta_T)
}{
L_{\mathrm{true}}(\delta_0)+10^{-12}
}.
\]

Also record per-step true loss growth:

\[
\frac{
L_{\mathrm{true}}(\delta_{k+1})
}{
L_{\mathrm{true}}(\delta_k)+10^{-12}
}.
\]

\subsubsection{Gradient Variables}

For PGD-type methods, record:

\[
\nabla_\delta L_{\mathrm{opt}}(\delta_k).
\]

Record:

\[
\|\nabla_\delta L_{\mathrm{opt}}(\delta_k)\|_2,
\]

\[
\|\nabla_\delta L_{\mathrm{opt}}(\delta_k)\|_\infty.
\]

Record cosine similarity between consecutive gradients:

\[
\cos(\nabla L_k,\nabla L_{k-1})
=
\frac{
\langle \nabla L_k,\nabla L_{k-1}\rangle
}{
\|\nabla L_k\|_2\|\nabla L_{k-1}\|_2+10^{-12}
}.
\]

This measures whether the attack direction is stable or oscillatory.

\subsubsection{Direction Variables}

Record normalized perturbation direction:

\[
u_k
=
\frac{\delta_k}{\|\delta_k\|_2+10^{-12}}.
\]

Record normalized gradient direction:

\[
r_k
=
\frac{
\nabla_\delta L_{\mathrm{opt}}(\delta_k)
}{
\|\nabla_\delta L_{\mathrm{opt}}(\delta_k)\|_2+10^{-12}
}.
\]

Record:

\[
\cos(u_k,r_k)
=
\frac{
\langle u_k,r_k\rangle
}{
\|u_k\|_2\|r_k\|_2+10^{-12}
}.
\]

For power iteration, record the normalized power direction:

\[
v_k.
\]

Record direction alignment:

\[
\cos(u_k,v_k)
=
\frac{
\langle u_k,v_k\rangle
}{
\|u_k\|_2\|v_k\|_2+10^{-12}
}.
\]

Also record:

\[
|\cos(u_k,v_k)|.
\]

The absolute value is important because singular vectors are sign-invariant.

\subsubsection{Singular-Vector and Spectral Diagnostics}

For Loss 1, record implicit spectral quantities.

For a power direction \(v_k\), compute \(J_f(x)v_k\) by JVP.

Record:

\[
\|J_f(x)v_k\|_2.
\]

Record estimated singular value:

\[
\hat{\sigma}_k
=
\frac{
\|J_f(x)v_k\|_2
}{
\|v_k\|_2+10^{-12}
}.
\]

Record Rayleigh quotient:

\[
R_k
=
\frac{
\|J_f(x)v_k\|_2^2
}{
\|v_k\|_2^2+10^{-12}
}.
\]

Record:

\[
\sqrt{R_k}.
\]

Also record:

\[
\|J_f(x)^\top J_f(x)v_k\|_2.
\]

For PGD versus singular-vector comparison, estimate the leading singular vector \(v_{\max}\), then compare:

\[
u_{\mathrm{PGD}}
=
\frac{
\delta_{\mathrm{PGD}}
}{
\|\delta_{\mathrm{PGD}}\|_2+10^{-12}
}.
\]

Record:

\[
\cos(u_{\mathrm{PGD}},v_{\max}),
\]

and

\[
|\cos(u_{\mathrm{PGD}},v_{\max})|.
\]

Also compare the same perturbation budget applied to the singular-vector direction:

\[
\delta_{\mathrm{sv}}
=
\epsilon v_{\max}.
\]

Evaluate:

\[
L_1(\delta_{\mathrm{PGD}})
\quad
\text{vs.}
\quad
L_1(\delta_{\mathrm{sv}}),
\]

and

\[
L_{\mathrm{true}}(\delta_{\mathrm{PGD}})
\quad
\text{vs.}
\quad
L_{\mathrm{true}}(\delta_{\mathrm{sv}}).
\]

This determines whether the singular-vector direction is only strong for model self-sensitivity, or also strong for true solver-model discrepancy.

\subsubsection{Timing Variables}

Record:

\[
t_{\mathrm{forward},k},
\qquad
t_{\mathrm{backward},k},
\qquad
t_{\mathrm{step},k},
\]

\[
t_{\mathrm{solver},k},
\qquad
t_{\mathrm{model},k}.
\]

Also record total attack time:

\[
T_{\mathrm{total}}.
\]

This is important because \(L_3\) may produce stronger attacks but may also be more expensive.

\subsection{Output Files}

Each experiment directory should save:

\begin{verbatim}
config.json
metrics.csv
summary.json

optimized_loss.npy
true_loss.npy
loss1.npy
loss2.npy
loss2_stopgrad.npy
loss3.npy

optimized_loss_squared.npy
true_loss_squared.npy

delta_norm_l2.npy
delta_norm_linf.npy
delta_budget_ratio.npy

gradient_norm_l2.npy
gradient_norm_linf.npy
gradient_cosine.npy
delta_gradient_cosine.npy

step_time.npy
forward_time.npy
backward_time.npy
solver_time.npy
model_time.npy

x_clean.npy
x_adv.npy
delta_final.npy

model_output_clean.npy
solver_output_clean.npy
model_output_adv.npy
solver_output_adv.npy
difference_adv.npy
\end{verbatim}

For power iteration / power-like methods, also save:

\begin{verbatim}
estimated_singular_value.npy
rayleigh_quotient.npy
sqrt_rayleigh_quotient.npy
power_direction_norm.npy
power_direction_final.npy
jvp_norm.npy
vjp_norm.npy
jtjv_norm.npy
\end{verbatim}

For PGD versus singular-vector comparison, also save:

\begin{verbatim}
pgd_direction_final.npy
singular_vector_direction.npy
cos_pgd_singular_vector.npy
abs_cos_pgd_singular_vector.npy
loss1_on_pgd_direction.npy
loss1_on_singular_vector.npy
true_loss_on_pgd_direction.npy
true_loss_on_singular_vector.npy
\end{verbatim}

For 2D recurrent frame-mode experiments, also save:

\begin{verbatim}
frame_mode_string.txt
frame_mode_per_step.json
all_a_metrics.csv
all_d_metrics.csv
all_w_metrics.csv
\end{verbatim}

\subsection{Summary Table}

Each run should produce one row in a final summary table.

The table should include:

\begin{verbatim}
case
loss_type
attack_method
norm
epsilon
alpha
steps
index
frame_mode
uses_solver_forward
uses_solver_backward

initial_optimized_loss
final_optimized_loss
optimized_loss_growth
optimized_loss_growth_ratio

initial_true_loss
final_true_loss
true_loss_growth
true_loss_growth_ratio

final_delta_l2
final_delta_linf
final_gradient_l2
final_gradient_linf

total_time
mean_step_time

estimated_sigma_final
rayleigh_quotient_final
cos_pgd_singular_vector
abs_cos_pgd_singular_vector
\end{verbatim}

The final table should make it possible to answer:

\begin{enumerate}
    \item Which method gives the largest optimized loss?
    \item Which method gives the largest true loss?
    \item How many times did the optimized loss increase?
    \item How many times did the true loss increase?
    \item Does PGD align with the leading singular-vector direction?
    \item Does detached solver forward help?
    \item Does full solver backward help more?
    \item How much extra time does each method cost?
\end{enumerate}

\subsection{Visualization Requirements}

\subsubsection{Single-Run Diagnostic Plots}

For every run, plot:

\begin{enumerate}
    \item optimized loss curve;
    \item true loss curve;
    \item \(L_1\), \(L_2\), \(L_{2,\mathrm{sg}}\), and \(L_3\) curves together;
    \item optimized loss growth ratio;
    \item true loss growth ratio;
    \item perturbation \(L_2\) norm;
    \item perturbation \(L_\infty\) norm;
    \item perturbation budget ratio;
    \item gradient \(L_2\) norm;
    \item gradient \(L_\infty\) norm;
    \item gradient cosine similarity;
    \item delta-gradient cosine similarity;
    \item step time;
    \item forward time and backward time.
\end{enumerate}

\subsubsection{Experiment A Plots}

For each PDE case and each loss, compare PGD and power-like results.

Plot:

\begin{enumerate}
    \item optimized loss: PGD vs power-like;
    \item true loss: PGD vs power-like;
    \item optimized loss growth ratio: PGD vs power-like;
    \item true loss growth ratio: PGD vs power-like;
    \item perturbation norm: PGD vs power-like;
    \item gradient norm: PGD vs power-like;
    \item step time: PGD vs power-like;
    \item final adversarial output comparison;
    \item final solver-model difference field:
    \[
    f(x+\delta)-g(x+\delta).
    \]
\end{enumerate}

For Loss 1, additionally plot:

\begin{enumerate}
    \item estimated singular value curve;
    \item Rayleigh quotient curve;
    \item \( \|J_f v_k\|_2 \) curve;
    \item \( \|J_f^\top J_f v_k\|_2 \) curve;
    \item cosine similarity between PGD direction and power direction;
    \item final PGD perturbation versus leading singular-vector perturbation.
\end{enumerate}

\subsubsection{Experiment B Plots}

For Experiment B, place four PGD runs in the same plot:

\[
\text{PGD optimizing } L_1,
\]

\[
\text{PGD optimizing } L_2,
\]

\[
\text{PGD optimizing } L_{2,\mathrm{sg}},
\]

\[
\text{PGD optimizing } L_3.
\]

The most important plot is:

\[
L_{\mathrm{true}}(\delta_k)
=
\|f(x+\delta_k)-g(x+\delta_k)\|_2.
\]

Also plot:

\[
L_{\mathrm{opt}}(\delta_k).
\]

Therefore, Experiment B must always include two families of curves:

\[
\text{surrogate / optimized / fake loss curves},
\]

and

\[
\text{true solver-model loss curves}.
\]

The true loss plot is the main evaluation plot. The optimized loss plot explains what the attack was actually optimizing.

\subsubsection{Frame-Mode Plots}

For the 2D recurrent case, compare:

\[
\texttt{aaaaaaaaaa},
\qquad
\texttt{dddddddddd},
\qquad
\texttt{wwwwwwwwww}.
\]

Plot:

\begin{enumerate}
    \item optimized loss curves;
    \item true loss curves;
    \item true loss growth ratio curves;
    \item solver time curves;
    \item backward time curves;
    \item final adversarial output fields;
    \item final solver-model difference fields.
\end{enumerate}

This should show the difference between:

\[
\text{no solver forward},
\]

\[
\text{solver forward but detached backward},
\]

and

\[
\text{solver forward with solver backward}.
\]

\subsubsection{Spatial and Temporal Visualizations}

For 1D Burgers, plot:

\begin{enumerate}
    \item clean input \(x\);
    \item adversarial input \(x+\delta\);
    \item perturbation \(\delta\);
    \item clean model output \(f(x)\);
    \item clean solver output \(g(x)\);
    \item adversarial model output \(f(x+\delta)\);
    \item adversarial solver output \(g(x+\delta)\);
    \item adversarial difference:
    \[
    f(x+\delta)-g(x+\delta).
    \]
\end{enumerate}

For 2D NS / Stokes, plot:

\begin{enumerate}
    \item clean input field;
    \item perturbation field;
    \item adversarial input field;
    \item model output field;
    \item solver output field;
    \item model-solver difference field;
    \item absolute difference heatmap.
\end{enumerate}

If the output has multiple time frames, make GIFs for:

\begin{enumerate}
    \item clean solver trajectory;
    \item adversarial solver trajectory;
    \item clean model trajectory;
    \item adversarial model trajectory;
    \item adversarial difference trajectory.
\end{enumerate}

When generating GIFs, temporary PNG frames may be saved first. After the GIF is successfully created, delete those temporary PNG files.

The PNG frames are only intermediate files. The GIF should be kept as the final visualization.

\subsection{Script Names and Command-Line Interface}

The main script should be:

\begin{verbatim}
run_loss_comparison_attack.py
\end{verbatim}

Recommended command-line arguments:

\begin{verbatim}
--case burgers / ns
--solver_backend jax
--model_backend torch
--loss_type loss1 / loss2 / loss2_stopgrad / loss3
--attack_method pgd / power / generalized_power
--norm 2 / inf
--epsilon
--alpha
--steps
--index
--frame_mode
--output_dir
--save_every
--make_gif
--cleanup_png
\end{verbatim}

The plotting and summary scripts should be:

\begin{verbatim}
plot_single_run_diagnostics.py
plot_loss_method_comparison.py
plot_true_loss_comparison.py
plot_frame_mode_comparison.py
make_attack_summary_table.py
\end{verbatim}

\subsection{Recommended Commands}

\subsubsection{1D Burgers: Experiment A}

\begin{verbatim}
python run_loss_comparison_attack.py \
  --case burgers \
  --solver_backend jax \
  --model_backend torch \
  --loss_type loss1 \
  --attack_method pgd \
  --norm 2 \
  --epsilon 20 \
  --alpha 0.06 \
  --steps 100 \
  --index 0 \
  --output_dir results/loss_compare/burgers/loss1_pgd

python run_loss_comparison_attack.py \
  --case burgers \
  --solver_backend jax \
  --model_backend torch \
  --loss_type loss1 \
  --attack_method power \
  --norm 2 \
  --epsilon 20 \
  --alpha 0.06 \
  --steps 100 \
  --index 0 \
  --output_dir results/loss_compare/burgers/loss1_power

python run_loss_comparison_attack.py \
  --case burgers \
  --solver_backend jax \
  --model_backend torch \
  --loss_type loss2 \
  --attack_method pgd \
  --norm 2 \
  --epsilon 20 \
  --alpha 0.06 \
  --steps 100 \
  --index 0 \
  --output_dir results/loss_compare/burgers/loss2_pgd

python run_loss_comparison_attack.py \
  --case burgers \
  --solver_backend jax \
  --model_backend torch \
  --loss_type loss2 \
  --attack_method power \
  --norm 2 \
  --epsilon 20 \
  --alpha 0.06 \
  --steps 100 \
  --index 0 \
  --output_dir results/loss_compare/burgers/loss2_power

python run_loss_comparison_attack.py \
  --case burgers \
  --solver_backend jax \
  --model_backend torch \
  --loss_type loss2_stopgrad \
  --attack_method pgd \
  --norm 2 \
  --epsilon 20 \
  --alpha 0.06 \
  --steps 100 \
  --index 0 \
  --output_dir results/loss_compare/burgers/loss2_stopgrad_pgd

python run_loss_comparison_attack.py \
  --case burgers \
  --solver_backend jax \
  --model_backend torch \
  --loss_type loss2_stopgrad \
  --attack_method power \
  --norm 2 \
  --epsilon 20 \
  --alpha 0.06 \
  --steps 100 \
  --index 0 \
  --output_dir results/loss_compare/burgers/loss2_stopgrad_power

python run_loss_comparison_attack.py \
  --case burgers \
  --solver_backend jax \
  --model_backend torch \
  --loss_type loss3 \
  --attack_method pgd \
  --norm 2 \
  --epsilon 20 \
  --alpha 0.06 \
  --steps 100 \
  --index 0 \
  --output_dir results/loss_compare/burgers/loss3_pgd

python run_loss_comparison_attack.py \
  --case burgers \
  --solver_backend jax \
  --model_backend torch \
  --loss_type loss3 \
  --attack_method power \
  --norm 2 \
  --epsilon 20 \
  --alpha 0.06 \
  --steps 100 \
  --index 0 \
  --output_dir results/loss_compare/burgers/loss3_power
\end{verbatim}

\subsubsection{2D NS / Stokes: Experiment A}

\begin{verbatim}
python run_loss_comparison_attack.py \
  --case ns \
  --solver_backend jax \
  --model_backend torch \
  --loss_type loss1 \
  --attack_method pgd \
  --norm 2 \
  --epsilon 32768 \
  --alpha 5 \
  --steps 100 \
  --index 1 \
  --output_dir results/loss_compare/ns/loss1_pgd

python run_loss_comparison_attack.py \
  --case ns \
  --solver_backend jax \
  --model_backend torch \
  --loss_type loss1 \
  --attack_method power \
  --norm 2 \
  --epsilon 32768 \
  --alpha 5 \
  --steps 100 \
  --index 1 \
  --output_dir results/loss_compare/ns/loss1_power

python run_loss_comparison_attack.py \
  --case ns \
  --solver_backend jax \
  --model_backend torch \
  --loss_type loss2 \
  --attack_method pgd \
  --norm 2 \
  --epsilon 32768 \
  --alpha 5 \
  --steps 100 \
  --index 1 \
  --output_dir results/loss_compare/ns/loss2_pgd

python run_loss_comparison_attack.py \
  --case ns \
  --solver_backend jax \
  --model_backend torch \
  --loss_type loss2 \
  --attack_method power \
  --norm 2 \
  --epsilon 32768 \
  --alpha 5 \
  --steps 100 \
  --index 1 \
  --output_dir results/loss_compare/ns/loss2_power

python run_loss_comparison_attack.py \
  --case ns \
  --solver_backend jax \
  --model_backend torch \
  --loss_type loss2_stopgrad \
  --attack_method pgd \
  --norm 2 \
  --epsilon 32768 \
  --alpha 5 \
  --steps 100 \
  --index 1 \
  --output_dir results/loss_compare/ns/loss2_stopgrad_pgd

python run_loss_comparison_attack.py \
  --case ns \
  --solver_backend jax \
  --model_backend torch \
  --loss_type loss2_stopgrad \
  --attack_method power \
  --norm 2 \
  --epsilon 32768 \
  --alpha 5 \
  --steps 100 \
  --index 1 \
  --output_dir results/loss_compare/ns/loss2_stopgrad_power

python run_loss_comparison_attack.py \
  --case ns \
  --solver_backend jax \
  --model_backend torch \
  --loss_type loss3 \
  --attack_method pgd \
  --norm 2 \
  --epsilon 32768 \
  --alpha 5 \
  --steps 100 \
  --index 1 \
  --output_dir results/loss_compare/ns/loss3_pgd

python run_loss_comparison_attack.py \
  --case ns \
  --solver_backend jax \
  --model_backend torch \
  --loss_type loss3 \
  --attack_method power \
  --norm 2 \
  --epsilon 32768 \
  --alpha 5 \
  --steps 100 \
  --index 1 \
  --output_dir results/loss_compare/ns/loss3_power
\end{verbatim}

\subsubsection{Experiment B: True Loss Comparison with PGD}

For Burgers:

\begin{verbatim}
python run_loss_comparison_attack.py \
  --case burgers \
  --solver_backend jax \
  --model_backend torch \
  --loss_type loss1 \
  --attack_method pgd \
  --norm 2 \
  --epsilon 20 \
  --alpha 0.06 \
  --steps 100 \
  --index 0 \
  --output_dir results/true_loss_compare/burgers/loss1_pgd

python run_loss_comparison_attack.py \
  --case burgers \
  --solver_backend jax \
  --model_backend torch \
  --loss_type loss2 \
  --attack_method pgd \
  --norm 2 \
  --epsilon 20 \
  --alpha 0.06 \
  --steps 100 \
  --index 0 \
  --output_dir results/true_loss_compare/burgers/loss2_pgd

python run_loss_comparison_attack.py \
  --case burgers \
  --solver_backend jax \
  --model_backend torch \
  --loss_type loss2_stopgrad \
  --attack_method pgd \
  --norm 2 \
  --epsilon 20 \
  --alpha 0.06 \
  --steps 100 \
  --index 0 \
  --output_dir results/true_loss_compare/burgers/loss2_stopgrad_pgd

python run_loss_comparison_attack.py \
  --case burgers \
  --solver_backend jax \
  --model_backend torch \
  --loss_type loss3 \
  --attack_method pgd \
  --norm 2 \
  --epsilon 20 \
  --alpha 0.06 \
  --steps 100 \
  --index 0 \
  --output_dir results/true_loss_compare/burgers/loss3_pgd
\end{verbatim}

For NS / Stokes:

\begin{verbatim}
python run_loss_comparison_attack.py \
  --case ns \
  --solver_backend jax \
  --model_backend torch \
  --loss_type loss1 \
  --attack_method pgd \
  --norm 2 \
  --epsilon 32768 \
  --alpha 5 \
  --steps 100 \
  --index 1 \
  --output_dir results/true_loss_compare/ns/loss1_pgd

python run_loss_comparison_attack.py \
  --case ns \
  --solver_backend jax \
  --model_backend torch \
  --loss_type loss2 \
  --attack_method pgd \
  --norm 2 \
  --epsilon 32768 \
  --alpha 5 \
  --steps 100 \
  --index 1 \
  --output_dir results/true_loss_compare/ns/loss2_pgd

python run_loss_comparison_attack.py \
  --case ns \
  --solver_backend jax \
  --model_backend torch \
  --loss_type loss2_stopgrad \
  --attack_method pgd \
  --norm 2 \
  --epsilon 32768 \
  --alpha 5 \
  --steps 100 \
  --index 1 \
  --output_dir results/true_loss_compare/ns/loss2_stopgrad_pgd

python run_loss_comparison_attack.py \
  --case ns \
  --solver_backend jax \
  --model_backend torch \
  --loss_type loss3 \
  --attack_method pgd \
  --norm 2 \
  --epsilon 32768 \
  --alpha 5 \
  --steps 100 \
  --index 1 \
  --output_dir results/true_loss_compare/ns/loss3_pgd
\end{verbatim}

\subsubsection{2D Frame-Mode Runs}

\begin{verbatim}
python run_loss_comparison_attack.py \
  --case ns \
  --solver_backend jax \
  --model_backend torch \
  --loss_type loss2_stopgrad \
  --attack_method pgd \
  --norm 2 \
  --epsilon 32768 \
  --alpha 5 \
  --steps 100 \
  --index 1 \
  --frame_mode aaaaaaaaaa \
  --output_dir results/frame_mode_compare/ns/all_a

python run_loss_comparison_attack.py \
  --case ns \
  --solver_backend jax \
  --model_backend torch \
  --loss_type loss2_stopgrad \
  --attack_method pgd \
  --norm 2 \
  --epsilon 32768 \
  --alpha 5 \
  --steps 100 \
  --index 1 \
  --frame_mode dddddddddd \
  --output_dir results/frame_mode_compare/ns/all_d

python run_loss_comparison_attack.py \
  --case ns \
  --solver_backend jax \
  --model_backend torch \
  --loss_type loss3 \
  --attack_method pgd \
  --norm 2 \
  --epsilon 32768 \
  --alpha 5 \
  --steps 100 \
  --index 1 \
  --frame_mode wwwwwwwwww \
  --output_dir results/frame_mode_compare/ns/all_w
\end{verbatim}

\subsection{Plotting Commands}

\begin{verbatim}
python plot_loss_method_comparison.py \
  --case burgers \
  --results_root results/loss_compare/burgers \
  --output_dir figures/loss_compare/burgers

python plot_loss_method_comparison.py \
  --case ns \
  --results_root results/loss_compare/ns \
  --output_dir figures/loss_compare/ns

python plot_true_loss_comparison.py \
  --case burgers \
  --results_root results/true_loss_compare/burgers \
  --output_dir figures/true_loss_compare/burgers

python plot_true_loss_comparison.py \
  --case ns \
  --results_root results/true_loss_compare/ns \
  --output_dir figures/true_loss_compare/ns

python plot_frame_mode_comparison.py \
  --case ns \
  --results_root results/frame_mode_compare/ns \
  --output_dir figures/frame_mode_compare/ns

python make_attack_summary_table.py \
  --results_roots \
    results/loss_compare/burgers \
    results/loss_compare/ns \
    results/true_loss_compare/burgers \
    results/true_loss_compare/ns \
    results/frame_mode_compare/ns \
  --output_csv results/attack_summary_table.csv \
  --output_tex results/attack_summary_table.tex
\end{verbatim}

\subsection{Final Claims to Test}

\subsubsection{Claim 1: Loss 1 is Closest to a Singular-Vector Problem}

Because

\[
L_1^2(\delta)
\approx
\delta^\top J_f(x)^\top J_f(x)\delta,
\]

Loss 1 should be closely related to the leading right singular vector of \(J_f(x)\).

Therefore, power iteration may be more suitable than PGD for Loss 1.

This should be tested using:

\[
\hat{\sigma}_k,
\qquad
R_k,
\qquad
|\cos(\delta_{\mathrm{PGD}},v_{\max})|,
\]

and the final comparison between:

\[
L_1(\delta_{\mathrm{PGD}})
\quad
\text{and}
\quad
L_1(\epsilon v_{\max}).
\]

\subsubsection{Claim 2: PGD May Be More Suitable for Loss 2, Stop-Gradient Loss 2, and Loss 3}

Loss 2 and Loss 3 contain linear terms in their local squared-loss expansions.

The stop-gradient loss has a moving solver target.

Therefore, these losses are not pure Rayleigh quotient problems.

For these losses, PGD may be more suitable than power iteration / power-like update.

This should be tested using:

\[
L_{\mathrm{opt}}(\delta_k),
\qquad
L_{\mathrm{true}}(\delta_k),
\qquad
\text{final true loss},
\qquad
\text{true loss growth ratio}.
\]

\subsubsection{Claim 3: The Stop-Gradient Loss Tests the Value of Solver Forward without Solver Backward}

The loss

\[
L_{2,\mathrm{sg}}(\delta)
=
\|f(x+\delta)-\operatorname{stopgrad}(g(x+\delta))\|_2
\]

tests whether using the perturbed solver output in the forward pass helps even if no gradient is propagated through the solver.

The expected ordering may be:

\[
L_2
\quad
<
\quad
L_{2,\mathrm{sg}}
\quad
<
\quad
L_3
\]

in true loss growth, but this must be tested empirically.

\subsubsection{Claim 4: Directly Optimizing Loss 3 Should Usually Produce the Largest True Loss Growth}

The true evaluation target is always

\[
L_{\mathrm{true}}(\delta)
=
L_3(\delta)
=
\|f(x+\delta)-g(x+\delta)\|_2.
\]

Therefore, directly optimizing \(L_3\) should usually produce the largest final true loss and the largest true loss growth ratio.

This is the main question of Experiment B.

\subsubsection{Claim 5: Frame-Level Solver Participation Matters in Recurrent 2D Models}

The all-\(\texttt{a}\), all-\(\texttt{d}\), and all-\(\texttt{w}\) experiments test how much attack strength comes from:

\begin{itemize}
    \item no solver frame usage;
    \item solver forward values without solver backward;
    \item full solver forward and backward.
\end{itemize}

The comparison

\[
\texttt{aaaaaaaaaa}
\quad
\text{vs.}
\quad
\texttt{dddddddddd}
\quad
\text{vs.}
\quad
\texttt{wwwwwwwwww}
\]

should be used as the first recurrent-frame experiment before trying any mixed frame-mode combinations.

\section{PGD, Gradient Updates, and Generalized Power Iteration}

This section clarifies the relationship between projected-gradient-type attack methods and generalized power iteration. The key point is that both methods can be interpreted as gradient-driven constrained update methods, but they differ in two essential places: first, how the driving direction is computed, and second, how the constraint is enforced after the driving direction is obtained.

\subsection{Generalized \(p\)-to-\(q\) Operator-Norm Objective}

Consider a local linear residual map

\[
r(v)=Jv,
\]

where \(J\) is the Jacobian of the chosen residual map with respect to the input. The generalized \(p\)-to-\(q\) operator-norm problem is

\[
\max_{\|v\|_p\le 1}\|Jv\|_q.
\]

Equivalently, one may use the \(q\)-power objective

\[
\Phi(v)
=
\frac{1}{q}\|Jv\|_q^q.
\]

Define the nonlinear output-space map

\[
\phi_q(z)
=
|z|^{q-2}\odot z
=
\operatorname{sign}(z)\odot |z|^{q-1}.
\]

Then, by the chain rule,

\[
\nabla_v \Phi(v)
=
J^\top \phi_q(Jv).
\]

Thus the core driving direction in generalized power iteration is

\[
w_k
=
J^\top \phi_q(Jv_k).
\]

This expression can be computed in two equivalent ways. One may directly differentiate

\[
\Phi(v)=\frac{1}{q}\|Jv\|_q^q
\]

with respect to \(v\), or one may explicitly compute

\[
y_k=Jv_k,
\]

\[
u_k=\phi_q(y_k),
\]

\[
w_k=J^\top u_k.
\]

The second form corresponds to a JVP--nonlinear-map--VJP computation:

\[
v_k
\longmapsto
Jv_k
\longmapsto
\phi_q(Jv_k)
\longmapsto
J^\top\phi_q(Jv_k).
\]

Therefore, for the pure local operator-norm objective, direct gradient computation and the explicit JVP/VJP formula produce the same driving direction.

\subsection{The Outer \(p\)-Constraint Update}

The vector

\[
w_k=J^\top \phi_q(Jv_k)
\]

is only the gradient-like driving direction. Generalized power iteration does not stop at \(w_k\). It then converts \(w_k\) into a feasible direction in the \(p\)-unit ball by solving

\[
v_{k+1}
=
\arg\max_{\|s\|_p\le 1}
\langle w_k,s\rangle.
\]

This step is not another gradient computation. It is a constrained linear maximization step. It selects the point in the \(p\)-ball that is most aligned with \(w_k\).

Let \(p^\ast\) be the dual exponent of \(p\):

\[
p^\ast=\frac{p}{p-1}.
\]

For \(1<p<\infty\), the solution is

\[
v_{k+1,i}
=
\frac{
\operatorname{sign}(w_{k,i})|w_{k,i}|^{p^\ast-1}
}{
\|w_k\|_{p^\ast}^{p^\ast-1}
}.
\]

Equivalently, if we define the \(p\)-constraint map

\[
\psi_p(w)
=
\arg\max_{\|s\|_p\le 1}\langle w,s\rangle,
\]

then generalized power iteration can be written compactly as

\[
v_{k+1}
=
\psi_p
\left(
J^\top\phi_q(Jv_k)
\right).
\]

For special values of \(p\),

\[
p=2:
\qquad
\psi_2(w)=\frac{w}{\|w\|_2},
\]

and

\[
p=\infty:
\qquad
\psi_\infty(w)=\operatorname{sign}(w).
\]

Therefore, when \(p=\infty\), the outer constrained direction selection in generalized power iteration has the same geometric form as the sign step in FGSM:

\[
\arg\max_{\|s\|_\infty\le 1}
\langle w,s\rangle
=
\operatorname{sign}(w).
\]

\subsection{Comparison with PGD-Type Updates}

A PGD-type update has the form

\[
v_{k+1}
=
\Pi_{\|v\|_p\le 1}
\left(
v_k+\alpha w_k
\right),
\]

where \(w_k\) is usually the gradient of the selected attack loss and \(\Pi\) is a projection or clipping operator.

Thus both PGD and generalized power iteration can be described by the same high-level template:

\[
\text{compute a driving direction } w_k,
\qquad
\text{then enforce a } p\text{-norm constraint}.
\]

However, they differ in two essential ways.

First, PGD keeps the current iterate and takes a finite step:

\[
v_k+\alpha w_k.
\]

Generalized power iteration does not add the gradient to the old iterate. Instead, it directly replaces the current direction by the most aligned feasible direction:

\[
v_{k+1}
=
\psi_p(w_k).
\]

Second, the constraint operation is different. PGD uses projection:

\[
\Pi_{\mathcal{B}_p}(z)
=
\arg\min_{s\in\mathcal{B}_p}
\|s-z\|_2^2,
\]

where

\[
\mathcal{B}_p=\{s:\|s\|_p\le 1\}.
\]

This projection finds the feasible point closest to \(z\). In contrast, generalized power iteration uses

\[
\arg\max_{\|s\|_p\le1}
\langle w_k,s\rangle,
\]

which finds the feasible point most aligned with \(w_k\).

Therefore, PGD projection and generalized-power constrained maximization solve different optimization problems:

\[
\text{PGD projection:}
\qquad
\arg\min_{\|s\|_p\le1}\|s-z\|_2^2,
\]

whereas

\[
\text{generalized power outer step:}
\qquad
\arg\max_{\|s\|_p\le1}\langle w,s\rangle.
\]

For \(p=2\), if \(w\) lies outside the unit ball, both operations may produce the same boundary direction:

\[
\Pi_{\|s\|_2\le1}(w)
=
\frac{w}{\|w\|_2},
\qquad
\arg\max_{\|s\|_2\le1}\langle w,s\rangle
=
\frac{w}{\|w\|_2}.
\]

If \(w\) is inside the unit ball, projection returns \(w\), while the constrained maximization returns \(w/\|w\|_2\). Thus even for \(p=2\), projection and constrained maximization are not identical operations.

For \(p=\infty\), the distinction is even clearer:

\[
\Pi_{\|s\|_\infty\le1}(w)
=
\operatorname{clip}(w,-1,1),
\]

whereas

\[
\arg\max_{\|s\|_\infty\le1}\langle w,s\rangle
=
\operatorname{sign}(w).
\]

They coincide only when every coordinate is clipped to the boundary, for example when \(|w_i|\ge1\) for all \(i\). Otherwise, they may even point in different directions.

\subsection{Three Losses and Their Local \(q\)-Power Gradients}

Now consider three attack losses. We use the \(q\)-power objective

\[
\Phi(\delta)=\frac{1}{q}L(\delta)^q,
\]

because its gradient has a clean analytic form.

\subsubsection{Loss 1: Surrogate Self-Change}

The first loss is

\[
L_1(\delta)
=
\|f(x+\delta)-f(x)\|_q.
\]

The corresponding \(q\)-power objective is

\[
\Phi_1(\delta)
=
\frac{1}{q}
\|f(x+\delta)-f(x)\|_q^q.
\]

The exact nonlinear gradient is

\[
\nabla_\delta \Phi_1(\delta)
=
J_f(x+\delta)^\top
\phi_q
\left(
f(x+\delta)-f(x)
\right).
\]

Near \(\delta=0\), linearize

\[
f(x+\delta)
\approx
f(x)+J_f(x)\delta.
\]

Then

\[
f(x+\delta)-f(x)
\approx
J_f\delta.
\]

Thus the local \(q\)-power objective is

\[
\Phi_1(\delta)
\approx
\frac{1}{q}\|J_f\delta\|_q^q,
\]

and the local gradient is

\[
\nabla_\delta \Phi_1(\delta)
\approx
J_f^\top\phi_q(J_f\delta).
\]

This is exactly the generalized power iteration driving direction. Therefore, for Loss 1, the direct gradient of the local \(q\)-power objective agrees with the explicit JVP--\(\phi_q\)--VJP computation.

\subsubsection{Loss 2: Fixed Target Loss}

The second loss is

\[
L_2(\delta)
=
\|f(x+\delta)-y\|_q,
\]

where \(y\) is a fixed target. Its \(q\)-power objective is

\[
\Phi_2(\delta)
=
\frac{1}{q}
\|f(x+\delta)-y\|_q^q.
\]

The exact nonlinear gradient is

\[
\nabla_\delta \Phi_2(\delta)
=
J_f(x+\delta)^\top
\phi_q
\left(
f(x+\delta)-y
\right).
\]

Linearizing \(f(x+\delta)\) near \(\delta=0\),

\[
f(x+\delta)
\approx
f(x)+J_f\delta.
\]

Hence

\[
f(x+\delta)-y
\approx
f(x)-y+J_f\delta.
\]

Define

\[
b_2=f(x)-y.
\]

Then the local objective is

\[
\Phi_2(\delta)
\approx
\frac{1}{q}
\|b_2+J_f\delta\|_q^q,
\]

and its local gradient is

\[
\nabla_\delta \Phi_2(\delta)
\approx
J_f^\top
\phi_q
\left(
b_2+J_f\delta
\right).
\]

Equivalently,

\[
\nabla_\delta \Phi_2(\delta)
\approx
J_f^\top
\left(
|b_2+J_f\delta|^{q-2}
\odot
(b_2+J_f\delta)
\right).
\]

This differs from the pure generalized-power direction

\[
J_f^\top\phi_q(J_f\delta)
\]

because the true local residual contains the offset \(b_2=f(x)-y\).

\subsubsection{Loss 3: Solver--Model Residual Loss}

The third loss is

\[
L_3(\delta)
=
\|f(x+\delta)-g(x+\delta)\|_q.
\]

Define the residual map

\[
h(x)=f(x)-g(x).
\]

Then

\[
L_3(\delta)=\|h(x+\delta)\|_q.
\]

The \(q\)-power objective is

\[
\Phi_3(\delta)
=
\frac{1}{q}
\|f(x+\delta)-g(x+\delta)\|_q^q.
\]

The exact nonlinear gradient is

\[
\nabla_\delta \Phi_3(\delta)
=
\left[
J_f(x+\delta)-J_g(x+\delta)
\right]^\top
\phi_q
\left(
f(x+\delta)-g(x+\delta)
\right).
\]

Linearizing both \(f\) and \(g\) near \(\delta=0\),

\[
f(x+\delta)
\approx
f(x)+J_f\delta,
\]

\[
g(x+\delta)
\approx
g(x)+J_g\delta.
\]

Thus

\[
f(x+\delta)-g(x+\delta)
\approx
f(x)-g(x)+(J_f-J_g)\delta.
\]

Define

\[
b_3=f(x)-g(x),
\]

and

\[
J_h=J_f-J_g.
\]

Then the local objective is

\[
\Phi_3(\delta)
\approx
\frac{1}{q}
\|b_3+J_h\delta\|_q^q,
\]

and its local gradient is

\[
\nabla_\delta \Phi_3(\delta)
\approx
J_h^\top
\phi_q
\left(
b_3+J_h\delta
\right).
\]

Equivalently,

\[
\nabla_\delta \Phi_3(\delta)
\approx
(J_f-J_g)^\top
\phi_q
\left(
f(x)-g(x)+(J_f-J_g)\delta
\right).
\]

This also differs from the pure generalized-power direction

\[
J_h^\top\phi_q(J_h\delta)
\]

because the true local residual contains the offset

\[
b_3=f(x)-g(x).
\]

\subsection{Unified Local Gradient Template}

All three local \(q\)-power gradients can be written in the unified affine form

\[
\Phi(\delta)
=
\frac{1}{q}\|b+J\delta\|_q^q,
\]

with gradient

\[
\nabla_\delta \Phi(\delta)
=
J^\top \phi_q(b+J\delta).
\]

The only differences among the three losses are the offset \(b\) and the Jacobian \(J\):

\[
\begin{array}{c|c|c|c}
\text{Loss} & b & J & \nabla_\delta \Phi(\delta)\\
\hline
L_1 & 0 & J_f
&
J_f^\top\phi_q(J_f\delta)
\\
L_2 & f(x)-y & J_f
&
J_f^\top\phi_q(f(x)-y+J_f\delta)
\\
L_3 & f(x)-g(x) & J_f-J_g
&
(J_f-J_g)^\top
\phi_q
\left(
f(x)-g(x)+(J_f-J_g)\delta
\right)
\end{array}
\]

This table shows the fundamental distinction. For Loss 1, \(b=0\), so the true local gradient is exactly the same as the pure generalized-power JVP/VJP direction. For Loss 2 and Loss 3, \(b\neq0\), so the true local gradient becomes

\[
J^\top\phi_q(b+J\delta),
\]

whereas the pure generalized-power direction is

\[
J^\top\phi_q(J\delta).
\]

These are generally different.

\subsection{The Special Case \(p=q=2\)}

When \(q=2\),

\[
\phi_2(z)=z.
\]

Therefore the three local gradients simplify substantially.

For Loss 1,

\[
\nabla_\delta \Phi_1(\delta)
\approx
J_f^\top J_f\delta.
\]

This is exactly the ordinary power iteration direction for

\[
A_1=J_f^\top J_f.
\]

Thus ordinary power iteration is

\[
v_{k+1}
=
\frac{J_f^\top J_fv_k}
{\|J_f^\top J_fv_k\|_2}.
\]

For Loss 2,

\[
\nabla_\delta \Phi_2(\delta)
\approx
J_f^\top
\left(
f(x)-y+J_f\delta
\right),
\]

so

\[
\nabla_\delta \Phi_2(\delta)
\approx
J_f^\top(f(x)-y)
+
J_f^\top J_f\delta.
\]

For Loss 3,

\[
\nabla_\delta \Phi_3(\delta)
\approx
(J_f-J_g)^\top
\left(
f(x)-g(x)+(J_f-J_g)\delta
\right),
\]

so

\[
\nabla_\delta \Phi_3(\delta)
\approx
(J_f-J_g)^\top(f(x)-g(x))
+
(J_f-J_g)^\top(J_f-J_g)\delta.
\]

Thus, when \(p=q=2\), Loss 1 gives a pure quadratic form

\[
\Phi_1(\delta)
\approx
\frac{1}{2}
\delta^\top J_f^\top J_f\delta,
\]

whose gradient is

\[
J_f^\top J_f\delta.
\]

This is exactly the ordinary power iteration structure. In contrast, Loss 2 and Loss 3 have affine local residuals,

\[
b+J\delta,
\]

so their gradients contain an additional term

\[
J^\top b.
\]

This additional term is the reason why Loss 2 and Loss 3 are not strict power-iteration objectives, even in the Euclidean case.

\subsection{Two Independent Choices for Loss 2 and Loss 3}

For Loss 2 and Loss 3, there are two independent modeling choices.

First, one must choose the driving direction \(w_k\). One option is the pure Jacobian direction

\[
w_k^{\mathrm{jac}}
=
J^\top\phi_q(Jv_k),
\]

which corresponds to strict generalized power iteration and local Jacobian sensitivity. This ignores the offset \(b\). Another option is the true affine local gradient

\[
w_k^{\mathrm{affine}}
=
J^\top\phi_q(b+Jv_k),
\]

which is the actual gradient of the local affine residual objective

\[
\frac{1}{q}\|b+Jv\|_q^q.
\]

Second, once \(w_k\) is chosen, one must choose the update mechanism. A generalized-power-style update uses

\[
v_{k+1}
=
\psi_p(w_k)
=
\arg\max_{\|s\|_p\le1}
\langle w_k,s\rangle.
\]

A PGD-style update uses

\[
v_{k+1}
=
\Pi_{\|v\|_p\le1}
\left(
v_k+\alpha w_k
\right).
\]

Therefore, for Loss 2 and Loss 3, the ambiguity comes from two independent differences:

\[
\boxed{
\text{pure Jacobian direction}
\quad \text{versus} \quad
\text{true affine residual gradient}
}
\]

and

\[
\boxed{
\text{direct dual-map update}
\quad \text{versus} \quad
\text{additive PGD projection update}.
}
\]

These two distinctions explain why generalized power iteration and PGD agree cleanly for Loss 1 but split into several reasonable variants for Loss 2 and Loss 3.

\subsection{Practical Interpretation}

For Loss 1, the local residual is

\[
J_f\delta,
\]

so the direct gradient and the JVP/VJP generalized-power direction are the same. Therefore generalized power iteration has a strict operator-norm interpretation.

For Loss 2, the local residual is

\[
f(x)-y+J_f\delta.
\]

For Loss 3, the local residual is

\[
f(x)-g(x)+(J_f-J_g)\delta.
\]

Because of these offsets, the true local gradients are not the same as the pure Jacobian sensitivity directions. One may still use pure generalized power iteration as a local sensitivity baseline, but it should not be interpreted as the true gradient optimizer for the finite Loss 2 or Loss 3 objective. Conversely, one may use the true affine residual gradient, but then the method becomes an affine gradient-power or PGD-style method rather than a strict generalized power iteration.

\section{Initialization Effects, Stationary Points, and Spectral Components}

This section summarizes the role of initialization in the three attack losses. The main issue is that Loss 1 has zero residual at \(\delta=0\), and therefore its first-order gradient also vanishes at \(\delta=0\). In contrast, Loss 2 and Loss 3 usually contain a nonzero residual offset, so their gradients at \(\delta=0\) are generally nonzero. This difference strongly affects how PGD-type methods, gradient-power methods, and generalized power iteration behave at initialization.

\subsection{Unified Local Form}

As discussed before, after local linearization the three losses can be written in the unified affine residual form

\[
\Phi(\delta)
=
\frac{1}{q}\|b+J\delta\|_q^q,
\]

where

\[
\phi_q(z)
=
|z|^{q-2}\odot z
=
\operatorname{sign}(z)\odot |z|^{q-1}.
\]

The first derivative with respect to \(\delta\) is

\[
\nabla_\delta \Phi(\delta)
=
J^\top\phi_q(b+J\delta).
\]

The second derivative, for the affine local residual model and for points where the relevant derivatives exist, is

\[
\nabla_\delta^2 \Phi(\delta)
=
J^\top
\left[
(q-1)\operatorname{diag}\left(|b+J\delta|^{q-2}\right)
\right]
J.
\]

For the special Euclidean case \(q=2\), this simplifies to

\[
\nabla_\delta \Phi(\delta)
=
J^\top(b+J\delta)
=
J^\top b+J^\top J\delta,
\]

and

\[
\nabla_\delta^2\Phi(\delta)
=
J^\top J.
\]

Therefore, in the \(q=2\) case, the first-order term is determined by \(J^\top b\), while the local curvature is determined by \(J^\top J\).

\subsection{Loss 1: Zero Residual and Zero First-Order Gradient}

Loss 1 is

\[
L_1(\delta)
=
\|f(x+\delta)-f(x)\|_q.
\]

The corresponding \(q\)-power objective is

\[
\Phi_1(\delta)
=
\frac{1}{q}
\|f(x+\delta)-f(x)\|_q^q.
\]

Near \(\delta=0\),

\[
f(x+\delta)
\approx
f(x)+J_f\delta,
\]

so

\[
f(x+\delta)-f(x)
\approx
J_f\delta.
\]

Thus the local objective is

\[
\Phi_1(\delta)
\approx
\frac{1}{q}\|J_f\delta\|_q^q.
\]

In the unified notation,

\[
b_1=0,
\qquad
J=J_f.
\]

Therefore,

\[
\nabla_\delta \Phi_1(\delta)
=
J_f^\top\phi_q(J_f\delta).
\]

At \(\delta=0\),

\[
\nabla_\delta \Phi_1(0)
=
J_f^\top\phi_q(0)
=
0.
\]

Thus Loss 1 has a zero first-order gradient at the origin. If PGD or any purely first-order gradient method starts from

\[
\delta_0=0,
\]

then

\[
\nabla_\delta \Phi_1(\delta_0)=0,
\]

and the method cannot move away from the origin.

In the special case \(q=2\),

\[
\Phi_1(\delta)
=
\frac{1}{2}\|J_f\delta\|_2^2
=
\frac{1}{2}\delta^\top J_f^\top J_f\delta.
\]

Then

\[
\nabla_\delta \Phi_1(\delta)
=
J_f^\top J_f\delta,
\]

and

\[
\nabla_\delta^2\Phi_1(0)
=
J_f^\top J_f.
\]

Therefore, \(\delta=0\) is a stationary point with zero first-order information, but the second-order matrix

\[
A_1=J_f^\top J_f
\]

contains the local growth directions. From the perspective of attack maximization, \(\delta=0\) is a bad starting point because first-order methods do not see any ascent direction there. A nonzero initialization is needed to provide a seed direction.

\subsection{Spectral Interpretation of Loss 1}

When \(p=q=2\), the local Loss 1 objective is governed by the symmetric positive semidefinite matrix

\[
A_1=J_f^\top J_f.
\]

Let

\[
A_1u_i=\lambda_i u_i,
\]

where

\[
\lambda_1\ge \lambda_2\ge \cdots \ge 0.
\]

If the initial perturbation is decomposed as

\[
\delta_0=\sum_i c_i u_i,
\]

then repeated multiplication by \(A_1\) gives

\[
A_1^k\delta_0
=
\sum_i c_i\lambda_i^k u_i.
\]

If

\[
c_1\neq0,
\]

then the dominant direction \(u_1\) eventually dominates. If

\[
c_1=0,
\]

then the dominant eigendirection \(u_1\) never appears in the iteration. The method then converges, if it converges, to the largest eigendirection present in the initial subspace.

In particular, if the initial direction is exactly the eigenvector corresponding to the smallest positive eigenvalue,

\[
\delta_0=u_{\min},
\]

then

\[
A_1\delta_0
=
\lambda_{\min}\delta_0.
\]

After normalization, the direction remains \(u_{\min}\). Therefore, the method can remain trapped in a poor spectral direction. If \(\lambda_{\min}=0\), then

\[
A_1\delta_0=0,
\]

and the update direction disappears entirely.

Thus Loss 1 does not create missing spectral components by itself. It only amplifies components that are already present in the initial perturbation. This is why random nonzero initialization is important.

For continuous random initialization, the probability that the initial direction is exactly orthogonal to a fixed dominant eigendirection is mathematically zero. However, in finite-precision computation, special initialization schemes, symmetry, or nearly orthogonal random draws can still create poor starts or slow convergence. Hence multiple random restarts are practically useful.

\subsection{Small Nonzero Initialization versus Zero Initialization}

For Loss 1, zero initialization is problematic:

\[
\delta_0=0
\quad\Rightarrow\quad
\nabla_\delta \Phi_1(\delta_0)=0.
\]

However, a very small but nonzero initialization can move, provided that

\[
J_f\delta_0\neq0.
\]

For the \(q\)-power objective,

\[
\nabla_\delta \Phi_1(\delta)
=
J_f^\top\phi_q(J_f\delta).
\]

Because

\[
\phi_q(cz)=c^{q-1}\phi_q(z),
\qquad c>0,
\]

we have the scaling behavior

\[
\nabla_\delta\Phi_1(c\delta)
=
c^{q-1}\nabla_\delta\Phi_1(\delta)
\]

for the local linear model.

Therefore, for unnormalized gradient ascent,

\[
\delta_{k+1}
=
\delta_k+\alpha\nabla_\delta\Phi_1(\delta_k),
\]

a very small initialization produces a very small gradient and may lead to slow movement. Conversely, a larger initialization can produce a larger gradient, which may create a rolling or snowball-like growth effect unless projection or normalization is applied.

In contrast, power iteration and generalized power iteration remove much of this scale dependence. For ordinary power iteration,

\[
v_{k+1}
=
\frac{A_1v_k}{\|A_1v_k\|_2},
\]

if

\[
v_0=\epsilon u_0,
\qquad
\epsilon>0,
\]

then

\[
\frac{A_1v_0}{\|A_1v_0\|_2}
=
\frac{\epsilon A_1u_0}{\|\epsilon A_1u_0\|_2}
=
\frac{A_1u_0}{\|A_1u_0\|_2}.
\]

Thus the scale \(\epsilon\) cancels. The direction matters; the magnitude does not.

Similarly, generalized power iteration uses

\[
v_{k+1}
=
\psi_p
\left(
J^\top\phi_q(Jv_k)
\right),
\]

where

\[
\psi_p(w)
=
\arg\max_{\|s\|_p\le1}\langle w,s\rangle.
\]

For \(c>0\),

\[
\psi_p(cw)=\psi_p(w),
\]

so positive scaling of the driving direction does not change the resulting constrained direction. Hence generalized power iteration also mainly depends on the initial direction rather than the initial magnitude, as long as the initialization is nonzero.

\subsection{Loss 2: Nonzero Offset and Nonzero Initial Gradient}

Loss 2 is

\[
L_2(\delta)
=
\|f(x+\delta)-y\|_q,
\]

where \(y\) is fixed. Its \(q\)-power objective is

\[
\Phi_2(\delta)
=
\frac{1}{q}
\|f(x+\delta)-y\|_q^q.
\]

Near \(\delta=0\),

\[
f(x+\delta)
\approx
f(x)+J_f\delta.
\]

Therefore,

\[
f(x+\delta)-y
\approx
f(x)-y+J_f\delta.
\]

Define

\[
b_2=f(x)-y.
\]

Then

\[
\Phi_2(\delta)
\approx
\frac{1}{q}
\|b_2+J_f\delta\|_q^q.
\]

The local gradient is

\[
\nabla_\delta \Phi_2(\delta)
=
J_f^\top\phi_q(b_2+J_f\delta).
\]

At \(\delta=0\),

\[
\nabla_\delta \Phi_2(0)
=
J_f^\top\phi_q(b_2).
\]

Thus Loss 2 generally has nonzero first-order information at the origin, provided that

\[
b_2\neq0
\]

and

\[
J_f^\top\phi_q(b_2)\neq0.
\]

For \(q=2\),

\[
\nabla_\delta \Phi_2(0)
=
J_f^\top b_2,
\]

and

\[
\nabla_\delta^2\Phi_2(0)
=
J_f^\top J_f.
\]

Therefore, Loss 2 can often be optimized from \(\delta_0=0\) because the offset \(b_2=f(x)-y\) already produces a gradient direction. If \(b_2=0\), then Loss 2 locally reduces to Loss 1, and the zero-initialization problem returns.

\subsection{Loss 3: Solver--Model Residual Offset}

Loss 3 is

\[
L_3(\delta)
=
\|f(x+\delta)-g(x+\delta)\|_q.
\]

Define

\[
h(x)=f(x)-g(x),
\]

\[
J_h=J_f-J_g,
\]

and

\[
b_3=f(x)-g(x).
\]

The \(q\)-power objective is

\[
\Phi_3(\delta)
=
\frac{1}{q}
\|f(x+\delta)-g(x+\delta)\|_q^q.
\]

After local linearization,

\[
f(x+\delta)-g(x+\delta)
\approx
f(x)-g(x)+(J_f-J_g)\delta.
\]

Thus

\[
\Phi_3(\delta)
\approx
\frac{1}{q}
\|b_3+J_h\delta\|_q^q.
\]

Its local gradient is

\[
\nabla_\delta \Phi_3(\delta)
=
J_h^\top\phi_q(b_3+J_h\delta).
\]

At \(\delta=0\),

\[
\nabla_\delta \Phi_3(0)
=
J_h^\top\phi_q(b_3).
\]

For \(q=2\),

\[
\nabla_\delta \Phi_3(0)
=
J_h^\top b_3
=
(J_f-J_g)^\top(f(x)-g(x)),
\]

and

\[
\nabla_\delta^2\Phi_3(0)
=
J_h^\top J_h
=
(J_f-J_g)^\top(J_f-J_g).
\]

Therefore, Loss 3 behaves similarly to Loss 2. If the model and solver already disagree at the clean input,

\[
b_3=f(x)-g(x)\neq0,
\]

then the origin may already have a nonzero gradient. If

\[
b_3=0,
\]

then the local residual becomes

\[
J_h\delta,
\]

and Loss 3 locally behaves like a Loss 1 problem with Jacobian \(J_h\). In that case, zero initialization again gives zero gradient.

\subsection{First- and Second-Order Structure at the Origin}

The three losses can be summarized by their first- and second-order information at \(\delta=0\).

For the local affine model

\[
\Phi(\delta)
=
\frac{1}{q}\|b+J\delta\|_q^q,
\]

we have

\[
\nabla_\delta\Phi(0)
=
J^\top\phi_q(b),
\]

and

\[
\nabla_\delta^2\Phi(0)
=
J^\top
\left[
(q-1)\operatorname{diag}(|b|^{q-2})
\right]
J.
\]

Therefore,

\[
\begin{array}{c|c|c|c}
\text{Loss} & b & J & \nabla_\delta\Phi(0)\\
\hline
L_1 & 0 & J_f & 0\\
L_2 & f(x)-y & J_f & J_f^\top\phi_q(f(x)-y)\\
L_3 & f(x)-g(x) & J_f-J_g
&
(J_f-J_g)^\top\phi_q(f(x)-g(x))
\end{array}
\]

For \(q=2\), the second-order matrices are

\[
\begin{array}{c|c}
\text{Loss} & \nabla_\delta^2\Phi(0)\\
\hline
L_1 & J_f^\top J_f\\
L_2 & J_f^\top J_f\\
L_3 & (J_f-J_g)^\top(J_f-J_g)
\end{array}
\]

This explains the initialization behavior:

\[
\boxed{
\text{Loss 1 has zero first-order information at } \delta=0.
}
\]

\[
\boxed{
\text{Loss 2 and Loss 3 usually have nonzero first-order information because their offsets } b_2,b_3 \text{ may be nonzero.}
}
\]

\[
\boxed{
\text{When the offset } b \text{ is zero, the problem becomes a pure Jacobian-amplification problem, and initialization becomes critical.}
}
\]

\subsection{Geometric Interpretation}

A useful geometric analogy is to imagine optimizing a surface in the perturbation space. For Loss 1, the point \(\delta=0\) is exactly the zero-residual point. In the local \(q=2\) model,

\[
\Phi_1(\delta)
=
\frac{1}{2}\delta^\top J_f^\top J_f\delta.
\]

At \(\delta=0\), the gradient vanishes:

\[
\nabla\Phi_1(0)=0.
\]

However, the curvature matrix

\[
J_f^\top J_f
\]

contains the directions in which the function increases fastest. A first-order method initialized exactly at \(\delta=0\) cannot see these directions. A second-order method, or a power-iteration-type method initialized with a nonzero seed, can use the curvature information.

Thus Loss 1 is not best understood as having no direction at all. Rather, its direction is not visible to first-order information at the origin. The direction is encoded in the second-order or Jacobian structure.

Loss 2 and Loss 3 usually start away from the zero-residual point because their residual offsets are

\[
b_2=f(x)-y
\]

and

\[
b_3=f(x)-g(x).
\]

Therefore, they may already have a nonzero gradient at \(\delta=0\). In the geometric analogy, they are not exactly at the flat stationary point. They are already located on a part of the surface with a first-order slope.

\subsection{Implications for PGD and Generalized Power Iteration}

For PGD-style methods,

\[
\delta_{k+1}
=
\Pi_{\|\delta\|_p\le\varepsilon}
\left(
\delta_k+\alpha\nabla_\delta\Phi(\delta_k)
\right),
\]

zero initialization fails whenever

\[
\nabla_\delta\Phi(0)=0.
\]

Therefore, Loss 1 requires nonzero initialization. Loss 2 and Loss 3 may not require it if their offsets are nonzero, but random initialization is still useful because the objectives are nonlinear and nonconvex.

For generalized power iteration,

\[
v_{k+1}
=
\psi_p
\left(
J^\top\phi_q(Jv_k)
\right),
\]

the initialization must also be nonzero. If

\[
v_0=0,
\]

then

\[
Jv_0=0,
\]

so

\[
J^\top\phi_q(Jv_0)=0.
\]

Then

\[
\psi_p(0)
\]

is not uniquely defined, because every feasible direction has the same zero inner product with \(0\).

Thus generalized power iteration should be initialized with a random nonzero direction, typically normalized to the \(p\)-unit boundary:

\[
\|v_0\|_p=1.
\]

For finite-radius attack perturbations, one may use

\[
\delta_0=\varepsilon v_0,
\qquad
\|v_0\|_p=1.
\]

\subsection{Random Initialization and Multiple Restarts}

Random initialization is used to avoid zero-gradient starts and poor spectral subspaces. In the ideal continuous setting, a random direction is almost surely not exactly orthogonal to the dominant eigendirection. However, the coefficient along the dominant direction may still be small, and the spectral gap may also be small. In that case convergence can be slow.

Moreover, in nonlinear attack objectives, different random starts can lead to different local optima or different high-loss directions. Therefore, one should use multiple random restarts:

\[
v_0^{(r)}\sim \text{random},
\qquad
r=1,\dots,R,
\]

run the attack or power-style update from each start, and keep the perturbation that gives the largest target evaluation.

For proxy-loss experiments, it is important to record both the optimized loss and the true evaluation loss. In particular, when comparing Loss 1, Loss 2, and Loss 3 as proxy objectives, the true solver-model loss

\[
L_3(\delta)
=
\|f(x+\delta)-g(x+\delta)\|_q
\]

should always be recorded as the final evaluation metric.

\subsection{Final Summary}

The initialization issue can be summarized as follows.

First, Loss 1 has no residual offset:

\[
b_1=0.
\]

Therefore,

\[
\nabla_\delta\Phi_1(0)=0.
\]

Its ascent directions are encoded in the Jacobian or Hessian structure, not in the first-order gradient at the origin. A zero initialization causes first-order methods to stop immediately. A nonzero initialization provides spectral components that can be amplified.

Second, Loss 2 and Loss 3 contain residual offsets:

\[
b_2=f(x)-y,
\]

and

\[
b_3=f(x)-g(x).
\]

If these offsets are nonzero, then

\[
\nabla_\delta\Phi_2(0)=J_f^\top\phi_q(b_2),
\]

and

\[
\nabla_\delta\Phi_3(0)=(J_f-J_g)^\top\phi_q(b_3),
\]

which are generally nonzero. Therefore, these losses can often move from zero initialization.

Third, for the \(p=q=2\) Loss 1 case, power iteration amplifies only the spectral components already present in the initial direction. If the initial direction is exactly orthogonal to the dominant eigendirection, the dominant eigendirection will never appear. If the initialization is exactly a smaller-eigenvalue eigenvector, the method may remain in that poorer direction. Random initialization avoids exact orthogonality almost surely in theory, but multiple restarts remain important in practice.

Finally, the key distinction is:

\[
\boxed{
\text{Loss 1 needs a nonzero seed because its first-order gradient at the origin is zero.}
}
\]

\[
\boxed{
\text{Loss 2 and Loss 3 can have nonzero first-order gradients at the origin because they contain residual offsets.}
}
\]

\[
\boxed{
\text{For power-style methods, the magnitude of a nonzero seed is less important than its direction; for unnormalized gradient methods, the magnitude can affect the update size.}
}
\]

\section{Loss Definitions, Ratio Objectives, and Optimization Methods}

This section clarifies the distinction between the attack loss, the normalized gain objective, and the optimization method. The main point is that the choice of loss and the choice of optimizer are two separate issues. Some losses naturally lead to an operator-norm or generalized singular-vector problem, while others lead to an affine residual maximization problem.

\subsection{Basic Objects}

Let \(x\) denote the clean input and let \(\delta\) denote the input perturbation. The attacked input is

\[
x_\delta = x+\delta.
\]

Let \(f\) be the surrogate model, \(g\) be the solver or ground-truth operator, and \(y\) be a fixed target or ground-truth output. We consider the following quantities:

\[
f(x), \qquad f(x+\delta), \qquad g(x), \qquad g(x+\delta), \qquad y.
\]

The input perturbation is constrained by a \(p\)-norm budget,

\[
\|\delta\|_p \leq \epsilon,
\]

while the output residual is measured by a \(q\)-norm.

\subsection{Three Attack Losses}

We consider three losses:

\[
L_1(\delta)
=
\|f(x+\delta)-f(x)\|_q,
\]

\[
L_2(\delta)
=
\|f(x+\delta)-y\|_q,
\]

and

\[
L_3(\delta)
=
\|f(x+\delta)-g(x+\delta)\|_q.
\]

Their zero-perturbation values are different:

\[
L_1(0)=0,
\]

\[
L_2(0)=\|f(x)-y\|_q,
\]

and

\[
L_3(0)=\|f(x)-g(x)\|_q.
\]

Thus, \(L_1\) directly measures output change caused by the perturbation, while \(L_2\) and \(L_3\) already contain a nonzero baseline residual at \(\delta=0\).

\subsection{Local Linearized Forms}

Locally, the three losses have different structures.

For \(L_1\), linearizing \(f(x+\delta)\) near \(x\) gives

\[
f(x+\delta)-f(x)
\approx
J_f(x)\delta.
\]

Therefore,

\[
L_1(\delta)
\approx
\|J_f\delta\|_q.
\]

This is a homogeneous linear residual.

For \(L_2\),

\[
f(x+\delta)-y
\approx
f(x)-y+J_f\delta.
\]

Define

\[
b_2=f(x)-y.
\]

Then

\[
L_2(\delta)
\approx
\|b_2+J_f\delta\|_q.
\]

For \(L_3\), define

\[
b_3=f(x)-g(x),
\]

and

\[
J_3=J_f-J_g.
\]

Then

\[
f(x+\delta)-g(x+\delta)
\approx
b_3+J_3\delta,
\]

so

\[
L_3(\delta)
\approx
\|b_3+J_3\delta\|_q.
\]

Hence \(L_2\) and \(L_3\) have affine residual forms, not purely linear residual forms.

\subsection{Different Objective Choices}

For each loss \(L_i\), there are several possible objective definitions. These objectives are not equivalent in general.

\[
\begin{array}{c|c|c}
\text{Objective Type} & \text{Formula} & \text{Meaning} \\
\hline
\text{Absolute loss} 
&
L_i(\delta)
&
\text{final perturbed loss}
\\
\text{Constrained loss maximization}
&
\max_{\|\delta\|_p\le \epsilon} L_i(\delta)
&
\text{standard PGD-style attack objective}
\\
\text{Loss increment}
&
L_i(\delta)-L_i(0)
&
\text{loss increase caused by perturbation}
\\
\text{Constrained increment maximization}
&
\max_{\|\delta\|_p\le \epsilon}
\left[L_i(\delta)-L_i(0)\right]
&
\text{same argmax as absolute loss under fixed budget}
\\
\text{Total ratio}
&
\frac{L_i(\delta)}{\|\delta\|_p}
&
\text{total loss per unit perturbation}
\\
\text{Incremental ratio}
&
\frac{L_i(\delta)-L_i(0)}{\|\delta\|_p}
&
\text{loss increase per unit perturbation}
\end{array}
\]

For fixed-budget constrained maximization,

\[
\arg\max_{\|\delta\|_p\le \epsilon} L_i(\delta)
=
\arg\max_{\|\delta\|_p\le \epsilon}
\left[L_i(\delta)-L_i(0)\right],
\]

because \(L_i(0)\) is constant with respect to \(\delta\).

However,

\[
\max_{\|\delta\|_p\le \epsilon} L_i(\delta)
\]

and

\[
\max_{0<\|\delta\|_p\le \epsilon}
\frac{L_i(\delta)-L_i(0)}{\|\delta\|_p}
\]

are generally different optimization problems. The first maximizes the final loss under a perturbation budget, while the second maximizes the loss increase per unit perturbation.

\subsection{Operator Norm Interpretation}

The induced \(p\)-to-\(q\) operator norm of a linear map \(J\) is

\[
\|J\|_{p\to q}
=
\max_{\delta\neq 0}
\frac{\|J\delta\|_q}{\|\delta\|_p}
=
\max_{\|\delta\|_p=1}
\|J\delta\|_q.
\]

This definition applies to a linear map. Therefore, \(L_1\) has a direct local operator-norm interpretation:

\[
\frac{L_1(\delta)}{\|\delta\|_p}
\approx
\frac{\|J_f\delta\|_q}{\|\delta\|_p}.
\]

Thus,

\[
\max_{\delta\neq 0}
\frac{L_1(\delta)}{\|\delta\|_p}
\approx
\|J_f\|_{p\to q}.
\]

Equivalently,

\[
\max_{\|\delta\|_p\le \epsilon}
L_1(\delta)
\approx
\epsilon \|J_f\|_{p\to q}.
\]

Therefore, for \(L_1\), maximizing the loss under a fixed perturbation budget and maximizing the normalized ratio lead to the same local maximizing direction.

In contrast, \(L_2\) and \(L_3\) have local affine forms,

\[
L_2(\delta)
\approx
\|b_2+J_f\delta\|_q,
\]

and

\[
L_3(\delta)
\approx
\|b_3+J_3\delta\|_q.
\]

The ratios

\[
\frac{\|b_2+J_f\delta\|_q}{\|\delta\|_p},
\qquad
\frac{\|b_3+J_3\delta\|_q}{\|\delta\|_p}
\]

are not induced operator norms, because the offsets \(b_2\) and \(b_3\) break homogeneity. Moreover, since \(L_2(0)\) and \(L_3(0)\) are generally nonzero, the total ratios

\[
\frac{L_2(\delta)}{\|\delta\|_p},
\qquad
\frac{L_3(\delta)}{\|\delta\|_p}
\]

may become ill-posed as \(\delta\to 0\). A more meaningful normalized quantity is the incremental ratio,

\[
\frac{L_i(\delta)-L_i(0)}{\|\delta\|_p}.
\]

For \(L_2\), this is

\[
\frac{
\|f(x+\delta)-y\|_q-\|f(x)-y\|_q
}{
\|\delta\|_p
}.
\]

For \(L_3\), this is

\[
\frac{
\|f(x+\delta)-g(x+\delta)\|_q
-
\|f(x)-g(x)\|_q
}{
\|\delta\|_p
}.
\]

These quantities measure loss growth per unit input perturbation, but they are generally not operator norms.

\subsection{Gradient Structure}

For a general affine residual objective,

\[
\Phi(\delta)
=
\frac{1}{q}\|b+J\delta\|_q^q,
\]

define

\[
\phi_q(z)
=
|z|^{q-2}\odot z.
\]

Then

\[
\nabla_\delta \Phi(\delta)
=
J^\top \phi_q(b+J\delta).
\]

For \(L_1\), \(b=0\), so the local gradient is

\[
\nabla_\delta \Phi_1(\delta)
\approx
J_f^\top \phi_q(J_f\delta).
\]

For \(L_2\),

\[
\nabla_\delta \Phi_2(\delta)
\approx
J_f^\top \phi_q(b_2+J_f\delta).
\]

For \(L_3\),

\[
\nabla_\delta \Phi_3(\delta)
\approx
J_3^\top \phi_q(b_3+J_3\delta).
\]

Thus, \(L_2\) and \(L_3\) depend on both the Jacobian and the baseline residual. Their local ascent direction is not determined by \(J\) alone.

When \(p=q=2\), the affine squared objective becomes

\[
\frac12\|b+J\delta\|_2^2
=
\frac12\|b\|_2^2
+
b^\top J\delta
+
\frac12\delta^\top J^\top J\delta.
\]

The gradient is

\[
\nabla_\delta
\frac12\|b+J\delta\|_2^2
=
J^\top b+J^\top J\delta.
\]

For \(L_1\), \(b=0\), and the gradient reduces to

\[
J_f^\top J_f\delta.
\]

For \(L_2\) and \(L_3\), the extra term \(J^\top b\) remains. This is the key reason why these losses are not pure singular-vector or operator-norm problems.

\subsection{Optimization Methods}

The optimization method should be chosen according to the structure of the objective.

\[
\begin{array}{c|c|c|c|c}
\text{Loss} 
& \text{Local Form}
& \text{Natural Interpretation}
& \text{Recommended Method}
& \text{Important Note}
\\
\hline
L_1
&
\|J_f\delta\|_q
&
p\to q \text{ operator gain}
&
\text{generalized power iteration}
&
\text{PGD needs nonzero initialization}
\\
L_2
&
\|b_2+J_f\delta\|_q
&
\text{prediction-error maximization}
&
\text{PGD}
&
\text{depends on } b_2 \text{ and } J_f
\\
L_3
&
\|b_3+J_3\delta\|_q
&
\text{model-solver residual maximization}
&
\text{PGD}
&
\text{depends on } b_3 \text{ and } J_3
\end{array}
\]

PGD solves a constrained loss maximization problem,

\[
\max_{\|\delta\|_p\le \epsilon} L(\delta),
\]

using the update

\[
\delta_{k+1}
=
\Pi_{\|\delta\|_p\le\epsilon}
\left(
\delta_k+\alpha \nabla_\delta L(\delta_k)
\right).
\]

It is a general-purpose method for nonlinear or affine residual losses.

Power iteration and generalized power iteration solve operator-gain problems. For \(p=q=2\), ordinary power iteration solves

\[
\max_{\|v\|_2=1}\|Jv\|_2,
\]

with update

\[
v_{k+1}
=
\frac{J^\top Jv_k}{\|J^\top Jv_k\|_2}.
\]

For general \(p\) and \(q\), generalized power iteration solves

\[
\max_{\|v\|_p=1}\|Jv\|_q,
\]

using

\[
w_k
=
J^\top\phi_q(Jv_k),
\]

and

\[
v_{k+1}
=
\psi_p(w_k)
=
\arg\max_{\|s\|_p\le 1}
\langle w_k,s\rangle.
\]

This method is naturally matched to \(L_1\), because \(L_1\) locally has the homogeneous linear form \(\|J_f\delta\|_q\). It is not directly matched to \(L_2\) or \(L_3\), because those losses contain affine offsets.

\subsection{Summary Table}

\[
\begin{array}{c|c|c|c}
\text{Question}
& L_1
& L_2
& L_3
\\
\hline
\text{Loss}
&
\|f(x+\delta)-f(x)\|_q
&
\|f(x+\delta)-y\|_q
&
\|f(x+\delta)-g(x+\delta)\|_q
\\
\hline
L_i(0)
&
0
&
\|f(x)-y\|_q
&
\|f(x)-g(x)\|_q
\\
\hline
\text{Local form}
&
\|J_f\delta\|_q
&
\|b_2+J_f\delta\|_q
&
\|b_3+J_3\delta\|_q
\\
\hline
\text{Map type}
&
\text{linear}
&
\text{affine}
&
\text{affine}
\\
\hline
\text{Direct ratio } L_i(\delta)/\|\delta\|_p
&
\text{meaningful}
&
\text{not recommended}
&
\text{not recommended}
\\
\hline
\text{Incremental ratio}
&
\frac{L_1(\delta)}{\|\delta\|_p}
&
\frac{L_2(\delta)-L_2(0)}{\|\delta\|_p}
&
\frac{L_3(\delta)-L_3(0)}{\|\delta\|_p}
\\
\hline
\text{Operator norm?}
&
\text{yes, locally}
&
\text{no}
&
\text{no}
\\
\hline
\text{Main direction source}
&
J_f
&
b_2 \text{ and } J_f
&
b_3 \text{ and } J_3
\\
\hline
\text{Natural method}
&
\text{generalized power iteration}
&
\text{PGD}
&
\text{PGD}
\end{array}
\]

\subsection{Main Conclusion}

The key distinction is that \(L_1\) is a homogeneous output-change loss, while \(L_2\) and \(L_3\) are affine residual losses. Therefore,

\[
L_1(\delta)
\approx
\|J_f\delta\|_q
\]

naturally leads to the induced operator norm

\[
\|J_f\|_{p\to q}
=
\max_{\delta\neq0}
\frac{\|J_f\delta\|_q}{\|\delta\|_p}.
\]

Thus generalized power iteration is a natural method for \(L_1\).

In contrast,

\[
L_2(\delta)
\approx
\|b_2+J_f\delta\|_q,
\]

and

\[
L_3(\delta)
\approx
\|b_3+J_3\delta\|_q.
\]

The offsets \(b_2\) and \(b_3\) break homogeneity, so these losses are not operator-norm objectives. They are better understood as affine residual maximization problems. Therefore, PGD is the more natural optimization method for \(L_2\) and \(L_3\).

\clearpage

\begin{landscape}
\thispagestyle{empty}

\footnotesize
\renewcommand{\arraystretch}{1.03}
\setlength{\tabcolsep}{1.35pt}
\setlength{\LTleft}{0pt}
\setlength{\LTright}{0pt}
\sloppy

\begin{longtable}{@{}
p{0.105\linewidth}
p{0.150\linewidth}
p{0.135\linewidth}
p{0.135\linewidth}
p{0.135\linewidth}
p{0.130\linewidth}
p{0.150\linewidth}
@{}}

\caption{Comparison of original loss, baseline-subtracted loss, ratio objectives, and local linearized objectives for $L_1$, $L_2$, and $L_3$.}
\label{tab:loss_objective_comparison}
\\

\toprule
\textbf{Objective Form}
&
\textbf{Formula}
&
\textbf{For $L_1$}
&
\textbf{For $L_2$}
&
\textbf{For $L_3$}
&
\textbf{Meaning}
&
\textbf{Method / Note}
\\
\midrule
\endfirsthead

\toprule
\textbf{Objective Form}
&
\textbf{Formula}
&
\textbf{For $L_1$}
&
\textbf{For $L_2$}
&
\textbf{For $L_3$}
&
\textbf{Meaning}
&
\textbf{Method / Note}
\\
\midrule
\endhead

\midrule
\multicolumn{7}{r}{Continued on next page}
\\
\endfoot

\bottomrule
\endlastfoot

Original loss
&
$L_i(\delta)$
&
$\|f(x+\delta)-f(x)\|_q$
&
$\|f(x+\delta)-y\|_q$
&
$\|f(x+\delta)-g(x+\delta)\|_q$
&
How large the perturbed loss is
&
Standard PGD objective
\\

Fixed-budget original loss
&
$\displaystyle \max_{\|\delta\|_p\le \epsilon} L_i(\delta)$
&
Maximize surrogate output change
&
Maximize prediction error
&
Maximize model-solver residual
&
Worst-case loss inside the $\epsilon$-ball
&
For $L_1$, locally related to $\epsilon\|J_f\|_{p\to q}$; for $L_2,L_3$, not an operator norm
\\

Baseline-subtracted loss increment
&
$L_i(\delta)-L_i(0)$
&
$L_1(\delta)-0=L_1(\delta)$
&
\makecell[l]{
$\|f(x+\delta)-y\|_q$\\
$-\|f(x)-y\|_q$
}
&
\makecell[l]{
$\|f(x+\delta)-g(x+\delta)\|_q$\\
$-\|f(x)-g(x)\|_q$
}
&
How much the perturbation increases the loss
&
Subtracting a constant does not change the argmax under a fixed budget
\\

Fixed-budget loss increment
&
$\displaystyle \max_{\|\delta\|_p\le \epsilon}\bigl[L_i(\delta)-L_i(0)\bigr]$
&
Same optimizer as maximizing $L_1(\delta)$
&
Same optimizer as maximizing $L_2(\delta)$
&
Same optimizer as maximizing $L_3(\delta)$
&
Worst-case loss increase inside the $\epsilon$-ball
&
Since $L_i(0)$ is constant, the optimizer is unchanged
\\

Total ratio
&
$\displaystyle \frac{L_i(\delta)}{\|\delta\|_p}$
&
$\displaystyle \frac{\|f(x+\delta)-f(x)\|_q}{\|\delta\|_p}$
&
$\displaystyle \frac{\|f(x+\delta)-y\|_q}{\|\delta\|_p}$
&
$\displaystyle \frac{\|f(x+\delta)-g(x+\delta)\|_q}{\|\delta\|_p}$
&
Total loss per unit input perturbation
&
For $L_2,L_3$, this is usually not recommended because baseline loss can dominate as $\delta\to0$
\\

Increment ratio
&
$\displaystyle \frac{L_i(\delta)-L_i(0)}{\|\delta\|_p}$
&
$\displaystyle \frac{L_1(\delta)}{\|\delta\|_p}$
&
\makecell[l]{
$\|f(x+\delta)-y\|_q$\\
$-\|f(x)-y\|_q$\\
$/\|\delta\|_p$
}
&
\makecell[l]{
$\|f(x+\delta)-g(x+\delta)\|_q$\\
$-\|f(x)-g(x)\|_q$\\
$/\|\delta\|_p$
}
&
Loss increase per unit perturbation
&
Useful as a sensitivity or efficiency metric; not always the best attack objective
\\

Fixed-radius increment ratio
&
$\displaystyle \max_{\|\delta\|_p=\epsilon}
\frac{L_i(\delta)-L_i(0)}{\|\delta\|_p}$
&
Equivalent to maximizing $L_1(\delta)$
&
Equivalent to maximizing $L_2(\delta)$
&
Equivalent to maximizing $L_3(\delta)$
&
Compare directions on the fixed-radius sphere
&
Since the denominator is fixed as $\epsilon$, the ratio does not change the direction
\\

Variable-radius increment ratio
&
$\displaystyle \max_{0<\|\delta\|_p\le\epsilon}
\frac{L_i(\delta)-L_i(0)}{\|\delta\|_p}$
&
Finds the local Lipschitz / operator-gain direction
&
Finds the direction with fastest prediction-error increase per unit perturbation
&
Finds the direction with fastest residual increase per unit perturbation
&
Maximize perturbation efficiency
&
Different from standard PGD; may prefer small perturbations
\\

Local linearization
&
$L_i(\delta)\approx \cdots$
&
$\|J_f\delta\|_q$
&
$\|b_2+J_f\delta\|_q$
&
$\|b_3+J_3\delta\|_q$
&
Local structure of the loss
&
$b_2=f(x)-y$, $b_3=f(x)-g(x)$, $J_3=J_f-J_g$
\\

Local ratio
&
$\displaystyle \frac{\text{local loss}}{\|\delta\|_p}$
&
$\displaystyle \frac{\|J_f\delta\|_q}{\|\delta\|_p}$
&
$\displaystyle \frac{\|b_2+J_f\delta\|_q}{\|\delta\|_p}$
&
$\displaystyle \frac{\|b_3+J_3\delta\|_q}{\|\delta\|_p}$
&
Local gain per unit perturbation
&
Only $L_1$ gives the standard $p\to q$ operator norm
\\

Local increment ratio
&
$\displaystyle \frac{\|b+J\delta\|_q-\|b\|_q}{\|\delta\|_p}$
&
$\displaystyle \frac{\|J_f\delta\|_q}{\|\delta\|_p}$
&
$\displaystyle \frac{\|b_2+J_f\delta\|_q-\|b_2\|_q}{\|\delta\|_p}$
&
$\displaystyle \frac{\|b_3+J_3\delta\|_q-\|b_3\|_q}{\|\delta\|_p}$
&
Baseline-removed unit growth rate
&
More reasonable than total ratio, but still depends on the baseline residual $b$
\\

Squared local objective when $p=q=2$
&
$\displaystyle \frac12\|b+J\delta\|_2^2$
&
$\displaystyle \frac12\delta^\top J_f^\top J_f\delta$
&
$\displaystyle \frac12\|b_2\|^2+b_2^\top J_f\delta+\frac12\delta^\top J_f^\top J_f\delta$
&
$\displaystyle \frac12\|b_3\|^2+b_3^\top J_3\delta+\frac12\delta^\top J_3^\top J_3\delta$
&
Separate linear and quadratic effects
&
For $L_2,L_3$, the direction depends on both $J^\top b$ and $J^\top J\delta$
\\

Small-perturbation initial direction
&
$\displaystyle \nabla_\delta \frac1q\|b+J\delta\|_q^q\bigg|_{\delta=0}$
&
$0$ or requires nonzero initialization
&
$J_f^\top \phi_q(b_2)$
&
$J_3^\top \phi_q(b_3)$
&
Whether optimization can move from $\delta=0$
&
For $L_1$, squared-loss PGD can get stuck at zero
\\

Most natural optimization method
&
--
&
Generalized power iteration / power iteration; PGD needs nonzero initialization
&
PGD
&
PGD
&
Match the method to the objective
&
$L_1$ is an operator-gain problem; $L_2,L_3$ are affine residual maximization problems
\\

Core problem
&
--
&
Find the largest input-output amplification direction of $J_f$
&
Find perturbations that move $f(x+\delta)$ farther from $y$
&
Find perturbations that move $f(x+\delta)$ farther from $g(x+\delta)$
&
Essential difference among the three losses
&
Only $L_1$ is a pure sensitivity problem
\\

\end{longtable}

\end{landscape}

\clearpage




\begin{table}[htbp]
\centering
\small
\renewcommand{\arraystretch}{1.3}
\setlength{\tabcolsep}{6pt}
\begin{tabular}{p{0.38\textwidth} p{0.28\textwidth} p{0.28\textwidth}}
\hline
\textbf{Question}
&
\textbf{$L_1$}
&
\textbf{$L_2,L_3$}
\\
\hline

Is the loss zero at $\delta=0$?
&
Yes
&
Usually no
\\

Can we directly use $\displaystyle \frac{L(\delta)}{\|\delta\|_p}$?
&
Yes
&
Not recommended, because it can be dominated by the baseline residual
\\

Should we subtract $L(0)$?
&
It makes no difference
&
It is necessary if one wants to define a meaningful ratio
\\

After subtracting $L(0)$, is optimizing without division by $\|\delta\|_p$ the same as optimizing the original loss?
&
Yes
&
Yes
\\

After subtracting $L(0)$, is optimizing with division by $\|\delta\|_p$ the same as optimizing the original loss?
&
The same on a fixed-radius sphere; different when the radius is allowed to vary
&
Generally different
\\

Is it a $p\to q$ operator norm problem?
&
Locally yes
&
No
\\

Is there a generalized singular-vector direction?
&
Yes: 
$\displaystyle \arg\max_{\|v\|_p=1}\|J_fv\|_q$
&
No direction determined only by $J$
\\

Recommended method
&
Generalized power iteration / power iteration
&
PGD
\\

\hline
\end{tabular}
\caption{Simplified comparison between $L_1$ and $L_2,L_3$.}
\label{tab:simplified_loss_comparison}
\end{table}

\section{A Unified Optimization View for PGD, Steepest-Ascent PGD, and Generalized Power Iteration}

We consider a perturbation variable $\delta$ constrained inside an $L_p$ ball,
\[
\|\delta\|_p \le \epsilon.
\]
The local form of the three attack losses can be written uniformly as
\[
r(\delta)=b+J\delta,
\]
where $r(\delta)$ is the local residual, $b$ is the baseline residual, and $J$ is the local Jacobian operator.

For the three losses, this gives
\[
\begin{aligned}
L_1(\delta)&=\|f(x+\delta)-f(x)\|_q,
&
b&=0,
&
J&=J_f,
\\
L_2(\delta)&=\|f(x+\delta)-y\|_q,
&
b&=f(x)-y,
&
J&=J_f,
\\
L_3(\delta)&=\|f(x+\delta)-g(x+\delta)\|_q,
&
b&=f(x)-g(x),
&
J&=J_f-J_g.
\end{aligned}
\]

Thus, all three losses can be locally written as
\[
L(\delta)=\|b+J\delta\|_q.
\]
For smoother differentiation, it is often convenient to use the $q$-power objective
\[
\mathcal{L}(\delta)
=
\frac{1}{q}\|b+J\delta\|_q^q.
\]
Let
\[
r_k=b+J\delta_k.
\]
Then the gradient with respect to $\delta$ is
\[
g_k
=
\nabla_\delta \mathcal{L}(\delta_k)
=
J^\top\left(|r_k|^{q-2}r_k\right).
\]
Equivalently,
\[
g_k
=
J^\top\left(
|b+J\delta_k|^{q-2}(b+J\delta_k)
\right).
\]

This formula is the common core behind PGD, steepest-ascent PGD, and generalized power iteration. The difference among these methods is not the gradient itself, but how this gradient is used to update $\delta$.

\subsection{Special Case: $p=q=2$ and $b=0$}

When
\[
b=0,
\qquad
q=2,
\]
we have
\[
\mathcal{L}(\delta)
=
\frac{1}{2}\|J\delta\|_2^2.
\]
This becomes the quadratic form
\[
\mathcal{L}(\delta)
=
\frac{1}{2}\delta^\top J^\top J\delta.
\]
Therefore,
\[
g_k
=
\nabla_\delta \mathcal{L}(\delta_k)
=
J^\top J\delta_k.
\]

In this special case, the ordinary power iteration update is
\[
\delta_{k+1}
=
\epsilon
\frac{J^\top J\delta_k}
{\|J^\top J\delta_k\|_2}.
\]
Using the gradient notation, this is simply
\[
\delta_{k+1}
=
\epsilon
\frac{g_k}{\|g_k\|_2}.
\]
Therefore, ordinary power iteration can be understood as a special gradient-based direction replacement method for the quadratic operator-gain problem.

\subsection{Projected Gradient Ascent / PGD}

The standard projected gradient ascent update is
\[
\delta_{k+1}
=
\Pi_{\|\delta\|_p\le \epsilon}
\left(
\delta_k+\alpha g_k
\right),
\]
where
\[
g_k=\nabla_\delta \mathcal{L}(\delta_k).
\]

This update has the form
\[
\text{old perturbation}
+
\text{gradient step}
+
\text{projection}.
\]
Thus, PGD is an additive update:
\[
\delta_{k+1}
=
\Pi_{\|\delta\|_p\le \epsilon}
\left(
\delta_k+\alpha g_k
\right).
\]

Here the projection
\[
\Pi_{\|\delta\|_p\le \epsilon}
\]
only ensures that the updated perturbation stays inside the allowed $L_p$ ball. It is a feasibility correction step, not a direction-selection step.

\subsection{$L_p$-Steepest-Ascent Direction}

Instead of directly using $g_k$ as the update direction, one may first find the direction that maximizes the linearized loss increase under an $L_p$ constraint:
\[
s_k
=
\arg\max_{\|s\|_p\le 1}
g_k^\top s.
\]
This direction $s_k$ is the steepest ascent direction under the $L_p$ geometry.

The value of this maximization is given by the dual norm:
\[
\max_{\|s\|_p\le 1} g_k^\top s
=
\|g_k\|_{p^\ast},
\]
where $p^\ast$ is the dual exponent of $p$:
\[
\frac{1}{p}+\frac{1}{p^\ast}=1.
\]

For common choices of $p$, the steepest-ascent direction has simple forms.

\paragraph{Case $p=2$.}
The dual exponent is also $2$, and
\[
s_k
=
\frac{g_k}{\|g_k\|_2}.
\]

\paragraph{Case $p=\infty$.}
The dual exponent is $1$, and
\[
s_k
=
\operatorname{sign}(g_k).
\]
This is the direction used in FGSM and $L_\infty$ PGD-style attacks.

\paragraph{Case $p=1$.}
The dual exponent is $\infty$. The optimal direction concentrates on the coordinate with the largest absolute gradient:
\[
i^\star=\arg\max_i |(g_k)_i|,
\]
and
\[
s_k
=
\operatorname{sign}\bigl((g_k)_{i^\star}\bigr)e_{i^\star}.
\]
Thus, the $L_1$ steepest-ascent direction is sparse.

\subsection{Steepest-Ascent PGD}

After computing the $L_p$-steepest direction $s_k$, one can still use a PGD-style additive update:
\[
\delta_{k+1}
=
\Pi_{\|\delta\|_p\le \epsilon}
\left(
\delta_k+\alpha s_k
\right).
\]

This differs from ordinary PGD. Ordinary PGD uses
\[
\delta_k+\alpha g_k,
\]
while steepest-ascent PGD uses
\[
\delta_k+\alpha s_k,
\]
where
\[
s_k
=
\arg\max_{\|s\|_p\le 1} g_k^\top s.
\]

Thus, the difference is:
\[
\text{PGD uses the raw gradient } g_k,
\]
while
\[
\text{$L_p$-steepest PGD uses the best direction induced by } g_k
\text{ under the } L_p \text{ geometry}.
\]

\subsection{Generalized Power Iteration}

Generalized power iteration uses the same gradient core, but it updates $\delta$ differently. Instead of taking an additive step,
\[
\delta_k+\alpha g_k
\]
or
\[
\delta_k+\alpha s_k,
\]
it directly replaces the current perturbation by the steepest direction:
\[
\delta_{k+1}
=
\epsilon s_k,
\]
where
\[
s_k
=
\arg\max_{\|s\|_p\le 1}
g_k^\top s.
\]

Thus, generalized power iteration has the form
\[
\boxed{
\delta_{k+1}
=
\epsilon
\arg\max_{\|s\|_p\le 1}
g_k^\top s
}
\]
with
\[
g_k
=
J^\top\left(
|b+J\delta_k|^{q-2}(b+J\delta_k)
\right).
\]

Therefore, generalized power iteration can be understood as a replacement update:
\[
\text{new perturbation}
=
\text{current steepest direction}.
\]

In contrast, PGD is an additive update:
\[
\text{new perturbation}
=
\text{old perturbation}
+
\text{one step of ascent}
+
\text{projection}.
\]

\subsection{Unified Comparison}

The three update rules can be summarized as follows.

\[
\begin{array}{c|c|c|c}
\text{Method}
&
\text{Gradient Core}
&
\text{Direction Used}
&
\text{Update Rule}
\\
\hline
\text{PGD}
&
g_k=\nabla_\delta \mathcal{L}(\delta_k)
&
g_k
&
\delta_{k+1}
=
\Pi_{\|\delta\|_p\le\epsilon}
(\delta_k+\alpha g_k)
\\
\hline
L_p\text{-steepest PGD}
&
g_k=\nabla_\delta \mathcal{L}(\delta_k)
&
s_k=\arg\max_{\|s\|_p\le1}g_k^\top s
&
\delta_{k+1}
=
\Pi_{\|\delta\|_p\le\epsilon}
(\delta_k+\alpha s_k)
\\
\hline
\text{Generalized Power Iteration}
&
g_k=\nabla_\delta \mathcal{L}(\delta_k)
&
s_k=\arg\max_{\|s\|_p\le1}g_k^\top s
&
\delta_{k+1}
=
\epsilon s_k
\\
\hline
\text{Ordinary Power Iteration}
&
g_k=J^\top J\delta_k
&
s_k=g_k/\|g_k\|_2
&
\delta_{k+1}
=
\epsilon
\frac{J^\top J\delta_k}{\|J^\top J\delta_k\|_2}
\end{array}
\]

\subsection{Interpretation for the Three Losses}

For $L_1$, we have
\[
b=0,
\qquad
J=J_f.
\]
Therefore,
\[
g_k
=
J_f^\top
\left(
|J_f\delta_k|^{q-2}
J_f\delta_k
\right).
\]
When $q=2$, this reduces to
\[
g_k=J_f^\top J_f\delta_k.
\]
Thus, for $p=q=2$, $L_1$ becomes the standard operator-gain problem, and ordinary power iteration is recovered.

For $L_2$, we have
\[
b_2=f(x)-y,
\qquad
J=J_f.
\]
Therefore,
\[
g_k
=
J_f^\top
\left(
|b_2+J_f\delta_k|^{q-2}
(b_2+J_f\delta_k)
\right).
\]
Because $b_2$ is generally nonzero, this is an affine residual maximization problem rather than a pure operator-gain problem. PGD or steepest-ascent PGD is usually more natural than ordinary power iteration.

For $L_3$, we have
\[
b_3=f(x)-g(x),
\qquad
J_3=J_f-J_g.
\]
Therefore,
\[
g_k
=
J_3^\top
\left(
|b_3+J_3\delta_k|^{q-2}
(b_3+J_3\delta_k)
\right).
\]
Again, because $b_3$ is generally nonzero, $L_3$ is also an affine residual maximization problem. Its optimization is naturally handled by PGD, steepest-ascent PGD, Adam-style optimization on $\delta$, or related first-order constrained methods.

\subsection{Core Takeaway}

All methods can be unified as
\[
r_k=b+J\delta_k,
\]
\[
g_k
=
J^\top\left(|r_k|^{q-2}r_k\right),
\]
followed by different update rules.

PGD uses
\[
\delta_{k+1}
=
\Pi_{\|\delta\|_p\le\epsilon}
(\delta_k+\alpha g_k).
\]

$L_p$-steepest PGD uses
\[
s_k
=
\arg\max_{\|s\|_p\le1}g_k^\top s,
\]
\[
\delta_{k+1}
=
\Pi_{\|\delta\|_p\le\epsilon}
(\delta_k+\alpha s_k).
\]

Generalized power iteration uses
\[
s_k
=
\arg\max_{\|s\|_p\le1}g_k^\top s,
\]
\[
\delta_{k+1}
=
\epsilon s_k.
\]

Thus, the common part is the gradient
\[
g_k=\nabla_\delta \mathcal{L}(\delta_k),
\]
and the difference is whether the method performs an additive PGD step or a direction-replacement step.

In the special case
\[
b=0,
\qquad
p=q=2,
\]
the generalized update reduces to ordinary power iteration:
\[
\delta_{k+1}
=
\epsilon
\frac{J^\top J\delta_k}
{\|J^\top J\delta_k\|_2}.
\]

\subsection{Two Equivalent Views of PGD, Steepest-Ascent PGD, and Generalized Power Iteration}

The optimization problem is written as
\[
\max_{\|\delta\|_p\le \epsilon} \mathcal L(\delta).
\]
At iteration $k$, define the gradient
\[
g_k=\nabla_\delta \mathcal L(\delta_k).
\]

\paragraph{View 1: Gradient-only formulation.}
This formulation does not require writing the Jacobian explicitly. It only assumes that the gradient
\[
g_k=\nabla_\delta \mathcal L(\delta_k)
\]
can be computed, for example by automatic differentiation.

\[
\begin{array}{c|c|c|c}
\textbf{Method}
&
\textbf{Gradient}
&
\textbf{Direction}
&
\textbf{Update}
\\
\hline
\text{PGD}
&
g_k=\nabla_\delta \mathcal L(\delta_k)
&
g_k
&
\delta_{k+1}
=
\Pi_{\|\delta\|_p\le\epsilon}
\left(
\delta_k+\alpha g_k
\right)
\\
\hline
L_p\text{-steepest PGD}
&
g_k=\nabla_\delta \mathcal L(\delta_k)
&
s_k=
\arg\max_{\|s\|_p\le1}
g_k^\top s
&
\delta_{k+1}
=
\Pi_{\|\delta\|_p\le\epsilon}
\left(
\delta_k+\alpha s_k
\right)
\\
\hline
\text{Generalized Power Iteration}
&
g_k=\nabla_\delta \mathcal L(\delta_k)
&
s_k=
\arg\max_{\|s\|_p\le1}
g_k^\top s
&
\delta_{k+1}
=
\epsilon s_k
\\
\hline
\text{Ordinary Power Iteration}
&
g_k=\nabla_\delta \mathcal L(\delta_k)
&
s_k=\dfrac{g_k}{\|g_k\|_2}
&
\delta_{k+1}
=
\epsilon\dfrac{g_k}{\|g_k\|_2}
\quad
\text{when } p=2
\\
\end{array}
\]

In this table, everything is written only through the gradient $g_k$.
The Jacobian does not need to appear explicitly.
The distinction among the methods is only how they use the same gradient.

\[
\text{PGD: additive raw-gradient update.}
\]
\[
L_p\text{-steepest PGD: additive steepest-direction update.}
\]
\[
\text{Generalized power iteration: direction-replacement update.}
\]

\paragraph{The $L_p$ steepest direction.}
The direction
\[
s_k=
\arg\max_{\|s\|_p\le1}g_k^\top s
\]
has simple forms for common $p$:

\[
\begin{array}{c|c|c}
\textbf{Constraint}
&
\textbf{Steepest Direction}
&
\textbf{Interpretation}
\\
\hline
p=2
&
s_k=\dfrac{g_k}{\|g_k\|_2}
&
\text{normalized gradient direction}
\\
\hline
p=\infty
&
s_k=\operatorname{sign}(g_k)
&
\text{FGSM / } L_\infty \text{ steepest direction}
\\
\hline
p=1
&
s_k=
\operatorname{sign}\bigl((g_k)_{i^\star}\bigr)e_{i^\star},
\quad
i^\star=\arg\max_i |(g_k)_i|
&
\text{sparse coordinate direction}
\\
\end{array}
\]

\paragraph{View 2: Jacobian-expanded formulation.}
Now suppose the local residual has the affine form
\[
r(\delta)=b+J\delta,
\]
and the local objective is
\[
\mathcal L(\delta)
=
\frac{1}{q}\|b+J\delta\|_q^q.
\]
Then at iteration $k$,
\[
r_k=b+J\delta_k,
\]
and the gradient becomes
\[
g_k
=
\nabla_\delta \mathcal L(\delta_k)
=
J^\top\left(|r_k|^{q-2}r_k\right).
\]
Equivalently,
\[
g_k
=
J^\top
\left(
|b+J\delta_k|^{q-2}
(b+J\delta_k)
\right).
\]

Substituting this expression into the previous table gives the Jacobian-expanded version.

\[
\begin{array}{c|c|c|c}
\textbf{Method}
&
\textbf{Jacobian-expanded Gradient}
&
\textbf{Direction}
&
\textbf{Update}
\\
\hline
\text{PGD}
&
g_k=
J^\top
\left(
|b+J\delta_k|^{q-2}
(b+J\delta_k)
\right)
&
g_k
&
\delta_{k+1}
=
\Pi_{\|\delta\|_p\le\epsilon}
\left(
\delta_k+\alpha g_k
\right)
\\
\hline
L_p\text{-steepest PGD}
&
g_k=
J^\top
\left(
|b+J\delta_k|^{q-2}
(b+J\delta_k)
\right)
&
s_k=
\arg\max_{\|s\|_p\le1}
g_k^\top s
&
\delta_{k+1}
=
\Pi_{\|\delta\|_p\le\epsilon}
\left(
\delta_k+\alpha s_k
\right)
\\
\hline
\text{Generalized Power Iteration}
&
g_k=
J^\top
\left(
|b+J\delta_k|^{q-2}
(b+J\delta_k)
\right)
&
s_k=
\arg\max_{\|s\|_p\le1}
g_k^\top s
&
\delta_{k+1}
=
\epsilon s_k
\\
\hline
\text{Ordinary Power Iteration}
&
g_k=J^\top J\delta_k
\quad
\text{when } b=0,\ q=2
&
s_k=\dfrac{J^\top J\delta_k}{\|J^\top J\delta_k\|_2}
&
\delta_{k+1}
=
\epsilon
\dfrac{J^\top J\delta_k}
{\|J^\top J\delta_k\|_2}
\\
\end{array}
\]

\paragraph{Mapping to the three losses.}
For the three losses, the parameters $b$ and $J$ are:

\[
\begin{array}{c|c|c|c}
\textbf{Loss}
&
\textbf{Local Residual}
&
\textbf{Baseline } b
&
\textbf{Jacobian } J
\\
\hline
L_1(\delta)=\|f(x+\delta)-f(x)\|_q
&
r_1(\delta)\approx J_f\delta
&
0
&
J_f
\\
\hline
L_2(\delta)=\|f(x+\delta)-y\|_q
&
r_2(\delta)\approx b_2+J_f\delta
&
b_2=f(x)-y
&
J_f
\\
\hline
L_3(\delta)=\|f(x+\delta)-g(x+\delta)\|_q
&
r_3(\delta)\approx b_3+J_3\delta
&
b_3=f(x)-g(x)
&
J_3=J_f-J_g
\\
\end{array}
\]

Therefore, the gradients are:

\[
\begin{array}{c|c}
\textbf{Loss}
&
\textbf{Jacobian-expanded Gradient}
\\
\hline
L_1
&
g_k
=
J_f^\top
\left(
|J_f\delta_k|^{q-2}
J_f\delta_k
\right)
\\
\hline
L_2
&
g_k
=
J_f^\top
\left(
|b_2+J_f\delta_k|^{q-2}
(b_2+J_f\delta_k)
\right)
\\
\hline
L_3
&
g_k
=
J_3^\top
\left(
|b_3+J_3\delta_k|^{q-2}
(b_3+J_3\delta_k)
\right),
\qquad
J_3=J_f-J_g
\\
\end{array}
\]

\paragraph{Key interpretation.}
The gradient-only formulation is the most general one:
\[
g_k=\nabla_\delta \mathcal L(\delta_k).
\]
It does not require explicitly forming $J$ or $A$.

The Jacobian-expanded formulation is only a more explicit local expression:
\[
g_k=
J^\top
\left(
|b+J\delta_k|^{q-2}
(b+J\delta_k)
\right).
\]

In the special case
\[
b=0,
\qquad
q=2,
\]
this reduces to
\[
g_k=J^\top J\delta_k.
\]
Then, if also $p=2$, generalized power iteration becomes ordinary power iteration:
\[
\delta_{k+1}
=
\epsilon
\frac{J^\top J\delta_k}
{\|J^\top J\delta_k\|_2}.
\]

Thus, the hierarchy is:

\[
\text{general gradient form}
\quad
\Longrightarrow
\quad
\text{Jacobian-expanded local form}
\quad
\Longrightarrow
\quad
J^\top J\delta
\text{ in the } b=0,\ q=2 \text{ case}.
\]

The practical implementation can use the gradient-only form. The Jacobian form is mainly for analysis and for understanding when the update reduces to ordinary or generalized power iteration.

\subsection{$L_p$-Steepest PGD as a Bridge Between PGD and Generalized Power Iteration}

The previous discussion shows that PGD, $L_p$-steepest PGD, and generalized power iteration can be written under the same gradient-based framework. The key point is that generalized power iteration does not have to be written explicitly in terms of a matrix such as $A$ or $J^\top J$. In the most general form, it can be written directly using the gradient of the objective with respect to the perturbation variable.

Consider the constrained maximization problem
\[
\max_{\|\delta\|_p\le \epsilon}\mathcal L(\delta).
\]
At iteration $k$, define
\[
g_k=\nabla_\delta \mathcal L(\delta_k).
\]
This gradient-only notation is the most general form. It does not require explicitly forming the Jacobian $J$, the matrix $J^\top J$, or any large matrix $A$.

The ordinary PGD update is
\[
\delta_{k+1}
=
\Pi_{\|\delta\|_p\le\epsilon}
\left(
\delta_k+\alpha g_k
\right).
\]
This update uses the raw gradient $g_k$ directly. It is an additive update:
\[
\text{new perturbation}
=
\text{old perturbation}
+
\text{gradient step}
+
\text{projection}.
\]

However, under an $L_p$ constraint, the raw gradient $g_k$ is not always the steepest direction in the $L_p$ geometry. The steepest ascent direction is obtained by solving the linearized subproblem
\[
s_k
=
\arg\max_{\|s\|_p\le1}
g_k^\top s.
\]
This direction $s_k$ is the direction that gives the largest first-order increase of the objective under a unit $L_p$ constraint. It depends on the geometry of the constraint set.

For common choices of $p$, this gives
\[
p=2:
\qquad
s_k=\frac{g_k}{\|g_k\|_2},
\]
\[
p=\infty:
\qquad
s_k=\operatorname{sign}(g_k),
\]
and
\[
p=1:
\qquad
s_k=
\operatorname{sign}\bigl((g_k)_{i^\star}\bigr)e_{i^\star},
\qquad
i^\star=\arg\max_i |(g_k)_i|.
\]

Therefore, $L_p$-steepest PGD uses the generalized-power-style direction selection step, but still keeps the PGD-style additive update:
\[
\delta_{k+1}
=
\Pi_{\|\delta\|_p\le\epsilon}
\left(
\delta_k+\alpha s_k
\right).
\]
Thus, $L_p$-steepest PGD is a bridge between PGD and generalized power iteration:
\[
\boxed{
L_p\text{-steepest PGD}
=
\text{generalized-power-style direction selection}
+
\text{PGD-style additive update and projection}.
}
\]

Generalized power iteration uses the same direction $s_k$, but updates the perturbation by direct replacement:
\[
\delta_{k+1}
=
\epsilon s_k.
\]
Equivalently,
\[
\delta_{k+1}
=
\epsilon
\arg\max_{\|s\|_p\le1}
g_k^\top s.
\]
Hence, generalized power iteration can also be understood as a gradient-based method:
\[
\boxed{
\text{generalized power iteration}
=
\text{gradient-induced steepest direction replacement}.
}
\]

This makes the relationship among the three methods clear:
\[
\begin{array}{c|c|c}
\textbf{Method}
&
\textbf{Direction Selection}
&
\textbf{Update Rule}
\\
\hline
\text{PGD}
&
g_k
&
\delta_{k+1}
=
\Pi_{\|\delta\|_p\le\epsilon}
(\delta_k+\alpha g_k)
\\
\hline
L_p\text{-steepest PGD}
&
s_k=
\arg\max_{\|s\|_p\le1}g_k^\top s
&
\delta_{k+1}
=
\Pi_{\|\delta\|_p\le\epsilon}
(\delta_k+\alpha s_k)
\\
\hline
\text{Generalized Power Iteration}
&
s_k=
\arg\max_{\|s\|_p\le1}g_k^\top s
&
\delta_{k+1}
=
\epsilon s_k
\end{array}
\]

The difference is therefore not in the gradient itself, but in how the gradient is used.

PGD uses the raw gradient and performs an additive step.  
$L_p$-steepest PGD first converts the gradient into the steepest $L_p$ direction, and then performs an additive PGD-style step.  
Generalized power iteration also converts the gradient into the steepest $L_p$ direction, but directly replaces the old perturbation by this new direction.

In the special quadratic case
\[
\mathcal L(\delta)
=
\frac12\delta^\top A\delta,
\]
the gradient is
\[
g_k=A\delta_k.
\]
If
\[
A=J^\top J,
\]
then
\[
g_k=J^\top J\delta_k.
\]
Thus, the ordinary power iteration
\[
\delta_{k+1}
=
\epsilon
\frac{J^\top J\delta_k}
{\|J^\top J\delta_k\|_2}
\]
can be rewritten as
\[
\delta_{k+1}
=
\epsilon
\frac{g_k}{\|g_k\|_2}.
\]
This shows that even ordinary power iteration is a special case of the gradient-based replacement view.

Therefore, the matrix form
\[
J^\top J\delta_k
\]
is only a special explicit expression for the gradient in the case of a squared $L_2$ operator-gain objective. In the general $p,q$ setting, the unified formulation should be written through
\[
g_k=\nabla_\delta \mathcal L(\delta_k),
\]
while the Jacobian-expanded form is only an analytical specialization.

\section{Automatic Differentiation, Gradients, JVP, and VJP}

In the previous optimization framework, the attack problem was written as
\[
\max_{\|\delta\|_p\le \epsilon} \mathcal L(\delta),
\]
where $\delta$ is the perturbation variable. For PGD, $L_p$-steepest PGD, Adam-style optimization, or generalized power-iteration-style updates, the central quantity is the gradient
\[
g_k=\nabla_\delta \mathcal L(\delta_k).
\]

In practical deep learning frameworks, this gradient is usually computed directly by automatic differentiation. One does not need to explicitly construct the Jacobian $J$, the matrix $J^\top J$, or any large matrix $A$. The Jacobian notation is mainly useful for mathematical analysis.

\subsection{JVP and VJP}

Suppose
\[
r=f(\delta),
\]
where
\[
f:\mathbb R^n\to \mathbb R^m.
\]
The Jacobian of $f$ is
\[
J_f(\delta)=\frac{\partial f}{\partial \delta}.
\]

There are two important Jacobian-related operations.

\paragraph{Jacobian-vector product.}
Given an input-space direction $v\in\mathbb R^n$, the Jacobian-vector product is
\[
J_f(\delta)v.
\]
This maps an input perturbation direction into the output space. It describes the first-order output change caused by moving along $v$:
\[
f(\delta+\eta v)
\approx
f(\delta)+\eta J_f(\delta)v.
\]
This is usually associated with forward-mode automatic differentiation.

\paragraph{Vector-Jacobian product.}
Given an output-space vector $u\in\mathbb R^m$, the vector-Jacobian product is
\[
J_f(\delta)^\top u.
\]
This maps an output-space sensitivity back into the input space. It is the main operation behind reverse-mode automatic differentiation and backpropagation.

\subsection{Gradient of a Scalar Loss}

Now suppose the scalar loss has the composition form
\[
\mathcal L(\delta)=\ell(r(\delta)),
\]
where
\[
r(\delta)=f(\delta).
\]
By the chain rule,
\[
\nabla_\delta \mathcal L(\delta)
=
J_f(\delta)^\top \nabla_r \ell(r).
\]
Therefore, computing the gradient of a scalar loss with respect to $\delta$ is exactly a vector-Jacobian product:
\[
\boxed{
\nabla_\delta \mathcal L
=
J_f^\top
\left(
\nabla_r \ell
\right).
}
\]

This is the reason why ordinary backpropagation is naturally suited for computing attack gradients. For a scalar loss, reverse-mode automatic differentiation can compute
\[
\nabla_\delta \mathcal L(\delta)
\]
efficiently without explicitly forming the full Jacobian.

\subsection{Application to the Local Residual Loss}

The local form of the losses can be written as
\[
r(\delta)=b+J\delta,
\]
and a common smooth objective is
\[
\mathcal L(\delta)
=
\frac{1}{q}\|b+J\delta\|_q^q.
\]
Let
\[
r=b+J\delta.
\]
Then
\[
\nabla_r \mathcal L
=
|r|^{q-2}r.
\]
Using the chain rule,
\[
\nabla_\delta \mathcal L
=
J^\top
\left(
|r|^{q-2}r
\right).
\]
Thus,
\[
\boxed{
\nabla_\delta \mathcal L(\delta)
=
J^\top
\left(
|b+J\delta|^{q-2}(b+J\delta)
\right).
}
\]

This formula shows the relationship between the Jacobian-expanded expression and the gradient-only expression:
\[
g_k
=
\nabla_\delta \mathcal L(\delta_k).
\]
The automatic differentiation framework computes this gradient directly through the computational graph. The explicit Jacobian form is not required in implementation.

\subsection{Special Case: Squared $L_2$ Loss}

If
\[
q=2,
\]
then
\[
\mathcal L(\delta)
=
\frac12\|b+J\delta\|_2^2.
\]
The gradient becomes
\[
\nabla_\delta \mathcal L
=
J^\top(b+J\delta).
\]

If additionally
\[
b=0,
\]
then
\[
\mathcal L(\delta)
=
\frac12\|J\delta\|_2^2
=
\frac12\delta^\top J^\top J\delta,
\]
and
\[
\nabla_\delta \mathcal L
=
J^\top J\delta.
\]

Therefore,
\[
J^\top J\delta
\]
is not a separate object from the gradient. It is the gradient of the squared $L_2$ operator-gain objective:
\[
\boxed{
J^\top J\delta
=
\nabla_\delta
\left(
\frac12\|J\delta\|_2^2
\right).
}
\]

This is why ordinary power iteration can be written either in matrix form,
\[
\delta_{k+1}
=
\epsilon
\frac{J^\top J\delta_k}
{\|J^\top J\delta_k\|_2},
\]
or in gradient form,
\[
\delta_{k+1}
=
\epsilon
\frac{g_k}{\|g_k\|_2},
\qquad
g_k=\nabla_\delta \mathcal L(\delta_k).
\]

\subsection{Gradient Computation in Practice}

For practical attack optimization, the most direct implementation is usually to compute
\[
g_k=\nabla_\delta \mathcal L(\delta_k)
\]
using automatic differentiation.

For example, in PyTorch one typically writes
\[
\texttt{loss.backward()}
\]
and reads the gradient from
\[
\texttt{delta.grad}.
\]
In JAX, one typically writes
\[
\texttt{grad\_delta = jax.grad(attack\_loss)(delta)}.
\]

These commands compute the scalar loss gradient directly. Internally, the automatic differentiation system applies the chain rule through the computational graph. For scalar losses, this is mainly reverse-mode automatic differentiation, which corresponds to vector-Jacobian products.

Thus, for PGD, $L_p$-steepest PGD, and Adam-style optimization, the practical update only needs
\[
g_k=\nabla_\delta \mathcal L(\delta_k).
\]

\subsection{When JVP and VJP Are Useful}

Although direct gradient computation is usually the most convenient way to optimize a scalar attack loss, JVP and VJP are still useful in specific situations.

\paragraph{JVP is useful when one needs}
\[
Jv.
\]
This appears when analyzing the local output change caused by a specific input perturbation direction $v$.

\paragraph{VJP is useful when one needs}
\[
J^\top u.
\]
This appears when pulling an output-space sensitivity back to the input space.

\paragraph{JVP plus VJP is useful when one needs}
\[
J^\top Jv.
\]
This is required in matrix-free power iteration or Lanczos methods for estimating the dominant singular direction of a Jacobian.

Therefore, JVP and VJP are especially useful for local linearized operator analysis, such as
\[
\max_{\|v\|_2=1}\|Jv\|_2,
\]
or for estimating singular values and singular vectors without explicitly constructing $J$.

\subsection{Main Conclusion}

The practical and mathematical views are consistent.

The implementation view is
\[
\boxed{
g_k=\nabla_\delta \mathcal L(\delta_k).
}
\]
This is the quantity needed by PGD, $L_p$-steepest PGD, generalized-power-style updates, and Adam-style perturbation optimization.

The Jacobian-expanded view is
\[
\boxed{
g_k
=
J^\top
\left(
|b+J\delta_k|^{q-2}
(b+J\delta_k)
\right).
}
\]
This expression explains what the gradient contains, but it does not need to be explicitly implemented.

In the special squared $L_2$ case with $b=0$, this reduces to
\[
\boxed{
g_k=J^\top J\delta_k.
}
\]
Thus, the familiar matrix form in power iteration is simply a special explicit form of the gradient.

The general rule is:
\[
\boxed{
\text{for optimizing the true scalar loss, use automatic differentiation to compute }
\nabla_\delta \mathcal L;
}
\]
\[
\boxed{
\text{for analyzing local Jacobian operator directions, use JVP/VJP or matrix-free spectral methods.}
}
\]

\section{Induced Matrix Norms, Generalized Singular Vectors, and Affine Residual Gains}

This section summarizes the relationship among induced matrix norms, generalized power iteration, generalized singular vectors, and the affine residual gain objectives used in the attack losses. The main point is that the standard matrix norm problem is already an optimization problem, and the affine residual version is a natural but nonstandard extension of this variational viewpoint.

\subsection{Induced $p\to q$ Matrix Norm}

Let
\[
J:\mathbb R^n\to\mathbb R^m
\]
be a linear operator. The induced $p\to q$ norm of $J$ is defined by
\[
\|J\|_{p\to q}
=
\sup_{\delta\neq 0}
\frac{\|J\delta\|_q}{\|\delta\|_p}.
\]
Equivalently, by the homogeneity of the linear map $\delta\mapsto J\delta$ and the homogeneity of the norms,
\[
\|J\|_{p\to q}
=
\sup_{\|\delta\|_p=1}
\|J\delta\|_q.
\]

Thus, the induced $p\to q$ matrix norm is not merely a formula. It is itself a constrained maximization problem:
\[
\max_{\|\delta\|_p=1}
\|J\delta\|_q.
\]

In the special case
\[
p=q=2,
\]
this reduces to the spectral norm:
\[
\|J\|_{2\to2}
=
\sigma_{\max}(J),
\]
where $\sigma_{\max}(J)$ is the largest singular value of $J$. The maximizing direction is the standard largest right singular vector.

For general $p$ and $q$, the induced norm remains well-defined, but it usually does not have a simple closed-form formula. It is generally treated as an optimization problem.

\subsection{Right $(p,q)$-Singular Vectors}

The maximizing direction of the induced $p\to q$ norm can be defined as
\[
\delta^\star
\in
\arg\max_{\|\delta\|_p=1}
\|J\delta\|_q.
\]
Such a direction may be called a right $(p,q)$-singular vector of $J$, or more conservatively, a maximizing direction of the induced $p\to q$ norm.

The corresponding value is
\[
\sigma_{p\to q}(J)
=
\|J\|_{p\to q}
=
\|J\delta^\star\|_q.
\]

When
\[
p=q=2,
\]
this definition reduces to the standard largest singular value and standard largest right singular vector:
\[
J^\top J\delta^\star
=
\sigma_{\max}(J)^2 \delta^\star.
\]

For general $p,q$, the object is defined variationally through the norm maximization problem. There is usually no orthogonal SVD structure comparable to the Euclidean case.

\subsection{Generalized Power Iteration for the $p\to q$ Norm}

The induced $p\to q$ norm problem can be written as
\[
\max_{\|\delta\|_p=1}
\frac{1}{q}\|J\delta\|_q^q.
\]
Let
\[
\phi_q(z)=|z|^{q-2}z
\]
for $1<q<\infty$. The gradient of the objective with respect to $\delta$ is
\[
g(\delta)
=
J^\top \phi_q(J\delta).
\]

A generalized power iteration uses this gradient-like quantity to update the direction. One common form is
\[
g_k
=
J^\top \phi_q(J\delta_k),
\]
then
\[
\delta_{k+1}
=
\arg\max_{\|s\|_p\le 1}
g_k^\top s.
\]

The second step converts the current gradient into the steepest direction under the $L_p$ geometry.

For common values of $p$, this direction has simple forms:
\[
p=2:
\qquad
\delta_{k+1}
=
\frac{g_k}{\|g_k\|_2},
\]
\[
p=\infty:
\qquad
\delta_{k+1}
=
\operatorname{sign}(g_k),
\]
and
\[
p=1:
\qquad
\delta_{k+1}
=
\operatorname{sign}\bigl((g_k)_{i^\star}\bigr)e_{i^\star},
\qquad
i^\star=\arg\max_i |(g_k)_i|.
\]

When
\[
p=q=2,
\]
we have
\[
\phi_2(z)=z,
\]
and therefore
\[
g_k
=
J^\top J\delta_k.
\]
The update becomes
\[
\delta_{k+1}
=
\frac{J^\top J\delta_k}
{\|J^\top J\delta_k\|_2},
\]
which is the ordinary power iteration for the dominant right singular vector.

Thus, ordinary power iteration is the special case of generalized power iteration when
\[
p=q=2.
\]

\subsection{Generalized Power Iteration Is Not a Free Update Direction}

A right $(p,q)$-singular vector is itself the solution of the optimization problem
\[
\delta^\star
\in
\arg\max_{\|\delta\|_p=1}
\|J\delta\|_q.
\]
Therefore, the instruction
\[
\text{``move in the right $(p,q)$-singular vector direction''}
\]
is not a cheap primitive operation. It means that one must first solve, or approximately solve, an inner $p\to q$ operator norm maximization problem.

In particular, for a nonlinear model with a local Jacobian
\[
J_k=J(x+\delta_k),
\]
using the exact local right $(p,q)$-singular direction at every attack step would require solving
\[
u_k^\star
\in
\arg\max_{\|u\|_p=1}
\|J_k u\|_q
\]
at every outer iteration. This is expensive and generally only possible approximately.

Therefore, practical methods usually use one of the following approximations:
\[
\text{one initial generalized singular direction, followed by PGD refinement,}
\]
\[
\text{a few inner generalized-power iterations per outer step,}
\]
or
\[
\text{direct first-order optimization of the true attack objective.}
\]

\subsection{Objective, Gradient, and Update Direction Are Different Objects}

In attack optimization, one must distinguish three objects:
\[
\text{objective: } \mathcal L(\delta),
\]
\[
\text{gradient: } g_k=\nabla_\delta\mathcal L(\delta_k),
\]
and
\[
\text{update direction: } d_k.
\]

In raw PGD, these are connected by
\[
d_k=g_k,
\]
and
\[
\delta_{k+1}
=
\Pi_{\|\delta\|_p\le\epsilon}
\left(
\delta_k+\alpha g_k
\right).
\]

In $L_p$-steepest PGD, the same objective gradient is first converted into the steepest $L_p$ direction:
\[
s_k
=
\arg\max_{\|s\|_p\le1}
g_k^\top s,
\]
and then the update is
\[
\delta_{k+1}
=
\Pi_{\|\delta\|_p\le\epsilon}
\left(
\delta_k+\alpha s_k
\right).
\]

In a generalized-power-style replacement update, the same steepest direction may replace the old perturbation:
\[
\delta_{k+1}
=
\epsilon s_k.
\]

Thus, PGD, $L_p$-steepest PGD, and generalized power iteration can be unified as follows:
\[
\begin{array}{c|c|c}
\textbf{Method}
&
\textbf{Direction}
&
\textbf{Update}
\\
\hline
\text{Raw PGD}
&
g_k=\nabla_\delta\mathcal L(\delta_k)
&
\delta_{k+1}
=
\Pi_{\|\delta\|_p\le\epsilon}
(\delta_k+\alpha g_k)
\\
\hline
L_p\text{-steepest PGD}
&
s_k=
\arg\max_{\|s\|_p\le1}g_k^\top s
&
\delta_{k+1}
=
\Pi_{\|\delta\|_p\le\epsilon}
(\delta_k+\alpha s_k)
\\
\hline
\text{Generalized power replacement}
&
s_k=
\arg\max_{\|s\|_p\le1}g_k^\top s
&
\delta_{k+1}
=
\epsilon s_k
\end{array}
\]

The difference is not necessarily the objective. The difference is how the gradient or gradient-induced direction is used to update $\delta$.

Thus, PGD, \(L_p\)-steepest PGD, and generalized power replacement can be written in a unified projected-proposal form:
\[
\delta_{k+1}
=
\Pi_{\|\delta\|_p\le\epsilon}(z_{k+1}).
\]
Their difference lies in how the proposal \(z_{k+1}\) is constructed:
\[
\begin{array}{c|c|c|c}
\textbf{Method}
&
\textbf{Direction}
&
\textbf{Proposal}
&
\textbf{Projected update}
\\
\hline
\text{Raw PGD}
&
g_k=\nabla_\delta\mathcal L(\delta_k)
&
z_{k+1}=\delta_k+\alpha g_k
&
\delta_{k+1}
=
\Pi_{\|\delta\|_p\le\epsilon}(z_{k+1})
\\
\hline
L_p\text{-steepest PGD}
&
s_k=
\arg\max_{\|s\|_p\le1}g_k^\top s
&
z_{k+1}=\delta_k+\alpha s_k
&
\delta_{k+1}
=
\Pi_{\|\delta\|_p\le\epsilon}(z_{k+1})
\\
\hline
\text{Generalized power replacement}
&
s_k=
\arg\max_{\|s\|_p\le1}g_k^\top s
&
z_{k+1}=\tau s_k,\quad \tau\ge\epsilon
&
\delta_{k+1}
=
\Pi_{\|\delta\|_p\le\epsilon}(z_{k+1})
=
\epsilon s_k
\end{array}
\]

For the generalized power replacement update, since \(s_k\) is chosen from the unit \(L_p\)-ball and is typically taken on its boundary,
\[
\|s_k\|_p=1.
\]
Therefore, for any \(\tau\ge\epsilon\),
\[
\|\tau s_k\|_p=\tau\ge\epsilon.
\]
Projecting \(\tau s_k\) back to the feasible \(L_p\)-ball gives the boundary point in the same direction:
\[
\Pi_{\|\delta\|_p\le\epsilon}(\tau s_k)
=
\epsilon s_k.
\]
Hence the generalized power replacement update
\[
\delta_{k+1}=\epsilon s_k
\]
can also be written as a projected update:
\[
\delta_{k+1}
=
\Pi_{\|\delta\|_p\le\epsilon}(\tau s_k),
\qquad
\tau\ge\epsilon.
\]
The key difference is that PGD-type methods use additive proposals of the form
\[
z_{k+1}=\delta_k+\alpha d_k,
\]
whereas generalized power replacement uses a replacement proposal
\[
z_{k+1}=\tau s_k.
\]
Thus, all three methods share the same projection operator, but they differ in how the pre-projection proposal is constructed.

All methods can be written in the unified projected-proposal form
\[
\delta_{k+1}
=
\Pi_{\|\delta\|_p\le\epsilon}(z_{k+1}).
\]
The proposal \(z_{k+1}\) is constructed from a direction \(d_k\). We distinguish two proposal rules:
\[
z_{k+1}^{add}
=
\delta_k+\alpha d_k,
\]
and
\[
z_{k+1}^{rep}
=
\tau d_k,
\qquad
\tau\ge \epsilon.
\]
The first one is an additive PGD-style proposal, while the second one is a replacement / boundary proposal.

The three direction rules are:
\[
d_k^{raw}
=
g_k
=
\nabla_\delta \mathcal L(\delta_k),
\]
\[
d_k^{lp}
=
s_k
=
\arg\max_{\|s\|_p\le1} g_k^\top s,
\]
and
\[
d_k^{pow}
=
u_k,
\]
where \(u_k\) denotes the generalized-power direction. In the simplest gradient-induced version, \(u_k=s_k\). In a quadratic or operator-norm-type problem, \(u_k\) is the normalized power-iteration direction induced by the current gradient or linearized operator.
\subsection{A Unified Direction--Proposal--Projection View}

We separate the update rule into two independent components: a \emph{direction rule} and a \emph{proposal rule}. 
All methods are written in the unified projected-proposal form
\[
\delta_{k+1}
=
\Pi_{\|\delta\|_p\le\epsilon}(z_{k+1}),
\]
where \(z_{k+1}\) is the pre-projection proposal.

Let
\[
g_k
=
\nabla_\delta \mathcal L(\delta_k)
\]
be the raw gradient of the attack objective with respect to the perturbation. The raw gradient \(g_k\) is not necessarily a unit \(L_p\)-direction, and its norm can vary across iterations. When a unit raw-gradient direction is needed, we define
\[
\widehat g_k
=
\frac{g_k}{\|g_k\|_p}.
\]

We also define the \(L_p\)-steepest direction induced by \(g_k\) as
\[
s_k
=
\arg\max_{\|s\|_p\le 1} g_k^\top s.
\]
Unlike \(g_k\), the direction \(s_k\) is constrained by construction. When \(g_k\neq 0\), the maximizer is typically attained on the boundary of the unit \(L_p\)-ball, so that
\[
\|s_k\|_p=1.
\]

Finally, let
\[
u_k
\]
denote the generalized-power direction. In the simplest gradient-induced implementation, one may take
\[
u_k=s_k.
\]
In an operator-norm or quadratic local approximation, \(u_k\) may instead represent the normalized power-iteration direction induced by the corresponding local linearized operator.

For any chosen direction \(d_k\), we consider two proposal rules. The first is the additive PGD-style proposal
\[
z_{k+1}^{add}
=
\delta_k+\alpha d_k,
\]
where the previous perturbation \(\delta_k\) is preserved and the new direction is added with step size \(\alpha\). The second is the replacement proposal
\[
z_{k+1}^{rep}
=
\tau d_k,
\qquad
\tau\ge \epsilon,
\]
where the previous perturbation \(\delta_k\) is discarded and the proposal is placed in the current direction \(d_k\). If \(\|d_k\|_p=1\), then
\[
\Pi_{\|\delta\|_p\le\epsilon}(\tau d_k)
=
\epsilon d_k,
\qquad
\tau\ge\epsilon.
\]
Thus, replacement proposals directly place the perturbation on the boundary of the feasible \(L_p\)-ball.

The resulting six direction--proposal combinations are summarized as follows:
\[
\begin{array}{c|c|c|c}
\textbf{Direction rule}
&
\textbf{Proposal rule}
&
\textbf{Proposal}
&
\textbf{Projected update}
\\
\hline
\text{Raw gradient}
&
\text{Additive}
&
z_{k+1}
=
\delta_k+\alpha g_k
&
\delta_{k+1}
=
\Pi_{\|\delta\|_p\le\epsilon}
(\delta_k+\alpha g_k)
\\
\hline
\text{Raw gradient}
&
\text{Replacement}
&
z_{k+1}
=
\tau \widehat g_k,
\quad
\tau\ge\epsilon,
\quad
\widehat g_k=
\dfrac{g_k}{\|g_k\|_p}
&
\delta_{k+1}
=
\Pi_{\|\delta\|_p\le\epsilon}
(\tau \widehat g_k)
=
\epsilon \widehat g_k
\\
\hline
L_p\text{-steepest}
&
\text{Additive}
&
z_{k+1}
=
\delta_k+\alpha s_k
&
\delta_{k+1}
=
\Pi_{\|\delta\|_p\le\epsilon}
(\delta_k+\alpha s_k)
\\
\hline
L_p\text{-steepest}
&
\text{Replacement}
&
z_{k+1}
=
\tau s_k,
\quad
\tau\ge\epsilon
&
\delta_{k+1}
=
\Pi_{\|\delta\|_p\le\epsilon}
(\tau s_k)
=
\epsilon s_k
\\
\hline
\text{Generalized power}
&
\text{Additive}
&
z_{k+1}
=
\delta_k+\alpha u_k
&
\delta_{k+1}
=
\Pi_{\|\delta\|_p\le\epsilon}
(\delta_k+\alpha u_k)
\\
\hline
\text{Generalized power}
&
\text{Replacement}
&
z_{k+1}
=
\tau u_k,
\quad
\tau\ge\epsilon
&
\delta_{k+1}
=
\Pi_{\|\delta\|_p\le\epsilon}
(\tau u_k)
=
\epsilon u_k
\end{array}
\]

The first row is the standard raw PGD update. It uses the raw gradient \(g_k\) and an additive proposal:
\[
\delta_{k+1}
=
\Pi_{\|\delta\|_p\le\epsilon}
(\delta_k+\alpha g_k).
\]
The second row is the boundary-replacement version of raw PGD. Since the raw gradient \(g_k\) is not generally a unit \(L_p\)-direction, it must first be normalized:
\[
\widehat g_k
=
\frac{g_k}{\|g_k\|_p}.
\]
Then the update
\[
\delta_{k+1}
=
\epsilon \widehat g_k
\]
places the perturbation directly on the boundary of the feasible ball.

The third row is the \(L_p\)-steepest PGD update. It uses the \(L_p\)-steepest direction \(s_k\), but still uses an additive PGD-style proposal:
\[
\delta_{k+1}
=
\Pi_{\|\delta\|_p\le\epsilon}
(\delta_k+\alpha s_k).
\]
Thus, compared with raw PGD, the direction rule is changed from \(g_k\) to \(s_k\), while the additive proposal rule remains the same.

The fourth row is the \(L_p\)-steepest replacement update. It uses the same direction \(s_k\), but replaces the additive proposal by a boundary replacement proposal:
\[
z_{k+1}
=
\tau s_k,
\qquad
\tau\ge\epsilon.
\]
Because \(\|s_k\|_p=1\), projection gives
\[
\delta_{k+1}
=
\Pi_{\|\delta\|_p\le\epsilon}(\tau s_k)
=
\epsilon s_k.
\]

The fifth row is the additive version of generalized power iteration. It uses the generalized-power direction \(u_k\), but updates it in a PGD-like additive manner:
\[
\delta_{k+1}
=
\Pi_{\|\delta\|_p\le\epsilon}
(\delta_k+\alpha u_k).
\]
This row is useful as an ablation because it keeps the generalized-power direction while changing the proposal rule from replacement to additive.

The sixth row is the generalized power replacement update:
\[
\delta_{k+1}
=
\Pi_{\|\delta\|_p\le\epsilon}
(\tau u_k)
=
\epsilon u_k.
\]
This method discards the previous perturbation \(\delta_k\) and directly replaces it with the current boundary-scaled generalized-power direction.

This decomposition shows that the difference between PGD, \(L_p\)-steepest PGD, and generalized power replacement is not only the direction rule. It is also the proposal rule. PGD-type methods use additive proposals of the form
\[
z_{k+1}
=
\delta_k+\alpha d_k,
\]
while replacement-type methods use boundary proposals of the form
\[
z_{k+1}
=
\tau d_k,
\qquad
\tau\ge\epsilon.
\]
Therefore, replacement methods reach the boundary immediately, whereas additive PGD methods may require several iterations to reach the boundary when \(\alpha\) is small.

If the generalized-power direction is implemented as the same \(L_p\)-steepest direction,
\[
u_k=s_k,
\]
then some rows in the table become identical. In particular,
\[
\Pi_{\|\delta\|_p\le\epsilon}
(\delta_k+\alpha u_k)
=
\Pi_{\|\delta\|_p\le\epsilon}
(\delta_k+\alpha s_k),
\]
so the generalized-power additive update becomes the same as \(L_p\)-steepest PGD. Similarly,
\[
\epsilon u_k
=
\epsilon s_k,
\]
so the generalized-power replacement update becomes the same as the \(L_p\)-steepest replacement update.

Thus, the six-row table should be understood as a conceptual direction--proposal decomposition. In a specific implementation, some rows may coincide depending on how the generalized-power direction \(u_k\) is defined. The key distinction is:
\[
\text{additive proposal:}
\qquad
z_{k+1}
=
\delta_k+\alpha d_k,
\]
versus
\[
\text{replacement proposal:}
\qquad
z_{k+1}
=
\tau d_k,
\qquad
\tau\ge\epsilon.
\]
The additive proposal preserves the previous perturbation and moves locally, while the replacement proposal discards the previous perturbation and directly places the new perturbation on the boundary.














\subsection{Affine Residual Maps}

For the second and third attack losses, the local residual is not purely linear in $\delta$. Instead, it has the affine form
\[
r(\delta)=b+J\delta.
\]
Here $b$ is the baseline residual and $J$ is the local Jacobian of the residual map.

For example, if
\[
L_2(\delta)
=
\|f(x+\delta)-g(x)\|_q,
\]
then
\[
b_2=f(x)-g(x),
\qquad
J=J_f,
\]
and locally
\[
f(x+\delta)-g(x)
\approx
b_2+J_f\delta.
\]

For
\[
L_3(\delta)
=
\|f(x+\delta)-g(x+\delta)\|_q,
\]
we have
\[
b_3=f(x)-g(x),
\qquad
J_3=J_f-J_g,
\]
and locally
\[
f(x+\delta)-g(x+\delta)
\approx
b_3+J_3\delta.
\]

If the solver output is stopped from backpropagation, then the forward residual may still use
\[
g(x+\delta),
\]
but the backward gradient does not include the solver Jacobian $J_g$.

\subsection{Why the Direct Affine Ratio Is Unstable}

A direct affine ratio
\[
\frac{\|b+J\delta\|_q}{\|\delta\|_p}
\]
is generally unstable when $b\neq0$. As
\[
\delta\to0,
\]
the numerator satisfies
\[
\|b+J\delta\|_q\to\|b\|_q,
\]
while the denominator satisfies
\[
\|\delta\|_p\to0.
\]
Therefore,
\[
\frac{\|b+J\delta\|_q}{\|\delta\|_p}
\to\infty
\]
whenever $b\neq0$.

Thus, for affine residual losses, the meaningful ratio is not the total ratio. The meaningful quantity is the baseline-subtracted increment ratio:
\[
\frac{
\|b+J\delta\|_q-\|b\|_q
}{
\|\delta\|_p
}.
\]

This removes the static baseline residual and measures the increase in residual norm per unit input perturbation.

\subsection{Finite-Radius Affine Residual Gain}

Motivated by the induced $p\to q$ norm, we define a finite-radius affine residual gain.

\begin{definition}[Finite-radius affine residual gain]
Let
\[
J:\mathbb R^n\to\mathbb R^m
\]
be a linear operator, let
\[
b\in\mathbb R^m
\]
be a baseline residual, and let $\|\cdot\|_p$ and $\|\cdot\|_q$ be input and output norms. For $\epsilon>0$, define
\[
\mathcal G_{p\to q}(J,b;\epsilon)
=
\sup_{0<\|\delta\|_p\le\epsilon}
\frac{
\|b+J\delta\|_q-\|b\|_q
}{
\|\delta\|_p
}.
\]
We call $\mathcal G_{p\to q}(J,b;\epsilon)$ the finite-radius affine residual gain.
\end{definition}

This quantity measures the maximum per-unit increase of the residual norm around the baseline residual $b$ within an input perturbation radius $\epsilon$.

When
\[
b=0,
\]
the definition reduces to
\[
\mathcal G_{p\to q}(J,0;\epsilon)
=
\sup_{0<\|\delta\|_p\le\epsilon}
\frac{\|J\delta\|_q}{\|\delta\|_p}
=
\|J\|_{p\to q}.
\]
Thus, the ordinary induced $p\to q$ matrix norm is the special case of the affine residual gain when the baseline residual vanishes.

\subsection{Affine Residual Gain Maximizing Direction}

The maximizer of the finite-radius affine residual gain is defined by
\[
\delta^\star
\in
\arg\max_{0<\|\delta\|_p\le\epsilon}
\frac{
\|b+J\delta\|_q-\|b\|_q
}{
\|\delta\|_p
}.
\]
We call such a direction an affine residual gain maximizing direction.

When
\[
b=0,
\]
this reduces to
\[
\delta^\star
\in
\arg\max_{\|\delta\|_p=1}
\|J\delta\|_q,
\]
which is a right $(p,q)$-singular vector of $J$.

When
\[
b\neq0,
\]
the maximizing direction depends on both $J$ and $b$, and generally also on the radius $\epsilon$. Therefore, it should not be called a standard right $(p,q)$-singular vector of $J$. A more accurate name is
\[
\text{affine residual gain maximizing direction}.
\]

If one wants terminology closer to the singular-vector language, one may define it as a right affine $(p,q)$-gain direction, but this should be explicitly introduced as a new definition.

\subsection{Affine Residual Gain Is Not a Standard Matrix Norm}

The standard induced matrix norm is defined for a linear map:
\[
\delta\mapsto J\delta.
\]
The affine residual map
\[
\delta\mapsto b+J\delta
\]
is not linear when $b\neq0$. Therefore, the finite-radius affine residual gain is not a matrix norm in the standard sense.

The issue is not that the norm $\|\cdot\|_q$ loses homogeneity. The norm remains homogeneous:
\[
\|\alpha z\|_q=|\alpha|\|z\|_q.
\]
The issue is that the composite function
\[
\delta\mapsto \|b+J\delta\|_q
\]
is generally not homogeneous in $\delta$:
\[
\|b+J(\alpha\delta)\|_q
=
\|b+\alpha J\delta\|_q
\]
is generally not equal to
\[
|\alpha|\|b+J\delta\|_q.
\]

Similarly, the baseline-subtracted increment
\[
\|b+J\delta\|_q-\|b\|_q
\]
is generally not homogeneous in $\delta$. Therefore,
\[
\frac{
\|b+J\delta\|_q-\|b\|_q
}{
\|\delta\|_p
}
\]
usually depends not only on the direction of $\delta$, but also on its radius.

This is why the affine residual gain must be written as a finite-radius quantity:
\[
\mathcal G_{p\to q}(J,b;\epsilon).
\]

\subsection{Relation to Local Slope and Calmness-Type Quantities}

Define the scalar function
\[
F(\delta)=\|b+J\delta\|_q.
\]
Then
\[
F(0)=\|b\|_q.
\]
The affine residual gain is a finite-radius, one-sided difference quotient:
\[
\mathcal G_{p\to q}(J,b;\epsilon)
=
\sup_{0<\|\delta\|_p\le\epsilon}
\frac{F(\delta)-F(0)}{\|\delta\|_p}.
\]

Thus, it is closely related in spirit to local Lipschitz slopes, calmness moduli, and metric slope concepts from variational analysis. The difference is that the affine residual gain is finite-radius and one-sided: it measures increase of the residual norm rather than absolute variation.

A local version may be defined as
\[
\mathcal G_{p\to q}^{\mathrm{loc}}(J,b)
=
\limsup_{\delta\to0,\ \delta\neq0}
\frac{
\|b+J\delta\|_q-\|b\|_q
}{
\|\delta\|_p
}.
\]

If $q>1$ and $b\neq0$, then $F$ is differentiable at $\delta=0$, and
\[
\nabla_\delta F(0)
=
J^\top
\frac{|b|^{q-2}b}{\|b\|_q^{q-1}}.
\]
Therefore,
\[
\mathcal G_{p\to q}^{\mathrm{loc}}(J,b)
=
\sup_{\|s\|_p\le1}
\left\langle
J^\top
\frac{|b|^{q-2}b}{\|b\|_q^{q-1}},
s
\right\rangle.
\]
By duality of norms,
\[
\mathcal G_{p\to q}^{\mathrm{loc}}(J,b)
=
\left\|
J^\top
\frac{|b|^{q-2}b}{\|b\|_q^{q-1}}
\right\|_{p^\ast},
\]
where $p^\ast$ is the dual exponent of $p$:
\[
\frac{1}{p}+\frac{1}{p^\ast}=1.
\]

Thus, the local affine gain direction is the $L_p$-steepest direction of the scalar function
\[
F(\delta)=\|b+J\delta\|_q
\]
at $\delta=0$.

\subsection{Finite-Radius Versus Fixed-Radius Objectives}

The finite-radius affine residual gain is
\[
\sup_{0<\|\delta\|_p\le\epsilon}
\frac{
\|b+J\delta\|_q-\|b\|_q
}{
\|\delta\|_p
}.
\]
Because the numerator is generally not homogeneous in $\delta$, the optimizer may occur at a radius smaller than $\epsilon$.

If instead one fixes the radius
\[
\|\delta\|_p=\epsilon,
\]
then the denominator becomes constant:
\[
\frac{
\|b+J\delta\|_q-\|b\|_q
}{
\epsilon
}.
\]
On the fixed-radius sphere, maximizing the ratio is equivalent to maximizing
\[
\|b+J\delta\|_q.
\]
Therefore, the ratio objective differs from the original objective only when the radius is allowed to vary:
\[
0<\|\delta\|_p\le\epsilon.
\]

\subsection{Application to the Three Losses}

The first loss is
\[
L_1(\delta)
=
\|f(x+\delta)-f(x)\|_q.
\]
Locally,
\[
f(x+\delta)-f(x)
\approx
J_f\delta.
\]
Therefore,
\[
\frac{L_1(\delta)}{\|\delta\|_p}
\approx
\frac{\|J_f\delta\|_q}{\|\delta\|_p}.
\]
This is the standard induced $p\to q$ operator norm problem for $J_f$:
\[
\|J_f\|_{p\to q}.
\]
Its maximizing direction is a right $(p,q)$-singular vector of $J_f$.

The second loss is
\[
L_2(\delta)
=
\|f(x+\delta)-g(x)\|_q.
\]
Let
\[
b_2=f(x)-g(x).
\]
Then locally,
\[
f(x+\delta)-g(x)
\approx
b_2+J_f\delta.
\]
The increment ratio is
\[
\frac{
L_2(\delta)-L_2(0)
}{
\|\delta\|_p
}
\approx
\frac{
\|b_2+J_f\delta\|_q-\|b_2\|_q
}{
\|\delta\|_p
}.
\]
This is an affine residual gain problem, not a standard matrix norm problem.

The third loss is
\[
L_3(\delta)
=
\|f(x+\delta)-g(x+\delta)\|_q.
\]
Let
\[
b_3=f(x)-g(x),
\]
and
\[
J_3=J_f-J_g.
\]
Then locally,
\[
f(x+\delta)-g(x+\delta)
\approx
b_3+J_3\delta.
\]
The increment ratio is
\[
\frac{
L_3(\delta)-L_3(0)
}{
\|\delta\|_p
}
\approx
\frac{
\|b_3+J_3\delta\|_q-\|b_3\|_q
}{
\|\delta\|_p
}.
\]
This is also an affine residual gain problem.

\subsection{Optimization Consequences}

For $L_1$, the local ratio is the standard induced matrix norm:
\[
\frac{\|J_f\delta\|_q}{\|\delta\|_p}.
\]
Therefore, generalized power iteration is a natural method for estimating the corresponding right $(p,q)$-singular direction.

For $L_2$ and $L_3$, the local ratio has the affine increment form:
\[
\frac{
\|b+J\delta\|_q-\|b\|_q
}{
\|\delta\|_p
}.
\]
This is not a standard matrix norm and does not have a standard singular-vector interpretation. It should be optimized as a scalar objective using first-order methods such as PGD, $L_p$-steepest PGD, Adam-style optimization, or generalized-power-style heuristics.

The most direct practical method is to optimize the true scalar objective
\[
R(\delta)
=
\frac{
L_i(\delta)-L_i(0)
}{
\|\delta\|_p+\eta
}
\]
by automatic differentiation:
\[
g_k=\nabla_\delta R(\delta_k).
\]
Then one may update by raw PGD:
\[
\delta_{k+1}
=
\Pi_{\|\delta\|_p\le\epsilon}
(\delta_k+\alpha g_k),
\]
or by $L_p$-steepest PGD:
\[
s_k=
\arg\max_{\|s\|_p\le1}g_k^\top s,
\]
\[
\delta_{k+1}
=
\Pi_{\|\delta\|_p\le\epsilon}
(\delta_k+\alpha s_k).
\]

A generalized-power-style replacement update may also be used:
\[
\delta_{k+1}
=
\rho_k s_k,
\]
but for affine residual gains it should be interpreted as a heuristic direction-replacement method rather than an exact matrix singular-vector computation.

\subsection{Main Takeaways}

The standard induced $p\to q$ norm
\[
\|J\|_{p\to q}
\]
is a matrix/operator norm because it is defined for the linear map
\[
\delta\mapsto J\delta.
\]

Its maximizing direction
\[
\delta^\star
\in
\arg\max_{\|\delta\|_p=1}\|J\delta\|_q
\]
can be called a right $(p,q)$-singular vector.

For $p=q=2$, this reduces to the standard largest right singular vector.

For general $p,q$, the right $(p,q)$-singular vector is defined through a norm maximization problem and typically must be approximated by iterative optimization, such as generalized power iteration.

For affine residual maps
\[
\delta\mapsto b+J\delta,
\]
the meaningful gain is the baseline-subtracted quantity
\[
\frac{
\|b+J\delta\|_q-\|b\|_q
}{
\|\delta\|_p
}.
\]
This is not a standard matrix norm, because it depends on the baseline residual $b$ and generally on the radius $\epsilon$.

A natural extension is the finite-radius affine residual gain:
\[
\mathcal G_{p\to q}(J,b;\epsilon)
=
\sup_{0<\|\delta\|_p\le\epsilon}
\frac{
\|b+J\delta\|_q-\|b\|_q
}{
\|\delta\|_p
}.
\]

Its maximizer is an affine residual gain maximizing direction. When $b=0$, this reduces to a standard right $(p,q)$-singular vector. When $b\neq0$, it is a new affine gain direction rather than a standard matrix singular vector.

\section{Zero-Initialization Degeneracy of Loss-1 Objectives}

For the first attack loss, the residual is defined as
\[
r_1(\delta) = f(x+\delta)-f(x),
\]
and the corresponding loss is
\[
L_1(\delta)=\|r_1(\delta)\|_q.
\]
This loss measures the output variation of the surrogate model itself under an input perturbation. In particular, it does not involve the ground-truth solver or any external target.

When
\[
\delta = 0,
\]
we have
\[
r_1(0)=f(x)-f(x)=0,
\]
and therefore
\[
L_1(0)=0.
\]

For \(r_1(\delta)\neq 0\), the gradient of \(L_1\) with respect to \(\delta\) is
\[
\nabla_\delta L_1(\delta)
=
J_f(x+\delta)^\top
\frac{
|r_1(\delta)|^{q-2}r_1(\delta)
}{
\|r_1(\delta)\|_q^{q-1}
},
\]
where
\[
J_f(x+\delta)=\frac{\partial f(x+\delta)}{\partial \delta}.
\]
However, at
\[
r_1(\delta)=0,
\]
the norm is generally non-smooth in the strict mathematical sense. In practical PyTorch/autograd computation, the gradient selected at this point is the zero subgradient:
\[
\nabla_{r_1}\|r_1\|_q\big|_{r_1=0}=0.
\]
As a result,
\[
\nabla_\delta L_1(0)
=
J_f(x)^\top 0
=
0.
\]

This creates a zero-initialization degeneracy for all three Loss-1 objectives considered below.

\subsection{Loss-1 Original Objective}

The original objective is
\[
O_{\mathrm{orig}}(\delta)
=
L_1(\delta)
=
\|f(x+\delta)-f(x)\|_q.
\]
Since
\[
\nabla_\delta L_1(0)=0,
\]
we immediately obtain
\[
\nabla_\delta O_{\mathrm{orig}}(0)=0.
\]

Therefore, if the optimization is initialized with
\[
\delta_0=0,
\]
then the first gradient is
\[
g_0=\nabla_\delta O_{\mathrm{orig}}(\delta_0)=0.
\]
Any first-order update rule based on this gradient will fail to move away from the origin.

\subsection{Loss-1 Increment-Ratio Objective}

The increment-ratio objective is defined as
\[
O_{\mathrm{inc}}(\delta)
=
\frac{
L_1(\delta)-L_1(0)
}{
\|\delta\|_p+\eta
},
\]
where \(\eta>0\) is a small numerical stabilizer.

Since
\[
L_1(0)=0,
\]
this objective becomes
\[
O_{\mathrm{inc}}(\delta)
=
\frac{
L_1(\delta)
}{
\|\delta\|_p+\eta
}.
\]

Its gradient is
\[
\nabla_\delta O_{\mathrm{inc}}(\delta)
=
\frac{
(\|\delta\|_p+\eta)\nabla_\delta L_1(\delta)
-
L_1(\delta)\nabla_\delta \|\delta\|_p
}{
(\|\delta\|_p+\eta)^2
}.
\]

At
\[
\delta=0,
\]
we have
\[
L_1(0)=0,
\]
\[
\nabla_\delta L_1(0)=0,
\]
and, under PyTorch/autograd,
\[
\nabla_\delta \|\delta\|_p\big|_{\delta=0}=0.
\]

Therefore,
\[
\nabla_\delta O_{\mathrm{inc}}(0)
=
\frac{
\eta\cdot 0 - 0\cdot 0
}{
\eta^2
}
=
0.
\]

Thus, the increment-ratio objective also produces a zero gradient when initialized exactly at \(\delta=0\).

\subsection{Loss-1 Regularized Objective}

The regularized objective is defined as
\[
O_{\mathrm{reg}}(\delta)
=
L_1(\delta)-c\|\delta\|_p,
\]
where \(c>0\) controls the strength of the perturbation penalty.

The gradient is
\[
\nabla_\delta O_{\mathrm{reg}}(\delta)
=
\nabla_\delta L_1(\delta)
-
c\nabla_\delta \|\delta\|_p.
\]

At
\[
\delta=0,
\]
we have
\[
\nabla_\delta L_1(0)=0,
\]
and PyTorch/autograd also gives
\[
\nabla_\delta \|\delta\|_p\big|_{\delta=0}=0.
\]

Therefore,
\[
\nabla_\delta O_{\mathrm{reg}}(0)
=
0-c\cdot 0
=
0.
\]

Hence, the regularized objective is also stuck at the origin if the perturbation is initialized exactly as zero.

\subsection{Consequence for PGD, \(L_p\)-Steepest PGD, and GPI}

All three optimization methods considered here depend on the gradient
\[
g_k=\nabla_\delta O(\delta_k).
\]

For standard PGD, the update has the form
\[
\delta_{k+1}
=
\Pi\left(\delta_k+\alpha g_k\right),
\]
where \(\Pi\) denotes projection onto the feasible perturbation set.

For \(L_p\)-steepest PGD, the update direction is obtained from
\[
s_k
=
\arg\max_{\|s\|_p\le 1} g_k^\top s.
\]

For generalized power iteration, the direction is similarly determined by
\[
s_k
=
\arg\max_{\|s\|_p\le 1} g_k^\top s.
\]

Therefore, if
\[
g_k=0,
\]
then standard PGD has no nonzero update direction, \(L_p\)-steepest PGD has no meaningful steepest direction, and GPI also has no effective direction to follow.

Combining the three Loss-1 objectives, we have
\[
\nabla_\delta O_{\mathrm{orig}}(0)=0,
\]
\[
\nabla_\delta O_{\mathrm{inc}}(0)=0,
\]
and
\[
\nabla_\delta O_{\mathrm{reg}}(0)=0.
\]

Thus, for Loss-1, if the initialization is exactly
\[
\delta_0=0,
\]
then all three objective definitions produce
\[
g_0=\nabla_\delta O(\delta_0)=0
\]
in PyTorch/autograd. Consequently, all nine combinations formed by the three Loss-1 objectives and the three gradient-based optimization methods are expected to get stuck at initialization.

\subsection{Practical Implication: Nonzero Random Start}

To avoid this degeneracy, Loss-1 attacks should use a small nonzero random initialization:
\[
\delta_0\neq 0.
\]
A typical choice is to sample \(\delta_0\) from a small random distribution and then project it into the allowed perturbation set:
\[
\delta_0
=
\Pi\left(\epsilon_{\mathrm{init}}\xi\right),
\]
where \(\xi\) is a random tensor and \(\epsilon_{\mathrm{init}}\) is a small initialization scale.

This small random start breaks the exact cancellation
\[
f(x+\delta_0)-f(x)=0,
\]
so that
\[
r_1(\delta_0)\neq 0
\]
in general. Once the residual becomes nonzero, PyTorch/autograd can produce a nonzero gradient, allowing PGD, \(L_p\)-steepest PGD, and GPI to start the optimization process.

In summary, for Loss-1 objectives, the zero initialization
\[
\delta_0=0
\]
is numerically degenerate under PyTorch/autograd. A small nonzero random start is necessary for all three objectives:
\[
O_{\mathrm{orig}}(\delta),
\qquad
O_{\mathrm{inc}}(\delta),
\qquad
O_{\mathrm{reg}}(\delta).
\]
This applies equally to standard PGD, \(L_p\)-steepest PGD, and GPI.

\section{Local Lipschitz Continuity, Induced Norms, and Jacobian Alignment}

Let
\[
f:\mathcal X\subseteq \mathbb R^n\to \mathbb R^m
\]
be the learned regression model, and let
\[
j:\mathcal X\subseteq \mathbb R^n\to \mathbb R^m
\]
be the ground-truth function, oracle, or numerical solver. Define the error field by
\[
e(x)=f(x)-j(x).
\]
We measure input perturbations using the \(p\)-norm and output errors using the \(q\)-norm.

For a fixed input \(x\) and perturbation radius \(\varepsilon>0\), define the local perturbation set
\[
\Delta_p(x,\varepsilon)
=
\{\delta:\|\delta\|_p\le \varepsilon,\ x+\delta\in \mathcal X\},
\]
and the corresponding local input neighborhood
\[
B_p(x,\varepsilon)
=
\{z\in\mathcal X:\|z-x\|_p\le \varepsilon\}.
\]
Equivalently,
\[
x+\Delta_p(x,\varepsilon)=B_p(x,\varepsilon).
\]

\subsection{Local Lipschitz Constants of \(f\), \(j\), and the Error Field}

The local Lipschitz constant of \(f\) on \(B_p(x,\varepsilon)\) is
\[
L_f^{\mathrm{loc}}(x,\varepsilon)
=
\sup_{\substack{z_1,z_2\in B_p(x,\varepsilon)\\ z_1\neq z_2}}
\frac{
\|f(z_1)-f(z_2)\|_q
}{
\|z_1-z_2\|_p
}.
\]
If \(f\) is differentiable, then this local Lipschitz constant is controlled by the induced norm of the Jacobian:
\[
L_f^{\mathrm{loc}}(x,\varepsilon)
\le
\sup_{z\in B_p(x,\varepsilon)}
\|J_f(z)\|_{p\to q},
\]
where
\[
\|J_f(z)\|_{p\to q}
=
\sup_{v\neq 0}
\frac{\|J_f(z)v\|_q}{\|v\|_p}.
\]
Thus, under a local linearization at \(x\),
\[
L_f^{\mathrm{lin}}(x)
=
\|J_f(x)\|_{p\to q}.
\]

Similarly, the local Lipschitz constant of \(j\) is
\[
L_j^{\mathrm{loc}}(x,\varepsilon)
=
\sup_{\substack{z_1,z_2\in B_p(x,\varepsilon)\\ z_1\neq z_2}}
\frac{
\|j(z_1)-j(z_2)\|_q
}{
\|z_1-z_2\|_p
}.
\]
If \(j\) is differentiable, then
\[
L_j^{\mathrm{loc}}(x,\varepsilon)
\le
\sup_{z\in B_p(x,\varepsilon)}
\|J_j(z)\|_{p\to q},
\]
where
\[
\|J_j(z)\|_{p\to q}
=
\sup_{v\neq 0}
\frac{\|J_j(z)v\|_q}{\|v\|_p}.
\]
Under a local linearization at \(x\),
\[
L_j^{\mathrm{lin}}(x)
=
\|J_j(x)\|_{p\to q}.
\]

For the error field
\[
e(x)=f(x)-j(x),
\]
the local Lipschitz constant is
\[
L_e^{\mathrm{loc}}(x,\varepsilon)
=
\sup_{\substack{z_1,z_2\in B_p(x,\varepsilon)\\ z_1\neq z_2}}
\frac{
\|e(z_1)-e(z_2)\|_q
}{
\|z_1-z_2\|_p
}.
\]
Substituting \(e=f-j\), we obtain
\[
L_e^{\mathrm{loc}}(x,\varepsilon)
=
\sup_{\substack{z_1,z_2\in B_p(x,\varepsilon)\\ z_1\neq z_2}}
\frac{
\|[f(z_1)-j(z_1)]-[f(z_2)-j(z_2)]\|_q
}{
\|z_1-z_2\|_p
}.
\]
Equivalently,
\[
L_e^{\mathrm{loc}}(x,\varepsilon)
=
\sup_{\substack{z_1,z_2\in B_p(x,\varepsilon)\\ z_1\neq z_2}}
\frac{
\|[f(z_1)-f(z_2)]-[j(z_1)-j(z_2)]\|_q
}{
\|z_1-z_2\|_p
}.
\]

If \(f\) and \(j\) are differentiable, then
\[
J_e(z)=J_f(z)-J_j(z).
\]
Therefore,
\[
L_e^{\mathrm{loc}}(x,\varepsilon)
\le
\sup_{z\in B_p(x,\varepsilon)}
\|J_e(z)\|_{p\to q}
=
\sup_{z\in B_p(x,\varepsilon)}
\|J_f(z)-J_j(z)\|_{p\to q}.
\]
The induced norm appearing here is
\[
\|J_f(z)-J_j(z)\|_{p\to q}
=
\sup_{v\neq 0}
\frac{
\|(J_f(z)-J_j(z))v\|_q
}{
\|v\|_p
}.
\]
Thus, under a local linearization at \(x\),
\[
L_e^{\mathrm{lin}}(x)
=
\|J_f(x)-J_j(x)\|_{p\to q}.
\]

\subsection{Perturbation-Based Interpretation}

For a small perturbation \(\delta\), the first-order approximations are
\[
f(x+\delta)-f(x)\approx J_f(x)\delta,
\]
\[
j(x+\delta)-j(x)\approx J_j(x)\delta,
\]
and
\[
e(x+\delta)-e(x)
\approx
[J_f(x)-J_j(x)]\delta.
\]
Therefore,
\[
L_f^{\mathrm{lin}}(x)
=
\sup_{\delta\neq 0}
\frac{
\|J_f(x)\delta\|_q
}{
\|\delta\|_p
}
=
\|J_f(x)\|_{p\to q},
\]
\[
L_j^{\mathrm{lin}}(x)
=
\sup_{\delta\neq 0}
\frac{
\|J_j(x)\delta\|_q
}{
\|\delta\|_p
}
=
\|J_j(x)\|_{p\to q},
\]
and
\[
L_e^{\mathrm{lin}}(x)
=
\sup_{\delta\neq 0}
\frac{
\|[J_f(x)-J_j(x)]\delta\|_q
}{
\|\delta\|_p
}
=
\|J_f(x)-J_j(x)\|_{p\to q}.
\]
Hence, \(L_f^{\mathrm{lin}}\) measures the local sensitivity of the model itself, \(L_j^{\mathrm{lin}}\) measures the local sensitivity of the ground-truth function itself, and \(L_e^{\mathrm{lin}}\) measures the mismatch between the local linear response of the model and the local linear response of the ground truth.

\subsection{Norm Bounds Among the Three Quantities}

Let
\[
A=J_f(x),
\qquad
B=J_j(x).
\]
Then
\[
L_f^{\mathrm{lin}}(x)=\|A\|_{p\to q},
\]
\[
L_j^{\mathrm{lin}}(x)=\|B\|_{p\to q},
\]
and
\[
L_e^{\mathrm{lin}}(x)=\|A-B\|_{p\to q}.
\]
By the triangle inequality for induced norms,
\[
\|A-B\|_{p\to q}
\le
\|A\|_{p\to q}
+
\|B\|_{p\to q}.
\]
Therefore,
\[
L_e^{\mathrm{lin}}(x)
\le
L_f^{\mathrm{lin}}(x)
+
L_j^{\mathrm{lin}}(x).
\]
By the reverse triangle inequality,
\[
\|A-B\|_{p\to q}
\ge
\left|
\|A\|_{p\to q}
-
\|B\|_{p\to q}
\right|.
\]
Thus,
\[
\left|
L_f^{\mathrm{lin}}(x)
-
L_j^{\mathrm{lin}}(x)
\right|
\le
L_e^{\mathrm{lin}}(x)
\le
L_f^{\mathrm{lin}}(x)
+
L_j^{\mathrm{lin}}(x).
\]
Equivalently,
\[
\boxed{
\left|
\|J_f(x)\|_{p\to q}
-
\|J_j(x)\|_{p\to q}
\right|
\le
\|J_f(x)-J_j(x)\|_{p\to q}
\le
\|J_f(x)\|_{p\to q}
+
\|J_j(x)\|_{p\to q}
}.
\]

On a local region \(B_p(x,\varepsilon)\), define
\[
M_f(x,\varepsilon)
=
\sup_{z\in B_p(x,\varepsilon)}
\|J_f(z)\|_{p\to q},
\]
\[
M_j(x,\varepsilon)
=
\sup_{z\in B_p(x,\varepsilon)}
\|J_j(z)\|_{p\to q},
\]
and
\[
M_e(x,\varepsilon)
=
\sup_{z\in B_p(x,\varepsilon)}
\|J_f(z)-J_j(z)\|_{p\to q}.
\]
Then
\[
M_e(x,\varepsilon)
\le
M_f(x,\varepsilon)+M_j(x,\varepsilon),
\]
and
\[
L_e^{\mathrm{loc}}(x,\varepsilon)
\le
M_e(x,\varepsilon)
\le
M_f(x,\varepsilon)+M_j(x,\varepsilon).
\]

\subsection{Directional Alignment Between \(J_f\) and \(J_j\)}

The bounds
\[
\left|
L_f^{\mathrm{lin}}(x)
-
L_j^{\mathrm{lin}}(x)
\right|
\le
L_e^{\mathrm{lin}}(x)
\le
L_f^{\mathrm{lin}}(x)+L_j^{\mathrm{lin}}(x)
\]
only give a coarse range for the error-field Lipschitz constant. The actual value of
\[
L_e^{\mathrm{lin}}(x)
=
\|J_f(x)-J_j(x)\|_{p\to q}
\]
depends strongly on the directional alignment between \(J_f(x)\) and \(J_j(x)\).

If
\[
J_f(x)\approx J_j(x),
\]
then
\[
J_f(x)-J_j(x)\approx 0,
\]
and therefore
\[
L_e^{\mathrm{lin}}(x)\approx 0.
\]
In this case, even if \(L_f^{\mathrm{lin}}(x)\) and \(L_j^{\mathrm{lin}}(x)\) are both large, the error field can still be locally stable because the model and the ground-truth function change in almost the same way.

If
\[
J_f(x)\approx -J_j(x),
\]
then
\[
J_f(x)-J_j(x)\approx -2J_j(x),
\]
and therefore
\[
L_e^{\mathrm{lin}}(x)
\]
can be large, possibly close to the upper bound
\[
L_f^{\mathrm{lin}}(x)+L_j^{\mathrm{lin}}(x).
\]
This means that if the model changes in a direction opposite to the ground truth, the error field becomes highly sensitive to input perturbations.

In one dimension, this relationship is especially clear. Suppose
\[
J_f(x)=a,
\qquad
J_j(x)=b.
\]
Then
\[
L_f^{\mathrm{lin}}=|a|,
\qquad
L_j^{\mathrm{lin}}=|b|,
\qquad
L_e^{\mathrm{lin}}=|a-b|.
\]
If \(a\) and \(b\) have the same sign, then the two local responses have the same direction and
\[
L_e^{\mathrm{lin}}
\]
is close to the lower bound
\[
\left|
L_f^{\mathrm{lin}}-L_j^{\mathrm{lin}}
\right|.
\]
If \(a\) and \(b\) have opposite signs, then the two local responses point in opposite directions and
\[
L_e^{\mathrm{lin}}
=
L_f^{\mathrm{lin}}+L_j^{\mathrm{lin}}.
\]
Thus, in one dimension,
\[
\text{same direction}
\quad\Longrightarrow\quad
L_e^{\mathrm{lin}} \text{ is small},
\]
while
\[
\text{opposite direction}
\quad\Longrightarrow\quad
L_e^{\mathrm{lin}} \text{ is large}.
\]

\subsection{Worst-Case Direction in the Multidimensional Case}

In multiple dimensions, the induced norm
\[
\|J_f(x)-J_j(x)\|_{p\to q}
\]
is determined by the worst perturbation direction:
\[
\|J_f(x)-J_j(x)\|_{p\to q}
=
\sup_{\delta\neq 0}
\frac{
\|(J_f(x)-J_j(x))\delta\|_q
}{
\|\delta\|_p
}.
\]
Since
\[
(J_f(x)-J_j(x))\delta
=
J_f(x)\delta-J_j(x)\delta,
\]
the error-field Lipschitz constant measures whether the two output responses
\[
J_f(x)\delta
\quad\text{and}\quad
J_j(x)\delta
\]
are aligned for the worst input perturbation direction \(\delta\).

For a fixed direction \(\delta\), define
\[
u=J_f(x)\delta,
\qquad
v=J_j(x)\delta.
\]
When \(q=2\), the squared difference is
\[
\|u-v\|_2^2
=
\|u\|_2^2+\|v\|_2^2-2\langle u,v\rangle.
\]
If
\[
\cos\theta_\delta
=
\frac{
\langle J_f(x)\delta,J_j(x)\delta\rangle
}{
\|J_f(x)\delta\|_2
\|J_j(x)\delta\|_2
},
\]
then
\[
\|J_f(x)\delta-J_j(x)\delta\|_2^2
=
\|J_f(x)\delta\|_2^2
+
\|J_j(x)\delta\|_2^2
-
2
\|J_f(x)\delta\|_2
\|J_j(x)\delta\|_2
\cos\theta_\delta.
\]
Therefore, if
\[
\cos\theta_\delta\approx 1,
\]
then \(J_f(x)\delta\) and \(J_j(x)\delta\) point in nearly the same output direction, and the error response is small. If
\[
\cos\theta_\delta\approx -1,
\]
then the two responses point in opposite directions, and the error response is large. If
\[
\cos\theta_\delta\approx 0,
\]
then the two responses are nearly orthogonal, and the error response is still significant:
\[
\|J_f(x)\delta-J_j(x)\delta\|_2^2
\approx
\|J_f(x)\delta\|_2^2
+
\|J_j(x)\delta\|_2^2.
\]

Thus, in the multidimensional case, \(L_e^{\mathrm{lin}}\) is small only when the model and the ground-truth function have similar responses along the important perturbation directions. Since the induced norm takes a supremum, a single bad perturbation direction can make
\[
L_e^{\mathrm{lin}}(x)
=
\|J_f(x)-J_j(x)\|_{p\to q}
\]
large.

\subsection{Relation to Matrix Similarity}

A global matrix similarity measure can also be defined using the Frobenius inner product:
\[
\langle A,B\rangle_F
=
\operatorname{tr}(A^\top B),
\]
with the matrix cosine similarity
\[
\cos(A,B)
=
\frac{
\langle A,B\rangle_F
}{
\|A\|_F\|B\|_F
}.
\]
Under the Frobenius norm,
\[
\|A-B\|_F^2
=
\|A\|_F^2+\|B\|_F^2
-
2\langle A,B\rangle_F,
\]
or equivalently,
\[
\|A-B\|_F^2
=
\|A\|_F^2+\|B\|_F^2
-
2\|A\|_F\|B\|_F\cos(A,B).
\]
This shows that if the Frobenius cosine similarity is close to \(1\), then \(\|A-B\|_F\) is small. If the similarity is close to \(-1\), then \(\|A-B\|_F\) is large.

However, the quantity directly controlling adversarial error amplification is not the Frobenius norm but the induced norm
\[
\|A-B\|_{p\to q}.
\]
In the special case \(p=q=2\),
\[
\|A-B\|_{2\to 2}
=
\sigma_{\max}(A-B),
\]
so
\[
L_e^{\mathrm{lin}}(x)
=
\sigma_{\max}(J_f(x)-J_j(x)).
\]
This quantity depends on the largest singular direction of the Jacobian mismatch. Therefore, even if two matrices look globally similar under a Frobenius similarity score, the adversarially relevant error-field Lipschitz constant may still be large if there exists a perturbation direction where \(J_f(x)\delta\) and \(J_j(x)\delta\) differ strongly.

\subsection{Interpretation for Regression Robustness}

For regression robustness with a moving ground truth \(j(x+\delta)\), the relevant adversarial error is
\[
\|f(x+\delta)-j(x+\delta)\|_q
=
\|e(x+\delta)\|_q.
\]
The local growth of this error is controlled by
\[
L_e^{\mathrm{lin}}(x)
=
\|J_f(x)-J_j(x)\|_{p\to q}.
\]
Thus, robustness is not determined only by whether \(f\) is smooth, nor only by whether \(j\) is smooth. It is determined by whether the local response of \(f\) matches the local response of \(j\). A small value of
\[
\|J_f(x)-J_j(x)\|_{p\to q}
\]
means that the model follows the ground-truth function under small perturbations. A large value means that there exists an input perturbation direction along which the model and the ground truth move differently, causing the prediction error to grow rapidly.

\section{Triangle-Interval Position and Directional Alignment of Local Jacobians}

Let
\[
f:\mathcal X\subseteq \mathbb R^n\to \mathbb R^m
\]
be the learned model, and let
\[
g:\mathcal X\subseteq \mathbb R^n\to \mathbb R^m
\]
be the ground-truth solver or oracle. At a fixed input \(x\), denote their local Jacobians by
\[
J_f(x)
\qquad\text{and}\qquad
J_g(x).
\]
The local error field is
\[
e(x)=f(x)-g(x),
\]
and its Jacobian is
\[
J_e(x)=J_f(x)-J_g(x).
\]

Using the \(p\)-norm for input perturbations and the \(q\)-norm for output responses, define the three local induced norms
\[
L_f(x)=\|J_f(x)\|_{p\to q},
\]
\[
L_g(x)=\|J_g(x)\|_{p\to q},
\]
and
\[
L_e(x)=\|J_f(x)-J_g(x)\|_{p\to q}.
\]
Here
\[
\|A\|_{p\to q}
=
\sup_{\|\delta\|_p=1}
\|A\delta\|_q.
\]

The quantity \(L_f(x)\) measures the local sensitivity of the model \(f\), \(L_g(x)\) measures the local sensitivity of the solver \(g\), and \(L_e(x)\) measures the local mismatch between the model response and the solver response. In particular,
\[
L_e(x)
=
\sup_{\|\delta\|_p=1}
\|(J_f(x)-J_g(x))\delta\|_q
=
\sup_{\|\delta\|_p=1}
\|J_f(x)\delta-J_g(x)\delta\|_q.
\]

\subsection{Triangle Inequality Bounds}

By the triangle inequality and the reverse triangle inequality for induced norms, we have
\[
\left|
\|J_f(x)\|_{p\to q}
-
\|J_g(x)\|_{p\to q}
\right|
\le
\|J_f(x)-J_g(x)\|_{p\to q}
\le
\|J_f(x)\|_{p\to q}
+
\|J_g(x)\|_{p\to q}.
\]
Equivalently,
\[
\boxed{
|L_f(x)-L_g(x)|
\le
L_e(x)
\le
L_f(x)+L_g(x)
}.
\]

Define the lower and upper triangle bounds by
\[
L_{\mathrm{low}}(x)
=
|L_f(x)-L_g(x)|,
\]
and
\[
L_{\mathrm{up}}(x)
=
L_f(x)+L_g(x).
\]
Then
\[
L_e(x)\in
\left[
L_{\mathrm{low}}(x),
L_{\mathrm{up}}(x)
\right].
\]

\subsection{Normalized Triangle-Interval Position}

To quantify where \(L_e(x)\) lies inside the triangle-inequality interval, define the normalized triangle-position index
\[
\rho_{\triangle}(x)
=
\frac{
L_e(x)-L_{\mathrm{low}}(x)
}{
L_{\mathrm{up}}(x)-L_{\mathrm{low}}(x)
}.
\]
Equivalently,
\[
\rho_{\triangle}(x)
=
\frac{
L_e(x)-|L_f(x)-L_g(x)|
}{
L_f(x)+L_g(x)-|L_f(x)-L_g(x)|
}.
\]

When the denominator is nonzero, this quantity satisfies
\[
0\le \rho_{\triangle}(x)\le 1.
\]
Since
\[
L_f(x)+L_g(x)-|L_f(x)-L_g(x)|
=
2\min\{L_f(x),L_g(x)\},
\]
the index degenerates when
\[
\min\{L_f(x),L_g(x)\}=0.
\]
In that case, the normalized position should be treated separately.

The interpretation is as follows:
\[
\rho_{\triangle}(x)\approx 0
\]
means that \(L_e(x)\) is close to the lower bound. This suggests that \(J_f(x)\) and \(J_g(x)\) have locally aligned responses, or at least that their worst-case induced-norm mismatch is close to the smallest possible value allowed by their individual norms.

On the other hand,
\[
\rho_{\triangle}(x)\approx 1
\]
means that \(L_e(x)\) is close to the upper bound. This suggests that the local responses of \(J_f(x)\) and \(J_g(x)\) are strongly mismatched, possibly close to opposite directions in the worst-case perturbation direction.

If
\[
\rho_{\triangle}(x)\approx 0.5,
\]
then \(L_e(x)\) lies in the middle of the triangle interval. In this case, the scalar position alone is not enough to determine the detailed geometry of the mismatch. One should inspect directional alignment more directly.

\subsection{Worst-Case Perturbation Direction}

The induced norm \(L_e(x)\) is a worst-case quantity. It is determined by the perturbation direction that maximizes the discrepancy between the local model response and the local solver response:
\[
\delta_e^\star
=
\arg\max_{\|\delta\|_p=1}
\|(J_f(x)-J_g(x))\delta\|_q.
\]
Equivalently,
\[
\delta_e^\star
=
\arg\max_{\|\delta\|_p=1}
\|J_f(x)\delta-J_g(x)\delta\|_q.
\]

After obtaining \(\delta_e^\star\), define the two output-space responses
\[
u_f
=
J_f(x)\delta_e^\star,
\]
and
\[
u_g
=
J_g(x)\delta_e^\star.
\]
Then
\[
(J_f(x)-J_g(x))\delta_e^\star
=
u_f-u_g.
\]

When \(q=2\), a direct measure of directional alignment is the cosine similarity
\[
\cos\theta_e(x)
=
\frac{
\langle J_f(x)\delta_e^\star,\ J_g(x)\delta_e^\star\rangle
}{
\|J_f(x)\delta_e^\star\|_2
\|J_g(x)\delta_e^\star\|_2
}.
\]
Equivalently,
\[
\cos\theta_e(x)
=
\frac{
\langle u_f,u_g\rangle
}{
\|u_f\|_2\|u_g\|_2
}.
\]

If
\[
\cos\theta_e(x)\approx 1,
\]
then the model and solver responses in the worst mismatch direction are nearly aligned. If
\[
\cos\theta_e(x)\approx -1,
\]
then the two responses are nearly opposite. If
\[
\cos\theta_e(x)\approx 0,
\]
then the two responses are nearly orthogonal.

For \(q=2\), the relation
\[
\|u_f-u_g\|_2^2
=
\|u_f\|_2^2+\|u_g\|_2^2
-
2\langle u_f,u_g\rangle
\]
shows explicitly how the magnitude of the mismatch depends on the response angle. Using the cosine similarity, this becomes
\[
\|u_f-u_g\|_2^2
=
\|u_f\|_2^2+\|u_g\|_2^2
-
2\|u_f\|_2\|u_g\|_2\cos\theta_e(x).
\]
Thus, a large positive cosine reduces the mismatch, while a negative cosine increases it.

\subsection{Random-Direction and Principal-Direction Diagnostics}

Although \(L_e(x)\) and \(\rho_{\triangle}(x)\) provide useful worst-case summaries, they do not prove that \(J_f(x)\) and \(J_g(x)\) are aligned in all directions. The induced norm only captures the largest mismatch:
\[
L_e(x)
=
\sup_{\|\delta\|_p=1}
\|(J_f(x)-J_g(x))\delta\|_q.
\]
Therefore, \(L_e(x)\) being close to the lower bound does not necessarily imply that the model and solver responses are similar in every perturbation direction.

To study the broader directional structure, one can sample multiple directions
\[
\delta_i,
\qquad
\|\delta_i\|_p=1,
\qquad
i=1,\dots,N,
\]
and compute the directional mismatch ratios
\[
r_i(x)
=
\|J_f(x)\delta_i-J_g(x)\delta_i\|_q.
\]
Since \(\|\delta_i\|_p=1\), each \(r_i(x)\) measures the local response mismatch along direction \(\delta_i\). More generally, without normalization,
\[
r_i(x)
=
\frac{
\|J_f(x)\delta_i-J_g(x)\delta_i\|_q
}{
\|\delta_i\|_p
}.
\]

When \(q=2\), one can also record the directional cosine similarities
\[
c_i(x)
=
\frac{
\langle J_f(x)\delta_i,\ J_g(x)\delta_i\rangle
}{
\|J_f(x)\delta_i\|_2
\|J_g(x)\delta_i\|_2
}.
\]

These statistics distinguish several different notions of mismatch:
\[
L_e(x)
\]
measures the worst-case local mismatch,
\[
\frac{1}{N}\sum_{i=1}^N r_i(x)
\]
estimates an average-direction mismatch, and
\[
c_i(x)
\]
measures whether the model and solver responses are aligned, orthogonal, or opposite along sampled directions.

\subsection{Recommended Experimental Records}

For each sample \(x\), the following quantities should be recorded:
\[
L_f(x)=\|J_f(x)\|_{p\to q},
\]
\[
L_g(x)=\|J_g(x)\|_{p\to q},
\]
\[
L_e(x)=\|J_f(x)-J_g(x)\|_{p\to q}.
\]

Then compute the triangle interval:
\[
L_{\mathrm{low}}(x)=|L_f(x)-L_g(x)|,
\]
\[
L_{\mathrm{up}}(x)=L_f(x)+L_g(x),
\]
and the normalized position:
\[
\rho_{\triangle}(x)
=
\frac{
L_e(x)-L_{\mathrm{low}}(x)
}{
L_{\mathrm{up}}(x)-L_{\mathrm{low}}(x)
}.
\]

When \(p=q=2\), one should additionally record the worst-direction cosine similarity:
\[
\cos\theta_e(x)
=
\frac{
\langle J_f(x)\delta_e^\star,\ J_g(x)\delta_e^\star\rangle
}{
\|J_f(x)\delta_e^\star\|_2
\|J_g(x)\delta_e^\star\|_2
}.
\]

For a more complete directional picture, one may also record sampled-direction statistics:
\[
r_i(x)
=
\frac{
\|J_f(x)\delta_i-J_g(x)\delta_i\|_q
}{
\|\delta_i\|_p
},
\]
and, when \(q=2\),
\[
c_i(x)
=
\frac{
\langle J_f(x)\delta_i,\ J_g(x)\delta_i\rangle
}{
\|J_f(x)\delta_i\|_2
\|J_g(x)\delta_i\|_2
}.
\]

\subsection{Interpretation}

The value of \(L_e(x)\) measures how quickly the error field
\[
e(x)=f(x)-g(x)
\]
can change under the worst local input perturbation. A small \(L_e(x)\) means that the local error field changes slowly, while a large \(L_e(x)\) means that the model and solver have a large local response mismatch.

The value of \(\rho_{\triangle}(x)\) indicates where the error-field norm lies inside the triangle-inequality interval. If \(\rho_{\triangle}(x)\) is close to \(0\), then the mismatch is close to the smallest possible mismatch allowed by the individual sensitivities of \(f\) and \(g\). This suggests stronger local response agreement. If \(\rho_{\triangle}(x)\) is close to \(1\), then the mismatch is close to the largest possible mismatch allowed by the triangle inequality. This suggests stronger local response disagreement, possibly including opposite response directions.

However, \(\rho_{\triangle}(x)\) is still a scalar summary. It should not be interpreted as a complete proof that \(J_f(x)\) and \(J_g(x)\) are aligned or misaligned in every direction. For a more precise geometric diagnosis, one should inspect the worst-case direction \(\delta_e^\star\) and compute the response similarity between
\[
J_f(x)\delta_e^\star
\qquad\text{and}\qquad
J_g(x)\delta_e^\star.
\]
One may also sample additional perturbation directions and examine the distribution of \(r_i(x)\) and \(c_i(x)\).

In summary, comparing
\[
\|J_f(x)\|_{p\to q},
\qquad
\|J_g(x)\|_{p\to q},
\qquad
\|J_f(x)-J_g(x)\|_{p\to q}
\]
reveals the local sensitivity and mismatch structure of the model and solver. The position of
\[
\|J_f(x)-J_g(x)\|_{p\to q}
\]
inside the interval
\[
\left[
\left|\|J_f(x)\|_{p\to q}-\|J_g(x)\|_{p\to q}\right|,
\;
\|J_f(x)\|_{p\to q}+\|J_g(x)\|_{p\to q}
\right]
\]
provides a useful summary of local response agreement. Values close to the lower bound suggest more aligned responses, while values close to the upper bound suggest stronger mismatch or opposite responses. For a concrete directional interpretation, the cosine similarity between \(J_f(x)\delta\) and \(J_g(x)\delta\) should be examined, especially along the worst-case direction \(\delta_e^\star\).

\section{Local SVD Interpretation of the Difference Operator}

Let \(f\) and \(g\) be two differentiable operators, and define their difference operator as
\[
h(x)=f(x)-g(x).
\]
Around a fixed point \(x_0\), the local behavior of these operators is described by their Jacobians:
\[
f(x_0+\delta)-f(x_0)\approx J_f(x_0)\delta,
\]
\[
g(x_0+\delta)-g(x_0)\approx J_g(x_0)\delta.
\]
Therefore, the local variation of the difference operator is
\[
h(x_0+\delta)-h(x_0)
\approx
\bigl(J_f(x_0)-J_g(x_0)\bigr)\delta.
\]
Thus, the SVD of the local difference operator is not obtained by subtracting the singular values of \(J_f\) and \(J_g\). Instead, one must first form the Jacobian difference
\[
J_h(x_0)=J_f(x_0)-J_g(x_0),
\]
and then perform SVD on this new matrix.

This distinction is important because singular values depend not only on the magnitudes of local stretching, but also on the alignment of the corresponding left and right singular vectors. If
\[
J_f\delta \approx J_g\delta
\]
for some perturbation direction \(\delta\), then the two operators respond similarly to that perturbation. In this case,
\[
\bigl(J_f-J_g\bigr)\delta \approx 0,
\]
so the perturbation may strongly affect both \(f\) and \(g\), while producing only a small change in their difference. This corresponds to a local cancellation effect.

On the other hand, if
\[
J_f\delta \approx -J_g\delta,
\]
then the responses of the two operators are locally opposed. In this case,
\[
\bigl(J_f-J_g\bigr)\delta
\]
can become large, because the two responses add in the difference. Therefore, the singular values of \(J_f-J_g\) measure the strongest directions in which the local responses of \(f\) and \(g\) disagree.

A basic bound follows from the triangle inequality:
\[
\left|
\sigma_{\max}(J_f)-\sigma_{\max}(J_g)
\right|
\le
\sigma_{\max}(J_f-J_g)
\le
\sigma_{\max}(J_f)+\sigma_{\max}(J_g).
\]
The upper bound is approached when the dominant responses of \(J_f\) and \(J_g\) are aligned in input direction but opposite in output direction. The lower bound is approached when the dominant responses are aligned both in input and output direction, so that the two local effects cancel.

To see this more explicitly, consider the rank-one case
\[
J_f=\sigma_f u_f v_f^\top,
\qquad
J_g=\sigma_g u_g v_g^\top.
\]
If
\[
u_f=u_g,\qquad v_f=v_g,
\]
then
\[
J_f-J_g=(\sigma_f-\sigma_g)uv^\top,
\]
and hence
\[
\sigma_{\max}(J_f-J_g)=|\sigma_f-\sigma_g|.
\]
This is the cancellation case. In contrast, if
\[
u_f=-u_g,\qquad v_f=v_g,
\]
then
\[
J_f-J_g=(\sigma_f+\sigma_g)uv^\top,
\]
and therefore
\[
\sigma_{\max}(J_f-J_g)=\sigma_f+\sigma_g.
\]
This is the reinforcement case.

Consequently, the top right singular vector of \(J_f-J_g\) identifies the input perturbation direction that maximally exposes the local disagreement between \(f\) and \(g\). The corresponding largest singular value
\[
\sigma_{\max}(J_f-J_g)
\]
quantifies the strongest local inconsistency between the two operators near \(x_0\). In this sense, the SVD of \(J_f-J_g\) is not measuring the sensitivity of \(f\) alone or \(g\) alone, but rather the sensitivity of their difference.

















\section{Related Literature on Adversarial Attacks, Universal Perturbations, Jacobian Singular Directions, and Spectral Regularization}

This section summarizes the main papers related to adversarial perturbations,
continuous-output attacks, universal singular-vector attacks, sparse perturbations,
matrix-free singular-vector computation, and Lipschitz or spectral-norm
regularization. The central theme is the local amplification problem
\[
    \max_{\|v\|_p = 1} \|J(x)v\|_q,
\]
where \(J(x)\) denotes either an input-output Jacobian or a hidden-layer
Jacobian. When \(p=q=2\), this reduces to the standard dominant right singular
vector problem. When \(p\) and \(q\) are general, the problem becomes a nonlinear
\((p,q)\)-singular vector problem and is usually handled by generalized power
iteration rather than ordinary SVD or Lanczos methods.

\subsection{Classical Adversarial Examples and Standard Attack Baselines}

\begin{landscape}
\begin{table}[p]
\centering
\footnotesize
\begin{tabular}{>{\raggedright\arraybackslash}p{2.4cm} >{\raggedright\arraybackslash}p{4.0cm} >{\raggedright\arraybackslash}p{3.4cm} >{\raggedright\arraybackslash}p{1.0cm} >{\raggedright\arraybackslash}p{3.2cm} >{\raggedright\arraybackslash}p{6.2cm}}
\hline
\textbf{Theme} & \textbf{Title} & \textbf{Authors} & \textbf{Year} & \textbf{Venue} & \textbf{Main role in this literature} \\
\hline
Early adversarial examples &
\textit{Intriguing Properties of Neural Networks} &
Christian Szegedy, Wojciech Zaremba, Ilya Sutskever, Joan Bruna, Dumitru Erhan, Ian Goodfellow, Rob Fergus &
2013/2014 &
ICLR 2014 &
One of the earliest papers showing that small, carefully optimized perturbations can cause deep networks to misclassify inputs. It establishes the basic phenomenon of adversarial examples. \\
\hline
Fast gradient attack &
\textit{Explaining and Harnessing Adversarial Examples} &
Ian J. Goodfellow, Jonathon Shlens, Christian Szegedy &
2014/2015 &
ICLR 2015 &
Introduces FGSM and gives the influential linearity-based explanation of adversarial vulnerability. It is the standard starting point for gradient-sign attacks. \\
\hline
Physical-world adversarial examples &
\textit{Adversarial Examples in the Physical World} &
Alexey Kurakin, Ian Goodfellow, Samy Bengio &
2016/2017 &
ICLR Workshop 2017 / arXiv &
Studies whether adversarial examples remain effective after being printed, photographed, or otherwise observed through real-world sensors. It is closely related to BIM/I-FGSM style iterative attacks. \\
\hline
Minimal perturbation attack &
\textit{DeepFool: A Simple and Accurate Method to Fool Deep Neural Networks} &
Seyed-Mohsen Moosavi-Dezfooli, Alhussein Fawzi, Pascal Frossard &
2015/2016 &
CVPR 2016 &
Computes small adversarial perturbations by iteratively linearizing the classifier boundary. It later becomes important for the original UAP construction. \\
\hline
Optimization-based attack &
\textit{Towards Evaluating the Robustness of Neural Networks} &
Nicholas Carlini, David Wagner &
2016/2017 &
IEEE Symposium on Security and Privacy 2017 &
Introduces the Carlini--Wagner attacks and provides a strong benchmark against many proposed defenses. \\
\hline
Robust optimization / PGD &
\textit{Towards Deep Learning Models Resistant to Adversarial Attacks} &
Aleksander Madry, Aleksandar Makelov, Ludwig Schmidt, Dimitris Tsipras, Adrian Vladu &
2017/2018 &
ICLR 2018 &
Frames adversarial robustness as a min--max robust optimization problem and popularizes PGD adversarial training. \\
\hline
\end{tabular}
\caption{Classical adversarial-example papers and standard attack baselines.}
\end{table}
\end{landscape}

\subsection{Universal Adversarial Perturbations and Singular-Vector Attacks}

\begin{landscape}
\begin{table}[p]
\centering
\footnotesize
\begin{tabular}{>{\raggedright\arraybackslash}p{2.4cm} >{\raggedright\arraybackslash}p{4.0cm} >{\raggedright\arraybackslash}p{3.4cm} >{\raggedright\arraybackslash}p{1.0cm} >{\raggedright\arraybackslash}p{3.2cm} >{\raggedright\arraybackslash}p{6.2cm}}
\hline
\textbf{Theme} & \textbf{Title} & \textbf{Authors} & \textbf{Year} & \textbf{Venue} & \textbf{Main role in this literature} \\
\hline
Universal adversarial perturbation &
\textit{Universal Adversarial Perturbations} &
Seyed-Mohsen Moosavi-Dezfooli, Alhussein Fawzi, Omar Fawzi, Pascal Frossard &
2016/2017 &
CVPR 2017 &
The classical UAP paper. It shows that a single image-agnostic perturbation can fool many natural images and often transfers across networks. \\
\hline
Geometry of UAP &
\textit{Analysis of Universal Adversarial Perturbations} &
Seyed-Mohsen Moosavi-Dezfooli, Alhussein Fawzi, Omar Fawzi, Pascal Frossard, Stefano Soatto &
2017 &
arXiv &
Analyzes why universal perturbations exist, focusing on geometric properties of decision boundaries. Useful as theoretical background for UAP. \\
\hline
Jacobian singular-vector UAP &
\textit{The Art of Singular Vectors and Universal Adversarial Perturbations} &
Valentin Khrulkov, Ivan Oseledets &
2017/2018 &
CVPR 2018 &
The most directly relevant paper for the problem
\(\max_{\|v\|_p=1}\|Jv\|_q\). It constructs UAPs using \((p,q)\)-singular vectors of hidden-layer Jacobians. \\
\hline
Sparse universal singular-vector attack &
\textit{Sparse and Transferable Universal Singular Vectors Attack} &
Kseniia Kuvshinova, Olga Tsymboi, Ivan Oseledets &
2024 &
arXiv 2024; ICOMP 2024 conference-paper version &
Extends the singular-vector UAP idea to sparse perturbations using truncated power iteration. It is directly relevant for sparse universal attack directions. \\
\hline
Adversarial Jacobian right singular vectors &
\textit{Adversarial Perturbations Are Formed by Iteratively Learning Linear Combinations of the Right Singular Vectors of the Adversarial Jacobian} &
Thomas Paniagua, Chinmay Savadikar, Tianfu ``Matt'' Wu &
2025 &
ICML 2025 / OpenReview &
Studies adversarial perturbations as linear combinations of right singular vectors of an adversarial Jacobian and introduces RisingAttacK based on sequential quadratic programming. \\
\hline
\end{tabular}
\caption{Universal perturbation papers and Jacobian singular-vector attack papers.}
\end{table}
\end{landscape}

\subsection{Sparse, Patch, and Structured Adversarial Attacks}

\begin{landscape}
\begin{table}[p]
\centering
\footnotesize
\begin{tabular}{>{\raggedright\arraybackslash}p{2.4cm} >{\raggedright\arraybackslash}p{4.0cm} >{\raggedright\arraybackslash}p{3.4cm} >{\raggedright\arraybackslash}p{1.0cm} >{\raggedright\arraybackslash}p{3.2cm} >{\raggedright\arraybackslash}p{6.2cm}}
\hline
\textbf{Theme} & \textbf{Title} & \textbf{Authors} & \textbf{Year} & \textbf{Venue} & \textbf{Main role in this literature} \\
\hline
Sparse pixel attack &
\textit{SparseFool: A Few Pixels Make a Big Difference} &
Apostolos Modas, Seyed-Mohsen Moosavi-Dezfooli, Pascal Frossard &
2018/2019 &
CVPR 2019 &
A classical sparse attack showing that changing only a few pixels can be sufficient to fool deep classifiers. \\
\hline
Universal patch attack &
\textit{Adversarial Patch} &
Tom B. Brown, Dandelion Mane, Aurko Roy, Martin Abadi, Justin Gilmer &
2017 &
arXiv / adversarial ML workshop-style publication &
Introduces universal, robust, targeted adversarial image patches that can work in real-world scenes. \\
\hline
Structured sparse attack &
\textit{Structured Adversarial Attack: Towards General Implementation and Better Interpretability} &
Kaidi Xu, Sijia Liu, Pu Zhao, Pin-Yu Chen, Huan Zhang, Quanfu Fan, Deniz Erdogmus, Yanzhi Wang, Xue Lin &
2018/2019 &
ICLR 2019 &
Uses group sparsity and ADMM to generate structured perturbations, giving more spatially interpretable attacks. \\
\hline
Stochastic sparse attack &
\textit{Stochastic Sparse Adversarial Attacks} &
Manon Cesaire, Lucas Schott, Hatem Hajri, Sylvain Lamprier, Patrick Gallinari &
2020/2021 &
IEEE ICTAI 2021 &
Introduces simple stochastic sparse attacks for targeted and untargeted classification settings. \\
\hline
Transferable structural sparse attack &
\textit{Transferable Structural Sparse Adversarial Attack Via Exact Group Sparsity Training} &
Di Ming, Peng Ren, Yunlong Wang, Xin Feng &
2024 &
CVPR 2024 &
Focuses on transferability and structural sparsity by training perturbation generators under exact group-sparsity constraints. \\
\hline
\end{tabular}
\caption{Sparse, patch-based, and structured adversarial attack papers.}
\end{table}
\end{landscape}

\subsection{Attacks for Regression, Continuous Outputs, Optical Flow, and Neural Operators}

\begin{landscape}
\begin{table}[p]
\centering
\footnotesize
\begin{tabular}{>{\raggedright\arraybackslash}p{2.4cm} >{\raggedright\arraybackslash}p{4.0cm} >{\raggedright\arraybackslash}p{3.4cm} >{\raggedright\arraybackslash}p{1.0cm} >{\raggedright\arraybackslash}p{3.2cm} >{\raggedright\arraybackslash}p{6.2cm}}
\hline
\textbf{Theme} & \textbf{Title} & \textbf{Authors} & \textbf{Year} & \textbf{Venue} & \textbf{Main role in this literature} \\
\hline
Autoencoder / VAE attack &
\textit{Adversarial Images for Variational Autoencoders} &
Pedro Tabacof, Julia Tavares, Eduardo Valle &
2016 &
NIPS 2016 Adversarial Training Workshop / arXiv &
Studies adversarial attacks on autoencoders and VAEs, where the output is a reconstructed image rather than a class label. \\
\hline
Classification-to-regression perturbation theory &
\textit{Perturbation Analysis of Learning Algorithms: Generation of Adversarial Examples From Classification to Regression} &
Emilio Rafael Balda, Arash Behboodi, Rudolf Mathar &
2018/2019 &
IEEE Transactions on Signal Processing, 2019 &
Provides a unified perturbation-analysis framework for generating adversarial examples in both classification and regression tasks. \\
\hline
Multivariate time-series regression attack &
\textit{Adversarial Examples in Deep Learning for Multivariate Time Series Regression} &
Gautam Raj Mode, Khaza Anuarul Hoque &
2020 &
IEEE Applied Imagery Pattern Recognition Workshop (AIPR) 2020 &
Applies adversarial attack methods to CNN, LSTM, and GRU models for multivariate time-series regression. \\
\hline
Steering-angle regression attack &
\textit{Are Self-Driving Cars Secure? Evasion Attacks Against Deep Neural Networks for Steering Angle Prediction} &
Alesia Chernikova, Alina Oprea, Cristina Nita-Rotaru, BaekGyu Kim &
2019 &
IEEE Workshop on the Internet of Safe Things, co-located with IEEE S\&P 2019 &
Studies evasion attacks against self-driving steering-angle prediction models, including regression-based steering outputs. \\
\hline
Physical-world autonomous driving attack &
\textit{DeepBillboard: Systematic Physical-World Testing of Autonomous Driving Systems} &
Husheng Zhou, Wei Li, Zelun Kong, Junfeng Guo, Yuqun Zhang, Bei Yu, Lingming Zhang, Cong Liu &
2018/2020 &
ICSE 2020 &
Constructs physical-world adversarial billboards to induce steering-angle errors in autonomous driving systems. \\
\hline
Optical-flow patch attack &
\textit{Attacking Optical Flow} &
Anurag Ranjan, Joel Janai, Andreas Geiger, Michael J. Black &
2019 &
ICCV 2019 &
Extends adversarial patch attacks to optical-flow networks, showing that a small patch can severely corrupt dense flow predictions. \\
\hline
Perturbation-constrained optical-flow attack &
\textit{A Perturbation-Constrained Adversarial Attack for Evaluating the Robustness of Optical Flow} &
Jenny Schmalfuss, Philipp Scholze, Andres Bruhn &
2022 &
ECCV 2022 &
Introduces PCFA, a global perturbation-constrained flow attack for robustness evaluation of optical-flow methods. \\
\hline
Neural operator / PDE surrogate attack &
\textit{Adversarial Vulnerabilities in Neural Operator Digital Twins: Gradient-Free Attacks on Nuclear Thermal-Hydraulic Surrogates} &
Samrendra Roy, Kazuma Kobayashi, Souvik Chakraborty, Rizwan-uddin, Syed Bahauddin Alam &
2026 &
arXiv 2026 &
Studies sparse gradient-free attacks on neural-operator digital twins for nuclear thermal-hydraulic surrogate modeling. Closest to PDE/neural-operator robustness. \\
\hline
\end{tabular}
\caption{Adversarial attacks for regression, continuous-output prediction, optical flow, autonomous driving, and neural operators.}
\end{table}
\end{landscape}

\subsection{Matrix-Free Singular-Vector and Eigenvector Computation}

\begin{landscape}
\begin{table}[p]
\centering
\footnotesize
\begin{tabular}{>{\raggedright\arraybackslash}p{2.4cm} >{\raggedright\arraybackslash}p{4.0cm} >{\raggedright\arraybackslash}p{3.4cm} >{\raggedright\arraybackslash}p{1.0cm} >{\raggedright\arraybackslash}p{3.2cm} >{\raggedright\arraybackslash}p{6.2cm}}
\hline
\textbf{Theme} & \textbf{Title} & \textbf{Authors} & \textbf{Year} & \textbf{Venue} & \textbf{Main role in this literature} \\
\hline
SVD / bidiagonalization &
\textit{Calculating the Singular Values and Pseudo-Inverse of a Matrix} &
Gene H. Golub, William Kahan &
1965 &
Journal of the Society for Industrial and Applied Mathematics, Series B: Numerical Analysis &
Classical source for Golub--Kahan bidiagonalization. It is the numerical linear algebra basis for matrix-free SVD using alternating products \(Av\) and \(A^\top u\). \\
\hline
Sparse eigenvector computation &
\textit{Truncated Power Method for Sparse Eigenvalue Problems} &
Xiao-Tong Yuan, Tong Zhang &
2011/2013 &
Journal of Machine Learning Research, 2013 &
Introduces truncated power iteration for sparse dominant eigenvector problems. This is the key algorithmic background for sparse singular-vector attacks. \\
\hline
Large-scale eigensolvers &
\textit{ARPACK Users' Guide: Solution of Large-Scale Eigenvalue Problems with Implicitly Restarted Arnoldi Methods} &
Richard B. Lehoucq, Danny C. Sorensen, Chao Yang &
1998 &
SIAM book &
Standard reference for matrix-free Arnoldi/Lanczos eigensolvers. Relevant when \(p=q=2\) and the problem is reduced to the dominant eigenvector of \(J^\top J\). \\
\hline
Power iteration &
\textit{Power Method / Power Iteration} &
Classical numerical linear algebra method &
Classical &
Textbook method &
Baseline algorithm for the dominant eigenvector or singular vector. For \(p=q=2\), one can iterate \(v \leftarrow J^\top Jv\) using JVP and VJP without explicitly forming \(J\). \\
\hline
Golub--Kahan Lanczos bidiagonalization &
\textit{Lanczos Bidiagonalization / Matrix-Free Partial SVD} &
Classical Krylov-subspace literature &
Classical &
Numerical linear algebra literature &
More efficient than naive power iteration for large matrix-free SVD problems. Directly uses \(Jv\) and \(J^\top u\), making it natural for Jacobian singular-vector computation when \(p=q=2\). \\
\hline
\end{tabular}
\caption{Matrix-free numerical linear algebra methods relevant to dominant Jacobian singular directions.}
\end{table}
\end{landscape}

\subsection{Jacobian Spectral Norm, Nuclear Norm, and Neural ODE Regularization}

\begin{landscape}
\begin{table}[p]
\centering
\footnotesize
\begin{tabular}{>{\raggedright\arraybackslash}p{2.4cm} >{\raggedright\arraybackslash}p{4.0cm} >{\raggedright\arraybackslash}p{3.4cm} >{\raggedright\arraybackslash}p{1.0cm} >{\raggedright\arraybackslash}p{3.2cm} >{\raggedright\arraybackslash}p{6.2cm}}
\hline
\textbf{Theme} & \textbf{Title} & \textbf{Authors} & \textbf{Year} & \textbf{Venue} & \textbf{Main role in this literature} \\
\hline
Exact Jacobian spectral norm regularization &
\textit{Exact Spectral Norm Regularization for Neural Networks} &
Anton Johansson, Claes Strannegard, Niklas Engsner, Petter Mostad &
2022 &
arXiv / later related versions from Chalmers group &
Directly targets the spectral norm of the input-output Jacobian. Closely related to \(\sigma_{\max}(J)\), but used for regularization rather than attack. \\
\hline
Improved Jacobian spectral norm regularization &
\textit{Improved Spectral Norm Regularization for Neural Networks} &
Anton Johansson, Niklas Engsner, Claes Strannegard, Petter Mostad &
2023 &
MDAI 2023 / Lecture Notes in Computer Science &
Improves the exact spectral norm regularization line and compares against other regularization methods. \\
\hline
Neural ODE Jacobian regularization &
\textit{How to Train Your Neural ODE: The World of Jacobian and Kinetic Regularization} &
Chris Finlay, Joern-Henrik Jacobsen, Levon Nurbekyan, Adam M. Oberman &
2020 &
ICML 2020 &
Uses Jacobian and kinetic regularization to simplify neural ODE dynamics and reduce solver cost. Related to stability and sensitivity control. \\
\hline
Neural ODE Jacobian norm and conditioning &
\textit{Jacobian Norm Regularisation and Conditioning in Neural ODEs} &
Shane Josias, Willie Brink &
2022 &
SACAIR 2022 / Springer proceedings &
Studies Jacobian norm and condition-number regularization in neural ODEs, especially in relation to solver behavior and robustness. \\
\hline
Jacobian nuclear norm regularization &
\textit{Nuclear Norm Regularization for Deep Learning} &
Christopher Scarvelis, Justin Solomon &
2024 &
NeurIPS 2024 &
Regularizes the nuclear norm of the function Jacobian, i.e.,
\(\|J\|_*=\sum_i \sigma_i(J)\). It focuses on the whole singular-value spectrum rather than only the largest singular value. \\
\hline
\end{tabular}
\caption{Papers on Jacobian spectral norm, Jacobian nuclear norm, and Neural ODE regularization.}
\end{table}
\end{landscape}

\subsection{Layer-Wise Spectral Norm and Lipschitz Regularization}

\begin{landscape}
\begin{table}[p]
\centering
\footnotesize
\begin{tabular}{>{\raggedright\arraybackslash}p{2.4cm} >{\raggedright\arraybackslash}p{4.0cm} >{\raggedright\arraybackslash}p{3.4cm} >{\raggedright\arraybackslash}p{1.0cm} >{\raggedright\arraybackslash}p{3.2cm} >{\raggedright\arraybackslash}p{6.2cm}}
\hline
\textbf{Theme} & \textbf{Title} & \textbf{Authors} & \textbf{Year} & \textbf{Venue} & \textbf{Main role in this literature} \\
\hline
Weight spectral norm regularization &
\textit{Spectral Norm Regularization for Improving the Generalizability of Deep Learning} &
Yuichi Yoshida, Takeru Miyato &
2017 &
arXiv 2017 &
Penalizes the spectral norms of weight matrices to reduce sensitivity to input perturbations and improve generalization. \\
\hline
Spectral normalization &
\textit{Spectral Normalization for Generative Adversarial Networks} &
Takeru Miyato, Toshiki Kataoka, Masanori Koyama, Yuichi Yoshida &
2018 &
ICLR 2018 &
Introduces spectral normalization for GAN discriminators. It uses power iteration to normalize layer-wise weight spectral norms. \\
\hline
Parseval / tight-frame regularization &
\textit{Parseval Networks: Improving Robustness to Adversarial Examples} &
Moustapha Cisse, Piotr Bojanowski, Edouard Grave, Yann Dauphin, Nicolas Usunier &
2017 &
ICML 2017 &
Constrains linear and convolutional layers to be approximately Parseval tight frames, thereby controlling layer-wise Lipschitz constants. \\
\hline
Fast approximate spectral normalization &
\textit{Fast Approximate Spectral Normalization for Robust Deep Neural Networks} &
Zhixin Pan, Prabhat Mishra &
2021 &
arXiv 2021 / related conference versions &
Accelerates spectral normalization using Fourier-transform and layer-separation ideas. Mainly concerns layer-wise spectral norms rather than full input-output Jacobians. \\
\hline
Toeplitz-based Lipschitz regularization &
\textit{On Lipschitz Regularization of Convolutional Layers using Toeplitz Matrix Theory} &
Alexandre Araujo, Benjamin Negrevergne, Yann Chevaleyre, Jamal Atif &
2021 &
AAAI 2021 &
Uses Toeplitz matrix theory to derive efficient upper bounds for convolutional layer Lipschitz constants. Relevant to robustness and spectral-norm control. \\
\hline
\end{tabular}
\caption{Layer-wise spectral normalization and Lipschitz regularization papers.}
\end{table}
\end{landscape}













\subsection{Summary of the Most Relevant Reading Path}

For the problem of finding a maximum local sensitivity direction without explicitly
forming the Jacobian, the most important path is
\[
\boxed{
\text{Khrulkov--Oseledets 2018}
\;\rightarrow\;
\text{Golub--Kahan / matrix-free SVD}
\;\rightarrow\;
\text{ARPACK / Lanczos}
\;\rightarrow\;
\text{Kuvshinova--Tsymboi--Oseledets 2024}
}
\]
when \(p=q=2\) or when the goal is close to a standard dominant singular-vector
problem. In that case, the key computational primitives are
\[
    v \mapsto Jv,
    \qquad
    u \mapsto J^\top u,
\]
which can be implemented by JVP and VJP without forming \(J\).

For the more general \((p,q)\)-singular vector problem,
\[
    \max_{\|v\|_p=1}\|Jv\|_q,
\]
standard SVD, Golub--Kahan, and Lanczos methods no longer directly solve the
original objective unless \(p=q=2\). The relevant algorithmic framework is instead
the generalized power iteration used in the singular-vector UAP literature. A
typical iteration has the form
\[
v_{k+1}
=
\frac{
\psi_{p'}
\left(
J^\top \psi_q(Jv_k)
\right)
}{
\left\|
\psi_{p'}
\left(
J^\top \psi_q(Jv_k)
\right)
\right\|_p
},
\]
where \(\psi_q\) and \(\psi_{p'}\) introduce nonlinear norm-dependent operations.

For sparse attacks, the most compact reading path is
\[
\boxed{
\text{UAP 2017}
\;\rightarrow\;
\text{Singular-vector UAP 2018}
\;\rightarrow\;
\text{Sparse universal singular vectors 2024}
\;\rightarrow\;
\text{Truncated power method 2013}
}
\]
with SparseFool, Adversarial Patch, and structured sparse attacks serving as
additional context.

For continuous-output models, regression models, optical flow, and neural
operators, the relevant line is
\[
\boxed{
\text{classification attack}
\;\rightarrow\;
\text{regression / continuous-output attack}
\;\rightarrow\;
\text{dense field attack}
\;\rightarrow\;
\text{neural-operator / PDE surrogate attack}.
}
\]
This line is especially relevant when the model output is a scalar, vector field,
time series, optical flow field, or PDE solution rather than a class label.





\section{The Basic Problem of Adversarial Attack}

The core problem of adversarial attack is: when the input is only allowed to undergo a very small perturbation, whether the model output will produce a clear error or a severe deviation.

Let the original input be
\[
x,
\]
the model be
\[
f,
\]
and the input perturbation be
\[
\delta.
\]
The attacked input is
\[
x_{\mathrm{adv}}=x+\delta.
\]

Usually, the perturbation is required to satisfy some norm constraint:
\[
\|\delta\|_p\leq \epsilon.
\]

Here, \(\epsilon\) is the attack budget. This quantity is very important. If the input itself is allowed to change a lot, then a large output change is normal and does not indicate that the model is fragile. The meaningful situation is:

\[
\text{The input perturbation is small, but the output change is large.}
\]

Therefore, for a continuous-output model, adversarial vulnerability cannot simply be described as:
\[
f(x+\delta)\neq f(x).
\]

A continuous function naturally changes when its input changes. A more meaningful question is whether
\[
\frac{\|f(x+\delta)-f(x)\|}{\|\delta\|}
\]
is large, or whether the post-attack error exceeds the task tolerance threshold.

\section{Adversarial Attack in Classification Tasks}

In a classification task, the model output is a class label. Let the true class be
\[
y.
\]

The success condition of an untargeted attack is usually:
\[
f(x+\delta)\neq y.
\]

The success condition of a targeted attack is:
\[
f(x+\delta)=t,
\]
where \(t\) is the target class specified by the attacker.

The special feature of classification problems is that they have natural discrete decision boundaries. The model often uses:
\[
\hat y=\arg\max_j f_j(x).
\]

Therefore, as long as the logit ranking changes, the output class will flip. Although the logits themselves change continuously, after applying \(\arg\max\), the class can suddenly change.

A common objective for classification attack is:
\[
\max_{\|\delta\|_p\leq \epsilon}
L(f(x+\delta),y).
\]

The form of FGSM is:
\[
x_{\mathrm{adv}}
=
x+
\epsilon\operatorname{sign}
\left(
\nabla_x L(f(x),y)
\right).
\]

The form of PGD is:
\[
x_{\mathrm{adv}}^{k+1}
=
\Pi_{B_\epsilon(x)}
\left(
x_{\mathrm{adv}}^k
+
\alpha
\operatorname{sign}
\left(
\nabla_x L(f(x_{\mathrm{adv}}^k),y)
\right)
\right),
\]
where
\[
B_\epsilon(x)=\{z:\|z-x\|_p\leq \epsilon\}.
\]

PGD is essentially multi-step projected gradient ascent inside the perturbation ball.

\section{Adversarial Attack in Regression Tasks}

In regression tasks, the output is continuous, for example:
\[
f(x)\in \mathbb{R}
\]
or
\[
f(x)\in \mathbb{R}^m.
\]

For tasks such as PDE surrogate modeling, optical flow, depth estimation, and neural operators, the output may even be a high-dimensional field:
\[
f(x)\in \mathbb{R}^{H\times W\times C}.
\]

In this case, there is no class flipping, so attack success cannot be directly defined by
\[
f(x+\delta)\neq f(x).
\]

If the true output \(y\) is available, we can define the clean error as:
\[
E_{\mathrm{clean}}
=
\|f(x)-y\|_q.
\]

The adversarial error is:
\[
E_{\mathrm{adv}}
=
\|f(x+\delta)-y\|_q.
\]

The error increase is:
\[
\Delta E
=
E_{\mathrm{adv}}-E_{\mathrm{clean}}.
\]

If a threshold \(\tau_{\mathrm{err}}\) is set, then
\[
\Delta E>\tau_{\mathrm{err}}
\]
can be defined as attack success.

If the true \(y\) is not convenient to obtain, or if we only care about the stability of the model itself, we can look at:
\[
D_{\mathrm{out}}
=
\|f(x+\delta)-f(x)\|_q.
\]

If a threshold \(\tau_{\mathrm{out}}\) is set, then
\[
D_{\mathrm{out}}>\tau_{\mathrm{out}}
\]
can be defined as attack success.

\section{The Ratio Between Input Change and Output Change}

For continuous-output models, a more fundamental quantity is the input-output amplification ratio:
\[
S(x,\delta)
=
\frac{
\|f(x+\delta)-f(x)\|_q
}{
\|\delta\|_p
}.
\]

If we take the local worst case:
\[
S_\epsilon(x)
=
\sup_{0<\|\delta\|_p\leq \epsilon}
\frac{
\|f(x+\delta)-f(x)\|_q
}{
\|\delta\|_p
}.
\]

This quantity measures the maximum output change caused by a unit input perturbation.

If \(f\) is differentiable, then we have the first-order approximation:
\[
f(x+\delta)-f(x)
\approx
J_f(x)\delta,
\]
where
\[
J_f(x)=\frac{\partial f(x)}{\partial x}
\]
is the Jacobian.

Therefore:
\[
\frac{
\|f(x+\delta)-f(x)\|_2
}{
\|\delta\|_2
}
\approx
\frac{
\|J_f(x)\delta\|_2
}{
\|\delta\|_2
}.
\]

Taking the maximum over all directions gives:
\[
\sup_{\delta\neq 0}
\frac{
\|J_f(x)\delta\|_2
}{
\|\delta\|_2
}
=
\|J_f(x)\|_2.
\]

And
\[
\|J_f(x)\|_2
=
\sigma_{\max}(J_f(x)).
\]

Therefore, under the local linear approximation from \(L_2\) to \(L_2\), the maximum input-output amplification factor is the largest singular value of the Jacobian.

\section{When the True Output Also Changes with the Input}

In PDEs, operator learning, and physical surrogate models, the true output is usually also a function of the input:
\[
y=y(x).
\]

If the input becomes:
\[
x+\delta,
\]
then the corresponding true output should also become:
\[
y(x+\delta).
\]

At this point, we should not compare with the old \(y(x)\). Instead, we should consider:
\[
E_{\mathrm{adv}}
=
\frac{
\|f(x+\delta)-y(x+\delta)\|_q
}{
\|y(x+\delta)\|_q
}.
\]

The clean error is:
\[
E_{\mathrm{clean}}
=
\frac{
\|f(x)-y(x)\|_q
}{
\|y(x)\|_q
}.
\]

The error increase is:
\[
\Delta E
=
E_{\mathrm{adv}}-E_{\mathrm{clean}}.
\]

That is:
\[
\Delta E
=
\frac{
\|f(x+\delta)-y(x+\delta)\|_q
}{
\|y(x+\delta)\|_q
}
-
\frac{
\|f(x)-y(x)\|_q
}{
\|y(x)\|_q
}.
\]

If we define a threshold:
\[
\Delta E>\tau_{\mathrm{inc}},
\]
then we can say that the attack causes an unacceptable increase in model error.

\section{The Relationship Between FGSM, BIM, and PGD}

FGSM is Fast Gradient Sign Method:
\[
\mathrm{FGSM}=\mathrm{Fast\ Gradient\ Sign\ Method}.
\]

It is a one-step attack:
\[
x_{\mathrm{adv}}
=
x+
\epsilon
\operatorname{sign}
\left(
\nabla_x L(f(x),y)
\right).
\]

BIM is Basic Iterative Method:
\[
\mathrm{BIM}=\mathrm{Basic\ Iterative\ Method}.
\]

It is also called I-FGSM:
\[
\mathrm{I\text{-}FGSM}=\mathrm{Iterative\ FGSM}.
\]

The form of BIM is:
\[
x_{\mathrm{adv}}^0=x,
\]
\[
x_{\mathrm{adv}}^{k+1}
=
\operatorname{Clip}_{x,\epsilon}
\left(
x_{\mathrm{adv}}^k
+
\alpha
\operatorname{sign}
\left(
\nabla_x L(f(x_{\mathrm{adv}}^k),y)
\right)
\right).
\]

Here,
\[
\operatorname{Clip}_{x,\epsilon}
\]
means clipping the attacked input back into the perturbation range around the original input \(x\).

PGD is Projected Gradient Descent, which in the attack context is usually projected gradient ascent:
\[
x_{\mathrm{adv}}^{k+1}
=
\Pi_{B_\epsilon(x)}
\left(
x_{\mathrm{adv}}^k
+
\alpha
\operatorname{sign}
\left(
\nabla_x L(f(x_{\mathrm{adv}}^k),y)
\right)
\right).
\]

The main difference between BIM and PGD is that PGD usually has a random start:
\[
x_{\mathrm{adv}}^0=x+\delta^0,
\qquad
\delta^0\sim \mathrm{Uniform}(-\epsilon,\epsilon),
\]
whereas BIM usually starts from:
\[
x_{\mathrm{adv}}^0=x.
\]

Therefore, we can say:
\[
\boxed{
\mathrm{BIM}\approx \mathrm{PGD\ without\ random\ start}.
}
\]

\section{Attacks Based on an Optimization Trade-off}

Some attacks do not take the FGSM, BIM, or PGD form. Instead, they put the input perturbation size and the output error into the same objective function.

For example, an attack on steering angle regression in autonomous driving can be written as:
\[
\min_{\sigma}
\|\sigma\|_2
-
c\left(F(x+\sigma)-y\right)^2.
\]

Here:
\[
\sigma
\]
is the input image perturbation,
\[
\|\sigma\|_2
\]
is encouraged to be as small as possible,
\[
\left(F(x+\sigma)-y\right)^2
\]
is encouraged to be as large as possible, and
\[
c
\]
controls the weight between the two terms.

The meaning of this objective is:
\[
\text{The input perturbation should be as small as possible, while the output error should be as large as possible.}
\]

This is different from PGD's projected update. PGD updates inside a perturbation ball:
\[
\delta^{k+1}
=
\Pi_{\|\delta\|\leq \epsilon}
\left(
\delta^k+\alpha \nabla_\delta L
\right).
\]

The trade-off form directly puts the perturbation norm into the optimization objective.

\section{Physical-World Attack in DeepBillboard}

DeepBillboard studies physical-world billboard attack. It is not ordinary full-image \(L_p\) noise. Instead, it optimizes the pattern on a roadside billboard so that the autonomous driving model produces wrong steering angles when the vehicle passes by.

Let the original image at frame \(i\) be:
\[
x_i.
\]

The image with the adversarial billboard is:
\[
x_i'.
\]

The model outputs steering angles:
\[
f(x_i),\qquad f(x_i').
\]

It defines the mean angle error as:
\[
M_0
=
\frac{1}{|\hat X|}
\sum_{i=1}^{|\hat X|}
\left(
f(x_i')-f(x_i)
\right).
\]

It also defines the proportion of frames exceeding a threshold:
\[
M_1
=
\frac{
\left|
\left\{
x_i:
f(x_i')-f(x_i)>\tau
\right\}
\right|
}{
|\hat X|
}.
\]

Here:
\[
\hat X
\]
is a set of video frames generated when the vehicle passes the billboard.

\[
M_1
\]
can be understood as the attack success rate for a continuous-output task.

The threshold \(\tau\) is not set arbitrarily. It is computed from physical driving behavior. For example:
\[
40\ \mathrm{MPH},
\qquad
0.2\ \mathrm{s},
\qquad
1\ \mathrm{m}
\]
of off-track target can be converted into:
\[
\tau=16.24^\circ.
\]

In the physical experiment, under:
\[
20\ \mathrm{MPH}
\]
and a time window of about:
\[
0.33\ \mathrm{s},
\]
the paper uses:
\[
\tau=19.8^\circ.
\]

This type of definition is more meaningful than simply reporting the output change, because it connects the output error with real physical risk.

\section{Optical Flow Attack: Patch Attack by Ranjan et al.}

The input of optical flow consists of two image frames:
\[
(I_t,I_{t+1}),
\]
and the output is a two-dimensional optical flow field:
\[
f(I_t,I_{t+1})\in \mathbb{R}^{M\times N\times 2}.
\]

Ranjan et al.'s \textit{Attacking Optical Flow} uses an adversarial patch. The patch is placed on both image frames, and the perturbed region is small, for example less than \(1\%\) or \(5\%\) of the image area.

Because natural videos often do not have complete dense ground-truth optical flow, they use the model prediction on the original image as pseudo ground truth.

Let the clean optical flow be:
\[
f_{\mathrm{clean}}.
\]

Let the attacked optical flow be:
\[
f_{\mathrm{adv}}.
\]

Their optimization objective is to minimize cosine similarity:
\[
\mathrm{CS}
=
\frac{
\langle f_{\mathrm{clean}}, f_{\mathrm{adv}}\rangle
}{
\|f_{\mathrm{clean}}\|_2
\|f_{\mathrm{adv}}\|_2
}.
\]

Minimizing this quantity means making the attacked optical flow direction deviate as much as possible from the original optical flow direction, even approaching the opposite direction.

The evaluation metric is endpoint error:
\[
\mathrm{EPE}
=
\frac{1}{MN}
\sum_{i=1}^{MN}
\|f_i-f_i^\ast\|_2.
\]

Relative degradation is:
\[
\mathrm{Relative\ Degradation}
=
\frac{
\mathrm{EPE}_{\mathrm{adv}}
-
\mathrm{EPE}_{\mathrm{clean}}
}{
\mathrm{EPE}_{\mathrm{clean}}
}
\times 100\%.
\]

This paper does not strictly define binary success or failure. It mainly compares the degradation degree of different optical flow architectures under the same patch attack.

Its conclusion is:
\[
\text{encoder-decoder optical flow networks are more easily damaged by patch attack,}
\]
while:
\[
\text{structures based on spatial pyramids, warping, and cost volumes are relatively more robust.}
\]

\section{The Core Idea of PCFA Optical Flow}

PCFA stands for:
\[
\textit{A Perturbation-Constrained Adversarial Attack for Evaluating the Robustness of Optical Flow}.
\]

It studies the following problem: optical flow methods were previously evaluated mainly by accuracy, but being accurate on clean inputs does not mean the output remains stable after small input perturbations.

Traditional accuracy looks at:
\[
\mathrm{AEE}(f,f_g),
\]
where \(f_g\) is the ground-truth flow.

Robustness looks at:
\[
\mathrm{AEE}(\breve f,f),
\]
where \(f\) is the predicted flow on the unattacked input, and \(\breve f\) is the adversarial flow on the attacked input.

PCFA argues that a robust method should not produce largely different flow prediction from slightly changed input frames.

\subsection{Input Perturbation Constraint in PCFA}

PCFA adds perturbations to two frames:
\[
\delta_t,\delta_{t+1}.
\]

The constraint is:
\[
\|\delta_t,\delta_{t+1}\|_2
\leq
\epsilon_2\sqrt{2I}.
\]

Here, \(I\) is the number of elements in a single image frame. Multiplying by:
\[
\sqrt{2I}
\]
makes the perturbation bound independent of image size.

When:
\[
\epsilon_2=0.01,
\]
it can be understood as an average distortion of about \(1\%\) of the color range per pixel.

\subsection{The Target Flow in PCFA}

PCFA is a targeted attack. It first specifies a target flow:
\[
f_t.
\]

The attacked output is:
\[
\breve f.
\]

The optimization objective is to make:
\[
\breve f
\]
as close as possible to:
\[
f_t.
\]

A common attack strength metric is:
\[
\mathrm{AEE}(\breve f,f_t)
=
\frac{1}{MN}
\sum_{i=1}^{MN}
\|\breve f_i-(f_t)_i\|_2.
\]

The smaller this quantity is, the closer the attacked optical flow is to the target flow, and the stronger the attack is.

\subsection{Several Targets in PCFA}

Negative flow target:
\[
f_t=-f.
\]

This means making the attacked optical flow as close as possible to the opposite direction of the original flow.

Zero flow target:
\[
f_t=0.
\]

This means making the model believe that there is no motion in the scene.

Other target:
\[
f_t=f_{\mathrm{other}}.
\]

This means making the output of the current image pair close to the flow of another scene.

\subsection{The Penalty Method in PCFA}

PCFA is not ordinary PGD. It writes the constraint into a penalty objective.

Let:
\[
\hat\delta=(\delta_t,\delta_{t+1}),
\]
\[
\hat\epsilon_2=\epsilon_2\sqrt{2I}.
\]

The constraint violation is:
\[
c(\hat\delta)
=
\max
\left(
0,
\|\hat\delta\|_2-\hat\epsilon_2
\right).
\]

If:
\[
\|\hat\delta\|_2\leq \hat\epsilon_2,
\]
then:
\[
c(\hat\delta)=0.
\]

If:
\[
\|\hat\delta\|_2>\hat\epsilon_2,
\]
then:
\[
c(\hat\delta)>0.
\]

PCFA optimizes:
\[
\phi(\hat\delta,\mu)
=
\mathcal{L}(\breve f,f_t)
+
\mu |c(\hat\delta)|.
\]

Therefore, PCFA does not always penalize larger perturbations. It only penalizes perturbations that exceed the budget.

\subsection{The Relationship Between PCFA and Lipschitz Continuity}

Lipschitz continuity says:
\[
\|f(x+\delta)-f(x)\|
\leq
L\|\delta\|.
\]

If \(L\) is small, then a small input change will not cause a large output change.

PCFA does not directly compute the Lipschitz constant, because the exact Lipschitz constant of a neural network is hard to obtain. It chooses to search for the maximum output change under a fixed input perturbation budget, which is equivalent to empirically evaluating:
\[
\frac{
\|f(x+\delta)-f(x)\|
}{
\|\delta\|
}.
\]

Therefore, PCFA operationalizes the Lipschitz idea: instead of explicitly solving for \(L\), it uses a strong attack to search for large output changes.

\section{Perturbation Analysis: A Jacobian Framework from Classification to Regression}

The core idea of this paper is that adversarial attack can be uniformly understood as a worst-case amplification problem where the input perturbation passes through the model's local sensitivity structure.

For a regression model:
\[
f:\mathbb{R}^d\to \mathbb{R}^m,
\]
if the true output \(y\) is available, the attack can be defined as:
\[
\max_{\|\eta\|_p\leq \epsilon}
\|y-f(x+\eta)\|_2^2.
\]

If the true \(y\) is unavailable, or if we only study the stability of the model itself, we can use the clean prediction \(f(x)\) as the reference:
\[
\max_{\|\eta\|_p\leq \epsilon}
\|f(x+\eta)-f(x)\|_2^2.
\]

When \(\eta\) is small:
\[
f(x+\eta)-f(x)
\approx
J_f(x)\eta.
\]

Therefore:
\[
\max_{\|\eta\|_p\leq \epsilon}
\|f(x+\eta)-f(x)\|_2^2
\approx
\max_{\|\eta\|_p\leq \epsilon}
\|J_f(x)\eta\|_2^2.
\]

This is the Jacobian amplification problem in regression attacks.

\subsection{The Largest Singular Value in the \(L_2\) Case}

When:
\[
p=2,
\]
we have:
\[
\max_{\|\eta\|_2\leq \epsilon}
\|J_f(x)\eta\|_2
=
\epsilon
\|J_f(x)\|_2.
\]

And:
\[
\|J_f(x)\|_2
=
\sigma_{\max}(J_f(x)).
\]

Therefore, the optimal perturbation direction is:
\[
\eta^\star=\epsilon v_1,
\]
where \(v_1\) is the largest right singular vector of \(J_f(x)\).

The maximum output change is:
\[
\epsilon\sigma_{\max}(J_f(x)).
\]

The squared form is:
\[
\max_{\|\eta\|_2\leq \epsilon}
\|J_f(x)\eta\|_2^2
=
\epsilon^2\sigma_{\max}(J_f(x))^2.
\]

\section{Why the Rayleigh Quotient Appears}

Start from:
\[
\|J\delta\|_2^2.
\]

Then:
\[
\|J\delta\|_2^2
=
(J\delta)^\top(J\delta)
=
\delta^\top J^\top J\delta.
\]

Let:
\[
A=J^\top J.
\]

The problem becomes:
\[
\max_{\|\delta\|_2\leq \epsilon}
\delta^\top A\delta.
\]

Let:
\[
\delta=\epsilon v,
\qquad
\|v\|_2=1.
\]

Then:
\[
\delta^\top A\delta
=
\epsilon^2 v^\top A v.
\]

Therefore:
\[
\max_{\|v\|_2=1}
v^\top A v
=
\lambda_{\max}(A).
\]

The corresponding direction satisfies:
\[
Av_1=\lambda_{\max}(A)v_1.
\]

Because:
\[
A=J^\top J,
\]
we have:
\[
\lambda_{\max}(A)=\sigma_{\max}(J)^2.
\]

Thus, \(v_1\) is the largest right singular vector of \(J\).

\section{Backpropagation, Power Iteration, and the Largest Singular Vector}

If we define:
\[
L(\delta)=\|J\delta\|_2^2,
\]
then:
\[
L(\delta)=\delta^\top A\delta.
\]

Its gradient is:
\[
\nabla_\delta L(\delta)=2A\delta=2J^\top J\delta.
\]

One step of gradient ascent:
\[
\delta^{k+1}
=
\delta^k+\alpha\nabla_\delta L(\delta^k)
\]
gives:
\[
\delta^{k+1}
=
\delta^k+2\alpha A\delta^k
=
(I+2\alpha A)\delta^k.
\]

If:
\[
Av_i=\lambda_i v_i,
\]
then:
\[
(I+2\alpha A)v_i
=
(1+2\alpha\lambda_i)v_i.
\]

Therefore, \(I+2\alpha A\) and \(A\) have the same eigenvectors, and the eigenvalues become:
\[
1+2\alpha\lambda_i.
\]

If the initial perturbation is expanded as:
\[
\delta^0
=
c_1v_1+c_2v_2+\cdots+c_dv_d,
\]
then after \(k\) steps:
\[
\delta^k
=
c_1(1+2\alpha\lambda_1)^k v_1
+
c_2(1+2\alpha\lambda_2)^k v_2
+
\cdots
+
c_d(1+2\alpha\lambda_d)^k v_d.
\]

If:
\[
\lambda_1>\lambda_2
\]
and:
\[
c_1\neq 0,
\]
then the direction gradually approaches:
\[
v_1.
\]

This is the mechanism of power iteration.

Therefore, backpropagation does not directly give the largest singular vector in one step. Instead, through multiple gradient ascent steps, it implicitly repeatedly applies:
\[
J^\top J,
\]
and gradually filters out the largest right singular vector direction.

\section{Why Starting from \(\delta=0\) Cannot Move}

If the objective is:
\[
L(\delta)=\|f(x+\delta)-f(x)\|_2^2,
\]
then:
\[
L(0)=0.
\]

Also:
\[
\nabla_\delta L(0)=0.
\]

Therefore, if gradient ascent starts from:
\[
\delta^0=0,
\]
it may not move.

Thus, in this kind of raw-output sensitivity attack, random start is usually needed:
\[
\delta^0\neq 0.
\]

As long as the random initial perturbation contains a nonzero component along the largest singular direction \(v_1\), the iteration will gradually amplify this direction.

\section{The Difference Between Accumulated Perturbation and Local Perturbation in PGD}

Here, two symbols must be distinguished.

The original input:
\[
x_0.
\]

The accumulated perturbation at PGD step \(k\):
\[
\Delta^k.
\]

The current attacked input:
\[
x_k=x_0+\Delta^k.
\]

If we stand at the current point \(x_k\) and consider taking another small local step, we should use another symbol:
\[
h.
\]

The local raw-output maximum change problem at the current point is:
\[
\max_{\|h\|_2\leq \alpha}
\|f(x_k+h)-f(x_k)\|_2.
\]

The first-order approximation is:
\[
f(x_k+h)-f(x_k)
\approx
J_f(x_k)h.
\]

Therefore:
\[
\max_{\|h\|_2\leq \alpha}
\|J_f(x_k)h\|_2.
\]

Its solution is:
\[
h^\star=\alpha v_1(x_k),
\]
where \(v_1(x_k)\) is the largest right singular vector of \(J_f(x_k)\).

Here:
\[
\Delta^k
\]
is the accumulated PGD perturbation, while:
\[
h
\]
is the local trial direction at the current point. These two should not be mixed.

\section{The Jacobian Changes in Nonlinear Networks}

The fixed-Jacobian power iteration explanation only holds under the local linear approximation.

In a real neural network:
\[
x_k=x_0+\Delta^k
\]
keeps changing, so:
\[
J_f(x_k)
\]
also changes.

If we maximize:
\[
L(\delta)
=
\|f(x+\delta)-f(x)\|_2^2,
\]
the true gradient is:
\[
\nabla_\delta L(\delta)
=
2J_f(x+\delta)^\top
\left(
f(x+\delta)-f(x)
\right).
\]

This is not the fixed expression:
\[
2J_f(x)^\top J_f(x)\delta.
\]

Only when:
\[
f(x+\delta)-f(x)\approx J_f(x)\delta
\]
and:
\[
J_f(x+\delta)\approx J_f(x)
\]
does it approximately become:
\[
\nabla_\delta L(\delta)
\approx
2J_f(x)^\top J_f(x)\delta.
\]

Therefore, true PGD can be understood as:
\[
\text{projected gradient ascent in a constantly changing local Jacobian field.}
\]

\section{The Relationship Between the PGD Gradient Direction and Singular Vectors}

If the maximized objective is a task loss:
\[
L(f(x),y),
\]
then by the chain rule:
\[
\nabla_x L
=
J_f(x)^\top \nabla_f L.
\]

Take the SVD of the Jacobian:
\[
J_f(x)=U\Sigma V^\top.
\]

Then:
\[
\nabla_x L
=
V\Sigma U^\top \nabla_f L.
\]

Expanding this gives:
\[
\nabla_x L
=
\sum_i
\sigma_i
\langle \nabla_f L,u_i\rangle
v_i.
\]

Therefore, the PGD direction is not simply the largest right singular vector. It is a weighted combination of many right singular vectors.

For a direction \(v_i\) to have a large weight in the PGD gradient, two conditions need to hold simultaneously:

\[
\sigma_i \text{ is large}
\]
and:
\[
\langle \nabla_f L,u_i\rangle \text{ is large}.
\]

That is, PGD favors:
\[
\text{directions that are both amplified by the Jacobian and able to increase the current loss.}
\]

Only when:
\[
\nabla_f L
\]
mainly points along the largest left singular vector \(u_1\), the PGD direction becomes close to the largest right singular vector \(v_1\).

\section{Why We Do Not Directly Solve for the Largest Singular Vector}

Theoretically, if the model is locally linear, the objective is to maximize the raw output's \(L_2\) change, and the constraint is also \(L_2\), then one can directly solve for the largest singular vector:
\[
\delta^\star=\epsilon v_1.
\]

However, in practice, this is usually not done, for the following reasons:

\begin{enumerate}
    \item The Jacobian \(J_f(x)\) has enormous dimensions, and explicitly constructing it is expensive.
    \item Full SVD or eigendecomposition is very time-consuming and memory-consuming.
    \item Neural networks are strongly nonlinear, and \(J_f(x+\delta)\) changes with position.
    \item Many attack objectives are not raw output change, but task loss.
    \item Many attacks use \(L_\infty\) constraints, where the largest singular vector is no longer a simple solution.
\end{enumerate}

The advantage of backpropagation is that it does not require explicitly constructing \(J\), nor does it require explicitly constructing \(J^\top J\). Through automatic differentiation, we can directly obtain:
\[
\nabla_\delta L.
\]

When:
\[
L(\delta)=\|f(x+\delta)-f(x)\|_2^2
\]
and the local linear approximation holds, backpropagation is equivalent to implicitly computing:
\[
2J^\top J\delta.
\]

Therefore, using an iterative method to approximate the most sensitive direction is usually more feasible than explicitly computing the Jacobian and SVD.

\section{The Difference Between Random Perturbation and Adversarial Perturbation}

Random perturbation:
\[
\delta_{\mathrm{rand}}
\]
usually does not deliberately align with the model's highly sensitive directions.

Adversarial perturbation intentionally searches for:
\[
\max_{\|\delta\|\leq \epsilon}
L(f(x+\delta),y)
\]
or:
\[
\max_{\|\delta\|\leq \epsilon}
\|f(x+\delta)-f(x)\|.
\]

If random perturbation and adversarial perturbation have similar effects, this means many directions are sensitive, and the model itself is locally unstable.

If adversarial perturbation is much stronger than random perturbation, this means the model has a small number of special fragile directions. In the language of Jacobians, it may mean:
\[
\sigma_1\gg \sigma_2,\sigma_3,\ldots.
\]

That is, the largest singular direction or a few highly sensitive directions dominate the output change.

\section{The Neural Operator Digital Twins Paper}

The paper:
\[
\textit{Adversarial Vulnerabilities in Neural Operator Digital Twins: Gradient-Free Attacks on Nuclear Thermal-Hydraulic Surrogates}.
\]

It studies nuclear thermal-hydraulic surrogate models / neural operator digital twins.

The input is:
\[
x\in\mathbb{R}^{102},
\]
including boundary conditions, sparse sensors, wall heat flux profiles, and so on.

The output is:
\[
f(x)\in\mathbb{R}^{3977\times 4},
\]
including pressure and three velocity components.

The attack uses sparse physical perturbation:
\[
\|\delta\|_0=k,
\]
where:
\[
k\in\{1,3,5,10\}.
\]

The attack objective is to maximize the relative field error:
\[
E_{\mathrm{rel}}(x+\delta)
=
\frac{
\|f(x+\delta)-y(x+\delta)\|_2
}{
\|y(x+\delta)\|_2
}.
\]

Note that here the true output also changes with the input:
\[
y=y(x+\delta).
\]

Attack success is defined as:
\[
E_{\mathrm{rel}}(x+\delta)>\tau.
\]

The thresholds include:
\[
10\%,\quad 20\%,\quad 30\%,\quad 40\%.
\]

The core conclusion of this paper is:
\[
\text{modifying less than }1\%\text{ of the input can push the relative }L_2\text{ field error from about }1.5\%\text{ to }37\%-63\%.
\]

It also proposes the effective perturbation dimension:
\[
d_{\mathrm{eff}}
=
\frac{
\left(
\sum_{i=1}^{d}
\|J_b[:,i]\|_2
\right)^2
}{
\sum_{i=1}^{d}
\|J_b[:,i]\|_2^2
}.
\]

Here:
\[
J_b[:,i]
\]
is the Jacobian column of the entire output field with respect to the \(i\)-th input coordinate.

This quantity measures whether sensitivity is concentrated on a small number of input coordinates.

If:
\[
d_{\mathrm{eff}}\approx 1,
\]
then the sensitivity is concentrated, and sparse attacks can easily hit the key coordinates.

If:
\[
d_{\mathrm{eff}}\approx d,
\]
then the sensitivity is dispersed, and sparse attacks are harder.

The importance of this paper is that it not only reports attack success, but also uses the structure of Jacobian column norms to explain why sparse attacks are effective.

\section{Active Learning and GAN-Style Neural-Operator Related Work}

Several neural-operator papers are adjacent to solver-integrated adversarial
training, but they should not be confused with it.

\subsection{Active Operator Learning with Predictive Uncertainty}

The paper
\[
\textit{Active Operator Learning with Predictive Uncertainty Quantification for Partial Differential Equations}
\]
by Winovich, Daneker, Lu, and Lin studies uncertainty-driven active operator
learning. Its main idea is to equip DeepONet and FNO-type models with
predictive uncertainty and use that uncertainty to choose which PDE examples
should be labeled next.

Its training loss is a Gaussian negative-log-likelihood style objective:
\[
\mathcal{L}_{\mathrm{NLL}}
=
\frac{1}{2}
\frac{(\mu-u)^2}{\sigma^2}
+
\frac{1}{2}\log(2\pi\sigma^2),
\]
where the model predicts both a mean \(\mu\) and an uncertainty \(\sigma\).
This loss calibrates predictive uncertainty to observed model error. It is not
an adversarial robustness loss.

The active-learning acquisition is model-uncertainty based:
\[
\text{select examples with large predicted uncertainty}.
\]
The solver is used to generate labels for selected examples, but the
acquisition objective is not the solver-model discrepancy:
\[
\|F_\theta(a+\delta)-G(a+\delta)\|.
\]

Its evaluation is therefore clean generalization, uncertainty calibration, and
data efficiency: MSE, MAE, relative \(L^1/L^2\) error, predicted uncertainty
versus observed error, coverage, and error versus training-set size. It does
not study PGD perturbations, epsilon budgets, Jacobian attack directions, or
solver-integrated adversarial training.

\subsection{Multi-Resolution Active Learning of FNOs}

The paper
\[
\textit{Multi-Resolution Active Learning of Fourier Neural Operators}
\]
studies multi-fidelity active learning for FNOs. The important distinction is
that it chooses both:
\[
\text{which input function to query}
\]
and:
\[
\text{which solver/simulation resolution to use}.
\]

Its acquisition score is a utility/cost rule based on probabilistic
multi-resolution FNO uncertainty. Thus it is solver-cost-aware and highly
relevant to data-efficient operator learning, but it is not an adversarial
attack method. It does not search for a perturbation \(\delta\) that maximizes
the solver-model discrepancy.

\subsection{Generative Adversarial Neural Operators}

The paper
\[
\textit{Generative Adversarial Neural Operators}
\]
uses adversarial learning in the GAN sense. It learns probability measures on
function spaces by pairing a generator neural operator with a discriminator
neural functional. The training objective is distribution matching:
\[
\text{generated function distribution}
\approx
\text{real function distribution}.
\]

This is different from adversarial robustness:
\[
\max_{\|\delta\|\leq\epsilon}
\|F_\theta(a+\delta)-G(a+\delta)\|.
\]

GANO does not use a numerical PDE solver as a model-solver reference inside
the objective. Its experiments study GRF data, mixtures of GRFs, and real
InSAR volcanic deformation functions. It is useful for showing that
adversarial objectives have been used with neural operators, but the
adversarial signal is generator-versus-discriminator, not perturbation
robustness.

\subsection{Karniadakis Group Turbulent-Flow adv-NO}

The paper
\[
\textit{Learning Turbulent Flows with Generative Models:
Super-resolution, Forecasting, and Sparse Flow Reconstruction}
\]
by Oommen, Khodakarami, Bora, Wang, and George Em Karniadakis introduces an
adversarially trained neural operator, often called adv-NO, for turbulent-flow
super-resolution and forecasting.

The main problem is spectral bias: ordinary \(L^2\)-trained neural operators
can recover low-wavenumber structure while oversmoothing high-frequency
turbulent structures. The adversarial loss is discriminator-based and is
combined with \(L^1\) and perceptual losses to produce sharper turbulent
fields and better energy spectra.

This is again not PGD-style solver-integrated robustness. It uses
adversarial/generative losses to improve high-frequency realism:
\[
\text{fake turbulent field}
\quad\text{vs}\quad
\text{real turbulent field}.
\]
It does not construct norm-bounded worst-case input perturbations and does not
define robustness through \(F_\theta(a+\delta)-G(a+\delta)\).

\section{Comparison of Different Papers}

\begin{center}\footnotesize\begin{tabular}{>{\raggedright\arraybackslash}p{2.8cm}>{\raggedright\arraybackslash}p{3.0cm}>{\raggedright\arraybackslash}p{3.0cm}>{\raggedright\arraybackslash}p{3.0cm}}
\hline
Paper & Input Perturbation Constraint & Output Damage Metric & Explicit Success Threshold \\
\hline
MTS Regression 2020 & \(\epsilon=0.2\) normalized perturbation & RMSE increase & No explicit threshold \\
Steering Angle 2019 & \(L_2\) image perturbation & MSE ratio & No explicit threshold in the regression part \\
DeepBillboard & Physical billboard region & Steering angle error & Yes, \(\tau\) comes from off-track risk \\
Attacking Optical Flow & Patch area constraint & EPE relative degradation & No binary threshold \\
PCFA & \(L_2\) bound, \(\epsilon_2=0.01\) has pixel-level meaning & AEE to target / AEE to original flow & No binary threshold \\
Perturbation Analysis & \(\|\eta\|_p\leq \epsilon\) & Jacobian amplification & Theoretical analysis, no binary threshold \\
Neural Operator 2026 & \(L_0\) sparse physical perturbation & Relative \(L_2\) field error & Yes, 10\%--40\% thresholds \\
Winovich active operator learning & No adversarial perturbation; sample acquisition by uncertainty & Clean relative \(L^1/L^2\), uncertainty calibration, data efficiency & No adversarial success threshold \\
MRA-FNO & No adversarial perturbation; choose input function and resolution & Relative error versus solver cost / data budget & No adversarial success threshold \\
GANO & No input attack; GAN distribution matching in function space & Function statistics, autocorrelation, histogram, circular statistics & No robustness threshold \\
Karniadakis adv-NO turbulence & No PGD input attack; discriminator/perceptual training for turbulent fields & Energy-spectrum error, sharp gradients, forecast/reconstruction quality & No adversarial robustness threshold \\
\hline
\end{tabular}
\end{center}

\section{Overall Summary}

For regression and continuous-output models, adversarial attack cannot simply mean:
\[
\text{the input changed, and the output also changed.}
\]

This is a normal property of continuous functions.

The meaningful situation is:
\[
\text{the input perturbation is small, but the output error is large.}
\]

More rigorously, one should report both:

\[
\|\delta\|_p\leq \epsilon,
\]
and:
\[
\|f(x+\delta)-f(x)\|_q
\]
or:
\[
\|f(x+\delta)-y(x+\delta)\|_q.
\]

If local amplification is being studied, one should report:
\[
\frac{
\|f(x+\delta)-f(x)\|_q
}{
\|\delta\|_p
}.
\]

After local linearization, this problem is controlled by the Jacobian:
\[
f(x+\delta)-f(x)\approx J_f(x)\delta.
\]

In the \(L_2\) case, the maximum local amplification direction is:
\[
\delta^\star=\epsilon v_1,
\]
where \(v_1\) is the largest right singular vector of \(J_f(x)\).

The maximum amplification factor is:
\[
\sigma_{\max}(J_f(x)).
\]

Under local linearity and a raw-output change objective, PGD and backpropagation can be viewed as approximate power iteration. However, in real nonlinear networks and task losses, the Jacobian changes with \(x+\delta\), and the PGD direction is not necessarily the largest singular vector direction.

Therefore, the most accurate understanding is:

\[
\boxed{
\text{Jacobian / singular values give the theory of local linear sensitivity, while PGD numerically searches for bad perturbations on a nonlinear model using gradients.}
}
\]

\section{Gradient-Based Perturbation Optimization, Rayleigh Quotients, and Jacobian-Free Spectral Methods}

\subsection{Problem Setup: Perturbing an Input}

Let
\[
F:\mathbb{R}^n\to \mathbb{R}^m
\]
be a vector-valued mapping, and let \(x\in \mathbb{R}^n\) be a fixed input. We consider a perturbation
\[
x \mapsto x+\delta,
\]
where \(\delta\in \mathbb{R}^n\). A common objective is to find a perturbation direction that maximizes the change in the output:
\[
L_1(\delta)
=
\|F(x+\delta)-F(x)\|_2^2.
\]
Define the residual
\[
r_1(\delta)=F(x+\delta)-F(x).
\]
Then
\[
L_1(\delta)=r_1(\delta)^T r_1(\delta).
\]

By the chain rule,
\[
\nabla_\delta L_1(\delta)
=
2J_{r_1}(\delta)^T r_1(\delta).
\]
Since
\[
r_1(\delta)=F(x+\delta)-F(x),
\]
and \(F(x)\) is constant with respect to \(\delta\), we have
\[
J_{r_1}(\delta)=J_F(x+\delta).
\]
Therefore,
\[
\boxed{
\nabla_\delta L_1(\delta)
=
2J_F(x+\delta)^T
\left[
F(x+\delta)-F(x)
\right].
}
\]

At \(\delta=0\),
\[
r_1(0)=F(x)-F(x)=0,
\]
so
\[
\nabla_\delta L_1(0)=0.
\]
Thus, if gradient ascent is initialized with \(\delta_0=0\), the iteration is stuck:
\[
\delta_{k+1}
=
\delta_k+\eta \nabla_\delta L_1(\delta_k),
\qquad
\delta_0=0
\quad\Longrightarrow\quad
\delta_k=0
\quad\text{for all }k.
\]

This does not mean that \(\delta=0\) is illegal. It only means that \(\delta=0\) is a stationary point for this objective, so first-order gradient ascent cannot move away from it.

\subsection{Local Linearization and the Rayleigh Quotient}

For small \(\delta\), we can linearize \(F\) around \(x\):
\[
F(x+\delta)-F(x)
\approx
J_F(x)\delta.
\]
Then the objective becomes
\[
L_1(\delta)
\approx
\|J_F(x)\delta\|_2^2.
\]
Expanding the squared norm gives
\[
\|J_F(x)\delta\|_2^2
=
\delta^T J_F(x)^T J_F(x)\delta.
\]
Let
\[
A=J_F(x)^T J_F(x).
\]
Then
\[
L_1(\delta)\approx \delta^T A\delta.
\]

If we impose the constraint
\[
\|\delta\|_2=\varepsilon,
\]
then the local maximization problem becomes
\[
\max_{\|\delta\|_2=\varepsilon}
\delta^T A\delta.
\]
Equivalently,
\[
\max_{\|\delta\|_2=1}
\delta^T A\delta.
\]
This is the classical Rayleigh quotient problem:
\[
R(\delta)=\frac{\delta^T A\delta}{\delta^T\delta}.
\]
For symmetric positive semidefinite \(A=J_F(x)^TJ_F(x)\),
\[
\max_{\|\delta\|_2=1}\delta^T A\delta
=
\lambda_{\max}(A).
\]
Since the eigenvalues of \(J_F(x)^TJ_F(x)\) are the squared singular values of \(J_F(x)\),
\[
\lambda_{\max}\left(J_F(x)^TJ_F(x)\right)
=
\sigma_{\max}(J_F(x))^2.
\]
Therefore, the optimal local perturbation direction is
\[
\boxed{
\delta^\star=\varepsilon v_{\max},
}
\]
where \(v_{\max}\) is the largest right singular vector of \(J_F(x)\), satisfying
\[
J_F(x)^TJ_F(x)v_{\max}
=
\sigma_{\max}(J_F(x))^2v_{\max}.
\]
The maximum output amplification is
\[
\boxed{
\|J_F(x)\delta^\star\|_2
=
\varepsilon \sigma_{\max}(J_F(x)).
}
\]

\subsection{Gradient Ascent versus the Singular Vector Solution}

For the local quadratic approximation
\[
L(\delta)=\delta^T A\delta,
\]
the gradient is
\[
\nabla_\delta L(\delta)=2A\delta.
\]
Gradient ascent gives
\[
\delta_{k+1}
=
\delta_k+\eta \nabla_\delta L(\delta_k)
=
\delta_k+2\eta A\delta_k,
\]
or
\[
\delta_{k+1}
=
(I+2\eta A)\delta_k.
\]

If \(A\) has eigen-decomposition
\[
Av_i=\lambda_i v_i,
\]
and the initial perturbation is expanded as
\[
\delta_0=\sum_i c_i v_i,
\]
then
\[
\delta_k
=
(I+2\eta A)^k\delta_0
=
\sum_i c_i(1+2\eta\lambda_i)^k v_i.
\]
Thus, directions associated with larger eigenvalues are amplified more strongly.

If the largest eigenvalue satisfies
\[
\lambda_1>\lambda_2\geq \cdots \geq \lambda_n,
\]
and if \(c_1\neq 0\), then repeated gradient ascent with normalization tends to emphasize the leading eigenvector direction:
\[
\delta_k \to v_1.
\]

Therefore, in a quadratic Rayleigh quotient problem, gradient ascent behaves similarly to a spectral iteration. However, one step of gradient ascent is generally not equal to the largest eigenvector direction. The gradient direction at \(\delta_k\) is
\[
2A\delta_k,
\]
which depends on the current iterate \(\delta_k\). It equals the largest eigenvector direction only if \(\delta_k\) is already aligned with that eigenvector.

Thus,
\[
\boxed{
\text{The singular vector solution gives the optimal direction directly, while gradient ascent approaches it iteratively.}
}
\]

\subsection{Power Iteration Interpretation}

The standard power iteration for finding the dominant eigenvector of \(A\) is
\[
\delta_{k+1}
=
\frac{A\delta_k}{\|A\delta_k\|_2}.
\]
If
\[
\delta_0=\sum_i c_i v_i,
\]
then
\[
A^k\delta_0
=
\sum_i c_i \lambda_i^k v_i.
\]
If
\[
|\lambda_1|>|\lambda_2|\geq \cdots,
\]
then
\[
\frac{A^k\delta_0}{\|A^k\delta_0\|_2}
\to
v_1,
\]
provided \(c_1\neq 0\).

Gradient ascent on the quadratic objective uses
\[
\delta_{k+1}
=
(I+2\eta A)\delta_k.
\]
This is not exactly the same as power iteration on \(A\), but it is power iteration on
\[
I+2\eta A.
\]
Since \(A\) and \(I+2\eta A\) have the same eigenvectors, gradient ascent can still converge toward the dominant eigenvector direction under suitable conditions.

Thus,
\[
\boxed{
\text{For a fixed quadratic objective, gradient ascent is closely related to power iteration.}
}
\]

\subsection{A Second Problem: Maximizing the Difference Between Two Functions}

Now consider two vector-valued functions
\[
F,G:\mathbb{R}^n\to \mathbb{R}^m.
\]
A different objective is
\[
L_2(\delta)
=
\|F(x+\delta)-G(x+\delta)\|_2^2.
\]
Define
\[
r_2(\delta)=F(x+\delta)-G(x+\delta).
\]
Then
\[
L_2(\delta)=r_2(\delta)^T r_2(\delta).
\]
By the chain rule,
\[
\nabla_\delta L_2(\delta)
=
2J_{r_2}(\delta)^T r_2(\delta).
\]
Since
\[
J_{r_2}(\delta)
=
J_F(x+\delta)-J_G(x+\delta),
\]
we obtain
\[
\boxed{
\nabla_\delta L_2(\delta)
=
2
\left[
J_F(x+\delta)-J_G(x+\delta)
\right]^T
\left[
F(x+\delta)-G(x+\delta)
\right].
}
\]

At \(\delta=0\),
\[
r_2(0)=F(x)-G(x),
\]
so
\[
\nabla_\delta L_2(0)
=
2
\left[
J_F(x)-J_G(x)
\right]^T
\left[
F(x)-G(x)
\right].
\]
This is generally nonzero. Therefore, unlike the previous objective, gradient ascent can often start from
\[
\delta_0=0.
\]

The key distinction is:
\[
F(x+0)-F(x)=0,
\]
but generally
\[
F(x)-G(x)\neq 0.
\]
Hence the first objective has zero residual at \(\delta=0\), while the second objective usually has nonzero residual at \(\delta=0\).

\subsection{Local Linearization of the Two-Function Objective}

Linearizing \(F\) and \(G\) around \(x\), we get
\[
F(x+\delta)-G(x+\delta)
\approx
F(x)-G(x)
+
\left[
J_F(x)-J_G(x)
\right]\delta.
\]
Let
\[
r_0=F(x)-G(x),
\]
and
\[
B=J_F(x)-J_G(x).
\]
Then
\[
L_2(\delta)
\approx
\|r_0+B\delta\|_2^2.
\]
Expanding gives
\[
L_2(\delta)
\approx
r_0^Tr_0
+
2r_0^TB\delta
+
\delta^TB^TB\delta.
\]
Unlike the first objective, this local approximation contains a linear term:
\[
2r_0^TB\delta.
\]
Therefore, at \(\delta=0\),
\[
\nabla_\delta L_2(0)
=
2B^Tr_0.
\]

If instead we study the local sensitivity of the difference function
\[
H(x)=F(x)-G(x),
\]
then we consider
\[
H(x+\delta)-H(x)
=
(F-G)(x+\delta)-(F-G)(x).
\]
Linearizing,
\[
H(x+\delta)-H(x)
\approx
J_H(x)\delta,
\]
where
\[
J_H(x)=J_F(x)-J_G(x).
\]
Then the Rayleigh quotient problem becomes
\[
\max_{\|\delta\|_2=1}
\left\|
\left[
J_F(x)-J_G(x)
\right]\delta
\right\|_2^2.
\]
Equivalently,
\[
\max_{\|\delta\|_2=1}
\delta^T
\left[
J_F(x)-J_G(x)
\right]^T
\left[
J_F(x)-J_G(x)
\right]
\delta.
\]
The optimal direction is the largest right singular vector of
\[
J_F(x)-J_G(x).
\]

\subsection{Zero Initialization and Stationary Points}

Gradient methods always require a current point. For gradient ascent,
\[
\delta_{k+1}
=
\delta_k+\eta \nabla_\delta L(\delta_k).
\]
Therefore, an initial guess \(\delta_0\) is always needed.

Whether \(\delta_0=0\) is a useful initialization depends on whether
\[
\nabla_\delta L(0)
\]
is zero.

If
\[
\nabla_\delta L(0)\neq 0,
\]
then zero initialization can move:
\[
\delta_1
=
0+\eta\nabla_\delta L(0).
\]

If
\[
\nabla_\delta L(0)=0,
\]
then zero initialization is stuck:
\[
\delta_1=0.
\]

A point satisfying
\[
\nabla L(\delta)=0
\]
is called a stationary point or critical point. It may be a local minimum, local maximum, saddle point, or degenerate flat point.

For
\[
L_1(\delta)=\|F(x+\delta)-F(x)\|_2^2,
\]
\(\delta=0\) is a stationary point because the residual vanishes:
\[
F(x)-F(x)=0.
\]
Since \(L_1(\delta)\geq 0\) and \(L_1(0)=0\), this point is typically a local minimum. If the goal is maximization, zero initialization is a poor starting point.

For
\[
L_2(\delta)=\|F(x+\delta)-G(x+\delta)\|_2^2,
\]
the residual at zero is generally
\[
F(x)-G(x)\neq 0,
\]
so the gradient at zero is usually nonzero.

Thus,
\[
\boxed{
\text{Gradient methods always need an initial point, but whether zero is a useful initial point depends on the objective.}
}
\]

\subsection{Hessian, Saddle Points, and Curvature Directions}

For a general smooth objective \(L(z)\), the local second-order approximation around a point \(z_0\) is
\[
L(z_0+s)
\approx
L(z_0)
+
\nabla L(z_0)^T s
+
\frac12 s^T H(z_0)s,
\]
where
\[
H(z_0)=\nabla^2 L(z_0)
\]
is the Hessian.

If
\[
\nabla L(z_0)=0,
\]
then the first-order term vanishes, and the local behavior is controlled by
\[
\frac12 s^TH(z_0)s.
\]
If \(H(z_0)\) has both positive and negative eigenvalues, then \(z_0\) is a saddle point. For maximization, the fastest second-order increasing direction is the eigenvector corresponding to the largest positive eigenvalue of \(H(z_0)\).

However, having positive and negative Hessian eigenvalues at a point does not imply that the gradient at that point is zero. For example,
\[
L(x,y)=x^2-y^2+x
\]
has Hessian
\[
H=
\begin{bmatrix}
2 & 0\\
0 & -2
\end{bmatrix},
\]
which has one positive and one negative eigenvalue, but
\[
\nabla L(0,0)=
\begin{bmatrix}
1\\
0
\end{bmatrix}
\neq 0.
\]
Thus, a saddle point requires both:
\[
\nabla L(z_0)=0
\]
and an indefinite Hessian.

\subsection{Why Explicit Jacobians Are Often Infeasible}

If
\[
F:\mathbb{R}^n\to \mathbb{R}^m,
\]
then
\[
J_F(x)\in \mathbb{R}^{m\times n}.
\]
For high-dimensional inputs and outputs, this matrix can be enormous.

For example, if
\[
n=10^6,
\]
then
\[
J_F(x)^TJ_F(x)\in \mathbb{R}^{10^6\times 10^6},
\]
which contains
\[
10^{12}
\]
entries. Explicitly storing or factorizing this matrix is usually impossible.

If the Jacobian \(J_F(x)\) can be explicitly formed, then one can compute its singular value decomposition:
\[
J_F(x)=U\Sigma V^T.
\]
The optimal local perturbation direction is the first right singular vector:
\[
v_{\max}=V_{:,1}.
\]
This is accurate and direct.

However, in large-scale neural networks, PDE solvers, image models, and other high-dimensional systems, forming \(J_F(x)\) or \(J_F(x)^TJ_F(x)\) explicitly is usually too expensive.

\subsection{Jacobian-Free and Matrix-Free Computation}

Instead of explicitly forming
\[
J_F(x)^TJ_F(x),
\]
one can compute its action on a vector:
\[
v\mapsto J_F(x)^TJ_F(x)v.
\]
This requires two operations.

First, compute the Jacobian-vector product:
\[
u=J_F(x)v.
\]
This is called a JVP.

Second, compute the vector-Jacobian product:
\[
J_F(x)^Tu.
\]
This is called a VJP.

Combining them gives
\[
J_F(x)^TJ_F(x)v.
\]

Therefore, one can perform spectral iteration using only matrix-vector products:
\[
v_{k+1}
=
\frac{J_F(x)^TJ_F(x)v_k}
{\|J_F(x)^TJ_F(x)v_k\|_2}.
\]
This is a Jacobian-free power iteration for approximating the largest right singular vector of \(J_F(x)\).

The same idea applies to
\[
J_F(x)-J_G(x).
\]
To find the most sensitive direction of the difference \(F-G\), one can apply power iteration to
\[
\left[J_F(x)-J_G(x)\right]^T
\left[J_F(x)-J_G(x)\right]
\]
without explicitly constructing either Jacobian.

\subsection{Gradient Methods, Power Iteration, and Large-Scale Spectral Problems}

For quadratic objectives,
\[
L(\delta)=\delta^TA\delta,
\]
gradient ascent gives
\[
\delta_{k+1}=(I+2\eta A)\delta_k.
\]
Power iteration gives
\[
\delta_{k+1}=\frac{A\delta_k}{\|A\delta_k\|_2}.
\]
Both methods repeatedly apply a fixed linear operator to the current vector, so both can be understood spectrally.

However, if the problem is explicitly a Rayleigh quotient or singular vector problem, then specialized spectral algorithms are usually more appropriate than vanilla gradient ascent. These include:
\[
\text{power iteration},
\qquad
\text{Lanczos iteration},
\qquad
\text{Arnoldi iteration},
\qquad
\text{SVD methods}.
\]
The reason is that these algorithms directly exploit the linear algebraic structure of the problem.

Thus,
\[
\boxed{
\text{If the matrix is available and small enough, use eigendecomposition or SVD.}
}
\]
If the matrix is too large but matrix-vector products are available, use matrix-free spectral methods:
\[
\boxed{
v\mapsto Av.
}
\]
If the objective is nonlinear and the Jacobian changes with \(\delta\), then the problem becomes a genuine nonlinear optimization problem rather than a fixed Rayleigh quotient problem.

\subsection{Summary}

The single-function perturbation objective
\[
\|F(x+\delta)-F(x)\|_2^2
\]
has zero residual at \(\delta=0\), so its gradient at zero is zero. Its local linearization leads to the Rayleigh quotient
\[
\max_{\|\delta\|_2=1}
\delta^T J_F(x)^T J_F(x)\delta,
\]
whose solution is the largest right singular vector of \(J_F(x)\).

The two-function discrepancy objective
\[
\|F(x+\delta)-G(x+\delta)\|_2^2
\]
usually has nonzero residual at \(\delta=0\), so gradient ascent can often start from zero. Its local gradient at zero is
\[
2
\left[
J_F(x)-J_G(x)
\right]^T
\left[
F(x)-G(x)
\right].
\]

If one studies the local sensitivity of the difference function \(F-G\), then the Rayleigh quotient involves
\[
\left[
J_F(x)-J_G(x)
\right]^T
\left[
J_F(x)-J_G(x)
\right].
\]

For fixed linear or quadratic problems, gradient ascent is closely related to power iteration. It repeatedly applies a linear operator to the current iterate and gradually amplifies dominant eigendirections. However, if the matrix or Jacobian is available, direct spectral methods are more efficient. In large-scale problems, explicit Jacobians are usually infeasible, so one uses Jacobian-free or matrix-free methods based on JVPs, VJPs, and iterative spectral algorithms.

\section{Jacobian Singular Directions, Matrix-Free Spectral Iteration, and Self-Attack Optimization}

\subsection{Problem Setup: Local Output Amplification}

Let
\[
F:\mathbb{R}^n\to \mathbb{R}^m
\]
be a neural network or a differentiable vector-valued mapping. Given an input
\[
x\in\mathbb{R}^n,
\]
we consider a small perturbation
\[
x\mapsto x+\delta.
\]
The local change in the network output is
\[
F(x+\delta)-F(x).
\]
For small \(\delta\), the first-order Taylor expansion gives
\[
F(x+\delta)-F(x)
\approx
J_F(x)\delta,
\]
where
\[
J_F(x)=\frac{\partial F(x)}{\partial x}\in\mathbb{R}^{m\times n}
\]
is the input-output Jacobian of the network at \(x\).

Therefore, the local self-attack problem
\[
\max_{\|\delta\|_2=1}
\|F(x+\delta)-F(x)\|_2
\]
is locally approximated by
\[
\max_{\|\delta\|_2=1}
\|J_F(x)\delta\|_2.
\]
Squaring the objective gives
\[
\max_{\|\delta\|_2=1}
\|J_F(x)\delta\|_2^2
=
\max_{\|\delta\|_2=1}
\delta^T J_F(x)^T J_F(x)\delta.
\]
Define
\[
B=J_F(x)^TJ_F(x).
\]
Then
\[
B=B^T,\qquad B\succeq 0.
\]
Thus the problem becomes the Rayleigh quotient maximization problem
\[
\max_{\|\delta\|_2=1}
\delta^T B\delta.
\]
The solution is the leading eigenvector of \(B\):
\[
Bv_{\max}=\lambda_{\max}v_{\max}.
\]
Since
\[
B=J_F(x)^TJ_F(x),
\]
we have
\[
\lambda_{\max}(B)=\sigma_{\max}(J_F(x))^2.
\]
Therefore,
\[
\boxed{
v_{\max}
=
\text{the leading right singular vector of }J_F(x)
}
\]
and
\[
\boxed{
\max_{\|\delta\|_2=1}\|J_F(x)\delta\|_2
=
\sigma_{\max}(J_F(x)).
}
\]

\subsection{The Exact Singular Vector Solution and Its Computational Difficulty}

If the Jacobian matrix \(J_F(x)\) were explicitly available, then one could compute its singular value decomposition:
\[
J_F(x)=U\Sigma V^T.
\]
The optimal local perturbation direction would be the first right singular vector:
\[
\delta^\star=v_1,
\]
where
\[
J_F(x)^TJ_F(x)v_1
=
\sigma_1^2 v_1.
\]

However, for neural networks, images, PDE states, or high-dimensional inputs, \(n\) and \(m\) can be very large. If
\[
J_F(x)\in\mathbb{R}^{m\times n},
\]
then
\[
B=J_F(x)^TJ_F(x)\in\mathbb{R}^{n\times n}.
\]
Explicitly forming \(J_F(x)\) is often infeasible, and explicitly forming \(B\) is even worse because its size scales like \(n^2\).

For example, if
\[
n=10^6,
\]
then
\[
B\in\mathbb{R}^{10^6\times 10^6},
\]
which contains
\[
10^{12}
\]
entries. This is usually impossible to store or factorize.

Therefore, the correct large-scale approach is not to form \(J_F(x)\) or \(J_F(x)^TJ_F(x)\). Instead, one uses a matrix-free method based only on the action
\[
v\mapsto J_F(x)^TJ_F(x)v.
\]

\subsection{Jacobian-Vector Products and Vector-Jacobian Products}

To compute
\[
Bv=J_F(x)^TJ_F(x)v
\]
without explicitly forming \(J_F(x)\), one uses two automatic differentiation primitives.

First, compute the Jacobian-vector product:
\[
u=J_F(x)v.
\]
This is called a JVP.

It represents the directional derivative of \(F\) at \(x\) along the direction \(v\):
\[
J_F(x)v
=
\left.
\frac{d}{d\epsilon}
F(x+\epsilon v)
\right|_{\epsilon=0}.
\]
Equivalently,
\[
F(x+\epsilon v)
\approx
F(x)+\epsilon J_F(x)v.
\]
A JVP is not an ordinary forward pass. A forward pass computes
\[
F(x),
\]
whereas a JVP computes the first-order change
\[
J_F(x)v.
\]

Second, compute the vector-Jacobian product:
\[
g=J_F(x)^Tu.
\]
This is called a VJP.

If
\[
\phi(x)=\langle F(x),u\rangle,
\]
then
\[
\nabla_x \phi(x)=J_F(x)^Tu.
\]
VJP is the fundamental operation behind reverse-mode automatic differentiation and backpropagation.

Combining JVP and VJP gives
\[
g
=
J_F(x)^T(J_F(x)v)
=
J_F(x)^TJ_F(x)v
=
Bv.
\]
Thus
\[
\boxed{
Bv
=
\operatorname{VJP}
\left(
\operatorname{JVP}(v)
\right).
}
\]

In PyTorch or JAX, this can be implemented conceptually as
\[
v
\longrightarrow
Jv
\longrightarrow
J^T(Jv).
\]

\subsection{Ordinary Forward Pass, JVP, and VJP}

It is useful to distinguish three operations.

An ordinary forward pass computes the function value:
\[
x\mapsto F(x).
\]

A JVP computes a directional derivative:
\[
v\mapsto J_F(x)v.
\]
It propagates a tangent vector forward through the computational graph.

A VJP computes a reverse-mode derivative:
\[
u\mapsto J_F(x)^Tu.
\]
It propagates an output-side cotangent vector backward through the computational graph.

In ordinary neural network training, one usually computes a scalar loss
\[
L(\theta)=\ell(F_\theta(x),y),
\]
then calls backpropagation to compute
\[
\nabla_\theta L.
\]
This is primarily a reverse-mode AD / VJP computation. It is not a JVP computation.

By contrast, the matrix-free singular vector problem requires
\[
J^TJv,
\]
which needs both:
\[
\boxed{
\text{JVP: }Jv
}
\]
and
\[
\boxed{
\text{VJP: }J^T(Jv).
}
\]

\subsection{Power Iteration for the Leading Singular Vector}

The simplest matrix-free method is power iteration on
\[
B=J^TJ.
\]
Initialize
\[
v_0\neq 0,\qquad \|v_0\|_2=1.
\]
Then iterate
\[
g_k=Bv_k=J^TJv_k,
\]
\[
v_{k+1}
=
\frac{g_k}{\|g_k\|_2}.
\]
Using JVP and VJP, each step becomes
\[
u_k=Jv_k,
\]
\[
g_k=J^Tu_k,
\]
\[
v_{k+1}
=
\frac{g_k}{\|g_k\|_2}.
\]
After convergence,
\[
v_k\approx v_{\max}.
\]
The largest singular value can be estimated by
\[
\sigma_{\max}(J)
\approx
\|Jv_k\|_2.
\]

This method is called
\[
\boxed{
\text{matrix-free power iteration}
}
\]
or, in this context,
\[
\boxed{
\text{Jacobian-free power iteration for the leading singular vector.}
}
\]

\subsection{Why Power Iteration Works}

Assume that
\[
Bv_i=\lambda_i v_i,
\]
with
\[
\lambda_1>\lambda_2\geq \cdots \geq \lambda_n\geq 0.
\]
Suppose
\[
v_0=c_1v_1+c_2v_2+\cdots+c_nv_n,
\]
with
\[
c_1\neq 0.
\]
Then
\[
B^k v_0
=
c_1\lambda_1^k v_1
+
c_2\lambda_2^k v_2
+\cdots+
c_n\lambda_n^k v_n.
\]
Factoring out \(\lambda_1^k\),
\[
B^k v_0
=
\lambda_1^k
\left[
c_1v_1
+
c_2\left(\frac{\lambda_2}{\lambda_1}\right)^k v_2
+
\cdots
+
c_n\left(\frac{\lambda_n}{\lambda_1}\right)^k v_n
\right].
\]
Since
\[
\left|\frac{\lambda_i}{\lambda_1}\right|<1
\qquad
(i\geq 2),
\]
all non-leading components decay relative to the leading component. Hence the normalized vector converges to
\[
v_1=v_{\max}.
\]

For
\[
B=J^TJ,
\]
the eigenvalues are
\[
\lambda_i=\sigma_i(J)^2.
\]
Therefore, the convergence rate of power iteration is controlled by
\[
\left(\frac{\lambda_2}{\lambda_1}\right)^k
=
\left(\frac{\sigma_2^2}{\sigma_1^2}\right)^k.
\]
If
\[
\sigma_2\approx \sigma_1,
\]
then convergence can be slow.

\subsection{Power Iteration versus Normalized Gradient Ascent}

Consider the quadratic objective
\[
L(v)=v^TBv.
\]
The gradient is
\[
\nabla_v L(v)=2Bv.
\]

Normalized gradient ascent updates
\[
v_{k+1}
=
\frac{v_k+\alpha\nabla L(v_k)}
{\|v_k+\alpha\nabla L(v_k)\|_2}.
\]
Since
\[
\nabla L(v_k)=2Bv_k,
\]
this becomes
\[
v_{k+1}
=
\frac{v_k+2\alpha Bv_k}
{\|v_k+2\alpha Bv_k\|_2}.
\]
Equivalently,
\[
v_{k+1}
=
\frac{(I+2\alpha B)v_k}
{\|(I+2\alpha B)v_k\|_2}.
\]
Thus normalized gradient ascent is power iteration on
\[
M=I+2\alpha B.
\]

If
\[
Bv_i=\lambda_i v_i,
\]
then
\[
Mv_i=(I+2\alpha B)v_i
=
(1+2\alpha\lambda_i)v_i.
\]
Therefore, \(B\) and \(M\) have the same eigenvectors, but different eigenvalues:
\[
\mu_i=1+2\alpha\lambda_i.
\]

Thus both standard power iteration and normalized gradient ascent converge to the same leading eigenvector \(v_1\), assuming suitable initialization. However, their convergence rates differ.

Power iteration has asymptotic ratio
\[
\rho_{\text{power}}
=
\frac{\lambda_2}{\lambda_1}.
\]
Normalized gradient ascent has asymptotic ratio
\[
\rho_{\text{grad}}
=
\frac{1+2\alpha\lambda_2}
{1+2\alpha\lambda_1}.
\]
For
\[
\lambda_1>\lambda_2\geq 0,
\qquad
\alpha>0,
\]
we have
\[
\frac{1+2\alpha\lambda_2}
{1+2\alpha\lambda_1}
>
\frac{\lambda_2}{\lambda_1}.
\]
Therefore,
\[
\boxed{
\rho_{\text{grad}}>\rho_{\text{power}}.
}
\]
Since a smaller convergence ratio means faster convergence,
\[
\boxed{
\text{standard power iteration is faster than normalized gradient ascent in this fixed quadratic setting.}
}
\]

Intuitively, power iteration directly replaces the current vector by
\[
Bv_k,
\]
whereas normalized gradient ascent uses
\[
v_k+2\alpha Bv_k.
\]
The identity term \(v_k\) preserves part of the old direction and slows down the spectral filtering process.

As
\[
\alpha\to\infty,
\]
we have
\[
I+2\alpha B\approx 2\alpha B,
\]
so normalized gradient ascent becomes closer to standard power iteration.

\subsection{Gradient Direction versus Global Singular Direction}

For
\[
L(v)=\|Jv\|_2^2=v^TJ^TJv,
\]
the gradient at the current point \(v_k\) is
\[
\nabla L(v_k)=2J^TJv_k.
\]
This is the local steepest ascent direction at the current \(v_k\).

However, the global optimal perturbation direction is the leading right singular vector
\[
v_{\max},
\]
which solves
\[
J^TJv_{\max}
=
\sigma_{\max}(J)^2v_{\max}.
\]

In general,
\[
J^TJv_k
\not\parallel
v_{\max}.
\]
They are parallel only when
\[
v_k\parallel v_{\max}.
\]
Therefore,
\[
\boxed{
\text{the gradient direction is a local update direction, while the leading singular vector is the global Rayleigh quotient optimizer.}
}
\]

Power iteration directly uses the current gradient direction
\[
J^TJv_k
\]
as the next direction after normalization:
\[
v_{k+1}
=
\frac{J^TJv_k}{\|J^TJv_k\|_2}.
\]
Thus, for \(p=q=2\), power iteration can be interpreted as repeatedly replacing the current direction by the normalized gradient direction of
\[
\|Jv\|_2^2.
\]

\subsection{Faster Matrix-Free Methods}

Power iteration is simple and memory-efficient, but it may be slow when
\[
\sigma_2\approx \sigma_1.
\]
Several faster methods are available.

\subsubsection{Lanczos Method on \(B=J^TJ\)}

Since
\[
B=J^TJ
\]
is symmetric positive semidefinite, one can use the Lanczos method.

Lanczos constructs the Krylov subspace
\[
\mathcal{K}_k(B,v_0)
=
\operatorname{span}\{v_0,Bv_0,B^2v_0,\dots,B^{k-1}v_0\}.
\]
Instead of keeping only the latest vector, as in power iteration, Lanczos keeps an orthonormal basis for this subspace and projects \(B\) onto a small tridiagonal matrix. The leading eigenvector of this small tridiagonal matrix gives a Ritz approximation to the leading eigenvector of \(B\).

Lanczos only requires matrix-vector products
\[
v\mapsto Bv=J^TJv.
\]
Therefore, it can be implemented matrix-free using JVP and VJP.

Compared with power iteration, Lanczos usually converges faster because it uses the entire Krylov subspace rather than only the latest iterate.

The memory cost is higher. If \(k\) Lanczos vectors are stored, the memory cost is approximately
\[
O(kn)
\]
instead of
\[
O(n)
\]
for power iteration. Restarted Lanczos can reduce memory usage by limiting the Krylov dimension.

\subsubsection{Golub--Kahan Lanczos Bidiagonalization}

Since the original problem is a singular value problem for \(J\), it is often better to work directly with \(J\) rather than with
\[
J^TJ.
\]
The matrix \(J^TJ\) squares the condition number:
\[
\kappa(J^TJ)=\kappa(J)^2.
\]
This can make the eigenvalue problem numerically harder.

Golub--Kahan Lanczos bidiagonalization directly uses
\[
v\mapsto Jv
\]
and
\[
u\mapsto J^Tu.
\]
It constructs left and right Krylov bases and a small bidiagonal matrix whose singular values approximate the singular values of \(J\).

A schematic form of the iteration is
\[
u_k
=
Jv_k-\beta_k u_{k-1},
\]
\[
\alpha_k=\|u_k\|_2,
\]
\[
u_k\leftarrow \frac{u_k}{\alpha_k},
\]
\[
v_{k+1}
=
J^Tu_k-\alpha_k v_k,
\]
\[
\beta_{k+1}=\|v_{k+1}\|_2,
\]
\[
v_{k+1}\leftarrow \frac{v_{k+1}}{\beta_{k+1}}.
\]
The scalars \(\alpha_k\) and \(\beta_k\) form a small bidiagonal matrix. The top singular vector of this small matrix is lifted back to the original input space to approximate the leading right singular vector of \(J\).

This method is more directly adapted to singular value problems than Lanczos on \(J^TJ\). It is typically the preferred matrix-free SVD method when one can compute both JVPs and VJPs.

\subsubsection{Randomized SVD}

Randomized SVD is useful when one wants several leading singular vectors. It uses random probe vectors to estimate the dominant range of \(J\).

A typical procedure is:
\[
\Omega\sim \text{random matrix},
\]
\[
Y=J\Omega,
\]
then orthonormalize \(Y\) to obtain a subspace basis \(Q\). A smaller matrix is then formed:
\[
C=Q^TJ,
\]
and an SVD is computed in the smaller space.

Randomized SVD can be fast when \(J\) is approximately low-rank, but it requires storing multiple probe vectors and output responses. Therefore, it is not always the most memory-efficient method if only one leading singular vector is needed.

\subsubsection{LOBPCG}

LOBPCG is a block method for large symmetric eigenvalue problems. It can be applied to
\[
B=J^TJ
\]
through matrix-free products
\[
v\mapsto Bv.
\]
It is especially effective when multiple eigenvectors are needed or when a good preconditioner is available.

However, for the single leading singular vector of a neural network Jacobian, LOBPCG may be more complex than necessary unless multiple directions are required.

\subsection{Recommended Algorithms for the Self-Attack Problem}

For the local self-attack problem
\[
\max_{\|v\|_2=1}\|J_F(x)v\|_2,
\]
under the constraint that \(J_F(x)\) cannot be explicitly constructed, the practical options are:

\paragraph{Most memory-efficient baseline.}
Use matrix-free power iteration:
\[
v_{k+1}
=
\frac{J^TJv_k}{\|J^TJv_k\|_2}.
\]
Each iteration needs one JVP and one VJP. This method stores only a few large vectors.

\paragraph{Faster convergence with moderate memory.}
Use restarted Lanczos on
\[
B=J^TJ.
\]
It still only requires the matrix-free operation
\[
v\mapsto J^TJv,
\]
but it uses Krylov subspace information and usually converges faster than power iteration.

\paragraph{Most appropriate singular-value solver.}
Use Golub--Kahan Lanczos bidiagonalization. This directly targets the SVD of \(J\), uses JVP and VJP, and avoids forming
\[
J^TJ
\]
explicitly. It is often the best choice when one wants a fast and matrix-free leading singular vector method.

Thus:
\[
\boxed{
\text{minimal memory: matrix-free power iteration;}
}
\]
\[
\boxed{
\text{better speed-memory tradeoff: restarted Lanczos;}
}
\]
\[
\boxed{
\text{most SVD-specific method: Golub--Kahan Lanczos bidiagonalization.}
}
\]

\subsection{Relation to Ordinary Neural Network Training}

The self-attack problem optimizes an input perturbation:
\[
x\mapsto x+\delta.
\]
The model parameters are fixed. The goal is to find the most sensitive input direction.

Ordinary neural network training optimizes model parameters:
\[
\theta\mapsto \theta-\eta\nabla_\theta L(\theta).
\]
The goal is to reduce a training loss, such as cross-entropy.

These are different problems.

The local self-attack problem has the structure
\[
\max_{\|v\|=1}\|J_xF_\theta(x)v\|,
\]
which is a singular vector problem.

Ordinary training solves
\[
\min_\theta L(\theta),
\]
which is a high-dimensional nonlinear optimization problem. It is generally not a singular vector problem.

Power iteration is therefore appropriate for spectral subproblems such as:
\[
\|J_xF_\theta(x)\|_2,
\]
\[
\sigma_{\max}(W),
\]
or
\[
\lambda_{\max}(\nabla_\theta^2 L).
\]
It is not a general replacement for gradient descent in neural network training.

\section{Related Literature}

\subsection{The Art of Singular Vectors and Universal Adversarial Perturbations}

\paragraph{Reference.}
Khrulkov and Oseledets, \emph{The Art of Singular Vectors and Universal Adversarial Perturbations}, arXiv 2017, CVPR 2018.

\paragraph{Main idea.}
This work constructs universal adversarial perturbations using singular vectors of hidden-layer Jacobians.

Let \(f_i(x)\) denote the feature map at layer \(i\). For small perturbations,
\[
f_i(x+\varepsilon)-f_i(x)
\approx
J_i(x)\varepsilon.
\]
The paper considers
\[
\max_{\|\varepsilon\|_p=L}
\|J_i(x)\varepsilon\|_q.
\]
This is called the \((p,q)\)-singular vector problem:
\[
\max_{\|v\|_p=1}\|Av\|_q.
\]

\paragraph{Case \(p=q=2\).}
When
\[
p=q=2,
\]
the problem becomes the ordinary largest singular vector problem:
\[
\max_{\|v\|_2=1}\|Av\|_2.
\]
The solution is the leading right singular vector of \(A\).

\paragraph{General \((p,q)\) setting.}
The paper also considers general \(p\) and \(q\). In the adversarial image setting, it mainly uses
\[
p=\infty,
\qquad
q=10.
\]
The \(p=\infty\) constraint controls the maximum per-pixel perturbation:
\[
\|\varepsilon\|_\infty\leq L.
\]
The choice \(q=10\) is used as a smooth approximation to \(q=\infty\), because
\[
\|z\|_q\to \|z\|_\infty
\quad
\text{as}
\quad
q\to\infty.
\]
The \(\infty\)-norm involves a maximum operation and is nonsmooth where the maximizing coordinate changes or where multiple coordinates tie. A finite large \(q\) gives a smoother surrogate.

\paragraph{Generalized power method.}
The paper gives a generalized power method for
\[
\max_{\|x\|_p=1}\|Ax\|_q.
\]
Define
\[
\psi_r(z)=\operatorname{sign}(z)|z|^{r-1},
\]
applied elementwise, and let \(p'\) be the conjugate exponent:
\[
\frac{1}{p}+\frac{1}{p'}=1.
\]
The update is
\[
Sx
=
\psi_{p'}\left(A^T\psi_q(Ax)\right),
\]
\[
x\leftarrow \frac{Sx}{\|Sx\|_p}.
\]
When
\[
p=q=2,
\]
we have
\[
\psi_2(z)=z,
\qquad p'=2,
\]
so the update becomes
\[
Sx=A^TAx,
\]
which is the standard power method.

\paragraph{Batch-level universal perturbation.}
For a batch
\[
X_b=\{x_1,\dots,x_b\},
\]
each sample has a Jacobian
\[
J_i(x_j)\in\mathbb{R}^{m\times n}.
\]
The paper stacks them vertically:
\[
J_i(X_b)
=
\begin{bmatrix}
J_i(x_1)\\
J_i(x_2)\\
\vdots\\
J_i(x_b)
\end{bmatrix}
\in\mathbb{R}^{bm\times n}.
\]
A single shared perturbation
\[
\varepsilon\in\mathbb{R}^n
\]
acts on all samples:
\[
J_i(X_b)\varepsilon
=
\begin{bmatrix}
J_i(x_1)\varepsilon\\
J_i(x_2)\varepsilon\\
\vdots\\
J_i(x_b)\varepsilon
\end{bmatrix}.
\]
Thus the batch objective is
\[
\max_{\|\varepsilon\|_p=L}
\|J_i(X_b)\varepsilon\|_q.
\]
For \(p=q=2\), this is equivalent to
\[
\max_{\|\varepsilon\|_2=1}
\varepsilon^T
\left(
\sum_{j=1}^b
J_i(x_j)^TJ_i(x_j)
\right)
\varepsilon.
\]
Therefore the universal perturbation is the leading right singular vector of the stacked Jacobian.

\paragraph{Stochastic power method.}
The full dataset objective is too large, so the paper samples a batch and applies the generalized power method to the batch-stacked Jacobian. This is called the stochastic power method.

The word ``stochastic'' refers to using a random batch rather than the whole dataset. The paper also notes that one could change the batch during power iterations, but fixed batches already work well empirically.

\paragraph{Fooling rate.}
The fooling rate measures the fraction of images whose predicted class changes after adding the universal perturbation:
\[
\text{fooling rate}
=
\frac{
\left|
\{x:\arg\max p(x)\neq \arg\max p(x+\varepsilon)\}
\right|
}{
|X|
}.
\]
This is a classification attack metric.

\subsection{Universal Adversarial Perturbations}

\paragraph{Reference.}
Moosavi-Dezfooli et al., \emph{Universal Adversarial Perturbations}, CVPR 2017.

\paragraph{Main idea.}
This is a foundational paper on universal adversarial perturbations. It seeks a single perturbation
\[
\varepsilon
\]
that fools many images:
\[
\arg\max p(x)
\neq
\arg\max p(x+\varepsilon).
\]
Its method is not primarily based on Jacobian singular vectors. Instead, it accumulates perturbations that push samples across decision boundaries.

\paragraph{Relevance.}
It is important as background for universal adversarial perturbations, but it is less directly connected to matrix-free Jacobian singular vector computation than Khrulkov and Oseledets.

\subsection{Sparse and Transferable Universal Singular Vectors Attack}

\paragraph{Reference.}
Kuvshinova et al., \emph{Sparse and Transferable Universal Singular Vectors Attack}, 2024.

\paragraph{Main idea.}
This work extends the singular-vector-based universal attack idea by introducing sparsity. It uses truncated or sparse variants of power iteration to obtain universal perturbations that affect fewer pixels.

\paragraph{Relevance.}
It is directly connected to the singular-vector attack line. It is especially relevant if one wants the perturbation to be sparse or more transferable.

\subsection{Exact Spectral Norm Regularization for Neural Networks}

\paragraph{Reference.}
\emph{Exact Spectral Norm Regularization for Neural Networks}, arXiv 2022.

\paragraph{Main idea.}
This work regularizes the input-output Jacobian spectral norm:
\[
\left\|
\frac{df_\theta(x)}{dx}
\right\|_2
=
\max_{\|v\|_2=1}
\left\|
\frac{df_\theta(x)}{dx}v
\right\|_2.
\]
The goal is not to attack the model, but to penalize local sensitivity and improve robustness.

\paragraph{Connection to ReLU networks.}
For a ReLU network, within a fixed activation region \(R\), the network is affine:
\[
f_\theta(x)=W_Rx+b_R.
\]
The local Jacobian is
\[
J_f(x)=W_R.
\]
The spectral norm regularizer is
\[
\|W_R\|_2.
\]

\paragraph{Method.}
The paper estimates the spectral norm using power iteration:
\[
v_{k+1}
=
\frac{W_R^TW_Rv_k}
{\|W_R^TW_Rv_k\|_2}.
\]
It avoids explicitly constructing \(W_R\) by applying forward-mode and backward-mode operations through the network with fixed activation masks.

\paragraph{Relevance.}
This paper is highly relevant to Jacobian spectral norm estimation. It is more on the regularization side than the attack side.

\subsection{Spectral Normalization for Generative Adversarial Networks}

\paragraph{Reference.}
Miyato et al., \emph{Spectral Normalization for Generative Adversarial Networks}, ICLR 2018.

\paragraph{Main idea.}
This paper normalizes neural network weight matrices by their spectral norm:
\[
\bar W=\frac{W}{\sigma_{\max}(W)}.
\]
The largest singular value is estimated by power iteration:
\[
v\leftarrow \frac{W^Tu}{\|W^Tu\|_2},
\]
\[
u\leftarrow \frac{Wv}{\|Wv\|_2}.
\]

\paragraph{Relevance.}
This is one of the most influential uses of power iteration in deep learning. However, it applies to layer weight matrices, not directly to the full input-output Jacobian.

\subsection{Spectral Norm Regularization for Improving the Generalizability of Deep Learning}

\paragraph{Reference.}
Yoshida and Miyato, 2017.

\paragraph{Main idea.}
This work regularizes the spectral norms of neural network weight matrices in order to improve generalization and robustness.

\paragraph{Relevance.}
It is related through spectral norm control, but it is layer-wise rather than directly based on the full input-output Jacobian.

\subsection{Parseval Networks}

\paragraph{Reference.}
Ciss\'e et al., \emph{Parseval Networks: Improving Robustness to Adversarial Examples}, ICML 2017.

\paragraph{Main idea.}
This paper encourages weight matrices to behave like Parseval tight frames, approximately satisfying
\[
W^TW\approx I.
\]
This controls layer Lipschitz constants and can improve robustness.

\paragraph{Relevance.}
It is related to controlling singular values and Lipschitz constants, but it does not directly solve the matrix-free Jacobian singular vector problem.

\subsection{Fast Approximate Spectral Normalization}

\paragraph{Reference.}
Pan and Mishra, \emph{Fast Approximate Spectral Normalization for Robust Deep Neural Networks}, 2021.

\paragraph{Main idea.}
This line of work seeks faster approximations to spectral normalization, often for layers such as convolutional layers.

\paragraph{Relevance.}
It is relevant to efficient spectral norm estimation, but it focuses more on layer-wise spectral normalization than on full input-output Jacobian singular vectors.

\subsection{Nuclear Norm Regularization for Deep Learning}

\paragraph{Reference.}
\emph{Nuclear Norm Regularization for Deep Learning}, NeurIPS 2024.

\paragraph{Main idea.}
Instead of penalizing only the largest singular value
\[
\sigma_{\max}(J),
\]
this line studies penalties involving the nuclear norm:
\[
\|J\|_*=\sum_i \sigma_i(J).
\]
This controls the whole singular value spectrum rather than only the top singular value.

\paragraph{Relevance.}
It is related to Jacobian singular values but less directly connected to finding the leading singular vector for self-attack.

\section{Relevance Ranking}

For the specific goal
\[
\text{``find the leading singular direction of a neural network Jacobian without explicitly forming the Jacobian,''}
\]
the most relevant works and methods are:

\begin{enumerate}
    \item \textbf{The Art of Singular Vectors and Universal Adversarial Perturbations.}
    
    This is the closest paper to the self-attack setting. It directly uses hidden-layer Jacobian singular vectors to construct universal perturbations.

    \item \textbf{Golub--Kahan Lanczos bidiagonalization.}
    
    This is the most appropriate numerical method if the goal is a fast matrix-free leading singular vector computation using \(Jv\) and \(J^Tu\).

    \item \textbf{Lanczos / ARPACK / Krylov eigensolvers.}
    
    These methods are important if one applies matrix-free eigensolvers to \(B=J^TJ\).

    \item \textbf{Exact Spectral Norm Regularization for Neural Networks.}
    
    This directly estimates and regularizes the input-output Jacobian spectral norm.

    \item \textbf{Spectral Normalization for GANs.}
    
    This is the classic deep learning application of power iteration for spectral norm estimation, although it focuses on layer weights rather than full input-output Jacobians.

    \item \textbf{Sparse and Transferable Universal Singular Vectors Attack.}
    
    This is a recent extension of the singular-vector universal attack idea.

    \item \textbf{Universal Adversarial Perturbations.}
    
    This is foundational background for universal attacks but less directly tied to singular vector computation.

    \item \textbf{Parseval Networks and layer-wise spectral norm regularization.}
    
    These are important for robustness and Lipschitz control but less directly tied to matrix-free Jacobian singular vectors.
\end{enumerate}

\section{Practical Recommendation}

For the self-attack problem
\[
\max_{\|\delta\|_2=1}
\|F(x+\delta)-F(x)\|_2
\approx
\max_{\|\delta\|_2=1}
\|J_F(x)\delta\|_2,
\]
the most practical route is:

\begin{enumerate}
    \item Implement a matrix-free product
    \[
    v\mapsto J^TJv
    \]
    using
    \[
    JVP + VJP.
    \]

    \item First test matrix-free power iteration:
    \[
    v_{k+1}
    =
    \frac{J^TJv_k}{\|J^TJv_k\|_2}.
    \]

    \item If convergence is too slow, use restarted Lanczos on
    \[
    B=J^TJ.
    \]

    \item If one wants the most SVD-specific method, use Golub--Kahan Lanczos bidiagonalization with the two operations:
    \[
    v\mapsto Jv,
    \qquad
    u\mapsto J^Tu.
    \]
\end{enumerate}

The final perturbation is
\[
\delta^\star=\epsilon v_{\max}.
\]
The corresponding local maximum output amplification is approximately
\[
\|F(x+\delta^\star)-F(x)\|_2
\approx
\epsilon\sigma_{\max}(J_F(x)).
\]

\section{Summary}

The local self-attack problem is fundamentally a leading singular vector problem for the neural network input-output Jacobian:
\[
\max_{\|\delta\|_2=1}
\|J_F(x)\delta\|_2.
\]
The exact solution is the leading right singular vector of \(J_F(x)\), but explicitly forming \(J_F(x)\) or \(J_F(x)^TJ_F(x)\) is usually infeasible.

The correct computational strategy is matrix-free:
\[
v\mapsto J_F(x)^TJ_F(x)v.
\]
This product can be computed by
\[
v
\xrightarrow{\text{JVP}}
Jv
\xrightarrow{\text{VJP}}
J^TJv.
\]

Power iteration is the simplest method:
\[
v_{k+1}
=
\frac{J^TJv_k}{\|J^TJv_k\|_2}.
\]
It is memory-efficient but can be slow when the top singular values are close.

Normalized gradient ascent is equivalent to power iteration on
\[
I+\alpha J^TJ,
\]
and is generally slower than standard power iteration on
\[
J^TJ
\]
for the fixed quadratic problem.

Faster matrix-free methods include Lanczos and Golub--Kahan Lanczos bidiagonalization. Among them, Golub--Kahan is most directly suited to singular value problems because it works with \(J\) and \(J^T\) directly rather than with \(J^TJ\).

The most directly related paper is \emph{The Art of Singular Vectors and Universal Adversarial Perturbations}, which uses hidden-layer Jacobian singular vectors to construct universal perturbations. For regularization, \emph{Exact Spectral Norm Regularization for Neural Networks} is highly relevant because it directly penalizes the neural network Jacobian spectral norm.

\section{Universal Adversarial Perturbations and Singular Directions}

Universal adversarial perturbations aim to find a single input perturbation \(v\) that fools a large fraction of samples from a data distribution. Given a classifier \(\hat{k}(\cdot)\), the goal is to construct a perturbation \(v\) such that
\[
\hat{k}(x+v)\neq \hat{k}(x)
\]
for many inputs \(x\), while keeping the perturbation small:
\[
\|v\|_p \leq \epsilon.
\]
Unlike image-dependent adversarial attacks, a universal perturbation is fixed and is applied to many different inputs.

The paper \emph{The Art of Singular Vectors and Universal Adversarial Perturbations} proposes a method called \emph{SingularFool}. Its core idea is to construct universal perturbations from the most amplified directions of hidden-layer Jacobians. Let \(f_i(x)\) denote the representation of input \(x\) at the \(i\)-th hidden layer. For a small perturbation \(\epsilon\), the first-order expansion gives
\[
f_i(x+\epsilon)-f_i(x)\approx J_i(x)\epsilon,
\]
where
\[
J_i(x)=\frac{\partial f_i(x)}{\partial x}
\]
is the Jacobian of the hidden-layer map with respect to the input.

Therefore, a natural way to find a sensitive input direction is to solve
\[
\max_{\|\epsilon\|_p=1}\|J_i(x)\epsilon\|_q.
\]
This direction is a generalized \((p,q)\)-singular vector of the Jacobian. For multiple samples \(x_1,\dots,x_N\), the universal direction can be obtained by maximizing
\[
\max_{\|v\|_p=1}\sum_{j=1}^N \|J_i(x_j)v\|_q^q.
\]
The resulting direction \(v\) is then scaled to the allowed perturbation magnitude and applied to all images:
\[
x \mapsto x+\epsilon v.
\]

\section{\((p,q)\)-Singular Vectors}

For a matrix \(A\in\mathbb{R}^{m\times n}\), the \((p,q)\)-operator norm is defined as
\[
\|A\|_{p\to q}
=
\max_{\|x\|_p=1}\|Ax\|_q.
\]
A maximizer of this problem is called a \((p,q)\)-singular vector. This generalizes the usual notion of a singular vector.

When
\[
p=q=2,
\]
we recover the usual spectral norm:
\[
\|A\|_{2\to 2}
=
\sigma_{\max}(A)
=
\max_{\|x\|_2=1}\|Ax\|_2.
\]
In this special case,
\[
\|Ax\|_2^2
=
(Ax)^\top(Ax)
=
x^\top A^\top A x,
\]
so the maximization becomes the Rayleigh quotient problem
\[
\max_{\|x\|_2=1}x^\top A^\top A x.
\]
The maximizer is the dominant eigenvector of \(A^\top A\), equivalently the dominant right singular vector of \(A\).

For general \(p,q\), however, the problem
\[
\max_{\|x\|_p=1}\|Ax\|_q
\]
cannot be reduced to an eigenvalue problem of \(A^\top A\). For example,
\[
\|Ax\|_q^q
=
\sum_i |(Ax)_i|^q,
\]
which is a quadratic form only when \(q=2\). Thus, for \(p,q\neq 2\), the problem is a nonlinear constrained optimization problem rather than a standard SVD problem.

\section{The Function \(\psi_r\) and Its Role}

The generalized power method used in the paper relies on the function
\[
\psi_r(z)
=
\operatorname{sign}(z)|z|^{r-1},
\]
applied componentwise to vectors.

This function naturally appears as the derivative of the \(r\)-power of the \(\ell_r\)-norm. Indeed, for
\[
\|z\|_r^r
=
\sum_i |z_i|^r,
\]
we have
\[
\nabla_z \|z\|_r^r
=
r\operatorname{sign}(z)|z|^{r-1}
=
r\psi_r(z).
\]
Thus, up to the constant factor \(r\), \(\psi_r(z)\) is the gradient direction of \(\|z\|_r^r\).

Several special cases are important:
\[
\psi_2(z)=z,
\]
so when \(r=2\), the map is linear. For \(r\neq 2\), \(\psi_r\) is nonlinear. For example,
\[
\psi_{10}(z)
=
\operatorname{sign}(z)|z|^9,
\]
which strongly emphasizes large components of \(z\). Also,
\[
\psi_1(z)=\operatorname{sign}(z),
\]
which appears in \(\ell_\infty\)-constrained perturbation problems.

\section{Derivation of the Generalized Power Method}

The goal is to solve
\[
\max_{\|x\|_p=1}\|Ax\|_q.
\]
Because the function \(t\mapsto t^q\) is monotone increasing for \(t\geq 0\), this is equivalent to maximizing
\[
F(x)=\|Ax\|_q^q.
\]
Expanding the \(q\)-norm gives
\[
F(x)
=
\sum_i |(Ax)_i|^q.
\]

Let
\[
z=Ax.
\]
Then
\[
F(x)=\|z\|_q^q.
\]
Using the chain rule,
\[
\nabla_x F(x)
=
A^\top \nabla_z \|z\|_q^q.
\]
Since
\[
\nabla_z \|z\|_q^q
=
q\psi_q(z),
\]
we get
\[
\nabla_x F(x)
=
qA^\top \psi_q(Ax).
\]
Ignoring the positive scalar \(q\), the essential gradient direction is
\[
g(x)
=
A^\top \psi_q(Ax).
\]

At the current iterate \(x_k\), define
\[
g_k=A^\top \psi_q(Ax_k).
\]
A first-order Taylor approximation gives
\[
F(y)
\approx
F(x_k)+\nabla F(x_k)^\top (y-x_k).
\]
Since the terms independent of \(y\) do not affect the maximizer, the next direction can be chosen by solving the linearized problem
\[
x_{k+1}
=
\arg\max_{\|y\|_p=1} g_k^\top y.
\]

This problem is solved using the duality between \(\ell_p\) and \(\ell_{p'}\), where the conjugate exponent \(p'\) satisfies
\[
\frac{1}{p}+\frac{1}{p'}=1,
\qquad
p'=\frac{p}{p-1}.
\]
By H\"older's inequality,
\[
g^\top y
\leq
\|g\|_{p'}\|y\|_p.
\]
Under the constraint \(\|y\|_p=1\), this gives
\[
g^\top y\leq \|g\|_{p'}.
\]
The maximizer is
\[
y
=
\frac{\psi_{p'}(g)}
{\|\psi_{p'}(g)\|_p}.
\]
Substituting
\[
g_k=A^\top\psi_q(Ax_k)
\]
gives the generalized power iteration
\[
x_{k+1}
=
\frac{
\psi_{p'}
\left(
A^\top \psi_q(Ax_k)
\right)
}{
\left\|
\psi_{p'}
\left(
A^\top \psi_q(Ax_k)
\right)
\right\|_p
}.
\]

Thus one iteration consists of the following steps:
\[
x_k
\longmapsto
Ax_k
\longmapsto
\psi_q(Ax_k)
\longmapsto
A^\top\psi_q(Ax_k)
\longmapsto
\psi_{p'}\left(A^\top\psi_q(Ax_k)\right)
\longmapsto
x_{k+1}.
\]

\section{Reduction to the Ordinary Power Method When \(p=q=2\)}

When
\[
p=q=2,
\]
the conjugate exponent satisfies
\[
p'=2,
\]
and
\[
\psi_2(z)=z.
\]
Therefore,
\[
\psi_q(Ax_k)=\psi_2(Ax_k)=Ax_k.
\]
Then
\[
A^\top\psi_q(Ax_k)
=
A^\top A x_k.
\]
Since
\[
\psi_{p'}=\psi_2
\]
is also the identity map, the generalized update becomes
\[
x_{k+1}
=
\frac{A^\top A x_k}{\|A^\top A x_k\|_2}.
\]
This is precisely the ordinary power iteration applied to \(A^\top A\).

Therefore,
\[
p=q=2
\]
is the special case in which the \((p,q)\)-singular vector problem becomes the usual singular vector problem. In all other cases, the method remains a generalized power method rather than ordinary SVD.

\section{Connection to FGSM, BIM, and PGD}

The generalized power update has a close conceptual relationship to adversarial attacks such as FGSM, BIM, and PGD, because all of them use gradient information under a norm constraint. However, the update mechanism is different.

FGSM solves the linearized adversarial problem under an \(\ell_\infty\) constraint:
\[
\max_{\|\delta\|_\infty\leq \epsilon}
\nabla_x \mathcal{L}(x,y)^\top \delta.
\]
The solution is
\[
\delta
=
\epsilon\operatorname{sign}(\nabla_x\mathcal{L}(x,y)).
\]
This is exactly the steepest ascent direction under an \(\ell_\infty\) constraint.

PGD and BIM use additive updates followed by projection:
\[
\delta_{k+1}
=
\Pi_{\|\delta\|_p\leq \epsilon}
\left(
\delta_k+\alpha d_k
\right),
\]
where \(d_k\) is a gradient-based ascent direction. These methods accumulate perturbations step by step.

The generalized power method differs because it does not update by
\[
x_k+\alpha g_k.
\]
Instead, it directly chooses the best direction for the linearized objective under the \(p\)-norm constraint:
\[
x_{k+1}
=
\arg\max_{\|y\|_p=1}g_k^\top y.
\]
This gives
\[
x_{k+1}
=
\frac{\psi_{p'}(g_k)}
{\|\psi_{p'}(g_k)\|_p}.
\]
Thus, it is better interpreted as a fixed-point or generalized power iteration rather than projected gradient ascent.

For example, when
\[
p=\infty,
\]
we have
\[
p'=1,
\]
so
\[
\psi_{p'}(g)=\psi_1(g)=\operatorname{sign}(g).
\]
Thus the update resembles a sign-gradient attack. However, the gradient is not the gradient of a classification loss. Instead, it is the gradient of a hidden-layer amplification objective:
\[
g_k=A^\top\psi_q(Ax_k).
\]

\section{SingularFool}

SingularFool is the method proposed in the paper for constructing universal adversarial perturbations. It is based on computing \((p,q)\)-singular vectors of hidden-layer Jacobians.

For a hidden layer \(f_i(x)\), the local feature change is approximated by
\[
f_i(x+v)-f_i(x)\approx J_i(x)v.
\]
The method searches for a direction \(v\) that is strongly amplified by this Jacobian. For a batch of samples \(x_1,\dots,x_N\), SingularFool solves the problem
\[
\max_{\|v\|_p=1}
\sum_{j=1}^N
\|J_i(x_j)v\|_q^q.
\]
This is done using the generalized power method. The computation is matrix-free: the full Jacobian does not need to be formed. Instead, the method only requires products of the form
\[
J_i(x)v
\]
and
\[
J_i(x)^\top u.
\]
In deep learning frameworks, these correspond to Jacobian-vector products and vector-Jacobian products.

After computing the direction \(v\), the universal perturbation is obtained by scaling it to the allowed perturbation radius:
\[
\delta=\epsilon v,
\]
and then applying the same perturbation to every image:
\[
x\mapsto x+\delta.
\]

\section{Comparison with the UAP Method}

The UAP method, introduced by Moosavi-Dezfooli et al., also constructs universal adversarial perturbations. It solves the same broad task as SingularFool:
\[
\text{find one perturbation that fools many inputs}.
\]
However, the construction mechanism is different.

UAP maintains a universal perturbation \(v\). It then iterates over training samples. If the current perturbation already fools the sample \(x_i\), the algorithm does nothing. If not, it computes an additional sample-specific correction \(\Delta v_i\), often using DeepFool, such that
\[
x_i+v+\Delta v_i
\]
crosses the decision boundary. Then it updates
\[
v\leftarrow v+\Delta v_i
\]
and projects back to the allowed perturbation set:
\[
v\leftarrow \Pi_{\|v\|_p\leq \epsilon}(v).
\]

Thus, UAP is a decision-boundary-based, sample-by-sample, cumulative perturbation method. It is indirectly gradient-based because DeepFool uses local linearization of the classifier and therefore uses gradient or Jacobian information. However, UAP is not a generalized power method.

SingularFool is different. It does not explicitly search for the closest decision boundary for each sample. Instead, it uses hidden-layer Jacobians and computes a universal direction that is maximally amplified in feature space:
\[
\max_{\|v\|_p=1}
\sum_j
\|J_i(x_j)v\|_q^q.
\]
Therefore, the main distinction is:
\[
\text{UAP: decision-boundary accumulation},
\]
whereas
\[
\text{SingularFool: hidden-Jacobian singular direction}.
\]

\section{Experimental Comparison Between SingularFool and UAP}

The paper compares SingularFool with the original UAP method at the attack-construction level. The comparison focuses on fooling rate, sample efficiency, and approximate running time.

The reported results show that SingularFool can obtain a high fooling rate using relatively few images. In particular, the paper reports that using only \(64\) images to construct the perturbation can produce a fooling rate above \(60\%\) on a validation set of \(50{,}000\) images.

In the comparison around Figure 9, the authors report that SingularFool reaches roughly \(60\%\) fooling rate in about one minute under their setting, using VGG-19, the block2 pool layer, batch size \(32\), and \(30\) power iterations. In contrast, under their hardware setting, UAP with batch size \(128\) takes about ten minutes and reaches only around \(20\%\) fooling rate. They also note that, according to the original UAP paper, reaching a fooling rate around \(60\%\) requires roughly \(3000\) images.

Therefore, the main empirical conclusion is that SingularFool is more sample-efficient and faster in the reported setting. It can find a strong universal perturbation from a small batch of images, whereas UAP requires more samples and more sequential boundary-correction steps.

However, this comparison should be interpreted carefully. The paper does not prove that generalized power iteration is globally optimal, nor does it systematically compare it against all possible numerical optimization methods such as PGD, BIM, Lanczos, or Adam. The comparison is mainly at the level of universal attack construction:
\[
\text{SingularFool vs. UAP}.
\]
The paper's main claim is that hidden-layer Jacobian singular directions provide an effective mechanism for constructing universal adversarial perturbations.

\section{Main Takeaways}

The central mathematical object in the paper is
\[
\max_{\|v\|_p=1}\|Jv\|_q,
\]
or, for multiple samples,
\[
\max_{\|v\|_p=1}
\sum_j
\|J(x_j)v\|_q^q.
\]
When \(p=q=2\), this reduces to the ordinary singular vector problem:
\[
v_{k+1}
=
\frac{J^\top Jv_k}{\|J^\top Jv_k\|_2}.
\]
For general \(p,q\), the update becomes
\[
v_{k+1}
=
\frac{
\psi_{p'}
\left(
J^\top \psi_q(Jv_k)
\right)
}{
\left\|
\psi_{p'}
\left(
J^\top \psi_q(Jv_k)
\right)
\right\|_p
}.
\]
This is the generalized power method.

The method is matrix-free, since it only requires
\[
Jv
\]
and
\[
J^\top u.
\]
This makes it suitable for neural networks, where explicitly forming the Jacobian is usually infeasible.

Conceptually, SingularFool shows that universal adversarial perturbations can be understood as directions that are strongly amplified by hidden-layer Jacobians. This provides a different explanation from purely decision-boundary-based methods such as UAP.

\section{Can Lanczos or Golub--Kahan Be Used for Matrix-Free Jacobian Singular Directions?}

Suppose we want to find a sensitive perturbation direction for a neural network Jacobian \(J\). The basic problem considered in the singular-vector attack literature is
\[
\max_{\|v\|_p=1}\|Jv\|_q.
\]
The answer to whether Lanczos or Golub--Kahan can be used depends crucially on whether
\[
p=q=2
\]
or not.

\subsection{The Case \(p=q=2\)}

When
\[
p=q=2,
\]
the problem becomes
\[
\max_{\|v\|_2=1}\|Jv\|_2.
\]
This is the standard largest singular value problem:
\[
\sigma_{\max}(J)
=
\max_{\|v\|_2=1}\|Jv\|_2.
\]
Equivalently,
\[
\|Jv\|_2^2
=
v^\top J^\top Jv,
\]
so the problem can be written as
\[
\max_{\|v\|_2=1}v^\top J^\top Jv.
\]
Therefore, the desired direction is the dominant eigenvector of
\[
J^\top J,
\]
which is the dominant right singular vector of \(J\).

In this case, ordinary power iteration is
\[
v_{k+1}
=
\frac{J^\top Jv_k}{\|J^\top Jv_k\|_2}.
\]
This can be computed without explicitly forming \(J\) or \(J^\top J\). One only needs the two matrix-free operations
\[
v\mapsto Jv
\]
and
\[
u\mapsto J^\top u.
\]
One iteration is therefore
\[
v_k
\longmapsto
Jv_k
\longmapsto
J^\top(Jv_k)
\longmapsto
v_{k+1}.
\]

Lanczos or Golub--Kahan methods can naturally replace ordinary power iteration in this setting. Lanczos can be applied to the symmetric positive semidefinite operator
\[
B=J^\top J,
\]
using only the matrix-free product
\[
Bv=J^\top(Jv).
\]
Golub--Kahan bidiagonalization is even more direct, because it works with the pair of operations
\[
Jv,\qquad J^\top u
\]
without explicitly forming \(J^\top J\).

Thus, for \(p=q=2\), one can use matrix-free Lanczos, Golub--Kahan, ARPACK-style partial SVD, or related Krylov methods. These methods do not require the explicit Jacobian matrix. They only require Jacobian-vector products and transposed-Jacobian-vector products.

\subsection{The Case \(p,q\neq 2\)}

For general \(p,q\), the problem is
\[
\max_{\|v\|_p=1}\|Jv\|_q.
\]
This is an \(\ell_p\to\ell_q\) operator norm maximization problem, not an ordinary singular value problem.

The generalized power method used in the \((p,q)\)-singular-vector setting has the form
\[
v_{k+1}
=
\frac{
\psi_{p'}
\left(
J^\top \psi_q(Jv_k)
\right)
}{
\left\|
\psi_{p'}
\left(
J^\top \psi_q(Jv_k)
\right)
\right\|_p
},
\]
where
\[
\psi_r(z)=\operatorname{sign}(z)|z|^{r-1},
\]
and
\[
\frac{1}{p}+\frac{1}{p'}=1.
\]

The key point is that, when \(p,q\neq 2\), the maps
\[
\psi_q(Jv_k)
\]
and
\[
\psi_{p'}(\cdot)
\]
are nonlinear. Therefore the update is not of the form
\[
v_{k+1}\propto Bv_k
\]
for some fixed linear operator \(B\). Instead, it is a nonlinear fixed-point iteration
\[
v_{k+1}=T(v_k),
\]
where
\[
T(v)
=
\frac{
\psi_{p'}
\left(
J^\top \psi_q(Jv)
\right)
}{
\left\|
\psi_{p'}
\left(
J^\top \psi_q(Jv)
\right)
\right\|_p
}.
\]

Standard Lanczos and Golub--Kahan methods require a fixed linear operator. For example, Lanczos constructs a Krylov subspace of the form
\[
\mathcal{K}_m(B,v_0)
=
\operatorname{span}
\{v_0,Bv_0,B^2v_0,\dots,B^{m-1}v_0\},
\]
which only makes sense when \(B\) is a fixed linear operator. In the general \((p,q)\) setting, the nonlinear maps \(\psi_q\) and \(\psi_{p'}\) destroy this fixed linear Krylov structure.

Therefore,
\[
\boxed{
\text{for }p,q\neq 2,\text{ standard Lanczos or Golub--Kahan cannot directly replace the generalized power method.}
}
\]

This is especially relevant for the setting used in the SingularFool paper, where the experiments often use
\[
p=\infty,\qquad q=10.
\]
In that case,
\[
p'=1,
\]
so
\[
\psi_{p'}=\psi_1=\operatorname{sign},
\]
and
\[
\psi_{10}(z)=\operatorname{sign}(z)|z|^9.
\]
The update becomes a nonlinear sign-type fixed-point iteration, not a standard linear SVD iteration.

\subsection{Matrix-Free SVD and the Role of SciPy \texttt{svds}}

If one works in the \(p=q=2\) setting, it is possible to use partial SVD methods without explicitly constructing the Jacobian. The correct interface is not a dense matrix \(J\), but a matrix-free linear operator.

Conceptually, one defines an operator \(A\) by specifying
\[
\operatorname{matvec}(v)=Jv
\]
and
\[
\operatorname{rmatvec}(u)=J^\top u.
\]
Then a partial SVD method can use these two operations to estimate the dominant singular value and singular vectors.

For example, in SciPy notation, this idea would look like
\[
A_{\mathrm{op}}
=
\operatorname{LinearOperator}
\left(
\text{shape}=(m,n),
\quad
\text{matvec}=v\mapsto Jv,
\quad
\text{rmatvec}=u\mapsto J^\top u
\right).
\]
Then a routine such as \texttt{svds} can, in principle, be applied to \(A_{\mathrm{op}}\).

However, this is only meaningful if \(Jv\) and \(J^\top u\) can be evaluated without forming \(J\). If the method first explicitly constructs the full Jacobian matrix and then passes it to an SVD routine, it violates the matrix-free requirement.

Therefore, the correct statement is:
\[
\boxed{
\texttt{svds} is only relevant here through a matrix-free \texttt{LinearOperator}; explicitly forming }J\text{ is not acceptable.}
\]

In neural network applications, there is also a practical engineering issue. SciPy routines usually operate on CPU-side NumPy arrays, while neural-network Jacobian-vector and vector-Jacobian products may live inside PyTorch or JAX on GPU. Passing data back and forth between GPU tensors and NumPy arrays can be inefficient. Therefore, although SciPy's \texttt{LinearOperator} interface is mathematically compatible with matrix-free SVD, it may not be the best engineering choice for large neural-network Jacobians. A native PyTorch or JAX implementation of matrix-free Lanczos or Golub--Kahan is often more appropriate.

\subsection{Practical Conclusion}

There are three distinct cases.

First, if the goal is the ordinary spectral direction
\[
\max_{\|v\|_2=1}\|Jv\|_2,
\]
then Lanczos, Golub--Kahan, ARPACK-style partial SVD, or related Krylov methods are appropriate. They can be implemented matrix-free using only
\[
Jv
\]
and
\[
J^\top u.
\]
No explicit Jacobian matrix is required.

Second, if the goal is the original generalized \((p,q)\)-singular-vector problem
\[
\max_{\|v\|_p=1}\|Jv\|_q
\]
with \(p,q\neq 2\), then standard Lanczos and Golub--Kahan cannot directly solve the problem, because the generalized power update contains nonlinear maps
\[
\psi_q
\quad\text{and}\quad
\psi_{p'}.
\]
The iteration is no longer a fixed linear Krylov method.

Third, one may still use a \(p=q=2\) spectral method as an initialization for the general \((p,q)\) problem. For example, one can first compute an approximate dominant \(\ell_2\to\ell_2\) singular direction by matrix-free Lanczos or Golub--Kahan, then normalize or transform it to satisfy the desired \(p\)-norm constraint, and finally continue with the generalized power method:
\[
v_{k+1}
=
\frac{
\psi_{p'}
\left(
J^\top \psi_q(Jv_k)
\right)
}{
\left\|
\psi_{p'}
\left(
J^\top \psi_q(Jv_k)
\right)
\right\|_p
}.
\]
In this hybrid approach, Lanczos is used only as a spectral initializer, not as a direct solver for the general \((p,q)\)-singular-vector problem.

The final conclusion is:
\[
\boxed{
p=q=2:
\text{ matrix-free Lanczos/Golub--Kahan is valid and often faster or more stable than ordinary power iteration.}
}
\]
\[
\boxed{
p,q\neq 2:
\text{ standard Lanczos/Golub--Kahan cannot directly replace the generalized power method.}
}
\]
\[
\boxed{
\text{Explicitly forming the Jacobian is unnecessary and should be avoided; all valid approaches must use }Jv\text{ and }J^\top u\text{ matrix-free.}
}
\]

\section{Sparse and Transferable Universal Singular Vectors Attack}

The paper \emph{Sparse and Transferable Universal Singular Vectors Attack} studies universal adversarial perturbations for image classification models. Its main goal is to construct a single perturbation vector \(v\) that can be added to many input images and cause the classifier output to change. Unlike dense universal perturbations, this paper focuses on sparse perturbations, where only a small number of pixels or patches are modified.

The setting is mainly image classification. The experiments are carried out on ImageNet pretrained classification models, including architectures such as ResNet, DenseNet, EfficientNet, Inception, ViT, and DEIT. The main evaluation metric is the fooling rate,
\[
    \mathrm{FR}(v)
    =
    \mathbb{P}_{x}
    \left[
        C(x+v)\neq C(x)
    \right],
\]
where \(C(x)\) is the original classifier prediction. The paper also reports attack success rate, which measures how often originally correctly classified samples become misclassified after adding the perturbation. Therefore, the concrete experimental setting of the paper is image classification adversarial attack.

The method is based on the earlier universal singular vector attack of Khrulkov and Oseledets. In that framework, one considers a hidden-layer representation
\[
    f_l(x),
\]
and studies how a small input perturbation \(v\) changes this representation:
\[
    f_l(x+v)-f_l(x).
\]
Using a first-order Taylor approximation,
\[
    f_l(x+v)-f_l(x)
    \approx
    J_l(x)v,
\]
where
\[
    J_l(x)
    =
    \frac{\partial f_l(x)}{\partial x}
\]
is the Jacobian of the hidden-layer feature map with respect to the input. The idea is that if a perturbation direction \(v\) is strongly amplified by \(J_l(x)\), then it can significantly change the internal representation of the network and eventually change the final classification result.

The dense universal singular vector attack can be written schematically as
\[
    \max_{\|v\|_p\leq 1}
    \frac{1}{N}
    \sum_{i=1}^{N}
    \|J_i v\|_q^q,
\]
where \(J_i = J_l(x_i)\) is the hidden-layer Jacobian for the input sample \(x_i\). The goal is to find a universal direction \(v\) that is amplified by many hidden-layer Jacobians simultaneously.

The sparse version studied in this paper adds an \(\ell_0\)-type sparsity constraint:
\[
    \max_{\|v\|_p\leq 1,\ \|v\|_0\leq k}
    \frac{1}{N}
    \sum_{i=1}^{N}
    \|J_i v\|_q^q.
\]
Here,
\[
    \|v\|_0\leq k
\]
means that the perturbation is allowed to have at most \(k\) nonzero entries, pixels, or patches. Thus, the paper is not only looking for a universal sensitive direction, but for a sparse universal sensitive direction.

The main algorithmic idea is to take the generalized power iteration used in the dense universal singular vector method and insert a truncation step. In this sense, the method can be understood as
\[
    \text{generalized power iteration}
    \quad+\quad
    \text{top-}k\text{ truncation}.
\]
After each power-type update, the algorithm keeps only the entries, pixels, or patches with the largest magnitudes and sets the rest to zero. A simplified form is
\[
    v_{t+1}
    =
    \operatorname{Normalize}_{p}
    \left(
        \operatorname{Truncate}_{k}
        \left(
            \widetilde v_{t+1}
        \right)
    \right),
\]
where \(\widetilde v_{t+1}\) is the vector obtained from the generalized power update before sparsity is enforced.

In the special case \(p=q=2\), the dense problem reduces to the standard largest singular vector problem:
\[
    \max_{\|v\|_2=1}
    \|Jv\|_2.
\]
Equivalently,
\[
    \max_{\|v\|_2=1}
    \|Jv\|_2^2
    =
    \max_{\|v\|_2=1}
    v^\top J^\top J v.
\]
The corresponding power iteration is
\[
    v_{t+1}
    \propto
    J^\top J v_t.
\]
With sparsity, this becomes
\[
    v_{t+1}
    =
    \operatorname{Normalize}_{2}
    \left(
        \operatorname{Truncate}_{k}
        \left(
            J^\top J v_t
        \right)
    \right).
\]
Therefore, for \(p=q=2\), the method is essentially ordinary power iteration with an additional top-\(k\) truncation step.

For general \(p\) and \(q\), however, the update is not simply
\[
    J^\top Jv.
\]
The expression \(J^\top Jv\) is specific to the Euclidean case \(p=q=2\). For a general \(q\)-norm objective, one must account for the nonlinear derivative of the \(q\)-norm. If
\[
    z = Jv,
\]
then
\[
    \nabla_z \|z\|_q^q
    =
    q\,\Phi_q(z),
\]
where, for \(1<q<\infty\),
\[
    \Phi_q(z)
    =
    |z|^{q-2}z
\]
entrywise. Therefore, the gradient-like term with respect to \(v\) is
\[
    J^\top \Phi_q(Jv),
\]
not simply
\[
    J^\top Jv.
\]
Only when \(q=2\) does
\[
    \Phi_2(z)=z
\]
hold, so that
\[
    J^\top \Phi_2(Jv)=J^\top Jv.
\]

Similarly, if \(p\neq 2\), the update of \(v\) is not just Euclidean normalization. Let \(p'\) denote the dual exponent of \(p\):
\[
    \frac{1}{p}+\frac{1}{p'}=1.
\]
If
\[
    g_t
    =
    J^\top \Phi_q(Jv_t),
\]
then the maximizer of the linearized objective over the \(\ell_p\) unit ball is related to the duality map associated with \(p'\). A typical generalized update can be written as
\[
    v_{t+1}
    \propto
    \Phi_{p'}(g_t),
\]
followed by normalization in the \(\ell_p\) norm. Entrywise, this has the form
\[
    \Phi_{p'}(g)
    =
    \operatorname{sign}(g)|g|^{p'-1}.
\]
Thus, for general \((p,q)\), the generalized power-type update has the schematic form
\[
    v_t
    \longmapsto
    Jv_t
    \longmapsto
    \Phi_q(Jv_t)
    \longmapsto
    J^\top \Phi_q(Jv_t)
    \longmapsto
    \Phi_{p'}
    \left(
        J^\top \Phi_q(Jv_t)
    \right).
\]
With sparsity, this becomes
\[
    v_{t+1}
    =
    \operatorname{Normalize}_{p}
    \left(
        \operatorname{Truncate}_{k}
        \left[
            \Phi_{p'}
            \left(
                J^\top \Phi_q(Jv_t)
            \right)
        \right]
    \right).
\]

The paper itself does not use the notation
\[
    \Phi_q,\qquad \Phi_{p'}
\]
as the main way of presenting the algorithm. Instead, it derives the update using the dual norm formulation and alternating maximization. The paper explicitly discusses \((p,q)\)-singular vectors and states that the method is formulated for arbitrary pairs \((p,q)\), but the original text does not literally say: first apply \(\Phi_q\), then multiply by \(J^\top\), then apply \(\Phi_{p'}\). That notation is a standard way to rewrite the same generalized power iteration idea.

The paper's derivation goes through the dual norm identity. For a norm \(\|\cdot\|_q\), one can write
\[
    \|a\|_q
    =
    \max_{\|u\|_{q'}\leq 1}
    \langle u,a\rangle,
\]
where
\[
    \frac{1}{q}+\frac{1}{q'}=1.
\]
Using this identity, the hidden-layer Jacobian objective can be rewritten as an alternating optimization problem involving the perturbation direction and dual variables. Fixing one variable gives a closed-form update for the other. The sparsity constraint is then enforced by introducing a binary mask or a truncation operator that keeps the largest-magnitude components.

Therefore, the most accurate interpretation is:
\[
    \boxed{
    \text{The paper formulates the sparse attack for general }(p,q)\text{-singular vectors.}
    }
\]
\[
    \boxed{
    \text{It uses dual norms and alternating maximization rather than explicitly writing }\Phi_q
    \text{ and }\Phi_{p'}.
    }
\]
\[
    \boxed{
    \text{For }p=q=2\text{, the generalized update reduces to }J^\top Jv.
    }
\]
\[
    \boxed{
    \text{For general }p,q\text{, the update is nonlinear and cannot be reduced to }J^\top Jv.
    }
\]

This also clarifies the relationship between this sparse method and classical numerical linear algebra methods. If the goal is only to compute the leading singular direction in the standard Euclidean setting,
\[
    \max_{\|v\|_2=1}\|Jv\|_2,
\]
then methods such as power iteration, Lanczos, randomized SVD, and Golub--Kahan bidiagonalization are the natural choices. These methods are designed for the standard \(p=q=2\) singular value problem and can operate through matrix-vector products such as \(Jv\) and \(J^\top u\), without explicitly forming the full Jacobian.

The method in \emph{Sparse and Transferable Universal Singular Vectors Attack} solves a different problem:
\[
    \max_{\|v\|_p\leq 1,\ \|v\|_0\leq k}
    \frac{1}{N}
    \sum_{i=1}^{N}
    \|J_i v\|_q^q.
\]
The additional sparsity constraint makes the problem nonconvex and combinatorial. Even in the case \(p=q=2\), the presence of
\[
    \|v\|_0\leq k
\]
turns the problem into a sparse singular vector or sparse PCA-type problem rather than a standard SVD problem. Therefore, the truncated power iteration used in the paper is better understood as a heuristic sparse generalized power method, not as a replacement for Golub--Kahan or Lanczos in the dense Euclidean case.

The practical importance of the paper depends on the goal. If the goal is to find the fastest dense maximum sensitivity direction, then the standard numerical linear algebra methods remain more fundamental. If the goal is to find a perturbation that is sparse, universal, and transferable across classification models, then the method of this paper is directly relevant.

In summary, the paper can be positioned as follows:
\[
    \text{Khrulkov--Oseledets universal singular vector attack}
    \quad+\quad
    \text{sparsity constraint}
    \quad+\quad
    \text{truncated power iteration}.
\]
Its central modification is simple but important: after the generalized power update, the perturbation is truncated so that only the largest-magnitude pixels or patches remain active. This makes the resulting universal perturbation sparse while preserving much of the attack strength and, according to the experiments, improving transferability in many image classification models.

\section{Adversarial Perturbation Objectives: PGD, Power Iteration, and Quadratic Loss Structures}

\subsection{Motivation}

We consider adversarial perturbations of an input \(x\). A model \(f\) maps the input to an output:
\[
    f: \mathbb{R}^d \to \mathbb{R}^m.
\]
Given a perturbation \(\delta\), the perturbed input is
\[
    x_\delta = x+\delta.
\]
The central question is how to define a meaningful adversarial loss and how the resulting optimization problem relates to gradient-based attacks such as PGD and spectral methods such as power iteration.

The key distinction is the following:

\[
\boxed{
\text{Power iteration is natural when the adversarial objective becomes a pure quadratic form.}
}
\]

In contrast,

\[
\boxed{
\text{PGD is natural for general nonlinear or affine-quadratic loss objectives.}
}
\]

More precisely, if the local adversarial objective has the form
\[
    \mathcal{L}(\delta)
    =
    \frac{1}{2}\delta^\top A\delta,
\]
then the maximization problem
\[
    \max_{\|\delta\|_2\le \varepsilon}
    \delta^\top A\delta
\]
is a Rayleigh quotient problem. The optimal direction is the top eigenvector of \(A\), and power iteration is directly applicable:
\[
    v_{k+1}
    =
    \frac{A v_k}{\|A v_k\|_2}.
\]

However, if the local adversarial objective has the form
\[
    \mathcal{L}(\delta)
    =
    b^\top \delta
    +
    \frac{1}{2}\delta^\top A\delta,
\]
then the objective is not a pure Rayleigh quotient. Its gradient is
\[
    \nabla_\delta \mathcal{L}
    =
    b+A\delta,
\]
and the optimal direction is generally not the top eigenvector of \(A\). In this case, a generic constrained gradient method such as PGD is more directly aligned with the true loss.

\subsection{Notation}

We use the following notation:
\[
    F_0 = f(x),
    \qquad
    F_\delta = f(x+\delta).
\]

If the ground truth is a fixed target \(y\), we write
\[
    Y = y.
\]

If the ground truth is itself a function or an oracle operator \(g\), then
\[
    G_0 = g(x),
    \qquad
    G_\delta = g(x+\delta).
\]

The Jacobian of the model is
\[
    J_f = \frac{\partial f(x)}{\partial x}.
\]

The Jacobian of the true function \(g\) is
\[
    J_g = \frac{\partial g(x)}{\partial x}.
\]

The local linearizations are
\[
    F_\delta
    =
    f(x+\delta)
    \approx
    F_0 + J_f\delta,
\]
and
\[
    G_\delta
    =
    g(x+\delta)
    \approx
    G_0 + J_g\delta.
\]

\subsection{PGD versus Power Iteration}

PGD-type methods optimize a loss directly. A standard projected gradient ascent step has the form
\[
    \delta_{k+1}
    =
    \Pi_{\mathcal{S}}
    \left(
        \delta_k + \alpha \nabla_\delta \mathcal{L}(\delta_k)
    \right),
\]
where \(\mathcal{S}\) is the allowed perturbation set, for example
\[
    \mathcal{S}
    =
    \{\delta:\|\delta\|_p\le \varepsilon\}.
\]

When the local loss is a pure quadratic form
\[
    \mathcal{L}(\delta)
    =
    \frac{1}{2}\delta^\top A\delta,
\]
the gradient is
\[
    \nabla_\delta \mathcal{L}
    =
    A\delta.
\]
In this case, the update direction is exactly the same operator used in power iteration:
\[
    \delta \mapsto A\delta.
\]

Power iteration uses
\[
    v_{k+1}
    =
    \frac{A v_k}{\|A v_k\|_2}.
\]

By contrast, a PGD-like normalized gradient step for the same quadratic objective may look like
\[
    v_{k+1}
    =
    \frac{v_k+\alpha A v_k}
    {\|v_k+\alpha A v_k\|_2}.
\]
This is equivalent to applying \(I+\alpha A\), not \(A\), so its convergence rate is governed by
\[
    \frac{1+\alpha\lambda_2}{1+\alpha\lambda_1},
\]
rather than
\[
    \frac{\lambda_2}{\lambda_1}.
\]
For \(\lambda_1>\lambda_2>0\), one generally has
\[
    \frac{1+\alpha\lambda_2}{1+\alpha\lambda_1}
    >
    \frac{\lambda_2}{\lambda_1}.
\]
Thus, for a pure fixed quadratic form, power iteration is usually more direct.

However, when the loss has an additional linear term,
\[
    \mathcal{L}(\delta)
    =
    b^\top \delta
    +
    \frac{1}{2}\delta^\top A\delta,
\]
the gradient becomes
\[
    \nabla_\delta \mathcal{L}
    =
    b+A\delta.
\]
The vector \(b\) changes the optimization direction. The top eigenvector of \(A\) is no longer the general optimal direction.

\subsection{Four Core Adversarial Losses}

After removing repeated or equivalent loss definitions, the meaningful losses discussed can be organized into four core types.

\begin{table}[h]
\centering
\begin{tabular}{p{0.8cm} p{4.4cm} p{4.8cm} p{4.2cm}}
\hline
No. & Loss & Meaning & Local structure \\
\hline
1 &
\(\displaystyle \frac{1}{2}\|F_\delta-F_0\|_2^2\)
&
Model output change
&
Pure quadratic:
\(\displaystyle \frac{1}{2}\delta^\top J_f^\top J_f\delta\)
\\[1.2em]

2 &
\(\displaystyle \frac{1}{2}\|F_\delta-Y\|_2^2\)
&
Prediction error against fixed target \(y\)
&
Linear plus quadratic:
\(\displaystyle e_y^\top J_f\delta+\frac{1}{2}\delta^\top J_f^\top J_f\delta\)
\\[1.2em]

3 &
\(\displaystyle \frac{1}{2}\|F_\delta-G_\delta\|_2^2\)
&
Prediction error against dynamic ground truth \(g(x+\delta)\)
&
Linear plus quadratic:
\(\displaystyle e_g^\top (J_f-J_g)\delta+\frac{1}{2}\delta^\top (J_f-J_g)^\top(J_f-J_g)\delta\)
\\[1.2em]

4 &
\(\displaystyle \frac{1}{2}\|(F_\delta-G_\delta)-(F_0-G_0)\|_2^2\)
&
Change of the error vector
&
Pure quadratic:
\(\displaystyle \frac{1}{2}\delta^\top (J_f-J_g)^\top(J_f-J_g)\delta\)
\\
\hline
\end{tabular}
\caption{Four core adversarial loss definitions and their local structures.}
\end{table}

Here
\[
    e_y = F_0-Y = f(x)-y,
\]
and
\[
    e_g = F_0-G_0 = f(x)-g(x).
\]

\subsection{Loss 1: Model Output Change}

The first loss is
\[
    \mathcal{L}_1(\delta)
    =
    \frac{1}{2}\|F_\delta-F_0\|_2^2
    =
    \frac{1}{2}\|f(x+\delta)-f(x)\|_2^2.
\]

This is the natural loss when no ground truth is available. It asks how much the model output changes under the input perturbation.

Using the linearization
\[
    f(x+\delta)-f(x)
    \approx
    J_f\delta,
\]
we obtain
\[
    \mathcal{L}_1(\delta)
    \approx
    \frac{1}{2}\|J_f\delta\|_2^2.
\]
Therefore,
\[
    \mathcal{L}_1(\delta)
    \approx
    \frac{1}{2}\delta^\top J_f^\top J_f\delta.
\]

Define
\[
    A_1 = J_f^\top J_f.
\]
Then
\[
    \mathcal{L}_1(\delta)
    \approx
    \frac{1}{2}\delta^\top A_1\delta.
\]

The gradient is
\[
    \nabla_\delta \mathcal{L}_1
    =
    A_1\delta
    =
    J_f^\top J_f\delta.
\]

Under the constraint
\[
    \|\delta\|_2\le \varepsilon,
\]
maximizing \(\mathcal{L}_1\) is equivalent to the Rayleigh quotient problem
\[
    \max_{\|\delta\|_2\le \varepsilon}
    \delta^\top J_f^\top J_f\delta.
\]

The optimal direction is the top right singular vector of \(J_f\). Therefore, this loss is naturally compatible with power iteration:
\[
    v_{k+1}
    =
    \frac{J_f^\top J_f v_k}
    {\|J_f^\top J_f v_k\|_2}.
\]

\subsection{Loss 2: Error Against a Fixed Ground Truth}

Suppose a fixed ground truth \(Y=y\) is available. The natural supervised regression attack loss is
\[
    \mathcal{L}_2(\delta)
    =
    \frac{1}{2}\|F_\delta-Y\|_2^2
    =
    \frac{1}{2}\|f(x+\delta)-y\|_2^2.
\]

Define the original residual
\[
    e_y = F_0-Y = f(x)-y.
\]

Using
\[
    F_\delta
    \approx
    F_0+J_f\delta,
\]
we obtain
\[
    F_\delta-Y
    \approx
    e_y+J_f\delta.
\]

Thus,
\[
    \mathcal{L}_2(\delta)
    \approx
    \frac{1}{2}\|e_y+J_f\delta\|_2^2.
\]

Expanding,
\[
    \mathcal{L}_2(\delta)
    \approx
    \frac{1}{2}e_y^\top e_y
    +
    e_y^\top J_f\delta
    +
    \frac{1}{2}\delta^\top J_f^\top J_f\delta.
\]

The constant term
\[
    \frac{1}{2}e_y^\top e_y
\]
does not affect the optimization over \(\delta\). Therefore,
\[
    \mathcal{L}_2(\delta)
    \sim
    e_y^\top J_f\delta
    +
    \frac{1}{2}\delta^\top J_f^\top J_f\delta.
\]

Equivalently,
\[
    \mathcal{L}_2(\delta)
    \sim
    (J_f^\top e_y)^\top \delta
    +
    \frac{1}{2}\delta^\top J_f^\top J_f\delta.
\]

Let
\[
    b_2 = J_f^\top e_y,
    \qquad
    A_2 = J_f^\top J_f.
\]
Then
\[
    \mathcal{L}_2(\delta)
    \sim
    b_2^\top \delta
    +
    \frac{1}{2}\delta^\top A_2\delta.
\]

The gradient is
\[
    \nabla_\delta \mathcal{L}_2
    =
    b_2 + A_2\delta.
\]
That is,
\[
    \nabla_\delta \mathcal{L}_2
    =
    J_f^\top e_y
    +
    J_f^\top J_f\delta.
\]

This is not generally a pure Rayleigh quotient problem. The residual term
\[
    J_f^\top e_y
\]
pulls the update direction toward the direction that increases the already existing prediction error. Therefore, the top right singular vector of \(J_f\) is not generally the optimal direction for this loss.

Only when
\[
    e_y \approx 0
\]
does the linear term become negligible. In that special case,
\[
    \mathcal{L}_2(\delta)
    \approx
    \frac{1}{2}\delta^\top J_f^\top J_f\delta,
\]
and the problem reduces approximately to Loss 1.

\subsection{Loss 3: Error Against a Dynamic Ground Truth Operator}

Now suppose the ground truth is itself a function or oracle operator \(g\). The natural loss is
\[
    \mathcal{L}_3(\delta)
    =
    \frac{1}{2}\|F_\delta-G_\delta\|_2^2
    =
    \frac{1}{2}\|f(x+\delta)-g(x+\delta)\|_2^2.
\]

Define the original error
\[
    e_g=F_0-G_0=f(x)-g(x).
\]

Using the linearizations
\[
    F_\delta\approx F_0+J_f\delta,
\]
and
\[
    G_\delta\approx G_0+J_g\delta,
\]
we get
\[
    F_\delta-G_\delta
    \approx
    e_g+(J_f-J_g)\delta.
\]

Let
\[
    K=J_f-J_g.
\]
Then
\[
    F_\delta-G_\delta
    \approx
    e_g+K\delta.
\]

Therefore,
\[
    \mathcal{L}_3(\delta)
    \approx
    \frac{1}{2}\|e_g+K\delta\|_2^2.
\]

Expanding,
\[
    \mathcal{L}_3(\delta)
    \approx
    \frac{1}{2}e_g^\top e_g
    +
    e_g^\top K\delta
    +
    \frac{1}{2}\delta^\top K^\top K\delta.
\]

Ignoring the constant term,
\[
    \mathcal{L}_3(\delta)
    \sim
    e_g^\top K\delta
    +
    \frac{1}{2}\delta^\top K^\top K\delta.
\]

Equivalently,
\[
    \mathcal{L}_3(\delta)
    \sim
    (K^\top e_g)^\top \delta
    +
    \frac{1}{2}\delta^\top K^\top K\delta.
\]

Let
\[
    b_3 = K^\top e_g
    =
    (J_f-J_g)^\top [f(x)-g(x)],
\]
and
\[
    A_3 = K^\top K
    =
    (J_f-J_g)^\top(J_f-J_g).
\]
Then
\[
    \mathcal{L}_3(\delta)
    \sim
    b_3^\top \delta
    +
    \frac{1}{2}\delta^\top A_3\delta.
\]

The gradient is
\[
    \nabla_\delta \mathcal{L}_3
    =
    b_3+A_3\delta.
\]
Explicitly,
\[
    \nabla_\delta \mathcal{L}_3
    =
    (J_f-J_g)^\top [f(x)-g(x)]
    +
    (J_f-J_g)^\top(J_f-J_g)\delta.
\]

Thus, Loss 3 is generally not a pure quadratic form. Its direction is affected by the original error vector
\[
    e_g=f(x)-g(x).
\]

Only when
\[
    e_g\approx 0
\]
does it reduce approximately to
\[
    \mathcal{L}_3(\delta)
    \approx
    \frac{1}{2}\delta^\top
    (J_f-J_g)^\top(J_f-J_g)\delta.
\]

In that special case, the top right singular vector of \(J_f-J_g\) becomes the relevant direction.

\subsection{Loss 4: Error-Change or Response-Mismatch Loss}

A different quantity is
\[
    \mathcal{L}_4(\delta)
    =
    \frac{1}{2}
    \|(F_\delta-G_\delta)-(F_0-G_0)\|_2^2.
\]

Equivalently,
\[
    \mathcal{L}_4(\delta)
    =
    \frac{1}{2}
    \|(F_\delta-F_0)-(G_\delta-G_0)\|_2^2.
\]

This loss measures how much the model error vector changes from \(x\) to \(x+\delta\). It is also the difference between the model's response change and the true operator's response change.

Using
\[
    F_\delta-F_0 \approx J_f\delta,
\]
and
\[
    G_\delta-G_0 \approx J_g\delta,
\]
we obtain
\[
    (F_\delta-F_0)-(G_\delta-G_0)
    \approx
    (J_f-J_g)\delta.
\]

Thus,
\[
    \mathcal{L}_4(\delta)
    \approx
    \frac{1}{2}\|(J_f-J_g)\delta\|_2^2.
\]

Therefore,
\[
    \mathcal{L}_4(\delta)
    \approx
    \frac{1}{2}\delta^\top
    (J_f-J_g)^\top(J_f-J_g)\delta.
\]

Let
\[
    A_4=(J_f-J_g)^\top(J_f-J_g).
\]
Then
\[
    \mathcal{L}_4(\delta)
    \approx
    \frac{1}{2}\delta^\top A_4\delta.
\]

The gradient is
\[
    \nabla_\delta \mathcal{L}_4
    =
    A_4\delta.
\]
That is,
\[
    \nabla_\delta \mathcal{L}_4
    =
    (J_f-J_g)^\top(J_f-J_g)\delta.
\]

This is a pure quadratic form. Therefore, if the goal is to maximize the change of the error vector, then power iteration is appropriate:
\[
    v_{k+1}
    =
    \frac{
    (J_f-J_g)^\top(J_f-J_g)v_k
    }
    {
    \|(J_f-J_g)^\top(J_f-J_g)v_k\|_2
    }.
\]

However, this loss should not be confused with Loss 3. Loss 3 measures the final error magnitude:
\[
    \|F_\delta-G_\delta\|.
\]
Loss 4 measures the change in the error vector:
\[
    \|(F_\delta-G_\delta)-(F_0-G_0)\|.
\]

Thus, Loss 4 is not the correct main objective if the goal is to make the final prediction error larger. It is appropriate only if the goal is to study the sensitivity of the error function itself.

\subsection{Geometric Interpretation}

Let
\[
    E(x)=f(x)-g(x)
\]
be the error function. Then
\[
    E_0=E(x),
    \qquad
    E_\delta=E(x+\delta).
\]

Locally,
\[
    E_\delta
    \approx
    E_0+\Delta E,
\]
where
\[
    \Delta E
    =
    (J_f-J_g)\delta.
\]

Loss 3 is
\[
    \mathcal{L}_3(\delta)
    =
    \frac{1}{2}\|E_\delta\|_2^2
    \approx
    \frac{1}{2}\|E_0+\Delta E\|_2^2.
\]

Loss 4 is
\[
    \mathcal{L}_4(\delta)
    =
    \frac{1}{2}\|E_\delta-E_0\|_2^2
    \approx
    \frac{1}{2}\|\Delta E\|_2^2.
\]

The difference is geometric. We have
\[
    \|E_0+\Delta E\|_2^2
    =
    \|E_0\|_2^2
    +
    2E_0^\top \Delta E
    +
    \|\Delta E\|_2^2.
\]

The term
\[
    2E_0^\top \Delta E
\]
depends on the angle between the original error \(E_0\) and the perturbation-induced change \(\Delta E\).

Let
\[
    \theta = \angle(E_0,\Delta E).
\]
Then
\[
    E_0^\top \Delta E
    =
    \|E_0\|_2\|\Delta E\|_2\cos\theta.
\]

Therefore,
\[
    \|E_0+\Delta E\|_2^2
    =
    \|E_0\|_2^2
    +
    2\|E_0\|_2\|\Delta E\|_2\cos\theta
    +
    \|\Delta E\|_2^2.
\]

If
\[
    \theta=0,
\]
then \(\Delta E\) is aligned with \(E_0\), and
\[
    \|E_0+\Delta E\|_2
    =
    \|E_0\|_2+\|\Delta E\|_2.
\]
In this case, making \(\|\Delta E\|\) large also increases the final error.

If
\[
    \theta=90^\circ,
\]
then
\[
    E_0^\top \Delta E=0,
\]
so \(\Delta E\) changes the error direction but gives no first-order increase in the error length.

If
\[
    \theta=180^\circ,
\]
then \(\Delta E\) points opposite to \(E_0\), and it may reduce the final error.

Thus,
\[
\boxed{
\max \|\Delta E\|
\neq
\max \|E_0+\Delta E\|
}
\]
in general.

\subsection{Why Loss 4 Can Be Misleading as an Attack Objective}

Loss 4 maximizes
\[
    \|E_\delta-E_0\|.
\]
This means it maximizes how far the error vector moves.

However, the final error magnitude is
\[
    \|E_\delta\|.
\]

It is possible that
\[
    \|E_\delta-E_0\|
\]
is large while
\[
    \|E_\delta\|
\]
does not increase.

For example, let
\[
    E_0=
    \begin{pmatrix}
    1\\
    0
    \end{pmatrix},
    \qquad
    E_\delta=
    \begin{pmatrix}
    -1\\
    0
    \end{pmatrix}.
\]
Then
\[
    \|E_0\|_2=1,
    \qquad
    \|E_\delta\|_2=1.
\]
The final error magnitude does not increase. But
\[
    \|E_\delta-E_0\|_2
    =
    2.
\]

Thus, Loss 4 can be large even when the final prediction error remains unchanged. It may simply rotate or move the error vector around a sphere of the same radius.

Therefore, Loss 4 is not the appropriate main loss if the goal is to make the model's final prediction error larger. It is appropriate if the goal is to measure sensitivity of the error function.

\subsection{Final Classification}

The losses can be classified as follows.

\begin{table}[h]
\centering
\begin{tabular}{p{0.9cm} p{4.2cm} p{4.4cm} p{3.2cm} p{2.4cm}}
\hline
No. & Loss & Main purpose & Local form & Power iteration? \\
\hline
1 &
\(\frac{1}{2}\|f(x+\delta)-f(x)\|^2\)
&
Model output sensitivity when no ground truth is available
&
Pure quadratic
&
Yes
\\[1.2em]

2 &
\(\frac{1}{2}\|f(x+\delta)-y\|^2\)
&
Increase error against fixed ground truth
&
Linear + quadratic
&
Generally no
\\[1.2em]

3 &
\(\frac{1}{2}\|f(x+\delta)-g(x+\delta)\|^2\)
&
Increase error against dynamic ground truth operator
&
Linear + quadratic
&
Generally no
\\[1.2em]

4 &
\(\frac{1}{2}\|[f(x+\delta)-g(x+\delta)]-[f(x)-g(x)]\|^2\)
&
Measure change of the error vector
&
Pure quadratic
&
Yes, but only for this sensitivity objective
\\
\hline
\end{tabular}
\caption{Classification of the four discussed losses.}
\end{table}

For the actual attack goal of making the final error larger, the appropriate losses are:
\[
    \mathcal{L}_1(\delta)
    =
    \frac{1}{2}\|f(x+\delta)-f(x)\|^2
\]
when no ground truth is available,
\[
    \mathcal{L}_2(\delta)
    =
    \frac{1}{2}\|f(x+\delta)-y\|^2
\]
when a fixed ground truth \(y\) is available, and
\[
    \mathcal{L}_3(\delta)
    =
    \frac{1}{2}\|f(x+\delta)-g(x+\delta)\|^2
\]
when the ground truth is a function \(g\).

Only \(\mathcal{L}_1\) is naturally a pure quadratic form under local linearization. The supervised losses \(\mathcal{L}_2\) and \(\mathcal{L}_3\) generally contain a residual-induced linear term, so they are not generally Rayleigh quotient problems.

\subsection{Main Conclusion}

The main conclusion is:

\[
\boxed{
\text{Power iteration is appropriate only when the adversarial objective reduces to a pure quadratic form.}
}
\]

For example,
\[
    \frac{1}{2}\|J_f\delta\|^2
    =
    \frac{1}{2}\delta^\top J_f^\top J_f\delta
\]
or
\[
    \frac{1}{2}\|(J_f-J_g)\delta\|^2
    =
    \frac{1}{2}\delta^\top (J_f-J_g)^\top(J_f-J_g)\delta.
\]

In such cases, the maximizer under an \(L_2\) constraint is the dominant eigenvector of the corresponding Gram matrix.

However, when the objective is
\[
    \frac{1}{2}\|e+J\delta\|^2,
\]
its local expansion is
\[
    e^\top J\delta
    +
    \frac{1}{2}\delta^\top J^\top J\delta.
\]
The gradient is
\[
    J^\top e + J^\top J\delta.
\]
The additional term \(J^\top e\) encodes the direction of the pre-existing residual. Therefore, the optimal direction is not determined only by the top singular vector of \(J\).

Thus:

\[
\boxed{
\text{PGD directly optimizes the actual loss, while power iteration optimizes a pure local amplification surrogate.}
}
\]

Consequently, power iteration can be faster and more natural for objectives based on output change or local sensitivity, but it does not generally replace PGD for supervised adversarial losses that include a residual term.

\section{Further Clarification: Gradients with Respect to the Input and the Perturbation}

\subsection{Gradient with Respect to \(a\) versus Gradient with Respect to \(\delta\)}

Let the input be written as
\[
    a = a_0+\delta,
\]
where \(a_0\) is the original input and \(\delta\) is the adversarial perturbation.

For the full oracle loss
\[
    \ell(a)
    =
    \|G(a)-g(a)\|_2^2,
\]
define the error function
\[
    e(a)=G(a)-g(a).
\]
Then
\[
    \ell(a)=\|e(a)\|_2^2.
\]

The Fréchet derivative or ordinary discrete gradient with respect to \(a\) is
\[
    \nabla_a \ell(a)
    =
    2\left(J_G(a)-J_g(a)\right)^\top e(a),
\]
where
\[
    J_G(a)=\frac{\partial G(a)}{\partial a},
    \qquad
    J_g(a)=\frac{\partial g(a)}{\partial a}.
\]

Since
\[
    a=a_0+\delta,
\]
we have
\[
    \frac{\partial a}{\partial \delta}=I.
\]
Therefore, differentiating with respect to \(a\) or \(\delta\) is not fundamentally different:
\[
    \nabla_\delta \ell(a_0+\delta)
    =
    \nabla_a \ell(a)\big|_{a=a_0+\delta}.
\]

Thus, the exact current-point gradient can be written as
\[
    \nabla_\delta \ell(\delta)
    =
    2\left(J_G(a_0+\delta)-J_g(a_0+\delta)\right)^\top
    \left(G(a_0+\delta)-g(a_0+\delta)\right).
\]

The formula differs from the local expression only because the local expression freezes the Jacobians at the base point \(a_0\).

\subsection{Deriving the Local Perturbation Gradient Step by Step}

At the base point \(a_0\), define
\[
    e_0=G(a_0)-g(a_0),
\]
and
\[
    K=J_G(a_0)-J_g(a_0).
\]

Linearize both \(G\) and \(g\) around \(a_0\):
\[
    G(a_0+\delta)
    \approx
    G(a_0)+J_G(a_0)\delta,
\]
and
\[
    g(a_0+\delta)
    \approx
    g(a_0)+J_g(a_0)\delta.
\]

Therefore,
\[
    e(a_0+\delta)
    =
    G(a_0+\delta)-g(a_0+\delta)
\]
is approximated by
\[
    e(a_0+\delta)
    \approx
    e_0+
    \left(J_G(a_0)-J_g(a_0)\right)\delta.
\]
Using the notation \(K=J_G(a_0)-J_g(a_0)\), this becomes
\[
    e(a_0+\delta)
    \approx
    e_0+K\delta.
\]

The loss as a function of the perturbation is
\[
    \ell(\delta)
    =
    \|e(a_0+\delta)\|_2^2.
\]
Using the linear approximation,
\[
    \ell(\delta)
    \approx
    \|e_0+K\delta\|_2^2.
\]

Expanding,
\[
    \ell(\delta)
    \approx
    e_0^\top e_0
    +
    2e_0^\top K\delta
    +
    \delta^\top K^\top K\delta.
\]

Taking the gradient with respect to \(\delta\),
\[
    \nabla_\delta \ell(\delta)
    \approx
    2K^\top e_0
    +
    2K^\top K\delta.
\]

Substituting
\[
    K=J_G(a_0)-J_g(a_0),
\]
we obtain
\[
    \boxed{
    \nabla_\delta \ell(\delta)
    \approx
    2\left(J_G-J_g\right)^\top e_0
    +
    2\left(J_G-J_g\right)^\top
    \left(J_G-J_g\right)\delta
    }.
\]

This is the local linearized form of the exact current-point gradient
\[
    \nabla_a \ell(a)
    =
    2\left(J_G(a)-J_g(a)\right)^\top e(a).
\]

The two expressions are therefore consistent. The first is exact at the current point \(a\), while the second is the local approximation around \(a_0\).

\section{Fixed Ground Truth versus Dynamic Ground Truth}

\subsection{Two Different Supervised Losses}

There are two important supervised settings.

\paragraph{Fixed target.}
The target is a fixed value \(y\), independent of \(a\):
\[
    \ell_y(a)
    =
    \|G(a)-y\|_2^2.
\]

\paragraph{Dynamic target.}
The target is a function or oracle \(g(a)\), which also changes with the input:
\[
    \ell_g(a)
    =
    \|G(a)-g(a)\|_2^2.
\]

These two losses may look similar, but their gradients are different because \(g(a)\) has its own derivative with respect to \(a\).

\subsection{Direct Gradients}

For the fixed-target loss, define
\[
    e_y(a)=G(a)-y.
\]
Since \(y\) is fixed,
\[
    \frac{\partial y}{\partial a}=0.
\]
Therefore,
\[
    \boxed{
    \nabla_a \ell_y(a)
    =
    2J_G(a)^\top [G(a)-y]
    }.
\]

For the dynamic-target loss, define
\[
    e_g(a)=G(a)-g(a).
\]
Since \(g(a)\) depends on \(a\),
\[
    \frac{\partial g(a)}{\partial a}=J_g(a).
\]
Therefore,
\[
    \boxed{
    \nabla_a \ell_g(a)
    =
    2\left(J_G(a)-J_g(a)\right)^\top [G(a)-g(a)]
    }.
\]

Thus the central difference is the appearance of the term \(J_g\).

If at a particular point \(a\) we set
\[
    y=g(a),
\]
then the forward residuals coincide:
\[
    G(a)-y
    =
    G(a)-g(a)
    =
    e(a).
\]
However, the gradients still differ:
\[
    \nabla_a \ell_y(a)
    =
    2J_G(a)^\top e(a),
\]
while
\[
    \nabla_a \ell_g(a)
    =
    2\left(J_G(a)-J_g(a)\right)^\top e(a).
\]

Their difference is
\[
    \nabla_a \ell_y(a)-\nabla_a \ell_g(a)
    =
    2J_g(a)^\top e(a).
\]

Equivalently,
\[
    \nabla_a \ell_g(a)
    =
    \nabla_a \ell_y(a)
    -
    2J_g(a)^\top e(a).
\]

\subsection{Local Perturbation Expansions}

Let
\[
    a=a_0+\delta.
\]

\subsubsection{Fixed Target \(y\)}

Define
\[
    e_{y,0}=G(a_0)-y.
\]
Linearizing \(G\) gives
\[
    G(a_0+\delta)
    \approx
    G(a_0)+J_G\delta.
\]
Thus
\[
    G(a_0+\delta)-y
    \approx
    e_{y,0}+J_G\delta.
\]

Therefore,
\[
    \ell_y(\delta)
    =
    \|G(a_0+\delta)-y\|_2^2
\]
has the local approximation
\[
    \ell_y(\delta)
    \approx
    \|e_{y,0}+J_G\delta\|_2^2.
\]

Expanding,
\[
    \ell_y(\delta)
    \approx
    e_{y,0}^\top e_{y,0}
    +
    2e_{y,0}^\top J_G\delta
    +
    \delta^\top J_G^\top J_G\delta.
\]

Ignoring the constant term,
\[
    \ell_y(\delta)
    \sim
    2e_{y,0}^\top J_G\delta
    +
    \delta^\top J_G^\top J_G\delta.
\]

The gradient is
\[
    \boxed{
    \nabla_\delta \ell_y(\delta)
    =
    2J_G^\top e_{y,0}
    +
    2J_G^\top J_G\delta
    }.
\]

Thus the local structure is
\[
    \text{linear term}+\text{quadratic term}.
\]

\subsubsection{Dynamic Target \(g(a)\)}

Define
\[
    e_{g,0}=G(a_0)-g(a_0).
\]
Linearizing both \(G\) and \(g\),
\[
    G(a_0+\delta)
    \approx
    G(a_0)+J_G\delta,
\]
and
\[
    g(a_0+\delta)
    \approx
    g(a_0)+J_g\delta.
\]

Thus
\[
    G(a_0+\delta)-g(a_0+\delta)
    \approx
    e_{g,0}+(J_G-J_g)\delta.
\]

Therefore,
\[
    \ell_g(\delta)
    =
    \|G(a_0+\delta)-g(a_0+\delta)\|_2^2
\]
has the local approximation
\[
    \ell_g(\delta)
    \approx
    \|e_{g,0}+(J_G-J_g)\delta\|_2^2.
\]

Expanding,
\[
    \ell_g(\delta)
    \approx
    e_{g,0}^\top e_{g,0}
    +
    2e_{g,0}^\top (J_G-J_g)\delta
    +
    \delta^\top
    (J_G-J_g)^\top(J_G-J_g)\delta.
\]

Ignoring the constant term,
\[
    \ell_g(\delta)
    \sim
    2e_{g,0}^\top (J_G-J_g)\delta
    +
    \delta^\top
    (J_G-J_g)^\top(J_G-J_g)\delta.
\]

The gradient is
\[
    \boxed{
    \nabla_\delta \ell_g(\delta)
    =
    2(J_G-J_g)^\top e_{g,0}
    +
    2(J_G-J_g)^\top(J_G-J_g)\delta
    }.
\]

This is also a linear-plus-quadratic objective in general.

\subsection{Direct Comparison When \(y=g(a_0)\)}

Suppose that the fixed target is chosen as
\[
    y=g(a_0).
\]
Then
\[
    e_{y,0}=G(a_0)-y
    =
    G(a_0)-g(a_0)
    =
    e_{g,0}.
\]
Denote this common residual by
\[
    e_0=G(a_0)-g(a_0).
\]

The fixed-target local loss becomes
\[
    \ell_y(\delta)
    \approx
    \|e_0+J_G\delta\|_2^2,
\]
with gradient
\[
    \nabla_\delta \ell_y(\delta)
    =
    2J_G^\top e_0
    +
    2J_G^\top J_G\delta.
\]

The dynamic-target local loss becomes
\[
    \ell_g(\delta)
    \approx
    \|e_0+(J_G-J_g)\delta\|_2^2,
\]
with gradient
\[
    \nabla_\delta \ell_g(\delta)
    =
    2(J_G-J_g)^\top e_0
    +
    2(J_G-J_g)^\top(J_G-J_g)\delta.
\]

The difference between the two gradients is
\[
    \nabla_\delta \ell_y(\delta)-\nabla_\delta \ell_g(\delta)
    =
    2J_g^\top e_0
    +
    2\left[
    J_G^\top J_G
    -
    (J_G-J_g)^\top(J_G-J_g)
    \right]\delta.
\]

Expanding the matrix difference,
\[
    (J_G-J_g)^\top(J_G-J_g)
    =
    J_G^\top J_G
    -
    J_G^\top J_g
    -
    J_g^\top J_G
    +
    J_g^\top J_g.
\]
Hence,
\[
    J_G^\top J_G
    -
    (J_G-J_g)^\top(J_G-J_g)
    =
    J_G^\top J_g
    +
    J_g^\top J_G
    -
    J_g^\top J_g.
\]

Therefore,
\[
    \boxed{
    \nabla_\delta \ell_y(\delta)-\nabla_\delta \ell_g(\delta)
    =
    2J_g^\top e_0
    +
    2\left[
    J_G^\top J_g
    +
    J_g^\top J_G
    -
    J_g^\top J_g
    \right]\delta
    }.
\]

This shows explicitly that the dynamic-target loss differs from the fixed-target loss because the target itself moves with \(a\).

\section{Detached and Proxy Versions}

\subsection{Full Gradient}

The full oracle loss is
\[
    \ell_{\mathrm{full}}(a)
    =
    \|G(a)-g(a)\|_2^2.
\]
Its gradient is
\[
    \boxed{
    \nabla_a \ell_{\mathrm{full}}(a)
    =
    2(J_G(a)-J_g(a))^\top [G(a)-g(a)]
    }.
\]

This is the mathematically correct gradient of the full loss.

\subsection{Detached Version}

The detached version uses the value \(g(a)\) in the forward pass but stops gradient through \(g\):
\[
    \ell_{\mathrm{det}}(a)
    =
    \|G(a)-\operatorname{sg}(g(a))\|_2^2,
\]
where \(\operatorname{sg}\) denotes stop-gradient.

The forward value is the same as the full loss:
\[
    \ell_{\mathrm{det}}(a)
    =
    \|G(a)-g(a)\|_2^2
\]
as a scalar value.

However, the derivative of \(\operatorname{sg}(g(a))\) is zero:
\[
    D[\operatorname{sg}(g(a))]=0.
\]

Thus the gradient becomes
\[
    \boxed{
    \nabla_a \ell_{\mathrm{det}}(a)
    =
    2J_G(a)^\top [G(a)-g(a)]
    }.
\]

This is not the true gradient of the full loss. It is the gradient of the stopped computational graph.

The difference between detached and full gradients is
\[
    \nabla_a \ell_{\mathrm{det}}(a)-\nabla_a \ell_{\mathrm{full}}(a)
    =
    2J_g(a)^\top [G(a)-g(a)].
\]

Thus,
\[
    \nabla_a \ell_{\mathrm{full}}(a)
    =
    \nabla_a \ell_{\mathrm{det}}(a)
    -
    2J_g(a)^\top [G(a)-g(a)].
\]

The detached version behaves like a fixed-target loss during backpropagation:
\[
    \text{forward target }=g(a),
    \qquad
    \text{backward target treated as constant}.
\]

\subsection{Proxy Version}

The proxy version defines a dictionary-based approximation:
\[
    \tilde a(a)
    =
    \arg\min_{b\in \mathrm{dict}}\|a-b\|,
\]
and
\[
    \tilde g(a)=g(\tilde a(a)).
\]

The proxy loss is
\[
    \tilde \ell(a)
    =
    \|G(a)-\tilde g(a)\|_2^2.
\]

If the dependence of \(\tilde g(a)\) on \(a\) is ignored, then
\[
    D\tilde g(a)\approx 0.
\]

Thus the approximate gradient is
\[
    \boxed{
    \nabla_a \tilde \ell(a)
    \approx
    2J_G(a)^\top [G(a)-\tilde g(a)]
    }.
\]

This differs from the full gradient in two ways:
\[
    \text{(i) the derivative through }g\text{ is ignored,}
\]
and
\[
    \text{(ii) the forward target is } \tilde g(a) \text{ rather than } g(a).
\]

Therefore the proxy gradient can deviate more strongly from the true full gradient.

\section{Why the Gradients Can Point in Very Different Directions}

\subsection{Moving Model Output versus Moving Target}

The fixed-target loss
\[
    \ell_y(a)=\|G(a)-y\|^2
\]
assumes that the target is stationary. Geometrically, only the model output moves.

By contrast, the dynamic-target loss
\[
    \ell_g(a)=\|G(a)-g(a)\|^2
\]
allows both the model output and the target to move. Locally,
\[
    \Delta G \approx J_G\delta,
\]
and
\[
    \Delta g \approx J_g\delta.
\]

The fixed-target loss measures the effect of
\[
    \Delta G.
\]
The dynamic-target loss measures the effect of
\[
    \Delta G-\Delta g.
\]

Therefore, the dynamic loss depends on
\[
    J_G-J_g,
\]
whereas the fixed-target loss depends only on
\[
    J_G.
\]

If
\[
    J_G\approx J_g,
\]
then
\[
    J_G-J_g\approx 0.
\]
In that case, the full dynamic gradient may be small, because the model and the true system respond similarly to the same input perturbation.

However, the detached or fixed-target gradient can still be large because it treats the target as fixed and ignores the true response \(\Delta g\).

\subsection{Residual Terms and Quadratic Terms}

Both fixed-target and dynamic-target losses generally have the local form
\[
    \ell(\delta)
    \sim
    2e_0^\top K\delta
    +
    \delta^\top K^\top K\delta,
\]
where \(K\) is either \(J_G\) or \(J_G-J_g\).

For the fixed-target loss,
\[
    K=J_G.
\]

For the dynamic-target loss,
\[
    K=J_G-J_g.
\]

The gradient is
\[
    \nabla_\delta \ell(\delta)
    =
    2K^\top e_0+2K^\top K\delta.
\]

The first term
\[
    2K^\top e_0
\]
is a residual-driven linear term.

The second term
\[
    2K^\top K\delta
\]
is the quadratic amplification term.

Only when
\[
    e_0\approx 0
\]
does the loss reduce approximately to a pure quadratic form:
\[
    \ell(\delta)
    \approx
    \delta^\top K^\top K\delta.
\]

In that special case, the dominant right singular vector of \(K\) becomes relevant.

\section{Connection Back to PGD and Power Iteration}

\subsection{Pure Quadratic Case}

If the objective is
\[
    \mathcal L(\delta)=\frac{1}{2}\delta^\top A\delta,
\]
then
\[
    \nabla_\delta \mathcal L=A\delta.
\]

The maximization under an \(L_2\) constraint,
\[
    \max_{\|\delta\|_2\le \varepsilon}\delta^\top A\delta,
\]
is a Rayleigh quotient problem. The optimal direction is the top eigenvector of \(A\).

Power iteration is natural:
\[
    v_{k+1}
    =
    \frac{A v_k}{\|A v_k\|_2}.
\]

\subsection{Linear plus Quadratic Case}

If the objective is
\[
    \mathcal L(\delta)
    =
    b^\top \delta
    +
    \frac{1}{2}\delta^\top A\delta,
\]
then
\[
    \nabla_\delta \mathcal L
    =
    b+A\delta.
\]

This is not a Rayleigh quotient problem. The optimal direction depends on both \(b\) and \(A\). The top eigenvector of \(A\) generally does not solve this problem.

For the fixed-target loss,
\[
    b=J_G^\top e_0,
    \qquad
    A=J_G^\top J_G.
\]

For the dynamic-target full loss,
\[
    b=(J_G-J_g)^\top e_0,
    \qquad
    A=(J_G-J_g)^\top(J_G-J_g).
\]

Therefore:

\[
\boxed{
\text{Power iteration is appropriate for pure quadratic sensitivity objectives.}
}
\]

But:

\[
\boxed{
\text{PGD is the direct method for the actual supervised loss when residual terms are present.}
}
\]

\subsection{Practical Implication}

The detached gradient
\[
    2J_G^\top(G-g)
\]
can be close to the full gradient
\[
    2(J_G-J_g)^\top(G-g)
\]
only under specific conditions, for example when
\[
    J_g^\top(G-g)
\]
is small, or when only the elementwise sign is used, as in an \(L_\infty\) PGD attack.

In \(L_\infty\) PGD, the update uses
\[
    \operatorname{sign}(\nabla_a \ell).
\]
Therefore, even if the magnitudes of the full and detached gradients differ, the update may be similar when their signs are mostly the same.

However, if
\[
    J_g^\top(G-g)
\]
is large or changes many coordinate signs, then the detached gradient can point in a significantly different direction from the full gradient.

The proxy version can deviate even more because it not only ignores \(Dg(a)\), but also replaces \(g(a)\) by \(\tilde g(a)\) in the forward target.

\section{Summary}

The fixed-target and dynamic-target losses differ because the target either remains stationary or moves with the input.

For a fixed target,
\[
    \ell_y(a)=\|G(a)-y\|^2,
\]
the gradient is
\[
    \nabla_a \ell_y(a)
    =
    2J_G^\top(G-y).
\]

For a dynamic target,
\[
    \ell_g(a)=\|G(a)-g(a)\|^2,
\]
the full gradient is
\[
    \nabla_a \ell_g(a)
    =
    2(J_G-J_g)^\top(G-g).
\]

The local perturbation forms are:
\[
    \ell_y(\delta)
    \approx
    \|e_0+J_G\delta\|^2,
\]
and
\[
    \ell_g(\delta)
    \approx
    \|e_0+(J_G-J_g)\delta\|^2.
\]

Both generally contain a residual-driven linear term and a quadratic term. Therefore, they are generally not pure Rayleigh quotient objectives.

Only when the residual is negligible do they reduce to pure quadratic forms:
\[
    \delta^\top J_G^\top J_G\delta
\]
or
\[
    \delta^\top (J_G-J_g)^\top(J_G-J_g)\delta.
\]

Thus, the appearance of \(J_g\) is the main mathematical difference between fixed-target and dynamic-target attacks. It can significantly change the gradient direction, especially when the true operator \(g\) is sensitive to the input or when \(J_G\) and \(J_g\) have similar responses.

\section{Power Iteration, PGD, and Spectral Convergence for Singular Vector Computation}

Let \(A\in\mathbb{R}^{m\times n}\). The problem of finding the dominant right singular vector of \(A\) can be written as

\[
v_1
=
\arg\max_{\|v\|_2=1}\|Av\|_2^2.
\]

Since

\[
\|Av\|_2^2
=
v^\top A^\top A v,
\]

this is equivalent to the Rayleigh quotient maximization problem

\[
v_1
=
\arg\max_{\|v\|_2=1} v^\top Bv,
\qquad
B=A^\top A.
\]

The matrix \(B\) is symmetric positive semidefinite. If

\[
Bv_i=\lambda_i v_i,
\qquad
\lambda_i=\sigma_i^2,
\]

where

\[
\lambda_1>\lambda_2\ge \lambda_3\ge\cdots,
\]

then \(v_1\) is the dominant right singular vector of \(A\).

\section{Power Iteration}

The standard power iteration for this problem is

\[
v_{k+1}
=
\frac{Bv_k}{\|Bv_k\|_2}
=
\frac{A^\top A v_k}{\|A^\top A v_k\|_2}.
\]

If the initial vector has a nonzero component in the dominant direction,

\[
v_0
=
c_1 v_1+c_2 v_2+\cdots+c_n v_n,
\qquad c_1\neq 0,
\]

then after \(k\) applications of \(B\),

\[
B^k v_0
=
c_1\lambda_1^k v_1
+
c_2\lambda_2^k v_2
+
\cdots
+
c_n\lambda_n^k v_n.
\]

After normalization, the non-dominant components are suppressed relative to the dominant component. The leading error term is controlled by

\[
\left(\frac{\lambda_2}{\lambda_1}\right)^k
=
\left(\frac{\sigma_2^2}{\sigma_1^2}\right)^k.
\]

Therefore, the direction convergence factor of power iteration is

\[
r_{\mathrm{power}}
=
\frac{\lambda_2}{\lambda_1}
=
\frac{\sigma_2^2}{\sigma_1^2}.
\]

A larger spectral gap, or equivalently a smaller ratio \(\lambda_2/\lambda_1\), gives faster convergence.

\section{Projected Gradient-Type Update}

A gradient ascent formulation of the same Rayleigh quotient problem is based on

\[
f(v)
=
\frac{1}{2}\|Av\|_2^2
=
\frac{1}{2}v^\top Bv.
\]

Its Euclidean gradient is

\[
\nabla f(v)
=
Bv.
\]

A simple normalized gradient ascent or PGD-like update is

\[
v_{k+1}
=
\frac{v_k+\alpha Bv_k}
{\|v_k+\alpha Bv_k\|_2}
=
\frac{(I+\alpha B)v_k}
{\|(I+\alpha B)v_k\|_2}.
\]

Thus, this update is not directly applying \(B\). It is applying the modified iteration matrix

\[
I+\alpha B.
\]

If \(Bv_i=\lambda_i v_i\), then

\[
(I+\alpha B)v_i
=
(1+\alpha\lambda_i)v_i.
\]

Therefore, the dominant direction is amplified by \(1+\alpha\lambda_1\), while the second direction is amplified by \(1+\alpha\lambda_2\). The corresponding convergence factor is

\[
r_{\mathrm{PGD}}
=
\frac{1+\alpha\lambda_2}
{1+\alpha\lambda_1}.
\]

For finite \(\alpha>0\), one usually has

\[
\frac{1+\alpha\lambda_2}
{1+\alpha\lambda_1}
>
\frac{\lambda_2}{\lambda_1},
\]

so the PGD-like update is generally slower than standard power iteration for the pure dominant singular vector problem.

As \(\alpha\to\infty\),

\[
\frac{1+\alpha\lambda_2}
{1+\alpha\lambda_1}
\to
\frac{\lambda_2}{\lambda_1}.
\]

Hence, a very large step size makes the normalized gradient update approach power iteration.

\section{Why Power Iteration Is Faster Than This PGD Update}

The power iteration update is

\[
v_{k+1}\propto Bv_k.
\]

Each eigendirection \(v_i\) is multiplied by \(\lambda_i\). Therefore, the dominant direction is favored over the second direction by the ratio

\[
\frac{\lambda_1}{\lambda_2}.
\]

The PGD-like update is

\[
v_{k+1}\propto v_k+\alpha Bv_k.
\]

Each eigendirection \(v_i\) is multiplied by

\[
1+\alpha\lambda_i.
\]

Thus, the dominant direction is favored over the second direction only by

\[
\frac{1+\alpha\lambda_1}
{1+\alpha\lambda_2}.
\]

Since

\[
\frac{1+\alpha\lambda_1}
{1+\alpha\lambda_2}
<
\frac{\lambda_1}{\lambda_2},
\]

the PGD-like update separates the dominant direction from the second direction more weakly than power iteration.

Therefore, for the pure Rayleigh quotient problem,

\[
\max_{\|v\|_2=1}v^\top Bv,
\]

standard power iteration is usually more efficient than a finite-step-size PGD update.

\section{Methods Faster Than Ordinary Power Iteration}

Although power iteration is faster than the simple PGD-like update above, it is not the fastest possible method for eigenvector or singular vector computation. Several constructions can have better convergence behavior.

\subsection{Higher-Power Filtering}

One direct construction is

\[
v_{k+1}
=
\frac{B^m v_k}
{\|B^m v_k\|_2}.
\]

The convergence factor per outer step becomes

\[
r_m
=
\left(\frac{\lambda_2}{\lambda_1}\right)^m.
\]

This is smaller than

\[
\frac{\lambda_2}{\lambda_1}
\]

when \(m>1\). However, this does not give a true computational acceleration if one counts matrix-vector products, because applying \(B^m\) requires \(m\) applications of \(B\).

\subsection{Chebyshev Polynomial Filtering}

Power iteration uses the polynomial

\[
p_k(\lambda)=\lambda^k.
\]

More generally, one can use

\[
p_k(B)v_0
\]

for a carefully chosen polynomial \(p_k\). The goal is to make

\[
p_k(\lambda_1)
\]

large while keeping

\[
p_k(\lambda_i),\qquad i\ge 2,
\]

small.

Assume that all non-dominant eigenvalues lie in

\[
[\lambda_{\min},\lambda_2],
\]

while \(\lambda_1\) lies outside this interval. Define

\[
x(\lambda)
=
\frac{2\lambda-(\lambda_2+\lambda_{\min})}
{\lambda_2-\lambda_{\min}}.
\]

Then

\[
\lambda\in[\lambda_{\min},\lambda_2]
\quad\Longrightarrow\quad
x(\lambda)\in[-1,1].
\]

Let \(T_k(x)\) be the Chebyshev polynomial of degree \(k\). It satisfies

\[
|T_k(x)|\le 1,
\qquad x\in[-1,1].
\]

For \(x>1\), \(T_k(x)\) grows rapidly. Therefore, one may define

\[
p_k(\lambda)
=
\frac{T_k(x(\lambda))}
{T_k(x(\lambda_1))}.
\]

Then

\[
p_k(\lambda_1)=1,
\]

while for every non-dominant eigenvalue,

\[
|p_k(\lambda_i)|
\le
\frac{1}{|T_k(x(\lambda_1))|}.
\]

Hence the Chebyshev-filtered error satisfies

\[
r_{\mathrm{Cheb},k}
\le
\frac{1}{|T_k(x(\lambda_1))|}.
\]

For \(x>1\),

\[
T_k(x)
=
\frac{1}{2}
\left[
\left(x+\sqrt{x^2-1}\right)^k
+
\left(x-\sqrt{x^2-1}\right)^k
\right].
\]

Therefore, approximately,

\[
r_{\mathrm{Cheb},k}
\approx
2
\left(
x(\lambda_1)+\sqrt{x(\lambda_1)^2-1}
\right)^{-k}.
\]

This can decay faster than the ordinary power iteration rate

\[
\left(\frac{\lambda_2}{\lambda_1}\right)^k.
\]

\subsection{Lanczos and Krylov Subspace Methods}

Lanczos methods use the Krylov subspace

\[
\mathcal{K}_k(B,v_0)
=
\operatorname{span}
\{v_0,Bv_0,B^2v_0,\ldots,B^{k-1}v_0\}.
\]

Power iteration only uses the final vector

\[
B^k v_0.
\]

Lanczos uses the entire Krylov subspace and finds an approximate eigenvector inside it. In this sense, Lanczos automatically constructs a better polynomial filter than the fixed monomial filter used by power iteration.

For symmetric matrices, Lanczos is one of the most important practical methods for computing extreme eigenvalues and eigenvectors. Since

\[
B=A^\top A
\]

is symmetric positive semidefinite, Lanczos is directly applicable to dominant singular vector computation.

\subsection{Shift-Invert Iteration}

Another much faster construction is shift-invert iteration. Choose a shift \(\mu\) close to the desired eigenvalue \(\lambda_1\), and iterate

\[
v_{k+1}
=
\frac{(B-\mu I)^{-1}v_k}
{\|(B-\mu I)^{-1}v_k\|_2}.
\]

If

\[
Bv_i=\lambda_i v_i,
\]

then

\[
(B-\mu I)^{-1}v_i
=
\frac{1}{\lambda_i-\mu}v_i.
\]

Thus the transformed eigenvalues are

\[
\theta_i
=
\frac{1}{\lambda_i-\mu}.
\]

The new convergence factor is

\[
r_{\mathrm{shift\text{-}invert}}
=
\left|
\frac{\lambda_1-\mu}
{\lambda_2-\mu}
\right|.
\]

If \(\mu\) is very close to \(\lambda_1\), then

\[
|\lambda_1-\mu|
\]

is small, and therefore

\[
r_{\mathrm{shift\text{-}invert}}
\ll
\frac{\lambda_2}{\lambda_1}.
\]

This can be much faster than power iteration. The main cost is that each step requires solving a linear system

\[
(B-\mu I)y=v_k.
\]

For very large implicit matrices, such as Jacobians in deep networks, this may be computationally expensive.

\subsection{Rayleigh Quotient Iteration}

Rayleigh quotient iteration is an adaptive version of shift-invert iteration. At step \(k\), define

\[
\mu_k
=
\frac{v_k^\top Bv_k}{v_k^\top v_k}.
\]

Then update

\[
v_{k+1}
=
\frac{(B-\mu_k I)^{-1}v_k}
{\|(B-\mu_k I)^{-1}v_k\|_2}.
\]

For symmetric matrices, once the iterate is sufficiently close to an eigenvector, Rayleigh quotient iteration has very fast local convergence. However, like shift-invert iteration, it requires solving a linear system at each step.

\section{The Role of the Rayleigh Quotient Structure}

The convergence analysis above depends heavily on the Rayleigh quotient structure

\[
f(v)=v^\top Bv,
\qquad \|v\|_2=1.
\]

The reason is that

\[
\nabla f(v)=2Bv.
\]

Therefore, the update can be written using a fixed linear operator:

\[
B,
\qquad
I+\alpha B,
\qquad
p(B),
\qquad
(B-\mu I)^{-1}.
\]

This allows us to decompose the initial vector as

\[
v_0=\sum_i c_i v_i
\]

and track the behavior of each eigendirection separately.

For example,

\[
B^k v_0
=
\sum_i c_i\lambda_i^k v_i.
\]

This is why one can obtain explicit convergence factors such as

\[
\frac{\lambda_2}{\lambda_1},
\qquad
\frac{1+\alpha\lambda_2}{1+\alpha\lambda_1},
\qquad
\frac{p(\lambda_2)}{p(\lambda_1)},
\qquad
\left|
\frac{\lambda_1-\mu}{\lambda_2-\mu}
\right|.
\]

These formulas rely on the existence of a fixed matrix \(B\) and its eigenvalue decomposition.

\section{When the Loss Is Not a Rayleigh Quotient}

For a general nonlinear loss

\[
\mathcal{L}(v),
\]

the gradient is generally not of the form

\[
\nabla \mathcal{L}(v)=Bv
\]

for a fixed matrix \(B\). Instead, the gradient depends nonlinearly on \(v\):

\[
\nabla \mathcal{L}(v)
=
\text{nonlinear function of }v.
\]

In that case, there is no fixed eigensystem

\[
Bv_i=\lambda_i v_i
\]

that globally controls the iteration. Therefore, the convergence factors

\[
\frac{\lambda_2}{\lambda_1},
\qquad
\frac{1+\alpha\lambda_2}{1+\alpha\lambda_1},
\]

do not directly apply.

For example, a general PGD adversarial attack has the form

\[
\delta_{k+1}
=
\Pi_{\|\delta\|\le \epsilon}
\left(
\delta_k+\alpha\nabla_\delta \mathcal{L}(x+\delta_k,y)
\right).
\]

Here the objective is

\[
\max_{\|\delta\|\le\epsilon}
\mathcal{L}(x+\delta,y),
\]

where \(\mathcal{L}\) is usually a nonlinear neural network loss. This is not generally a quadratic Rayleigh quotient problem.

Therefore, it is not correct to directly interpret such PGD iterations as power iterations for a fixed matrix. In this setting, PGD is a general constrained optimization method, not a spectral iteration method.

\section{Local Linearization and When Spectral Methods Become Relevant Again}

Although a general nonlinear loss is not globally a Rayleigh quotient, spectral methods can become relevant after local approximation.

For example, suppose one studies the local output change of a nonlinear model \(f\) around an input \(x\). For a small perturbation \(\delta\),

\[
f(x+\delta)-f(x)
\approx
J(x)\delta,
\]

where \(J(x)\) is the Jacobian of \(f\) at \(x\). Then

\[
\|f(x+\delta)-f(x)\|_2^2
\approx
\|J(x)\delta\|_2^2
=
\delta^\top J(x)^\top J(x)\delta.
\]

This gives a local Rayleigh quotient problem:

\[
\max_{\|\delta\|_2=1}
\delta^\top J(x)^\top J(x)\delta.
\]

The solution is the dominant right singular vector of \(J(x)\). In this local linearized setting, power iteration, Lanczos, Chebyshev filtering, and other spectral methods again have direct meaning.

Similarly, if a nonlinear loss is locally approximated by a second-order Taylor expansion,

\[
\mathcal{L}(x+\delta)
\approx
\mathcal{L}(x)
+
g^\top \delta
+
\frac{1}{2}\delta^\top H\delta,
\]

then the quadratic term involves the Hessian \(H\). If the quadratic term dominates, one may again study an eigenvalue problem such as

\[
\max_{\|\delta\|_2=1}\delta^\top H\delta.
\]

In that local quadratic approximation, spectral methods become meaningful again.

\section{Main Conclusion}

The comparison between PGD, power iteration, Chebyshev filtering, Lanczos, shift-invert iteration, and Rayleigh quotient iteration is meaningful when the target problem has a spectral structure, such as

\[
\max_{\|v\|_2=1}v^\top Bv
\]

or

\[
\max_{\|v\|_2=1}\|Av\|_2^2.
\]

In this setting,

\[
B=A^\top A,
\]

and the dominant right singular vector of \(A\) is the dominant eigenvector of \(B\).

For this problem, the ordinary power iteration has convergence factor

\[
r_{\mathrm{power}}
=
\frac{\lambda_2}{\lambda_1}.
\]

A PGD-like normalized gradient update has convergence factor

\[
r_{\mathrm{PGD}}
=
\frac{1+\alpha\lambda_2}{1+\alpha\lambda_1},
\]

which is usually larger and therefore slower for finite \(\alpha\).

More advanced spectral methods can improve on power iteration. Chebyshev filtering uses a better polynomial \(p(B)\), Lanczos uses the full Krylov subspace, and shift-invert methods transform the spectrum so that the target eigendirection becomes much more dominant.

However, if the objective is a general nonlinear loss rather than a Rayleigh quotient or a local quadratic approximation, then these spectral convergence factors no longer apply directly. In that case, PGD should be viewed as a general constrained optimization method, while power iteration and its accelerations should be viewed as methods for linear or locally linearized spectral problems.

\section{Why a Squared Norm Loss Is Not Always a Quadratic Form in the Perturbation}

A common source of confusion is that many adversarial losses have the form
\[
    \frac{1}{2}\|\cdot\|_2^2.
\]
At first glance, all such losses appear to be quadratic. However, whether a loss is a \emph{pure quadratic form in the perturbation} \(\delta\) depends on what appears inside the norm after local linearization.

The key criterion is:

\[
\boxed{
\text{A squared norm becomes a pure quadratic form in } \delta
\text{ only if the vector inside the norm is of the form } A\delta.
}
\]

If the vector inside the norm is instead
\[
    c + A\delta,
\]
then
\[
    \frac{1}{2}\|c+A\delta\|_2^2
    =
    \frac{1}{2}c^\top c
    +
    c^\top A\delta
    +
    \frac{1}{2}\delta^\top A^\top A\delta.
\]
After dropping the constant term, the loss is still
\[
    c^\top A\delta
    +
    \frac{1}{2}\delta^\top A^\top A\delta,
\]
which contains both a linear term and a quadratic term. It is therefore not a pure Rayleigh-quotient-type quadratic form.

\subsection{Notation}

Let
\[
    F_0=f(x),
    \qquad
    F_\delta=f(x+\delta),
\]
and, when a dynamic ground-truth operator is available,
\[
    G_0=g(x),
    \qquad
    G_\delta=g(x+\delta).
\]
For a fixed ground truth, write
\[
    Y=y.
\]

The local linearizations are
\[
    f(x+\delta)
    \approx
    f(x)+J_f\delta,
\]
and
\[
    g(x+\delta)
    \approx
    g(x)+J_g\delta.
\]

\subsection{Loss 1: Model Output Change}

The first loss is
\[
    \mathcal L_1(\delta)
    =
    \frac{1}{2}\|F_\delta-F_0\|_2^2
    =
    \frac{1}{2}\|f(x+\delta)-f(x)\|_2^2.
\]

Using the local linearization,
\[
    f(x+\delta)
    \approx
    f(x)+J_f\delta,
\]
we get
\[
    f(x+\delta)-f(x)
    \approx
    J_f\delta.
\]

Therefore,
\[
    \mathcal L_1(\delta)
    \approx
    \frac{1}{2}\|J_f\delta\|_2^2.
\]
Equivalently,
\[
    \mathcal L_1(\delta)
    \approx
    \frac{1}{2}\delta^\top J_f^\top J_f\delta.
\]

Thus, Loss 1 is a pure quadratic form in \(\delta\). It can be written as
\[
    \mathcal L_1(\delta)
    \approx
    \frac{1}{2}\delta^\top A_1\delta,
    \qquad
    A_1=J_f^\top J_f.
\]

Consequently, maximizing it under an \(L_2\) constraint leads to a Rayleigh quotient problem:
\[
    \max_{\|\delta\|_2\le \varepsilon}
    \delta^\top J_f^\top J_f\delta.
\]
The optimal direction is the dominant right singular vector of \(J_f\).

\subsection{Loss 2: Error Against a Fixed Target}

The second loss is
\[
    \mathcal L_2(\delta)
    =
    \frac{1}{2}\|F_\delta-Y\|_2^2
    =
    \frac{1}{2}\|f(x+\delta)-y\|_2^2.
\]

Using the local linearization,
\[
    f(x+\delta)
    \approx
    f(x)+J_f\delta,
\]
we obtain
\[
    f(x+\delta)-y
    \approx
    f(x)-y+J_f\delta.
\]

Define the original residual
\[
    e_y=f(x)-y.
\]
Then
\[
    f(x+\delta)-y
    \approx
    e_y+J_f\delta.
\]

Therefore,
\[
    \mathcal L_2(\delta)
    \approx
    \frac{1}{2}\|e_y+J_f\delta\|_2^2.
\]

Expanding,
\[
    \mathcal L_2(\delta)
    \approx
    \frac{1}{2}e_y^\top e_y
    +
    e_y^\top J_f\delta
    +
    \frac{1}{2}\delta^\top J_f^\top J_f\delta.
\]

The constant term
\[
    \frac{1}{2}e_y^\top e_y
\]
does not affect optimization over \(\delta\). However, the linear term
\[
    e_y^\top J_f\delta
\]
remains. Thus,
\[
    \mathcal L_2(\delta)
    \sim
    e_y^\top J_f\delta
    +
    \frac{1}{2}\delta^\top J_f^\top J_f\delta.
\]

Hence, Loss 2 is generally not a pure quadratic form in \(\delta\). It is a linear-plus-quadratic objective:
\[
    \mathcal L_2(\delta)
    \sim
    b_2^\top \delta
    +
    \frac{1}{2}\delta^\top A_2\delta,
\]
where
\[
    b_2=J_f^\top e_y,
    \qquad
    A_2=J_f^\top J_f.
\]

The reason is that \(f(x+\delta)-y\) does not subtract off the original model output \(f(x)\). The original residual
\[
    f(x)-y
\]
survives as a constant offset inside the norm, producing the linear term.

\subsection{Loss 3: Error Against a Dynamic Ground Truth}

The third loss is
\[
    \mathcal L_3(\delta)
    =
    \frac{1}{2}\|F_\delta-G_\delta\|_2^2
    =
    \frac{1}{2}\|f(x+\delta)-g(x+\delta)\|_2^2.
\]

Using the local linearizations,
\[
    f(x+\delta)
    \approx
    f(x)+J_f\delta,
\]
and
\[
    g(x+\delta)
    \approx
    g(x)+J_g\delta,
\]
we obtain
\[
    f(x+\delta)-g(x+\delta)
    \approx
    f(x)-g(x)+(J_f-J_g)\delta.
\]

Define the original model error
\[
    e_g=f(x)-g(x).
\]
Then
\[
    f(x+\delta)-g(x+\delta)
    \approx
    e_g+(J_f-J_g)\delta.
\]

Therefore,
\[
    \mathcal L_3(\delta)
    \approx
    \frac{1}{2}\|e_g+(J_f-J_g)\delta\|_2^2.
\]

Expanding,
\[
    \mathcal L_3(\delta)
    \approx
    \frac{1}{2}e_g^\top e_g
    +
    e_g^\top (J_f-J_g)\delta
    +
    \frac{1}{2}
    \delta^\top
    (J_f-J_g)^\top(J_f-J_g)
    \delta.
\]

Dropping the constant term,
\[
    \mathcal L_3(\delta)
    \sim
    e_g^\top (J_f-J_g)\delta
    +
    \frac{1}{2}
    \delta^\top
    (J_f-J_g)^\top(J_f-J_g)
    \delta.
\]

Thus, Loss 3 is generally not a pure quadratic form. It is again a linear-plus-quadratic objective:
\[
    \mathcal L_3(\delta)
    \sim
    b_3^\top \delta
    +
    \frac{1}{2}\delta^\top A_3\delta,
\]
where
\[
    b_3=(J_f-J_g)^\top e_g,
\]
and
\[
    A_3=(J_f-J_g)^\top(J_f-J_g).
\]

The reason is that Loss 3 measures the perturbed error itself,
\[
    f(x+\delta)-g(x+\delta),
\]
rather than the change in that error from \(x\) to \(x+\delta\). Unless
\[
    f(x)-g(x)=0,
\]
the original error vector remains inside the norm and produces a linear term.

\subsection{Loss 4: Error-Change Loss}

The fourth loss is
\[
    \mathcal L_4(\delta)
    =
    \frac{1}{2}
    \|(F_\delta-G_\delta)-(F_0-G_0)\|_2^2.
\]

Equivalently,
\[
    \mathcal L_4(\delta)
    =
    \frac{1}{2}
    \|
    [f(x+\delta)-g(x+\delta)]
    -
    [f(x)-g(x)]
    \|_2^2.
\]

This is the squared norm of the change in the error function
\[
    E(x)=f(x)-g(x).
\]
Indeed,
\[
    \mathcal L_4(\delta)
    =
    \frac{1}{2}\|E(x+\delta)-E(x)\|_2^2.
\]

Using the previous local expansion,
\[
    f(x+\delta)-g(x+\delta)
    \approx
    f(x)-g(x)+(J_f-J_g)\delta.
\]
Subtracting \(f(x)-g(x)\) gives
\[
    [f(x+\delta)-g(x+\delta)]-[f(x)-g(x)]
    \approx
    (J_f-J_g)\delta.
\]

Therefore,
\[
    \mathcal L_4(\delta)
    \approx
    \frac{1}{2}
    \|(J_f-J_g)\delta\|_2^2.
\]

Equivalently,
\[
    \mathcal L_4(\delta)
    \approx
    \frac{1}{2}
    \delta^\top
    (J_f-J_g)^\top(J_f-J_g)
    \delta.
\]

Thus, Loss 4 is a pure quadratic form in \(\delta\). It can be written as
\[
    \mathcal L_4(\delta)
    \approx
    \frac{1}{2}\delta^\top A_4\delta,
\]
where
\[
    A_4=(J_f-J_g)^\top(J_f-J_g).
\]

\subsection{Summary Table}

\begin{table}[h]
\centering
\begin{tabular}{p{0.8cm} p{4.2cm} p{4.2cm} p{4.8cm} p{2.2cm}}
\hline
No. & Loss & Linearized vector inside norm & Local loss structure & Pure quadratic? \\
\hline
1 &
\(\displaystyle \frac{1}{2}\|F_\delta-F_0\|^2\)
&
\(\displaystyle J_f\delta\)
&
\(\displaystyle \frac{1}{2}\delta^\top J_f^\top J_f\delta\)
&
Yes
\\[1.2em]

2 &
\(\displaystyle \frac{1}{2}\|F_\delta-Y\|^2\)
&
\(\displaystyle e_y+J_f\delta\)
&
\(\displaystyle e_y^\top J_f\delta+\frac{1}{2}\delta^\top J_f^\top J_f\delta\)
&
No
\\[1.2em]

3 &
\(\displaystyle \frac{1}{2}\|F_\delta-G_\delta\|^2\)
&
\(\displaystyle e_g+(J_f-J_g)\delta\)
&
\(\displaystyle e_g^\top(J_f-J_g)\delta+\frac{1}{2}\delta^\top(J_f-J_g)^\top(J_f-J_g)\delta\)
&
No
\\[1.2em]

4 &
\(\displaystyle \frac{1}{2}\|(F_\delta-G_\delta)-(F_0-G_0)\|^2\)
&
\(\displaystyle (J_f-J_g)\delta\)
&
\(\displaystyle \frac{1}{2}\delta^\top(J_f-J_g)^\top(J_f-J_g)\delta\)
&
Yes
\\
\hline
\end{tabular}
\caption{Why some squared norm losses become pure quadratic forms in \(\delta\), while others contain residual-driven linear terms.}
\end{table}

\subsection{General Rule}

The general rule is:

\[
\boxed{
\frac{1}{2}\|A\delta\|^2
=
\frac{1}{2}\delta^\top A^\top A\delta
\quad
\text{is a pure quadratic form.}
}
\]

But

\[
\boxed{
\frac{1}{2}\|c+A\delta\|^2
=
\frac{1}{2}c^\top c
+
c^\top A\delta
+
\frac{1}{2}\delta^\top A^\top A\delta
}
\]
contains a linear term unless \(c=0\).

Therefore, the decisive question is not whether the loss is a squared norm, but whether the vector inside the norm is purely proportional to the perturbation \(\delta\).

\[
\boxed{
\text{A squared norm loss is a pure quadratic form in } \delta
\text{ only if its linearized residual has no constant offset.}
}
\]

\subsection{Connection to Rayleigh Quotient and Power Iteration}

If the local loss has the form
\[
    \mathcal L(\delta)
    =
    \frac{1}{2}\delta^\top A\delta,
\]
then maximizing it under an \(L_2\) constraint,
\[
    \max_{\|\delta\|_2\le \varepsilon}
    \delta^\top A\delta,
\]
is a Rayleigh quotient problem. The maximizer is the dominant eigenvector of \(A\), and power iteration is appropriate:
\[
    v_{k+1}
    =
    \frac{Av_k}{\|Av_k\|}.
\]

This applies to Loss 1 and Loss 4.

However, if the local loss has the form
\[
    \mathcal L(\delta)
    =
    b^\top\delta
    +
    \frac{1}{2}\delta^\top A\delta,
\]
then it is not a Rayleigh quotient problem. The optimizer depends on both the residual-driven linear term \(b\) and the quadratic matrix \(A\).

This applies to Loss 2 and Loss 3.

Thus,

\[
\boxed{
\text{Although all four losses are squared norms, only Loss 1 and Loss 4 reduce to pure quadratic forms under local linearization.}
}
\]

\[
\boxed{
\text{Loss 2 and Loss 3 retain the original residual as a constant offset, which generates a linear term in } \delta.
}
\]


```latex
\section{Classification of Adversarial Losses by Comparison Target, Input Terms, and Spectral Structure}

This section summarizes a set of adversarial attack and robustness methods according to the structure of their objective functions. The central question is not only what quantity is optimized, but also what the perturbed output is compared against. In particular, we distinguish whether the loss compares \(f(x+\delta)\) with the original output \(f(x)\), with a label \(y\), with a target class \(t\), with a target output \(\tilde f\), or with a physical ground-truth simulator \(g(x+\delta)\). This distinction determines whether the local linearization of the loss leads to a pure quadratic form of the type
\[
    \|J_f(x)\delta\|_2^2
    =
    \delta^\top J_f(x)^\top J_f(x)\delta,
\]
which can be solved by singular-vector or spectral methods.

\subsection{Output-Change Losses and Pure Quadratic Forms}

The cleanest spectral case is an output-change loss:
\[
    \mathcal{L}_{\mathrm{out}}(\delta)
    =
    \|f(x+\delta)-f(x)\|_2^2.
\]
Using the local linearization
\[
    f(x+\delta)
    \approx
    f(x)+J_f(x)\delta,
\]
we obtain
\[
    f(x+\delta)-f(x)
    \approx
    J_f(x)\delta.
\]
Therefore,
\[
    \mathcal{L}_{\mathrm{out}}(\delta)
    \approx
    \|J_f(x)\delta\|_2^2.
\]
If the input perturbation is constrained by
\[
    \|\delta\|_2=1,
\]
then the optimization problem becomes
\[
    \max_{\|\delta\|_2=1}
    \|J_f(x)\delta\|_2^2
    =
    \max_{\|\delta\|_2=1}
    \delta^\top J_f(x)^\top J_f(x)\delta.
\]
This is a Rayleigh quotient problem. The optimizer is the dominant eigenvector of
\[
    J_f(x)^\top J_f(x),
\]
equivalently the dominant right singular vector of \(J_f(x)\). Thus, in this special setting, power iteration, Lanczos methods, and singular-vector algorithms are directly meaningful.

This type of objective appears explicitly in singular-vector-based universal perturbation methods, such as methods based on maximizing
\[
    \|J_\phi(x)v\|_q
\]
for a feature map \(\phi\), or in the multi-sample version
\[
    \max_{\|v\|_p=1}
    \sum_{i=1}^N
    \|J_\phi(x_i)v\|_q^q.
\]
When \(p=q=2\), this becomes
\[
    \max_{\|v\|_2=1}
    v^\top
    \left(
        \sum_{i=1}^N
        J_\phi(x_i)^\top J_\phi(x_i)
    \right)
    v.
\]
This is a genuine spectral or singular-vector attack.

\subsection{Label-Based and Target-Based Losses Are Usually Not Pure Quadratic}

Many adversarial attacks do not compare \(f(x+\delta)\) with \(f(x)\). Instead, they compare the perturbed output with a label, a target class, a target output, or a physical ground truth.

For example, a supervised regression loss has the form
\[
    \mathcal{L}_{\mathrm{reg}}(\delta)
    =
    \|f(x+\delta)-y\|_2^2.
\]
Linearizing \(f(x+\delta)\) gives
\[
    f(x+\delta)-y
    \approx
    f(x)-y+J_f(x)\delta.
\]
Hence,
\[
    \|f(x+\delta)-y\|_2^2
    \approx
    \|f(x)-y+J_f(x)\delta\|_2^2.
\]
Expanding this expression yields
\[
    \|f(x)-y\|_2^2
    +
    2(f(x)-y)^\top J_f(x)\delta
    +
    \delta^\top J_f(x)^\top J_f(x)\delta.
\]
The middle term
\[
    2(f(x)-y)^\top J_f(x)\delta
\]
is linear in \(\delta\). Therefore, this is not a pure quadratic Rayleigh quotient unless the residual is negligible:
\[
    f(x)\approx y.
\]
Only under such an approximation can the regression loss be reduced to the spectral form
\[
    \max_{\|\delta\|_2=1}
    \|J_f(x)\delta\|_2^2.
\]

Similarly, classification attacks based on cross-entropy,
\[
    \max_{\|\delta\|\le \epsilon}
    L(f(x+\delta),y),
\]
or target-class losses,
\[
    L(f(x+\delta),t),
\]
are not pure Jacobian spectral problems. Their gradients are task-loss gradients, not simply the action of \(J_f^\top J_f\) on a perturbation vector.

\subsection{Whether the Loss Contains an Input Term}

A separate distinction concerns whether the objective function itself contains an input perturbation penalty, such as
\[
    \|\delta\|_2^2,
    \qquad
    \|\delta\|_0,
    \qquad
    \Omega(\delta).
\]
This is different from having an input constraint such as
\[
    \|\delta\|_\infty\le \epsilon.
\]
If the perturbation size only appears in the constraint, then it should not be counted as an input term inside the loss.

For example, FGSM solves a problem of the form
\[
    \max_{\|\delta\|_\infty\le \epsilon}
    L(f(x+\delta),y).
\]
The objective is only
\[
    L(f(x+\delta),y).
\]
The perturbation constraint
\[
    \|\delta\|_\infty\le \epsilon
\]
is not part of the loss itself. Thus, the loss contains only an output term, while the input perturbation appears as a hard constraint.

By contrast, the Szegedy et al. L-BFGS attack uses an objective of the form
\[
    \min_\delta
    \|\delta\|_2^2
    +
    cL(f(x+\delta),t),
\]
with a box constraint
\[
    x+\delta\in[0,1]^d.
\]
Here the loss itself contains both an input term,
\[
    \|\delta\|_2^2,
\]
and an output term,
\[
    L(f(x+\delta),t).
\]

Carlini--Wagner attacks have the same structural property:
\[
    \min_\delta
    \|\delta\|_p^2
    +
    c\,h(f(x+\delta),t),
\]
where \(\|\delta\|_p^2\) penalizes the input perturbation and \(h(f(x+\delta),t)\) is a logit-margin attack loss.

DeepFool has a different structure:
\[
    \min_\delta \|\delta\|_2
    \quad
    \text{s.t.}
    \quad
    f(x+\delta)\neq f(x).
\]
The objective contains only the input perturbation size. The output condition appears as a constraint. SparseFool is analogous:
\[
    \min_\delta \|\delta\|_0
    \quad
    \text{s.t.}
    \quad
    f(x+\delta)\neq f(x).
\]

\subsection{Classification of Representative Methods}

\begin{table}[h]
\centering
\small
\begin{tabular}{p{3.2cm} p{3.4cm} p{4.8cm} p{2.2cm} p{2.0cm}}
\hline
\textbf{Method} & \textbf{Comparison Target} & \textbf{Objective / Loss Structure} & \textbf{Input Term in Loss?} & \textbf{Singular-Vector Attack?} \\
\hline
Szegedy et al. L-BFGS
&
\(f(x+\delta)\) vs. target class \(t\)
&
\(\min_\delta \|\delta\|_2^2+cL(f(x+\delta),t)\)
&
Yes
&
No
\\

FGSM
&
\(f(x+\delta)\) vs. label \(y\)
&
\(\max_{\|\delta\|_\infty\le \epsilon} L(f(x+\delta),y)\)
&
No; input only in constraint
&
No
\\

BIM / I-FGSM
&
\(f(x+\delta)\) vs. label \(y\)
&
Iterative maximization of \(L(f(x+\delta),y)\)
&
No; projection constraint
&
No
\\

Madry PGD
&
\(f(x+\delta)\) vs. label \(y\)
&
\(\max_{\delta\in S} L(f(x+\delta),y)\)
&
No; input only in \(S\)
&
No
\\

Carlini--Wagner
&
Logits vs. target class \(t\)
&
\(\min_\delta \|\delta\|_p^2+c\,h(f(x+\delta),t)\)
&
Yes
&
No
\\

DeepFool
&
Decision boundary
&
\(\min_\delta \|\delta\|_2\), subject to \(f(x+\delta)\neq f(x)\)
&
Yes; objective is input size
&
No
\\

SparseFool
&
Decision boundary
&
\(\min_\delta \|\delta\|_0\), subject to \(f(x+\delta)\neq f(x)\)
&
Yes; objective is sparsity
&
No
\\

Original UAP
&
\(f(x+v)\) vs. \(f(x)\) in classification decision
&
Maximize fooling rate, subject to \(\|v\|_p\le \xi\)
&
No; input only in constraint
&
No
\\

Art of Singular Vectors and UAP
&
\(\phi(x+v)\) vs. \(\phi(x)\)
&
\(\max_{\|v\|_p=1}\|J_\phi(x)v\|_q\)
&
No; input norm only in constraint
&
Yes
\\

Sparse Universal Singular Vectors
&
Hidden/output change
&
\(\max_v\sum_i\|J(x_i)v\|^q\), subject to norm and sparsity constraints
&
No; constraints only
&
Yes
\\

Adversarial Patch
&
Target-class probability
&
\(\max_p \mathbb{E}_{x,t}\log P(y_t\mid A(x,p,t))\)
&
Usually no explicit norm penalty
&
No
\\

Structured Adversarial Attack
&
Classification loss with structured sparsity
&
\(\mathcal{L}_{\mathrm{attack}}(x+\delta)+\lambda\Omega(\delta)\)
&
Yes
&
No
\\

VAE Adversarial Images
&
Target latent code or reconstruction target
&
\(\Delta(z(x+\delta),z(x_{\mathrm{target}}))+C\|\delta\|\)
&
Yes
&
No
\\

Time-series regression attacks
&
\(f(x+\delta)\) vs. \(y\)
&
\(\max_{\|\delta\|\le \epsilon}\|f(x+\delta)-y\|^2\)
&
No; input only in constraint
&
No
\\

Steering-angle attacks
&
\(f(x+\delta)\) vs. true steering angle \(y\)
&
\(\min_\delta \|\delta\|_2^2-c(f(x+\delta)-y)^2\)
&
Yes
&
No
\\

DeepBillboard
&
Steering trajectory deviation
&
Maximize cumulative steering error under physical constraints
&
Usually no simple norm penalty
&
No
\\

Optical-flow patch attacks
&
Perturbed flow vs. original/target flow
&
Maximize flow error or directional discrepancy via patch optimization
&
Usually no explicit norm penalty
&
No
\\

PCFA optical-flow attack
&
Predicted flow vs. target flow \(\tilde f\)
&
\(\min_{\|\delta\|_2\le \epsilon}\mathcal{L}(F(I+\delta),\tilde f)\)
&
No; input only in constraint
&
No
\\

Neural-operator digital-twin attack
&
\(\mathcal{G}_\theta(u+\delta)\) vs. \(\mathcal{G}_\theta(u)\)
&
\(\max_\delta \frac{\|\mathcal{G}_\theta(u+\delta)-\mathcal{G}_\theta(u)\|_2}{\|\mathcal{G}_\theta(u)\|_2}\)
&
No; sparsity and physical bounds are constraints
&
No
\\

Jacobian spectral norm regularization
&
Local output amplification
&
\(\sigma_{\max}(J_f)=\max_{\|v\|=1}\|J_fv\|\)
&
No concrete attack perturbation
&
Not an attack; spectral regularization
\\
\hline
\end{tabular}
\caption{Classification of adversarial and robustness methods by comparison target, input terms in the objective, and whether they are singular-vector attacks.}
\end{table}

\subsection{Original UAP Versus Singular-Vector UAP}

The original Universal Adversarial Perturbation method does not solve a singular-vector problem. Its objective is to find a universal perturbation \(v\) that fools many samples:
\[
    \frac{1}{N}
    \sum_{i=1}^N
    \mathbf{1}\{f(x_i+v)\neq f(x_i)\}.
\]
The perturbation is constrained by
\[
    \|v\|_p\le \xi.
\]
The original algorithm repeatedly scans the dataset. If the current \(v\) does not fool a sample \(x_i\), it computes an additional small perturbation \(\Delta v_i\), often using a DeepFool-like step, and then projects the updated universal perturbation back into the allowed norm ball:
\[
    v \leftarrow P_{p,\xi}(v+\Delta v_i).
\]
Therefore, the original UAP method is not power iteration and is not an SVD-based method.

The later singular-vector interpretation replaces the fooling-rate objective with a local feature-sensitivity objective:
\[
    \max_{\|v\|_p=1}
    \|J_\phi(x)v\|_q.
\]
This is the point where generalized singular vectors and generalized power iteration become relevant. Thus, the original UAP and singular-vector UAP should not be conflated.

\subsection{Regression Perturbation Theory}

In regression, one may start from a supervised loss:
\[
    \max_{\|\delta\|\le \epsilon}
    \|f(x+\delta)-y\|_2^2.
\]
This is not a pure quadratic form because the linearized expansion contains the residual \(f(x)-y\):
\[
    \|f(x+\delta)-y\|_2^2
    \approx
    \|f(x)-y\|_2^2
    +
    2(f(x)-y)^\top J_f(x)\delta
    +
    \delta^\top J_f(x)^\top J_f(x)\delta.
\]
However, if one assumes
\[
    y\approx f(x),
\]
then the residual term is negligible, and the regression attack reduces approximately to
\[
    \max_{\|\delta\|\le \epsilon}
    \|J_f(x)\delta\|_2^2.
\]
For an \(\ell_2\) constraint, this becomes a dominant singular-vector problem:
\[
    \delta^\star
    =
    \pm \epsilon v_{\max},
\]
where \(v_{\max}\) is the dominant eigenvector of
\[
    J_f(x)^\top J_f(x).
\]
Therefore, regression perturbation theory becomes spectral only under this approximation.

\subsection{Neural-Operator Digital-Twin Attack}

The neural-operator digital-twin attack is a useful example because it is output-change based but not a singular-vector attack. The attacked model is a neural operator
\[
    \mathcal{G}_\theta,
\]
which maps a physical input vector \(u\) to a high-dimensional predicted physical field:
\[
    \mathcal{G}_\theta(u).
\]
The attack compares the adversarial prediction and the clean prediction:
\[
    \mathcal{G}_\theta(u+\delta)
    -
    \mathcal{G}_\theta(u).
\]
The attack score is a relative output change:
\[
    \mathrm{Score}(u,\delta)
    =
    \frac{
        \|\mathcal{G}_\theta(u+\delta)-\mathcal{G}_\theta(u)\|_2
    }{
        \|\mathcal{G}_\theta(u)\|_2
    }.
\]
The denominator is the clean output norm, not the input perturbation norm. Hence this is not an output-change-over-input-change objective.

The perturbation is constrained to be sparse:
\[
    \|\delta\|_0\le k,
\]
and the modified input values must remain physically plausible. Thus, the input perturbation appears as a sparsity and physical-feasibility constraint, not as an input penalty inside the objective.

The paper optimizes this objective using gradient-free differential evolution. A candidate attack is encoded as a set of index-value pairs:
\[
    \theta
    =
    (i_1,a_1,i_2,a_2,\ldots,i_k,a_k),
\]
where \(i_j\) is the input coordinate to be modified and \(a_j\) is the replacement value. Each candidate is decoded into an adversarial input \(u_{\mathrm{adv}}\), evaluated by a forward pass,
\[
    \mathcal{G}_\theta(u_{\mathrm{adv}}),
\]
and assigned a fitness score:
\[
    \frac{
        \|\mathcal{G}_\theta(u_{\mathrm{adv}})-\mathcal{G}_\theta(u)\|_2
    }{
        \|\mathcal{G}_\theta(u)\|_2
    }.
\]
Differential evolution then updates a population of candidates by mutation, crossover, and greedy selection. This procedure requires only forward evaluations of the neural operator; it does not require gradients, backpropagation, Jacobian-vector products, or power iteration.

Therefore, this method is:
\[
    \text{output-change based},
\]
but it is not:
\[
    \text{a pure quadratic Jacobian singular-vector attack}.
\]
It is better described as a sparse, physically constrained, black-box evolutionary attack.

\subsection{Input-Output Amplification Versus Optimization Objective}

Another important distinction is whether a method optimizes an input-output amplification ratio:
\[
    \frac{\text{output change}}{\text{input change}}.
\]
Most attacks do not directly optimize this ratio. They either optimize an output loss under an input constraint, or they combine an input penalty and an output loss additively.

For example, FGSM, PGD, PCFA, time-series regression attacks, and the neural-operator attack use objectives containing only output terms, while input size appears in constraints. Szegedy, Carlini--Wagner, VAE attacks, structured attacks, and steering-angle attacks place the input perturbation penalty directly inside the objective. DeepFool and SparseFool minimize input perturbation size subject to an output-change condition.

The neural-operator attack does report an input-output amplification ratio as an analysis metric:
\[
    \mathrm{Amplification}
    =
    \frac{
        \text{output percentage change}
    }{
        \text{input percentage change}
    }.
\]
However, this ratio is used for interpretation and evaluation, not as the main optimization objective. The main optimization objective is the relative output change normalized by the clean output magnitude.

\subsection{Summary of the Main Distinctions}

The methods can be grouped into the following structural classes.

First, some methods have only an output term in the objective, while input magnitude is handled only through a constraint:
\[
    \max_{\delta\in S}
    \mathcal{L}_{\mathrm{out}}(x+\delta).
\]
This class includes FGSM, BIM, PGD, original UAP, PCFA, time-series regression attacks, and the neural-operator digital-twin attack. These methods should not be described as having an input penalty in the loss.

Second, some methods have both an input term and an output term in the objective:
\[
    \min_\delta
    \mathrm{InputPenalty}(\delta)
    +
    \lambda \mathrm{OutputLoss}(x+\delta).
\]
This class includes Szegedy et al. L-BFGS, Carlini--Wagner, VAE adversarial images, structured adversarial attacks, and some steering-angle attacks.

Third, some methods minimize input perturbation size subject to an output-change condition:
\[
    \min_\delta \mathrm{InputSize}(\delta)
    \quad
    \text{s.t.}
    \quad
    \text{output changes}.
\]
This class includes DeepFool and SparseFool.

Fourth, singular-vector attacks are a narrower class. They require a local output-change objective of the form
\[
    \max_{\|v\|=1}
    \|J_f(x)v\|,
\]
or its multi-sample version. These are the methods for which power iteration, Lanczos, and SVD-based reasoning are directly applicable. The main representatives are the singular-vector UAP methods and their sparse variants.

Finally, spectral or Jacobian regularization methods are not attacks. They use related mathematics, such as
\[
    \sigma_{\max}(J_f)
    =
    \max_{\|v\|=1}
    \|J_fv\|,
\]
but their goal is to reduce sensitivity during training rather than generate adversarial perturbations.

\subsection{Conclusion}

The most important classification principle is the following. If the loss compares the perturbed output to the original output,
\[
    f(x+\delta)-f(x),
\]
then its local approximation can become
\[
    J_f(x)\delta,
\]
and the squared loss becomes a pure quadratic form:
\[
    \delta^\top J_f(x)^\top J_f(x)\delta.
\]
This is the setting in which singular-vector attacks and spectral methods are natural.

If the loss compares the perturbed output to a label, target, ground truth, or simulator output, such as
\[
    L(f(x+\delta),y),
    \qquad
    \|f(x+\delta)-y\|^2,
    \qquad
    \|f(x+\delta)-\tilde f\|^2,
    \qquad
    \|G(u+\delta)-g(u+\delta)\|^2,
\]
then the local expansion generally contains residual or linear terms, and the problem is not a pure singular-vector problem.

Moreover, whether the input perturbation appears in the optimization variable or constraint is not enough. One must distinguish whether the objective itself contains an input penalty. A constraint such as
\[
    \|\delta\|_\infty\le \epsilon
\]
does not mean that the loss contains an input term. Only expressions such as
\[
    \|\delta\|^2,
    \qquad
    \|\delta\|_0,
    \qquad
    \Omega(\delta)
\]
inside the objective count as explicit input terms in the loss.
```

\section{Evaluation Criteria, Input Constraints, and Success Thresholds}

The previous section classified adversarial attacks according to their loss structure, especially whether the objective compares \(f(x+\delta)\) with \(f(x)\), with a label \(y\), with a target output, or with a physical ground truth. This section further refines the classification from the viewpoint of experimental evaluation. In particular, we distinguish three related but different questions:
\begin{enumerate}
    \item Does the objective function itself contain an explicit input-perturbation penalty, such as \(\|\delta\|^2\), \(\|\delta\|_0\), or \(\Omega(\delta)\)?
    \item Does the method impose an input perturbation constraint, such as \(\|\delta\|_\infty \leq \epsilon\), \(\|\delta\|_2\leq \epsilon\), \(\|v\|_p\leq \xi\), or \(\|\delta\|_0\leq k\)?
    \item Does the paper define a concrete output-side success threshold, such as a classification change, a target class condition, a steering-angle threshold, or a relative error threshold?
\end{enumerate}

These distinctions are important because the presence of an input constraint does not mean that the loss contains an input term. For example,
\[
    \max_{\|\delta\|_\infty\leq \epsilon}
    L(f(x+\delta),y)
\]
has an input constraint, but the objective itself is only the output loss
\[
    L(f(x+\delta),y).
\]
Thus, \(\|\delta\|_\infty\leq \epsilon\) should not be counted as an input term inside the loss.

By contrast, the objective
\[
    \min_\delta
    \|\delta\|_2^2
    +
    cL(f(x+\delta),t)
\]
contains both an explicit input penalty,
\[
    \|\delta\|_2^2,
\]
and an output attack loss,
\[
    L(f(x+\delta),t).
\]

\subsection{Input Terms in the Objective versus Input Constraints}

We separate attacks into three categories.

\paragraph{Output-only objective with input constraint.}
These methods optimize an output-side loss while controlling the perturbation through a constraint:
\[
    \max_{\delta\in S}
    \mathcal{L}_{\mathrm{out}}(x+\delta).
\]
Typical examples include FGSM, BIM/I-FGSM, PGD, original UAP, PCFA, time-series regression attacks, and neural-operator digital-twin attacks. In these cases, the input budget is part of the feasible set, not part of the objective.

\paragraph{Objective containing both input and output terms.}
These methods explicitly combine an input penalty and an output attack loss:
\[
    \min_\delta
    \mathrm{InputPenalty}(\delta)
    +
    \lambda \mathrm{OutputLoss}(x+\delta).
\]
Examples include Szegedy et al. L-BFGS, Carlini--Wagner, structured adversarial attacks, VAE adversarial images, and some steering-angle attacks.

\paragraph{Input-minimization objective with output constraint.}
These methods minimize the input perturbation size subject to a successful output change:
\[
    \min_\delta \mathrm{InputSize}(\delta)
    \quad
    \text{s.t.}
    \quad
    \text{output changes}.
\]
DeepFool and SparseFool are representative examples:
\[
    \min_\delta \|\delta\|_2
    \quad
    \text{s.t.}
    \quad
    f(x+\delta)\neq f(x),
\]
and
\[
    \min_\delta \|\delta\|_0
    \quad
    \text{s.t.}
    \quad
    f(x+\delta)\neq f(x).
\]

\subsection{Classification of Methods by Objective Structure and Success Criteria}

\begin{table}[h]
\centering
\small
\begin{tabular}{p{3.2cm} p{2.5cm} p{3.5cm} p{3.5cm} p{2.7cm}}
\hline
\textbf{Method} & \textbf{Task Type} & \textbf{Input-Side Control} & \textbf{Output-Side Success Criterion} & \textbf{Input Term in Objective?} \\
\hline
Szegedy et al. L-BFGS
&
Classification
&
Minimizes \(\|\delta\|_2^2\); box constraint \(x+\delta\in[0,1]^d\)
&
Target class or misclassification
&
Yes
\\

FGSM
&
Classification
&
\(\|\delta\|_\infty\leq \epsilon\)
&
Classification error
&
No
\\

BIM / I-FGSM
&
Classification
&
\(\|\delta\|_\infty\leq \epsilon\), step size \(\alpha\), projection
&
Classification error
&
No
\\

Madry PGD
&
Classification
&
\(\delta\in S\), often \(\|\delta\|_\infty\leq \epsilon\)
&
Classification error; robust accuracy
&
No
\\

Carlini--Wagner
&
Classification
&
Minimizes \(\|\delta\|_p\); box constraint
&
Target logit margin with confidence \(\kappa\)
&
Yes
\\

DeepFool
&
Classification
&
Minimizes \(\|\delta\|_2\)
&
Crosses decision boundary
&
Yes; objective is input size
\\

SparseFool
&
Classification
&
Minimizes or controls \(\|\delta\|_0\)
&
Classification changes
&
Yes; objective is sparsity
\\

Original UAP
&
Classification
&
\(\|v\|_p\leq \xi\)
&
Fooling rate: prediction label changes
&
No
\\

Art of Singular Vectors and UAP
&
Classification / UAP
&
\(\|v\|_p=1\), later scaled
&
Fooling rate
&
No
\\

Sparse Universal Singular Vectors
&
Classification / UAP
&
\(\|v\|_p=1\), \(\|v\|_0\leq k\)
&
Fooling rate
&
No
\\

Adversarial Patch
&
Classification
&
Patch size, location, pixel range
&
Target class is predicted
&
Usually no explicit norm penalty
\\

Structured Adversarial Attack
&
Classification
&
\(L_p\) distortion and structured sparsity
&
Targeted or untargeted attack success rate
&
Yes
\\

VAE Adversarial Images
&
Reconstruction / generative model
&
Input distortion and pixel box constraints
&
Target reconstruction or target latent similarity
&
Yes
\\

Time-series regression attacks
&
Regression
&
Usually \(\|\delta\|\leq \epsilon\)
&
MSE/RMSE increase; usually no hard binary threshold
&
No
\\

Steering-angle attacks
&
Regression / control
&
Input \(L_2\) distortion or image perturbation
&
MSE increase or classification success if discretized
&
Often yes
\\

DeepBillboard
&
Continuous steering control
&
Physical billboard constraints
&
Steering-angle error exceeds \(19.8^\circ\)
&
Usually no simple norm penalty
\\

Attacking Optical Flow
&
Dense regression
&
Patch smaller than about \(1\%\) of image area
&
Flow degradation / EPE; usually no hard binary threshold
&
Usually no explicit norm penalty
\\

PCFA optical-flow attack
&
Dense regression
&
\(\|\delta\|_2\leq \epsilon\)
&
Target-flow deviation / robustness ranking
&
No
\\

Neural-operator digital-twin attack
&
Continuous physical-field prediction
&
\(\|\delta\|_0\leq k\), physical bounds; fewer than \(1\%\) inputs
&
Relative output error exceeds threshold, e.g. \(30\%\)
&
No
\\

Jacobian spectral norm regularization
&
Regularization
&
No adversarial input threshold
&
No attack success definition
&
No concrete perturbation penalty
\\
\hline
\end{tabular}
\caption{Evaluation-oriented classification of adversarial and robustness methods. The key distinction is whether input control appears inside the objective or only as a constraint, and whether a concrete success threshold is defined.}
\end{table}

\subsection{Success Criteria in Classification Tasks}

For classification attacks, the success criterion is usually discrete and does not require an additional continuous threshold. For an untargeted attack, success means
\[
    \arg\max f(x+\delta)\neq y.
\]
For a targeted attack, success means
\[
    \arg\max f(x+\delta)=t,
\]
where \(t\) is the target class.

For universal perturbations, the success metric is usually the fooling rate:
\[
    \mathrm{FR}(v)
    =
    \frac{1}{N}
    \sum_{i=1}^{N}
    \mathbf{1}
    \left\{
        \arg\max f(x_i+v)
        \neq
        \arg\max f(x_i)
    \right\}.
\]
Thus, for original UAP, singular-vector UAP, and sparse universal singular-vector attacks, the per-sample success condition is still a label change, while the reported metric is the fraction of samples whose predictions are changed.

Carlini--Wagner attacks introduce a more refined logit-margin condition. For a targeted attack, one can express the success condition as
\[
    Z_t(x+\delta)
    \geq
    \max_{i\neq t} Z_i(x+\delta)+\kappa,
\]
where \(Z_i\) denotes the logit for class \(i\), and \(\kappa\) is a confidence parameter. When \(\kappa=0\), the target logit only needs to exceed all other logits. When \(\kappa>0\), the target class must win by a larger margin.

\subsection{Success Criteria in Regression and Continuous-Output Tasks}

For regression or continuous-output models, success is not automatically defined by a label change. The output may be a scalar, a vector, a trajectory, an optical-flow field, or a physical field. Therefore, the paper must either report continuous error metrics or define a hard threshold such as
\[
    \mathrm{error}>\tau.
\]

Examples of continuous-output error metrics include
\[
    \mathrm{MSE},
    \qquad
    \mathrm{RMSE},
    \qquad
    \mathrm{EPE},
    \qquad
    \text{relative }L_2\text{ error}.
\]
Some papers report these metrics as continuous quantities rather than converting them into binary success/failure outcomes.

\subsection{Regression and Continuous-Output Methods Extracted from the Full List}

If we restrict attention to purely regression or continuous-output settings, we obtain the following subset:
\[
\begin{array}{ll}
\text{Time-series regression attacks} 
& \text{continuous numerical prediction}, \\[1mm]
\text{Self-driving steering-angle attacks} 
& \text{continuous steering-angle regression}, \\[1mm]
\text{DeepBillboard} 
& \text{continuous steering trajectory / physical driving control}, \\[1mm]
\text{Attacking Optical Flow} 
& \text{dense continuous optical-flow field}, \\[1mm]
\text{PCFA optical-flow attack} 
& \text{dense continuous optical-flow field}, \\[1mm]
\text{Neural-operator digital-twin attack} 
& \text{continuous PDE / physical-field prediction}, \\[1mm]
\text{Regression perturbation theory} 
& \text{theoretical regression-output sensitivity analysis}.
\end{array}
\]

VAE adversarial images can also be viewed as a continuous-output reconstruction attack, but it is better treated separately as an autoencoder or generative-model attack rather than a standard supervised regression problem.

\subsection{Regression and Continuous-Output Methods: Detailed Table}

\begin{table}[h]
\centering
\small
\begin{tabular}{p{3.3cm} p{2.8cm} p{3.4cm} p{3.4cm} p{2.5cm}}
\hline
\textbf{Method} & \textbf{Output Type} & \textbf{Attack Objective} & \textbf{Input Threshold / Budget} & \textbf{Hard Output Success Threshold?} \\
\hline
Time-series regression attacks
&
Continuous scalar or multivariate time-series prediction
&
\(\max \|f(x+\delta)-y\|^2\)
&
Usually \(\|\delta\|\leq \epsilon\)
&
Usually no; report MSE/RMSE increase
\\

Steering-angle attacks
&
Continuous steering angle
&
\(\min \|\delta\|_2^2-c(f(x+\delta)-y)^2\)
&
Input \(L_2\) distortion or image perturbation
&
Regression part usually reports MSE increase; classification-discretized part uses label change
\\

DeepBillboard
&
Continuous steering trajectory
&
Maximize cumulative steering deviation
&
Billboard region, physical position, printability, viewpoint
&
Yes; steering-angle error threshold \(19.8^\circ\)
\\

Attacking Optical Flow
&
Dense optical-flow vector field
&
Maximize flow degradation, EPE, or directional discrepancy
&
Patch area less than about \(1\%\) of image area
&
Usually no; report EPE / flow degradation
\\

PCFA optical-flow attack
&
Dense optical-flow vector field
&
\(\min \mathcal{L}(F(I+\delta),\tilde f)\)
&
\(\|\delta\|_2\leq \epsilon\)
&
Usually no; used for robustness ranking
\\

Neural-operator digital-twin attack
&
High-dimensional physical field
&
\(\max \dfrac{\|\mathcal{G}(u+\delta)-\mathcal{G}(u)\|_2}{\|\mathcal{G}(u)\|_2}\)
&
\(\|\delta\|_0\leq k\), physical bounds, fewer than \(1\%\) inputs
&
Yes; e.g. relative \(L_2\) error \(>30\%\)
\\

Regression perturbation theory
&
General continuous output
&
Original: \(\max \|f(x+\delta)-y\|^2\); approximation: \(\max \|J_f(x)\delta\|^2\)
&
\(\|\delta\|\leq \epsilon\)
&
Theoretical; usually no experimental hard threshold
\\
\hline
\end{tabular}
\caption{Regression and continuous-output attacks extracted from the broader list. These methods differ from classification attacks because success is not automatically defined by label change.}
\end{table}

\subsection{Time-Series Regression Attacks}

Time-series regression attacks usually adapt FGSM, BIM, or PGD-style methods to regression models such as CNNs, LSTMs, or GRUs. A typical objective is
\[
    \max_{\|\delta\|\leq \epsilon}
    \|f(x+\delta)-y\|_2^2.
\]
The input budget is usually a norm constraint such as
\[
    \|\delta\|\leq \epsilon.
\]
The objective itself contains only the regression error. The perturbation norm appears in the constraint, not inside the loss.

The output is continuous, so success is normally not defined by a label flip. Instead, the paper typically reports the increase in
\[
    \mathrm{MSE},
    \qquad
    \mathrm{RMSE},
    \qquad
    \mathrm{MAE},
\]
or related regression metrics. In many such papers, there is no hard threshold
\[
    \mathrm{RMSE}>\tau
\]
that converts the attack into a binary success/failure event.

\subsection{Steering-Angle Regression Attacks}

Steering-angle attacks target models that output a continuous steering angle
\[
    f(x)=\hat{\theta}.
\]
A representative regression-style objective is
\[
    \min_\delta
    \|\delta\|_2^2
    -
    c(f(x+\delta)-y)^2,
\]
where \(y\) is the true steering angle. This objective contains both an input perturbation penalty,
\[
    \|\delta\|_2^2,
\]
and an output error term,
\[
    (f(x+\delta)-y)^2.
\]

If the steering angle is discretized into classes, then attack success can be defined by classification error. In the pure regression setting, however, the common evaluation is an increase in MSE or steering-angle deviation. Unlike DeepBillboard, ordinary steering-angle regression attacks do not always define a universal hard safety threshold such as
\[
    |\hat{\theta}_{\mathrm{adv}}-\theta|>\tau.
\]

\subsection{DeepBillboard}

DeepBillboard attacks a physical driving-control system using an adversarial billboard. The output is a continuous steering trajectory rather than a class label. The input is constrained by physical conditions such as billboard location, printable pattern, viewing angle, distance, and illumination.

A key feature of DeepBillboard is that it defines a concrete output-side threshold. The steering-angle error threshold is
\[
    19.8^\circ.
\]
Thus, a dangerous frame or event can be identified when
\[
    |\Delta\theta|>19.8^\circ.
\]
The paper also uses metrics such as \(M_0\) and \(M_1\), where \(M_0\) measures the magnitude of steering deviation and \(M_1\) measures the frequency or duration with which the deviation exceeds the safety threshold.

This is an important example of a continuous-output attack where a hard numerical success line is explicitly defined.

\subsection{Attacking Optical Flow}

Optical-flow attacks target a dense continuous vector field:
\[
    F(I_1,I_2)(p)=(u_p,v_p),
\]
where \(p\) indexes pixels. The attack may optimize a patch to cause noisy flow, erased motion, directional errors, or increased endpoint error. The input constraint is often expressed as a patch-size constraint, for example a patch smaller than about \(1\%\) of the image area.

The output evaluation is typically continuous:
\[
    \mathrm{EPE},
    \qquad
    \text{flow degradation},
    \qquad
    \text{visual corruption of the flow field}.
\]
There is usually no universal hard threshold
\[
    \mathrm{EPE}>\tau
\]
that defines binary attack success. Therefore, this line of work is better understood as measuring attack strength through continuous degradation metrics rather than through a classification-like success rate.

\subsection{PCFA Optical-Flow Attack}

PCFA attacks optical-flow models by pushing the predicted flow toward a specified target flow \(\tilde f\):
\[
    \min_{\|\delta\|_2\leq \epsilon}
    \mathcal{L}(F(I+\delta),\tilde f).
\]
The input threshold is explicit:
\[
    \|\delta\|_2\leq \epsilon.
\]
However, the output side is generally evaluated by how far the prediction is shifted toward the target flow, how much optical-flow error is induced, or how the model ranks under adversarial robustness evaluation. PCFA is usually not framed as a binary success/failure attack with a universal threshold
\[
    \mathcal{L}(F(I+\delta),\tilde f)>\tau.
\]
Instead, it is used to compare robustness across optical-flow methods under a fixed input perturbation budget.

\subsection{Neural-Operator Digital-Twin Attack}

The neural-operator digital-twin attack targets a neural operator
\[
    \mathcal{G}_\theta
\]
that maps a physical input vector \(u\) to a high-dimensional physical field:
\[
    \mathcal{G}_\theta(u).
\]
The attack compares the perturbed prediction and the clean prediction:
\[
    \mathcal{G}_\theta(u+\delta)
    -
    \mathcal{G}_\theta(u).
\]
The attack score is a relative output change:
\[
    \mathrm{Score}(u,\delta)
    =
    \frac{
        \|\mathcal{G}_\theta(u+\delta)-\mathcal{G}_\theta(u)\|_2
    }{
        \|\mathcal{G}_\theta(u)\|_2
    }.
\]
The denominator is the clean output norm, not the input perturbation norm. Thus, the optimization objective is not an input-output amplification ratio of the form
\[
    \frac{\text{output change}}{\text{input change}}.
\]

The input threshold is sparse and physical:
\[
    \|\delta\|_0\leq k,
\]
together with physical bounds on modified input coordinates. If the input dimension is \(102\), then \(k=1\) means
\[
    \frac{1}{102}\approx 0.98\%,
\]
so the attack modifies fewer than \(1\%\) of the input coordinates.

Unlike many continuous-output attacks, this paper uses a hard relative-error threshold to define attack success. A typical success condition can be written as
\[
    \frac{
        \|\mathcal{G}_\theta(u+\delta)-\mathcal{G}_\theta(u)\|_2
    }{
        \|\mathcal{G}_\theta(u)\|_2
    }
    >
    \tau.
\]
For example, one may use
\[
    \tau=30\%.
\]
The success rate is then the fraction of test cases for which the relative output change exceeds this threshold:
\[
    \mathrm{ASR}_{\tau}
    =
    \frac{
        \#\left\{
            u_i:
            \mathrm{Score}(u_i,\delta_i)>\tau
        \right\}
    }{
        \#\{\text{test inputs}\}
    }.
\]

The optimization method is gradient-free differential evolution. A candidate attack is encoded by index-value pairs:
\[
    (i_1,a_1,\ldots,i_k,a_k),
\]
where \(i_j\) is the input coordinate to modify and \(a_j\) is the replacement value. Each candidate is evaluated only by forward passes through the neural operator. Differential evolution updates the candidate population through mutation, crossover, and selection. It does not use gradients, Jacobian-vector products, SVD, or power iteration.

Therefore, the neural-operator attack is output-change based, sparse, physically constrained, and black-box. It is not a pure quadratic singular-vector attack.

\subsection{Regression Perturbation Theory}

For a general regression model, one may start from the supervised regression objective
\[
    \max_{\|\delta\|\leq \epsilon}
    \|f(x+\delta)-y\|_2^2.
\]
This is not a pure quadratic form in \(\delta\), because
\[
    f(x+\delta)-y
    \approx
    f(x)-y+J_f(x)\delta,
\]
so
\[
    \|f(x+\delta)-y\|_2^2
    \approx
    \|f(x)-y\|_2^2
    +
    2(f(x)-y)^\top J_f(x)\delta
    +
    \delta^\top J_f(x)^\top J_f(x)\delta.
\]
The linear term
\[
    2(f(x)-y)^\top J_f(x)\delta
\]
prevents this from being a pure Rayleigh quotient.

However, if the model is already accurate at \(x\), so that
\[
    f(x)\approx y,
\]
then
\[
    f(x+\delta)-y
    \approx
    f(x+\delta)-f(x)
    \approx
    J_f(x)\delta.
\]
The problem then reduces approximately to
\[
    \max_{\|\delta\|\leq \epsilon}
    \|J_f(x)\delta\|_2^2.
\]
Under an \(\ell_2\) constraint, the optimal perturbation direction is the dominant right singular vector of \(J_f(x)\):
\[
    \delta^\star
    =
    \pm \epsilon v_{\max},
\]
where \(v_{\max}\) is the dominant eigenvector of
\[
    J_f(x)^\top J_f(x).
\]
Thus, regression perturbation theory becomes a singular-vector problem only under the approximation \(f(x)\approx y\) and typically under an \(\ell_2\) perturbation constraint.

\subsection{Which Regression Methods Define a Hard Output Threshold?}

Among the continuous-output methods discussed above, the clearest cases with hard output-side thresholds are:
\[
    \text{DeepBillboard}
\]
and
\[
    \text{neural-operator digital-twin attacks}.
\]

For DeepBillboard, the threshold is a steering-angle error:
\[
    |\Delta\theta|>19.8^\circ.
\]
For the neural-operator attack, the threshold is a relative output error:
\[
    \frac{
        \|\mathcal{G}(u+\delta)-\mathcal{G}(u)\|_2
    }{
        \|\mathcal{G}(u)\|_2
    }
    >
    \tau,
\]
for example
\[
    \tau=30\%.
\]

By contrast, time-series regression attacks, ordinary steering-angle regression attacks, optical-flow patch attacks, and PCFA usually report continuous degradation metrics such as
\[
    \mathrm{MSE},
    \qquad
    \mathrm{RMSE},
    \qquad
    \mathrm{EPE},
    \qquad
    \text{relative error},
\]
rather than defining a universal binary success threshold.

\subsection{Final Summary}

For classification tasks, attack success is usually easy to define:
\[
    \arg\max f(x+\delta)\neq y
\]
for untargeted attacks, or
\[
    \arg\max f(x+\delta)=t
\]
for targeted attacks. The success rate is simply the fraction of samples satisfying the classification success condition.

For regression and continuous-output tasks, success is not automatic. The output is a number, a trajectory, a vector field, or a physical field. Therefore, a paper must either report continuous error metrics or define a numerical threshold:
\[
    \mathrm{error}>\tau.
\]
Some continuous-output papers, such as DeepBillboard and neural-operator digital-twin attacks, do define such thresholds. Others, such as time-series regression attacks, optical-flow attacks, and PCFA, mainly report continuous degradation metrics without converting them into a universal binary success/failure definition.

Finally, the presence of an input perturbation constraint does not mean the loss contains an input term. A constraint such as
\[
    \|\delta\|_\infty\leq \epsilon
\]
or
\[
    \|\delta\|_0\leq k
\]
controls the feasible perturbation set, but it is not an input penalty inside the loss. Only terms such as
\[
    \|\delta\|^2,
    \qquad
    \|\delta\|_0,
    \qquad
    \Omega(\delta)
\]
appearing directly inside the objective should be counted as explicit input terms in the loss.





\end{document}
